% ============================================================================
% Thesis: Confidence-Aware Episodic Memory for Hallucination Mitigation
% Author: Aksan Gony Alif
% CSE400, BRAC University, 2026
%
% This is the active thesis report. The earlier draft at
% pre thesis 1 report/ is preserved as a reference baseline; do not
% edit it during the rewrite.
%
% Build instructions: see BUILD.md.
% Style rules: see audit_phase_a.md and the per-chapter scope notes.
% ============================================================================

\documentclass[12pt,a4paper,oneside,openany]{report}

% --- Page geometry and typography -------------------------------------------
\usepackage[a4paper,width=150mm,top=25mm,bottom=25mm]{geometry}  % BRAC template
\usepackage[utf8]{inputenc}
\usepackage[T1]{fontenc}
\usepackage{lmodern}
\usepackage{microtype}
\usepackage{setspace}
% single spacing, as in the BRAC template (no \onehalfspacing)

% --- Math and symbols -------------------------------------------------------
\usepackage{amsmath}
\usepackage{amssymb}
\usepackage{amsthm}
\usepackage{bm}
\usepackage{mathtools}
\usepackage{nicefrac}

% --- Tables and figures -----------------------------------------------------
\usepackage{graphicx}
\graphicspath{{./}{./images/}{./figures/}}
\usepackage{booktabs}
\usepackage{multirow}
\usepackage{multicol}
\usepackage{array}
\usepackage{tabularx}
\usepackage{longtable}
\usepackage{enumitem}
\usepackage{rotating}
\usepackage{float}
\usepackage{placeins}
\usepackage{needspace}
\usepackage{caption}
\usepackage{subcaption}
\usepackage[table]{xcolor}
\usepackage{tcolorbox}
\usepackage{tikz}
\usetikzlibrary{positioning,arrows,arrows.meta,shapes,shapes.geometric,fit,backgrounds,calc,decorations.pathreplacing,shadows.blur}
\usepackage{pgfplots}
\pgfplotsset{compat=1.18}

% --- Algorithms -------------------------------------------------------------
\usepackage{algorithm}
\usepackage{algpseudocode}

% --- Cross-references and hyperlinks ----------------------------------------
\usepackage[colorlinks=true,linkcolor=black,filecolor=magenta,citecolor=black,urlcolor=black]{hyperref}  % BRAC template colours
\urlstyle{same}
\usepackage[noabbrev,capitalize]{cleveref}

% --- Bibliography -----------------------------------------------------------
\usepackage[backend=biber,style=ieee,sorting=ynt]{biblatex}  % BRAC template
\addbibresource{bibliography/references.bib}
% Render-time backstop: suppress any field that might still leak thesis-internal
% commentary or non-bibliographic clutter into the rendered bibliography.
\AtEveryBibitem{%
  \clearfield{note}%
  \clearfield{abstract}%
  \clearfield{annotation}%
  \clearfield{addendum}%
  \clearfield{keywords}%
  \clearfield{issn}%
  \clearfield{isbn}%
}

% --- Code-path typesetting --------------------------------------------------
% \codepath breaks long monospace identifiers (file paths, function names,
% config keys) at any character so unbreakable underscores and slashes do
% not cause overfull-hbox bleeds in narrow appendix columns and prose. The
% sloppy + emergencystretch combination is applied at the paragraph level
% via an \appendixoverflowsafe declaration loaded inside each appendix.
\usepackage{seqsplit}
\newcommand{\codepath}[1]{\texttt{\seqsplit{#1}}}
% Apply once per appendix paragraph that contains long monospace identifiers
% to license aggressive whitespace stretching as a last-resort fallback.
\newcommand{\appendixoverflowsafe}{%
  \sloppy
  \emergencystretch=5em
  \hbadness=10000
}

% --- Lists and indices ------------------------------------------------------
\usepackage[intoc]{nomencl}
\makenomenclature
\usepackage{tocloft}

% --- Custom commands --------------------------------------------------------
\newcommand{\caem}{\textsc{CAEM}}
\newcommand{\upre}{u_{\text{pre}}}
\newcommand{\ustored}{u_{\text{stored}}}
\newcommand{\taustore}{\tau_{\text{store}}}
\newcommand{\taudefer}{\tau_{\text{defer}}}
\newcommand{\tautrain}{\tau_{\text{train}}}
\newcommand{\alphastore}{\alpha_{\text{store}}}
\newcommand{\alphadefer}{\alpha_{\text{defer}}}

% --- Theorem environments ---------------------------------------------------
\theoremstyle{plain}
\newtheorem{theorem}{Theorem}[chapter]
\newtheorem{lemma}[theorem]{Lemma}
\newtheorem{corollary}[theorem]{Corollary}
\newtheorem{proposition}[theorem]{Proposition}
\theoremstyle{definition}
\newtheorem{definition}[theorem]{Definition}
\newtheorem{assumption}[theorem]{Assumption}
\theoremstyle{remark}
\newtheorem{remark}[theorem]{Remark}
\theoremstyle{definition}
\newtheorem{heuristic}[theorem]{Heuristic}
\crefname{heuristic}{Heuristic}{Heuristics}
\Crefname{heuristic}{Heuristic}{Heuristics}
% Argument sketches for heuristics end without the end-of-proof box, which is kept for proofs.
\newenvironment{sketch}{\begin{proof}[Argument sketch]\renewcommand{\qedsymbol}{}}{\end{proof}}

% --- BRAC template layout controls ------------------------------------------
\setlength{\parindent}{0em}
% The template sets no paragraph spacing; with zero indent, paragraphs need a visible gap.
\usepackage[skip=0.5\baselineskip plus 2pt]{parskip}
\pagestyle{plain}
\let\cleardoublepage=\clearpage

% ============================================================================
\begin{document}

% --- Front matter -----------------------------------------------------------
\thispagestyle{empty}
\begin{titlepage}
\renewcommand*{\thepage}{Title} % Change page number in PDF

    \begin{center} 
        \vspace*{3cm} % For creating Vertical Blank Space
        
        {\fontsize{16pt}{22pt}\selectfont{Confidence-Aware Episodic Memory with Self-Improvement for Hallucination Reduction in Large Language Models}
        } % "fontsize{font size}{line space}\selectfont{}" command to override font size and line space for the Title
        
        \vspace{1.5cm}
        
        \text{by}
        
        \vspace{0.5cm}
        
        	Md. Aksan Gony Alif\\
	        24341256\\
            Maliha Binte Shamim\\
	        23101466\\
	        Md. Mahmudur Rahman\\
	        22101355


        \vspace{1.5cm}
        
        	A thesis submitted to the Department of Computer Science and Engineering\\
            in partial fulfillment of the requirements for the degree of\\
            B.Sc. in Computer Science and Engineering

        
        \vspace{2.5cm}
        
    		Department of Computer Science and Engineering\\
            Brac University\\
            Fall 2026.
        
        \vspace{3cm}
        
    		\copyright\ 2026. Brac University\\
            All rights reserved.
    
    \end{center}

\end{titlepage}
\cleardoublepage
\pagenumbering{roman}
\phantomsection
\addcontentsline{toc}{chapter}{Declaration}
% Following command is used to create grouped signature line for Three Authors
\newcommand*\wildcard[2][6cm]{\vspace{2cm}\parbox{#1}{\hrulefill\par#2}}

% A "parbox{}{}" is a box whose contents are created in paragraph mode.
% "hrulefill{}" sets the underline thickness.

\section*{Declaration}

It is hereby declared that

\begin{enumerate} % begin{enumerate} function to create numbered list
  \item The thesis submitted is my/our own original work while completing degree at Brac University.
  \item The thesis does not contain material previously published or written by a third party, except where this is appropriately cited through full and accurate referencing.
  \item The thesis does not contain material which has been accepted, or submitted, for any other degree or diploma at a university or other institution.
  \item We have acknowledged all main sources of help.
\end{enumerate}

\vspace{1cm}
\textbf{Student's Full Name \& Signature:}

\begingroup

    \begin{center}
        \parbox{6.5cm}{\centering\includegraphics[width=6.5cm]{alif.png}\\\hrule\vspace{0.2cm}Md. Aksan Gony Alif\\ID: 24341256}
        \hspace{1cm}
        \parbox{5cm}{\centering\includegraphics[width=4cm]{{maliha sign}.png}\\\hrule\vspace{0.2cm}Maliha Binte Shamim\\ID: 23101466}
        \hspace{1cm}
        \parbox{5cm}{\centering\includegraphics[width=4cm]{mahmud.png}\\\hrule\vspace{0.2cm}Md. Mahmudur Rahman\\ID: 22101355}
    \end{center}

\endgroup


\pagebreak






\phantomsection
\addcontentsline{toc}{chapter}{Approval}
\section*{Approval}

The thesis/project titled “Confidence-Aware Episodic Memory with Self-Improvement
for Progressive Hallucination Reduction in Large Language
Models” submitted by 
\begin{enumerate}
  \item Md. Aksan Gony Alif (24341256)
  \item Maliha Binte Shamim (23101466)
  \item Md. Mahmudur Rahman (22101355) 
\end{enumerate}

of Fall, 2025 has been accepted as satisfactory in partial fulfillment of the requirement for the degree of B.Sc. in Computer Science in Spring 2026. 

\vspace{0.5cm}
\textbf{Examining Committee:}

\vspace{1cm}

\noindent Supervisor:\\
(Member)
\begin{center}
    \hspace{7cm} \wildcard{\centerline{Dr. Md. Golam Rabiul Alam} ~\\ \centerline{Professor}~\\ \centerline{Department of Computer Science \& Engineering}~\\ \centerline{Brac University} } \hspace{1cm}
\end{center}

\noindent Co-Supervisor:\\
(Member)
\begin{center}
    \hspace{7cm} \wildcard{\centerline{Md. Tanzim Reza} ~\\ \centerline{Senior Lecturer}~\\ \centerline{Department of Computer Science \& Engineering}~\\ \centerline{Brac University} } \hspace{1cm}
\end{center}


\noindent Thesis Coordinator:\\
(Member)
\begin{center}
    \hspace{7cm} \wildcard{\centerline{Dr. Md. Golam Rabiul Alam} ~\\ \centerline{Professor}~\\ \centerline{Department of Computer Science \& Engineering}~\\ \centerline{Brac University} } \hspace{1cm}
\end{center}

\noindent Head of Department:\\
(Chair)
\begin{center}
    \hspace{7cm} \wildcard{\centerline{Dr. Sadia Hamid Kazi} ~\\ \centerline{Associate Professor \& Chairperson}~\\ \centerline{Department of Computer Science \& Engineering}~\\ \centerline{Brac University} } \hspace{1cm}
\end{center}

\pagebreak
\input{core/ethics_statement}      % optional page; empty, so no TOC entry
\phantomsection
\addcontentsline{toc}{chapter}{Abstract}
% ============================================================================
% Abstract — rewritten 2026-10-19 on the final numbers (baseline re-run under
% the comparison protocol, one scoring rule, TruthfulQA judge v2).
% Every number traces to outputs/analysis/chm_judged.json,
% outputs/analysis/paired_baseline_tests_final.json and the chapter-five tables
% built by scripts/make_final_results.sh.
% ============================================================================

\section*{Abstract}

Large language models give confident but false answers, and a user cannot tell from the answer which ones to trust: asked for supporting literature, chatbots have returned references that do not exist \cite{alkaissi2023artificial}, for up to $91.4$ percent of the references they generate in systematic-review searches \cite{info:doi/10.2196/53164}. This thesis proposes \emph{Confidence-Aware Episodic Memory with Self-Improvement (CAEM)}, a system built around a small open generator. A router sends each question to a verified memory, to zero-shot generation or to retrieval-augmented generation, with a classifier that decides when retrieval is likely to help. A verifier combines nine signals of internal confidence, agreement across sampled answers, support from retrieved passages and question match into a calibrated probability of correctness, and a four-outcome decision stores, defers, abstains on or discards each answer. Every answer reaches the user with that probability and a status tag (verified, provisional, conflicting or insufficient evidence). Verified answers are kept in memory, re-checked at every cycle, and used to fine-tune a low-rank adapter, with a retention guard that rolls the adapter back when held-out accuracy falls. Two results are proven: storage makes the stored pool purer than the generator's raw output exactly when the verifier's balanced accuracy exceeds one half, and the pool's precision is bounded below by the calibration precision when new questions resemble the calibration data.

CAEM was run for six cycles with Qwen-2.5-3B-Instruct on five benchmarks (FEVER, TriviaQA and CommonsenseQA for training; TruthfulQA and StrategyQA held out) and compared with eight baselines under one protocol: the same questions, the same prompt, one scoring rule and the same hallucination instrument. The architecture alone added almost no accuracy (pooled exact match $0.517$ against $0.511$ for the zero-shot baseline); self-improvement raised it to $0.535$ after two cycles, after which it eroded to $0.508$, and no difference from the zero-shot baseline is significant on any benchmark. CAEM is level with strong prompting baselines, significantly better than the retrieval baselines on several benchmarks, and level on accuracy with a fine-tuned baseline that used about eleven times as many gold labels, which reached a lower hallucination metric. Its measured contribution is selectivity: the status tag predicts correctness, with verified answers right $65$ to $75$ percent of the time against $51$ percent for the zero-shot baseline. The gains are concentrated where the verifier can check answers against retrieved passages, and after self-improvement the verifier scored more of the wrong answers highly (confident confabulation $0.305$ against $0.208$). Of four pre-specified hypotheses, two are supported, one partially and one not.

\vspace{0.5cm}
\noindent\textbf{Keywords:} hallucination; episodic memory; calibrated confidence; selective prediction; verifier; self-improvement; retrieval-augmented generation; low-rank adaptation; large language models.
\pagebreak

\input{core/dedication}            % optional page; empty, so no TOC entry
\phantomsection
\addcontentsline{toc}{chapter}{Acknowledgment}
% DRAFT (2026-10-16): written for the authors to personalise; replace or extend freely.
\section*{Acknowledgement}
We are deeply grateful to our supervisor, Dr. Md. Golam Rabiul Alam, Professor, Department of Computer Science and Engineering, Brac University, for his guidance, encouragement, and critical feedback throughout this thesis. We also thank our co-supervisor, Md. Tanzim Reza, Senior Lecturer, Department of Computer Science and Engineering, Brac University, for his support and advice at every stage of the work.

We thank the Department of Computer Science and Engineering, Brac University, for the opportunity and resources to carry out this research.

Finally, we thank our families for their constant support and patience during our studies.
\pagebreak


\renewcommand{\contentsname}{Table of Contents}
\cleardoublepage
\phantomsection
\addcontentsline{toc}{chapter}{Table of Contents}
\tableofcontents
\cleardoublepage
\phantomsection
\addcontentsline{toc}{chapter}{\listfigurename}
\listoffigures
\cleardoublepage
\phantomsection
\addcontentsline{toc}{chapter}{\listtablename}
\listoftables
\cleardoublepage
\phantomsection
\addcontentsline{toc}{chapter}{List of Algorithms}
\listofalgorithms
% ============================================================================
% List of Acronyms — single table, two-column, footnotesize.
% ============================================================================

\chapter*{List of Acronyms}
\addcontentsline{toc}{chapter}{List of Acronyms}

{\footnotesize
\setlength{\tabcolsep}{4pt}
\renewcommand{\arraystretch}{1.05}
\begin{longtable}{@{}p{1.7cm}p{5.9cm}@{\hspace{1.2em}}p{1.7cm}p{5.9cm}@{}}
\multicolumn{4}{@{}l@{}}{\textbf{\normalsize Architecture and methods}}\\
\textbf{CAEM}    & Confidence-Aware Episodic Memory with Self-Improvement & \textbf{SIL}     & Self-Improvement Loop \\
\textbf{RUC}     & Retrieval-Utility Classifier (Stage 3b) & \textbf{NLI}     & Natural Language Inference \\
\textbf{RAG}     & Retrieval-Augmented Generation & \textbf{DPR}     & Dense Passage Retrieval \\
\textbf{CoT}     & Chain-of-Thought prompting & \textbf{EWC}     & Elastic Weight Consolidation \\
\textbf{LoRA}    & Low-Rank Adaptation & \textbf{MC}      & Monte Carlo (MC-Dropout) \\
\textbf{KLE}     & Kernel Language Entropy & \textbf{STaR}    & Self-Taught Reasoner (\textbf{V-STaR}: verifier) \\
\textbf{NTM}     & Neural Turing Machine & \textbf{DNC}     & Differentiable Neural Computer \\
\textbf{A-MEM}   & Agentic Memory & \textbf{FActScore} & Atomic-fact factuality score \\
\textbf{FLARE}   & Forward-Looking Active Retrieval & \textbf{Self-RAG} & Self-reflective RAG \\
\textbf{SI / MAS} & Synaptic Intelligence; Memory-Aware Synapses & \textbf{CP}      & Conformal Prediction (\textbf{split-CP}; registered then replaced by fixed-threshold storage rule) \\
\\[-0.5em]
\multicolumn{4}{@{}l@{}}{\textbf{\normalsize Models and judges}}\\
\textbf{LLM}     & Large Language Model & \textbf{SBERT}   & Sentence-BERT (cosine encoder) \\
\textbf{MiniCheck} & Flan-T5-Large NLI judge for claim-support & \textbf{MNLI / SNLI} & Multi-genre / Stanford NLI corpora \\
\\[-0.5em]
\multicolumn{4}{@{}l@{}}{\textbf{\normalsize Benchmarks}}\\
\textbf{FEVER}   & Fact Extraction and VERification & \textbf{NQ}      & Natural Questions (dropped at cycle zero; RUC training only) \\
\textbf{TQA}     & TriviaQA & \textbf{TruthfulQA} & Truthful QA (Haiku-judged EM) \\
\textbf{CSQA}    & CommonsenseQA & \textbf{StrategyQA} & Multi-hop strategy QA \\
\textbf{HotpotQA} & Multi-hop QA (dropped at cycle zero) & & \\
\textbf{HaluEval-QA} & Hallucination-Evaluation QA (RUC training only) & \textbf{OpenBookQA} & Open-Book QA (RUC training only) \\
\textbf{MMLU}    & Massive Multitask Language Understanding (retention probe) & \textbf{HaluEval} & Hallucination Evaluation benchmark \\
\textbf{BIG-bench} & Beyond the Imitation Game & & \\
\textbf{SQuAD}   & Stanford Question Answering Dataset & & \\
\\[-0.5em]
\multicolumn{4}{@{}l@{}}{\textbf{\normalsize Metrics, statistics, infrastructure}}\\
\textbf{EM}      & Exact Match & \textbf{F$_1$}   & Token-level F$_1$ \\
\textbf{ROUGE}   & Recall-Oriented Understudy for Gisting Evaluation & \textbf{CHM}     & Composite Hallucination Metric \\
\textbf{CES}     & Composite Evaluation Score (geometric mean of ACC, EPI, RET, CAL, VER; \textbf{EPI} $=1-\mathrm{CHM}$) & \textbf{ECE}     & Expected Calibration Error \\
\textbf{AUROC}   & Area Under ROC curve & \textbf{TPR / FPR / TNR} & True / False / True-Negative Rate \\
\textbf{TPR}     & True-Positive Rate & \textbf{CI}      & Confidence Interval \\
\textbf{ID / OOD} & In- / Out-of-Distribution & \textbf{EMA}     & Exponential Moving Average \\
\textbf{FPR}     & False-Positive Rate & \textbf{FAISS}   & Facebook AI Similarity Search (\textbf{IVF}: inverted-file) \\
\textbf{TTL}     & Time To Live (deferred-buffer expiry) & \textbf{QA}      & Question Answering \\
\end{longtable}
}

% ============================================================================
% List of Symbols — single table, two-column, footnotesize.
% ============================================================================

\chapter*{List of Symbols}
\addcontentsline{toc}{chapter}{List of Symbols}

{\footnotesize
\setlength{\tabcolsep}{3pt}
\renewcommand{\arraystretch}{1.05}
\begin{longtable}{@{}>{\raggedright\arraybackslash}p{0.17\textwidth}>{\raggedright\arraybackslash}p{0.30\textwidth}@{\hspace{0.5em}}>{\raggedright\arraybackslash}p{0.17\textwidth}>{\raggedright\arraybackslash}p{0.30\textwidth}@{}}
\multicolumn{4}{@{}l@{}}{\textbf{\normalsize Confidence, verifier signals, thresholds}}\\
$\upre$ & Pre-routing confidence & $\ustored$ & Calibrated post-generation composite \\
$u_{\text{token}}$ & Mean token log-probability & $u_{\text{dropout}}$ & Dropout disagreement share, $K{=}5$ \\
$u_{\text{internal}}$ & $0.5u_{\text{tok}}{+}0.5(1{-}u_{\text{drop}})$ & $s_{\text{avg}}$ & Mean pairwise cosine, $M{=}3$ \\
$p_{\text{entail}}$ & Chain-to-answer NLI entailment & $h_{\text{norm}}$ & Normalised semantic entropy (dropped), $K{=}10$ \\
$p_{\text{ground,max}}$ & Top-$1$ passage entailment & $p_{\text{ground,mean}}$ & Mean entailment over top-$3$ \\
$p_{\text{ground,atomic}}$ & Length-gated atomic-fact entailment & $q_{\text{match}}$ & Question-answer match score \\
\multicolumn{4}{@{}p{\textwidth}@{}}{\footnotesize $h_{\text{norm}}$ (normalised semantic entropy, $K{=}10$) was dropped at cycle zero (learned weight $\approx -5\times10^{-4}$, largest verifier cost); the deployed verifier uses the other nine signals.} \\
$\taustore,\taudefer$ & Fixed store / defer cut-points on $\ustored$ ($0.60$ / $0.45$) & $\tautrain$ & Strict training-pool threshold ($0.70$) \\
$\tau_{\text{retro}}$ & Retroactive prune threshold ($0.50$) & $\tau_{\text{ttl}}$ & Deferred-buffer time-to-live ($2$ passes) \\
$\tau_{\text{cons}}$ & Memory consolidation cosine ($0.92$) & $\phi_{\text{pg}}$ & Abstention floor on top-passage entailment ($0.20$) \\
$\phi_{\text{safe}}$ & Pre-routing safety floor on $\upre$ ($0.38$; per-benchmark refit) & $\theta_{1},\theta_{\text{sim}}$ & Tier 1 combined-score / Tier 2 similarity thresholds ($0.90$ / $0.75$) \\
$\eta_u,\eta_g$ & Early-exit thresholds & $\tau_{\text{RUC}}$ & Retrieval-utility classifier threshold ($0.375$) \\
\\[-0.5em]
\multicolumn{4}{@{}l@{}}{\textbf{\normalsize Calibrated probability composite}}\\
\multicolumn{4}{@{}p{\textwidth}@{}}{\footnotesize $\alphastore,\alphadefer$ (split-CP miscoverage targets) belonged to the conformal storage gate and were retired when the gate was replaced by the fixed cut-points above (\Cref{sec:conformal-gate}). The composite refit follows the drift of the signals by re-scoring the same calibration fold at every cycle; $\alpha_{\text{ema}} = 0.7$ smooths only the pre-routing temperature.} \\
$n_{\min}$ & Minimum top-$k$ for fit ($20$) & $n_{\text{cal}},n_{\text{eval}}$ & Calibration / evaluation fold sizes \\
$\pi_{\text{store}},\pi_{\text{defer}}$ & Calibration-fold empirical precision & $\pi_{\text{retro}}$ & Retroactive-survivor precision \\
$\pi_t,\pi_\infty$ & Pool purity at cycle $t$, asymptote & $P_c,P_c^{\text{theory}}$ & Empirical / Bayes pool purity \\
$P_j,\mathrm{iso}_j$ & Per-signal isotonic probability & $L_{i,j}$ & Calibration log-odds matrix \\
$\mathbf{w},b,w_j$ & Boost weights, intercept & $C$ & Boost L$_2$ strength ($0.01$) \\
$T$ & Logit-space temperature & $\sigma(\cdot),\ell$ & Logistic; composite log-odds \\
$\alpha_{\text{shrink}}$ & Per-benchmark shrinkage weight ($0.6$) & $c_{\text{conv}}$ & Layer-convergence ratio in $\upre$ \\
\\[-0.5em]
\multicolumn{4}{@{}l@{}}{\textbf{\normalsize Self-improvement, router, memory}}\\
$\boldsymbol{\theta},\boldsymbol{\theta}^{(t)}$ & Generator parameters at cycle $t$ & $\boldsymbol{\theta}^\star$ & Pre-retention-guard argmin \\
$\lambda_{\text{anc}}$ & L$_2$ anchor strength ($0$ on the adapter path) & $\mathcal{L}_{\mathrm{CE}},\mathcal{L}$ & CE loss; total per-cycle loss \\
$\rho,\rho_{\min}$ & Retention ratio; floor ($0.93$) & $\rho_{\text{pre/post}}$ & Pre / post fine-tune probe accuracy \\
$\mathcal{T}$ & Training pool above $\tautrain$ & $\mathcal{M},\mathcal{M}_t$ & Episodic memory; pool at cycle $t$ \\
$\mathcal{M}_t^{\text{kept}},\mathcal{N}_{t+1}$ & Survivors; new admissions & $\mathcal{B},K_{\text{def}}$ & Deferred buffer; capacity \\
$\lambda$ & Routing similarity weight ($0.70$) & $s,\mathcal{S}$ & Cosine; dispatch score $\lambda s{+}(1{-}\lambda)\ustored$ \\
$\tau_1(t)$ & Tier-$1$ share of queries & $H_{1}(t)$ & Tier-$1$ error rate \\
$H_{\text{mem}}(t),H_{\text{gen}}(t)$,\newline $H_t$ & Memory / generated / overall error rate & $M,K$ & Chain count ($3$); answer count ($10$) \\
$p_{\text{RAG}}$ & RUC retrieval-helpfulness score & $p_{\text{IK}}$ & Self-knowledge probe (Kadavath-style) \\
\\[-0.5em]
\multicolumn{4}{@{}l@{}}{\textbf{\normalsize Theorem trajectory and statistics}}\\
$d,d_t$ & Cohen's $d$ at cycle $t$ & $d_{\text{id}},d_{\text{transfer}}$ & ID and transfer Cohen's $d$ \\
$\mu_+,\mu_-,\varsigma$ & Within-class means; shared std & $p_+,p_-$ & Correct / incorrect base rates \\
$\Phi(\cdot),\phi(\cdot)$ & Std-normal CDF; density & $\sigma_t$ & Empirical storage rate at cycle $t$ \\
$r$ & Geometric contraction rate & $G_t$ & Cycle-$t$ gap $|\pi_t{-}\pi_\infty|$ \\
$A,k,c^{*}(\varepsilon)$ & Decay fit; time-to-equilibrium & $p_c,\alpha_c$ & Base correctness; verifier balanced acc. \\
$\mathrm{TPR}_c,\mathrm{FPR}_c$,\newline $\mathrm{TNR}_c$ & Storage-gate rates at cycle $c$ & $f(p)$ & Recurrence $p\alpha/(p\alpha{+}(1{-}p)(1{-}\alpha))$ \\
$K(C)$ & Questions supportable by corpus $C$ & $\chi^{2}_{\text{McN}},b,c$ & McNemar; discordant pair counts \\
$B$ & Bootstrap replicates ($2{,}000$) & $\alpha$ & Significance level (overloaded) \\
$\mathrm{EM},\mathrm{F}_1$ & Per-task correctness & $\mathrm{CHM},\mathrm{CES},\mathrm{ECE}$ & Composite hallucination / eval / cal. error \\
$N,T_{\text{cycles}}$ & Cycle index; registered cap ($10$); realised closure at $6$ & $\lvert\theta\rvert$ & Backbone parameters ($3.086{\times}10^{9}$) \\
\end{longtable}
}


% --- Main matter ------------------------------------------------------------
\cleardoublepage
\pagenumbering{arabic}

% Chapter list mandated by the BRAC University CSE400 thesis rubric. Each
% chapter title and section number is fixed by the rubric; depth is added
% through subsections inside the prescribed sections. See
% images/thesis_rubric.png for the full Course Outcome to chapter mapping.
\chapter{Introduction}
\label{ch:introduction}

This chapter introduces the problem the thesis addresses, motivates the choice to tackle it through architectural rather than incremental means, formalises the problem statement, declares the objectives, summarises the proposed system at a high level, and registers the scope within which the work is bounded.

\section{Background}
\label{sec:background}

\subsection{Capability and the persistent failure mode}
\label{sec:cap-failure}

Large language models trained on broad text corpora and instruction-tuned for assistant-style interaction now perform competitively on a wide range of natural language tasks, including factoid question answering, multi-step reasoning, summarisation, and program synthesis. Frontier systems such as GPT-4, Claude, and PaLM reach or exceed reported human accuracy on standardised academic benchmarks, and their accessibility through chat interfaces has made them part of everyday research, education, and professional workflows \cite{openai2023gpt4, anthropic2023claude, anil2023palm}. The same systems, however, generate confident but factually incorrect statements at deployment-relevant rates. The accepted name for this failure mode is hallucination, although the term covers heterogeneous causes ranging from retrieval limitations to confabulation under genuine uncertainty.

A practical consequence is that the same fluency that makes large language models useful also makes their hallucinations dangerous. A false answer phrased with the same confidence as a true one is harder to detect than a false answer that sounds uncertain. Users without independent verification, particularly in high-stakes domains, take the model's confident phrasing as a signal of reliability and then act on incorrect information. This thesis takes the position that hallucination is therefore not merely a benchmark-accuracy problem but a deployment-safety problem, and that addressing it requires architectural change rather than incremental decoding adjustments.

\subsection{The hallucination problem at scale}
\label{sec:halluc-scale}

The empirical record on hallucination is consistent across benchmarks and domains. The TruthfulQA benchmark of Lin et al.~\cite{lin-etal-2022-truthfulqa} isolates questions on which humans themselves are prone to hold confident misconceptions, and reports that the largest frontier model evaluated achieved only $58$ percent on the truthfulness axis and $21$ percent on the combined truthful-and-informative axis, against a human baseline of $94$ and $87$ percent respectively. The benchmark also documented an inverse-scaling relationship in which larger models were more, not less, prone to echo human falsehoods, indicating that scale alone does not solve the problem.

Long-form generation shows the same pattern at finer granularity. The atomic-fact decomposition protocol of Min et al.~\cite{min-etal-2023-factscore}, evaluating chat-style models on biographical generation, found that only $58$ percent of the atomic facts in ChatGPT's biographies were supported by the reference knowledge source. Domain-specific evaluations are still more concerning. When Alkaissi and McFarlane~\cite{alkaissi2023artificial} asked ChatGPT to support a clinical explanation with references, it returned five citations of which none existed, and a study of systematic-review searches by Chelli et al.~\cite{info:doi/10.2196/53164} found that $39.6$, $28.6$ and $91.4$ percent of the references produced by GPT-3.5, GPT-4 and Bard respectively were hallucinated, far above any threshold a responsible deployment would tolerate.

These results indicate that hallucination is neither an artefact of small models nor a residual that scale alone will eliminate; the survey literature describes it as an artefact of neural generation that persists in large, instruction-tuned and preference-tuned models \cite{ji2023survey,huang2025survey}. This motivates changing the structure of the system around the model rather than relying on further training alone.

\subsection{Three reinforcing shortcomings}
\label{sec:stateless-gap}

Three structural shortcomings of the current inference paradigm shape every other limitation a deployed model exhibits.

\textbf{No memory:} Standard inference is stateless: every query is processed independently, no representation of past successful answers is retained across sessions, and identical questions can produce inconsistent answers from one interaction to the next. This wastes computation on queries the system has effectively already solved and prevents the accumulation of domain-specific knowledge through use.

\textbf{No reliable confidence:} Modern neural networks are systematically overconfident. Guo et al.~\cite{pmlr-v70-guo17a} found that modern neural classifiers assign confidence that exceeds their actual accuracy, with expected calibration errors above $10$ percentage points on standard image benchmarks, and Kadavath et al.~\cite{kadavath2022language} found that language models fine-tuned with reinforcement learning from human feedback appear badly miscalibrated until their outputs are rescaled. The confidence a model states in words discriminates correct from incorrect answers only weakly: in the study of Xiong et al.~\cite{xiong2023can}, GPT-4's verbalised confidence averages an area under the receiver operating curve of about $0.63$, close to the $0.5$ of random guessing, and no single elicitation method is consistently best. A user therefore cannot reliably tell which answers are likely to be wrong from the model's stated confidence alone.

\textbf{No learning after deployment:} Deployed models are static. Once trained they neither learn from their mistakes nor adapt to the patterns of an institution that uses them, because naive online fine-tuning incurs catastrophic forgetting on the broad capabilities the base training instilled, as Kirkpatrick et al.~\cite{doi:10.1073/pnas.1611835114} documented. The implication is that verification effort spent today produces no compounding benefit tomorrow: a thousand answers verified by humans this morning will not change a single answer the model gives this afternoon.

These three shortcomings reinforce one another. Without memory the model cannot reuse verified knowledge. Without calibrated confidence the system cannot recognise when memory is missing. Without self-improvement the verifications that humans or external systems do produce decay into one-shot interventions. The remainder of this thesis develops a system that breaks the cycle along all three axes simultaneously. The next section argues why a combined intervention at the architectural level is the appropriate response.

\section{Rationale of the Study or Motivation}
\label{sec:rationale}

\subsection{The deployment gap}
\label{sec:deployment-gap}

The capability gap that motivates this thesis is the distance between what large language models do well and what reliable deployment requires of them. Fluency, breadth of topic coverage, and on-demand reasoning are already present in current frontier models. What is not present is the discipline of distinguishing facts that the model has correctly internalised from confabulations that emerge under genuine uncertainty \cite{farquhar2024semantic}, and the discipline of accumulating verified knowledge across interactions instead of regenerating answers from scratch every time. The gap is not aesthetic. In educational deployments, students who learn from confidently incorrect tutoring risk acquiring misconceptions. In healthcare information services, a hallucinated drug interaction or treatment dosage would produce decisions whose costs are measured in patient outcomes rather than in benchmark points. In enterprise knowledge workflows, fabricated citations and missing provenance erode the trust that the entire deployment rests on.

A second, less visible cost is computational. The stateless protocol of a typical large-model deployment regenerates every answer from first principles, including answers whose substance has already been verified in a previous interaction. As deployment scale grows, the cost of unnecessary regeneration accumulates against any computational budget. A system that can recognise a familiar query and serve a previously verified answer at near-zero cost converts the same compute envelope into more delivered value, and that conversion is itself part of the rationale.

\subsection{Why prior pipelines do not close the gap}
\label{sec:prior-pipelines}

Three families of prior systems have attempted to reduce hallucination, and each has a structural limitation that prevents it from closing the gap on its own.

\textbf{Retrieval-augmented generation:} Retrieval-augmented generation, introduced by Lewis et al.~\cite{NEURIPS2020_6b493230}, grounds the model in retrieved evidence, which raises factual recall when the retrieval succeeds but is silent when retrieval misses or returns evidence the model summarises incorrectly. The output remains a single-pass generation whose quality is bounded by retrieval quality, and the system has no representation of which past answers were correct.

\textbf{Prompting interventions:} Prompting interventions, including chain-of-thought reasoning (Wei et al.~\cite{wei2022chain}) and self-consistency over sampled chains (Wang et al.~\cite{wang2023selfconsistency}), expose internal reasoning and exploit agreement across resampled answers, but they provide no persistence: each session begins with no memory of prior verified answers, and the same question asked tomorrow triggers the same reasoning effort as today.

\textbf{Post-generation verification:} Post-generation verifiers such as SelfCheckGPT (Manakul et al.~\cite{manakul-etal-2023-selfcheckgpt}) and semantic entropy (Farquhar et al.~\cite{farquhar2024semantic}) detect hallucinations after they occur and support refusal or revision, but they do not feed back into model weights or memory, so the same hallucination patterns recur in the next session.

The implication is that hallucination reduction at deployment-relevant rates is unlikely to come from incremental improvement of any one family. Stronger retrievers, longer chains of thought, and sharper post-generation verifiers each remain bounded by the structural limitation of their family when used in isolation. What is missing is a system in which verification feeds memory, memory feeds routing, routing feeds generation, and improvement compounds across cycles instead of evaporating at session boundaries. That system is an architectural intervention, and arguing for it is the work this thesis attempts.

\subsection{The architectural opportunity}
\label{sec:arch-opportunity}

Three capabilities, taken together, would close the gap that the existing three families leave open. The first is a verified episodic memory whose admission rule is finer than a single binary threshold. An answer that is plausible but has not earned the precision required for retrieval should not be discarded, because the same query may be answered more confidently after the model has been updated; it should instead enter a deferred buffer and be reconsidered at the next cycle boundary. A four-outcome decision (store, defer, abstain, discard) gives the system somewhere to put borderline outputs without contaminating the high-precision pool. The second capability is confidence-aware routing that reads pre-generation uncertainty as an independent safety signal alongside the combined memory similarity and stored quality. It adds a per-query retrieval-utility decision on the queries that would otherwise fall to the retrieval-augmented tier, so that retrieval is consulted only where it is expected to help rather than unconditionally. When the pre-generation uncertainty is high, the routing layer should force grounded retrieval before generation has started, rather than rely on post-generation review to catch an answer that should never have been produced. The third capability is a constrained self-improvement loop that fine-tunes a small low-rank adapter on a frozen backbone using a verified high-quality pool, so that the bounded adapter plays the role that the parameter anchor plays in elastic weight consolidation \cite{doi:10.1073/pnas.1611835114}, and that rolls the adapter back if a held-out retention probe indicates general capability has eroded. Together the three capabilities convert verification from an event that fires once per query into a discipline that compounds across cycles.

A practical consequence of the joint architecture, named the confidence-aware episodic memory with self-improvement (CAEM), is that the three serving tiers are intended to act not as static cost categories but as a compute ladder that reorganises itself under the cycle loop, and that two distinct statelessness problems of the standard inference paradigm are converted into two stateful mechanisms by two different architectural elements. The first statelessness problem is that retrieval-augmented generation pays the full retrieval-and-generation cost in every session because the model never absorbs what retrieval just taught it; CAEM closes this by feeding verified retrieval-augmented outputs into the self-improvement loop's training pool, so each cycle's low-rank-adapter update folds the previously retrieval-dependent reasoning into the generator's parametric capacity. The design expectation is that the next cycle answers part of the previously retrieval-dependent queries from the zero-shot tier alone, without paying the retrieval cost. The second statelessness problem is that even when the model already knows the answer, the standard protocol re-generates it from scratch on every recurrence and offers no guarantee that two sessions produce the same answer; CAEM closes this through the verified episodic memory, where exact and near-duplicate queries are served directly from the stored answer at embedding-lookup cost, and where the retroactive re-verification pass at every cycle boundary refreshes the stored score of every cached entry against the updated generator, so that the cached answer's quality tracks the model's current capability rather than freezing at the moment of admission. In this design the zero-shot tier absorbs queries that the self-improvement loop has taught the generator to answer, the memory-direct tier absorbs queries that recur and whose cached answer survives the retroactive pass, and the retrieval-augmented tier retains the novel queries that neither mechanism yet covers. On the realised six-cycle trajectory this redistribution is muted on the evaluation fold: the retrieval-augmented tier served $51.3$ percent of evaluation queries at cycle zero and between $44.7$ and $55.1$ percent at later cycles, and the memory-direct tier served at most $0.1$ percent, because the evaluation questions are content-disjoint from the stored episodes by construction (\Cref{tab:tier-trajectory}). The memory-direct tier is built for repeated and paraphrased questions: on a separate probe of $100$ questions built from episodes stored by cycle five, it fired on $42$ of $50$ exact repeats and $19$ of $50$ paraphrases, and its answers were correct in $41$ of the $42$ and all $19$ cases (\Cref{tab:tier1-probe}).

The user-facing surface mirrors the same logic. Every served answer carries the calibrated probability of correctness as a continuous percentage alongside the four-class decision tag, so a user sees both the categorical reliability label (verified, provisional, conflicting, or insufficient) and the underlying probability that the answer is right. When the calibrated probability and the grounding evidence both fall short of the registered floors, the system abstains explicitly rather than serving a guessed answer with an inflated confidence label. The architecture therefore targets three deployment-time properties at once: hallucination that falls across cycles through the verifier-gated self-improvement loop, per-query cost that falls as queries migrate from the retrieval-augmented tier into the cheaper zero-shot and memory-direct tiers, and answers that always carry the system's own calibrated estimate of their reliability.

There is also a theoretical case for why this combined intervention can succeed even when verification is itself imperfect. By Bayes' rule, the stored pool is purer than the generator's raw output whenever the storage gate is more likely to admit a correct answer than an incorrect one, however imperfect the gate is otherwise. The formal statement is \Cref{thm:bayes-purity}; how the property is expected to carry across cycles is stated in \Cref{ch:methodology} as a design heuristic rather than a proven result. For the rationale here, the relevant point is that the architecture does not require a perfect verifier; it requires a verifier whose balanced accuracy exceeds one half on the distribution of claims actually scored, and \Cref{ch:evaluation} reports the evidence on whether that condition holds on the five-benchmark panel.

The scope of the intervention is bounded so that it can be argued empirically within an undergraduate research budget: a single small open generator (Qwen-2.5-3B-Instruct), five public factual question-answering benchmarks, and a single consumer-grade graphics processor, so that the result is reproducible without specialised infrastructure; \Cref{sec:scope} sets out the full scope. The next section formalises the problem the intervention is designed to solve.

\section{Problem Statement}
\label{sec:problem}

\subsection{Formal problem}
\label{sec:formal-problem}

Let a deployed question-answering system be specified by a generator distribution and a retrieval substrate, both held fixed at deployment time. A query arrives and produces an answer through the system's per-query pipeline, after which the system must decide whether to commit the produced answer into a persistent state that will influence future answers. The thesis addresses the following problem. Given a stream of queries drawn from a knowledge-intensive distribution, design a deployable system whose committed answer pool is dominated by correct answers, whose user-facing answers carry a calibrated probability of being correct that is refitted on labelled data at every cycle, so that its precision can be checked against held-out data, whose computational cost per query falls as the system processes more of the stream, and whose general capability does not erode as it accumulates domain-specific knowledge. The system is required to satisfy these properties without per-domain calibration at deployment, without external service calls, and on a single small instruction-tuned generator.

The objects under control are the storage decision applied to each generated answer, the calibration of the confidence used to make that decision, the routing of incoming queries among a memory tier and one or more generation tiers, and the cycle-boundary update of the model parameters. The objects not under control are the generator's prior over factual content, the retrieval corpus, and the eventual stream distribution. Within the controlled set the design must produce a falsifiable empirical claim against a panel of external baselines under matched protocol.

\subsection{Shortcomings addressed and exclusions}
\label{sec:three-shortcomings}

The system must address the three shortcomings set out in \Cref{sec:stateless-gap} together: stateless operation that discards verified answers, stated confidence that does not separate correct from incorrect answers, and a static model that cannot learn from verified use. Because the three reinforce one another, addressing any one in isolation yields a system that is better on one axis and unchanged on the other two: a memory without calibrated confidence cannot tell when it is worth consulting, and calibration without self-improvement leaves the underlying error rate unchanged. The thesis therefore commits to a joint treatment in which memory feeds routing, routing feeds generation, and verification feeds both memory and the cycle-boundary update.

Two classes of approach are explicitly excluded from the present scope. The first is changes to the generator's pre-training corpus or pre-training objective. The thesis treats the generator as fixed and instruments the system around it, so that the architectural claim is separable from any pre-training advantage. The second is reliance on external commercial verification services. Every verification signal the system uses must be reproducible from the open weights of a small specialised model and a publicly available retrieval corpus, so that the result is reproducible without paid infrastructure. The boundary of the problem the thesis solves is therefore an architecture that operates on a fixed open generator, with open verification machinery, on standard hardware, and that is designed to satisfy, and evaluated against, the four properties registered in \Cref{sec:formal-problem} on the five-benchmark factual question-answering panel of \Cref{sec:setup}. The next section converts this problem statement into the specific objectives that organise the rest of the work.

\section{Objective}
\label{sec:objective}

\subsection{Primary objective}
\label{sec:primary-objective}

The primary objective of this thesis is to develop a deployable architecture that reduces confident factual hallucination in knowledge-intensive question answering on a single small instruction-tuned generator, while preserving the generator's general capability across self-improvement cycles. The architecture, CAEM, addresses the three reinforcing shortcomings registered in \Cref{sec:problem} through a verified episodic memory paired with a calibrated probability composite, a four-outcome storage decision, a confidence-aware three-tier router with an independent safety override, retroactive cycle-boundary review of stored content, and a constrained self-improvement loop guarded by a multi-modal retention probe, a panel of held-out probes that spans three task formats (general-knowledge multiple choice, open-text factoid, and commonsense multiple choice). The success of the architecture is measured against a panel of external baselines under matched protocol on a five-benchmark factual question-answering panel partitioned into three training benchmarks and two transfer benchmarks, using a composite hallucination metric that aggregates eight measurable subtypes of confident-error failure mode, together with per-tier latency and the per-cycle retention ratio as auxiliary diagnostic axes.

\subsection{Specific objectives}
\label{sec:specific-objectives}

The primary objective decomposes into seven specific objectives. Each is a separate engineering or scientific deliverable that contributes to the joint claim and that admits an independent empirical check in the later chapters.

\begin{enumerate}
\item \textbf{Verified episodic memory with a four-outcome decision:} Design an episodic memory whose admission rule is finer than a single binary threshold. A generated answer enters one of four mutually exclusive outcomes determined by a calibrated probability composite, a confabulation early-exit rule that fires when high internal confidence coincides with low external grounding, and a deferral path that holds borderline answers for cycle-boundary reconsideration. Storage uses a similarity-indexed encoder over query and reasoning content with a novelty filter and a value-based pruning policy that combines recency, retrieval frequency, and stored-answer quality.

\item \textbf{Pre-generation confidence, three-tier routing with safety override, and a retrieval-utility classifier:} Compute a pre-generation confidence estimate from a short greedy decode over the query that fuses the token-probability aggregate of the first thirty-two decoded tokens with a layer-convergence ratio read from the hidden states over the query, calibrate it via per-benchmark temperature scaling on a disjoint fold, and use it independently from the memory similarity score and the stored-answer quality. The router combines query-to-memory similarity with stored-answer quality into a single dispatch score that selects between memory-direct retrieval, zero-shot generation, and retrieval-augmented generation. An independent safety override forces the conservative tier whenever pre-generation confidence falls below a fixed floor, so that grounded retrieval is available before generation begins and is not entangled with the quality of memory the query happens to land near. A retrieval-utility classifier (Stage 3b) then intercepts the safety-override dispatch and the score-formula fall-through to the retrieval-augmented tier and refines the default into either the zero-shot tier or the retrieval-augmented tier on a per-query basis (\Cref{sec:ruc}); its contribution is measured by the on/off ablation in \Cref{ch:evaluation}.

\item \textbf{Nine-signal post-generation verifier with a calibrated probability composite:} Combine nine complementary signals into a per-answer record across four families: internal calibration of the generator, sample-set agreement across resampled chains, external grounding against retrieved evidence under three aggregation rules, and question-answer relevance. The signals enter a calibrated probability composite that fits per-signal isotonic regression both on the pooled training-panel fold and on each per-benchmark slice, blends each per-benchmark child curve with the pooled fit through a James-Stein style shrinkage prior, and aggregates the per-signal calibrated probabilities through a regularised logistic regression layer fitted on the same calibration fold. The learned aggregation follows the multi-scoring approach of Valentin et al.~\cite{valentin2024hallucination}, which combines the logits of several hallucination scores through logistic regression, in place of a fixed weighted sum.

\item \textbf{Fixed-threshold storage gate with two registered cut-points:} Convert the calibrated probability composite into the storage and deferral decisions through two cut-points and one abstention floor, all fixed before the main run and held constant across cycles. The storage cut-point is chosen on the cycle-zero calibration fold and sets how sure the system must be before an answer enters memory; the deferral cut-point, set lower, sends a wider band of borderline answers to the deferred buffer for reconsideration at the next cycle boundary.

\item \textbf{Retroactive review and deferred reconsideration at cycle boundaries:} Because the verifier's internal-calibration signals read the generator's hidden state directly, stored composites are model-relative snapshots that drift after every fine-tune. A retroactive re-verification pass at every cycle boundary rescores stored episodes under the updated model, overwrites each stored composite with its new value, raised or lowered, and prunes any entry whose composite has fallen below a prune threshold set beneath the storage threshold. A companion pass over the deferred buffer promotes entries that have crossed the storage threshold and discards entries whose age has exceeded a fixed time-to-live.

\item \textbf{Constrained self-improvement with a parameter-efficient adapter and a multi-modal retention guard:} Curate the fine-tuning pool from stored episodes that exceed a strict quality threshold set above the storage cut-point and pass the curated pool through a benchmark-balancing reweighting layer that prevents single-benchmark dominance. Fine-tune the generator through a low-rank adaptation layer on the attention and feed-forward projections rather than through full-parameter fine-tuning \cite{hu2022lora}, with the trainable footprint at roughly $2$ percent of the backbone parameter count and the underlying backbone frozen across the cycle. The bounded adapter parameter budget, layered on a frozen backbone, acts as the implicit cycle-to-cycle drift bound in the spirit of elastic weight consolidation \cite{doi:10.1073/pnas.1611835114}, without the hand-tuned quadratic anchor coefficient that a full-parameter fine-tune would require. Probe a registered set of held-out distributions before and after each fine-tune and abort the cycle with an adapter rollback when the worst probe ratio falls below a fixed floor; the memory is preserved across rollbacks, so accumulated knowledge survives abort events even when the adapter does not.

\item \textbf{Empirical validation against external baselines and architectural ablations:} Evaluate the architecture on the five-benchmark panel and the three-probe retention panel described in \Cref{sec:scope-research}. Report the composite hallucination metric, the per-tier behaviour, the retention ratio, and the calibration trajectory cycle by cycle. Compare against eight external baselines under matched protocol using paired statistical tests with multiple-comparison correction within each benchmark family, and run the pre-specified retrieval-utility-classifier on/off ablation, with its acceptance gates declared in advance and its bundled-confound scope disclosed.
\end{enumerate}

\subsection{Pre-specified hypotheses}
\label{sec:hypotheses}

The thesis registers five falsifiable hypotheses ahead of the empirical chapters. Each hypothesis names a specific quantity, a deciding direction, and a deciding statistic, so that the empirical chapters either confirm or refute the hypothesis without ambiguity. The first four hypotheses pair with the formal statements of \Cref{ch:methodology}, which comprise two proven results, design heuristics, and remarks on idealised models; the fifth tests the operational claim about deployed cost. Below each hypothesis, a short note says what it tests and how it is decided.

\begin{enumerate}
\item \textbf{Verifier balanced accuracy exceeds one half on the calibration-matched distribution at every cycle:} Storing answers makes the pool purer than the generator's raw output only if the gate admits correct answers more often than incorrect ones, which is the same as a balanced accuracy above one half (\Cref{thm:bayes-purity}). The deciding statistic is the gate's balanced accuracy at the storage threshold, measured on each cycle's calibration fold and reported per benchmark. The hypothesis is falsified by any cycle in which the condition fails on the pooled training panel.

\item \textbf{Pool purity is non-decreasing across cycles whenever the precondition is satisfied:} If the condition of H1 holds and the verifier does not change, an idealised model of repeated filtering (\Cref{thm:bayes-convergence}) predicts that each cycle's pool is at least as pure as the last. The deciding statistic is the pool purity at each cycle. The hypothesis is falsified by any cycle in which the condition holds and pool purity falls. \Cref{ch:methodology} explains why the realised system, which refits its composite at every cycle, may fluctuate around this prediction rather than follow it exactly.

\item \textbf{Late-cycle pool-purity increments enter a bounded stationary band on the committed-cycle subsequence:} The convergence heuristic (\Cref{thm:convergence}) expects pool purity to settle. A convergence rate cannot be estimated from six cycles, so the test is whether purity stays inside a narrow band over the committed cycles, that is, the cycles whose fine-tune was kept rather than rolled back. The deciding statistic is the purity on those cycles and the width of the band. The sharper test of a geometric convergence rate needs at least fifteen committed cycles and is left to future work.

\item \textbf{The composite hallucination metric at the realised cycle horizon approaches a parameter-bounded asymptote:} The full-pipeline hallucination decomposition (\Cref{thm:asymptotic-elim}) splits the user-facing hallucination rate into a part served from memory and a part generated afresh, and predicts a floor set by the memory-tier hit rate and by the correctness ceiling of the generator. The deciding statistic is the gap between the composite hallucination metric at the final cycle and that predicted floor.

\item \textbf{Per-cycle Tier 1 hit rate grows monotonically as memory accumulates, producing a measurable per-query cost reduction:} Serving a query from memory costs only an embedding lookup, so a growing memory share should lower the average cost per query. The deciding statistic is the share of queries served by the memory tier at each cycle and the corresponding pooled latency.
\end{enumerate}

The empirical chapters report each hypothesis with its deciding statistic and tag each as supported, partially supported, or unsupported at the trajectory's end. Together the seven specific objectives and the five hypotheses constitute the complete falsifiable content of the architectural claim, each anchored to a measurement reported in \Cref{ch:evaluation}. \Cref{tab:hyp-registration} summarises the registration in compact form.

\begin{table}[!htbp]
\centering
\small
\begin{tabular}{@{}p{0.6cm}>{\raggedright\arraybackslash}p{6.0cm}>{\raggedright\arraybackslash}p{4.3cm}>{\raggedright\arraybackslash}p{2.6cm}@{}}
\toprule
\textbf{ID} & \textbf{Hypothesis} & \textbf{Deciding statistic} & \textbf{Formal claim} \\
\midrule
H1 & Verifier balanced accuracy exceeds one half on the calibration-matched distribution at every cycle. & Per-cycle verifier balanced accuracy at the fitted storage threshold, pooled training panel. & \Cref{thm:bayes-purity} precondition \\
H2 & Pool purity is non-decreasing across cycles whenever the precondition is satisfied. & Per-cycle pool-purity trajectory under the precondition-satisfied regime. & \Cref{thm:bayes-convergence} \\
H3 & Late-cycle pool-purity increments enter a bounded stationary band on the committed-cycle subsequence. & Per-cycle pool-purity reading on the committed-cycle subsequence and realised band width. & \Cref{thm:convergence} \\
H4 & The composite hallucination metric at the realised cycle horizon approaches a parameter-bounded asymptote. & Gap between final-cycle composite hallucination metric and the predicted asymptote. & \Cref{thm:asymptotic-elim} \\
H5 & Per-cycle Tier 1 hit rate grows monotonically as memory accumulates, producing a measurable per-query cost reduction. & Per-cycle Tier 1 share and per-cycle pooled latency. & Operational claim (\Cref{sec:econ-deployment}) \\
\bottomrule
\end{tabular}
\caption[The five pre-specified hypotheses with deciding statistic and the theoretical or operational claim each tests.]{The five pre-specified hypotheses with deciding statistic and the theoretical or operational claim each tests. Verdicts are reported in \Cref{tab:hypothesis-verdicts}.}
\label{tab:hyp-registration}
\end{table}

\section{Methodology in Brief}
\label{sec:methodology-brief}

\subsection{Architecture overview}
\label{sec:arch-overview}

The architecture organises a single query into an eight-stage pipeline that is traversed once per query and wrapped in a per-cycle update loop. The first stage produces a pre-routing confidence from a short greedy decode over the query, fusing the decode's token-probability aggregate with a layer-convergence ratio. The second stage encodes the query with a sentence-embedding encoder and searches a similarity-indexed episodic memory for the nearest stored entry, returning a similarity score and the neighbour's stored-quality score. The third stage combines the similarity score and the stored-quality score into a single dispatch score and routes the query to the memory-direct tier when the score is high, the zero-shot tier in a middle band, and the retrieval-augmented tier when it is low, with an independent safety override that forces the retrieval-augmented default whenever the pre-routing confidence falls below a fixed floor. A retrieval-utility classifier then intercepts both the safety-override dispatch and the score-formula fall-through to the retrieval-augmented tier and refines that default into either the zero-shot tier or the retrieval-augmented tier on a per-query basis, so that retrieval is consulted only on queries it is expected to help; the classifier is specified in \Cref{sec:ruc}. The fourth stage executes the chosen tier: the memory-direct tier serves the neighbour's stored answer; the zero-shot tier runs a generation on the fine-tuned generator; the retrieval-augmented tier retrieves passages from a public corpus and runs a grounded generation. The fifth stage passes every generated answer through a nine-signal post-generation verifier whose output is a calibrated probability that the answer is correct. The sixth stage applies a four-outcome decision (store, defer, abstain, discard) to that probability, with a confabulation early-exit rule that takes precedence over the thresholds. Both stages are summarised below. The seventh stage executes the implied memory and user-facing actions. The eighth stage runs at cycle boundaries and is detailed below. \Cref{fig:architecture} gives the high-level view; the full per-stage pipeline appears in \Cref{ch:methodology}.

\begin{figure}[!htbp]
\centering
\resizebox{\textwidth}{!}{%
\begin{tikzpicture}[
    >=stealth,
    node distance=8mm and 10mm,
    box/.style={rectangle, rounded corners=4pt, draw=#1!70!black, fill=#1!12,
                line width=0.9pt, align=center, text width=3.0cm, minimum height=1.6cm,
                inner sep=4pt, font=\normalsize\bfseries},
    store/.style={rectangle, rounded corners=4pt, draw=teal!70!black, fill=teal!10,
                  line width=0.9pt, align=center, text width=3.0cm, minimum height=1.6cm,
                  inner sep=4pt, font=\normalsize\bfseries},
    arrow/.style={->, line width=1.1pt, draw=#1!70!black},
    darrow/.style={->, line width=1.0pt, draw=#1!70!black, dashed},
    elabel/.style={font=\small, fill=white, inner sep=2pt, text=black!75}
]
%% Top row
\node[box=blue]                                 (qry)  {Query};
\node[box=cyan,   right=of qry]                 (pre)  {Stage 1\\Pre-route};
\node[box=teal,   right=of pre]                 (enc)  {Stage 2\\Encode};
\node[box=orange, right=of enc]                 (route){Stage 3\\Route};

%% Memory above encode
\node[store, above=of enc]                      (mem)  {Memory};

%% RUC (Stage 3b) between Route and Tier band
\node[box=yellow!60!orange, right=8mm of route, fill=yellow!18, text width=2.4cm, minimum height=1.4cm]
                                                (ruc)  {Stage 3b\\RUC};

%% Tier band right of RUC
\node[box=green,  right=10mm of ruc, yshift=18mm]  (t1) {Tier 1\\memory};
\node[box=yellow, right=10mm of ruc]                (t2) {Tier 2\\zero-shot};
\node[box=red,    right=10mm of ruc, yshift=-18mm] (t3) {Tier 3\\RAG};

%% External corpus
\node[box=blue, dashed, below=8mm of t3, fill=blue!4] (kb) {Public\\corpus};

%% Verify, Decide, Output
\node[box=violet, right=22mm of t2]   (ver) {Stage 5\\Verify};
\node[box=purple, below=of ver]       (dec) {Stage 6\\Decide};
\node[box=green,  right=of ver]       (out) {Stage 7\\Output};

%% Generator below route
\node[store, below=12mm of route]     (gen) {Generator};

%% Arrows -- ingress
\draw[arrow=blue]   (qry)   -- (pre);
\draw[arrow=cyan]   (pre)   -- (enc);
\draw[arrow=teal]   (enc)   -- (route);

%% Memory probe
\draw[darrow=teal]  (enc) -- (mem);
\draw[darrow=teal]  (mem) -| (route);

%% Router to T1 direct; T2/T3 fall-through via RUC
\draw[arrow=green, rounded corners=4pt]  (route.east) -- ++(0.3,0) |- (t1.west);
\draw[arrow=orange] (route) -- (ruc);
\draw[arrow=orange] (ruc.east) to[out=20, in=180]  (t2.west);
\draw[arrow=red]    (ruc.east) to[out=-20, in=180] (t3.west);

%% Safety override (pre-route low confidence -> RUC, which then decides T2 vs T3)
\draw[darrow=red, rounded corners=4pt]   (pre.south) -- ++(0,-0.8) -| node[elabel, pos=0.25, anchor=north] {safety} (ruc.south);

%% Memory to T1, Generator to T2
\draw[arrow=green]  (mem.east) to[out=0, in=120] (t1.north);
\draw[arrow=orange] (gen.east) to[out=0, in=-120] (t2.south);

%% External corpus to T3
\draw[darrow=blue]  (kb) -- (t3);

%% Tiers to verify
\draw[arrow=green]  (t1) -| (ver.north);
\draw[arrow=orange] (t2) -- (ver);
\draw[arrow=red]    (t3) -| (ver.south);

%% Verify -> Decide -> Output
\draw[arrow=violet] (ver) -- (dec);
\draw[arrow=green]  (ver) -- (out);

%% Decide -> Memory (commit), Memory -> Generator (training pool)
\coordinate (commitUp) at (0, 4.2);
\draw[darrow=purple, rounded corners=4pt] (dec.east) -- ++(4.5, 0) |- node[elabel, pos=0.25] {commit} (mem.north |- commitUp) -- (mem.north);
\draw[darrow=purple, rounded corners=4pt] (mem.south east) -- ++(0.4,0) |- node[elabel, pos=0.75] {train} (gen.west);

%% Stage 8 frame
\node[draw=purple!65, dashed, line width=0.7pt, rounded corners=4pt,
      fit={(mem) (gen) (route)}, inner sep=4pt,
      label={[font=\scriptsize\bfseries, text=purple!75!black]above:Stage 8 \ Self-improvement (cycle boundary)}] {};
\end{tikzpicture}
}

\caption[Architecture overview]{\small Eight-stage pipeline: pre-route, encode, route, Stage 3b retrieval-utility classifier on the score-formula and safety-override fall-through to the retrieval-augmented tier, tier execution, verify, decide, output, and a cycle-boundary self-improvement block; full specification in \Cref{ch:methodology}.}
\label{fig:architecture}
\end{figure}

\subsection{Verification and the storage decision}
\label{sec:per-query-pipeline}

The verifier evaluates nine complementary signals across four families. Internal calibration signals read the generator's own state and include token-level uncertainty, dropout-based disagreement, and an internal-state composite. Sample-set signals measure agreement across resampled chains, including average pairwise similarity and chain-to-answer entailment. External grounding signals measure entailment between retrieved passages and the answer under three aggregation rules: top-one entailment, pooled mean entailment, and a weakest-link atomic-fact decomposition that fires only when the answer is long enough to admit decomposition. A question-answer relevance signal closes the off-topic-but-grounded failure mode by checking whether the answer addresses the question. The signals enter a calibrated probability composite that maps each per-signal score to a calibrated probability via isotonic regression and aggregates them through a regularised logistic regression layer fitted on a disjoint calibration fold; the calibrated output is the storage score. The confabulation early-exit rule discards an answer outright when high internal confidence coincides with low external grounding. Otherwise the four-outcome decision maps the storage score to one of four mutually exclusive actions: commit to memory, defer to a bounded buffer with a time-to-live, abstain explicitly, or discard. The user-facing action that follows depends on the outcome, so that the user receives a stored answer, a generated answer, or a calibrated refusal rather than a silently uncertain response.

\subsection{Calibration discipline}
\label{sec:calibration-discipline}

The storage score is converted into the storage and deferral decisions by two fixed cut-points, both chosen before the main run, the storage cut-point on the cycle-zero calibration fold, and held constant across cycles (\Cref{sec:conformal-gate} explains why fixed cut-points replaced an earlier conformal formulation). What changes from cycle to cycle is the calibrated probability composite: at every cycle boundary a labelled calibration fold refits it, per benchmark under a shrinkage prior toward the pooled fit, while the verifier-curated answers produced by the pipeline supply the training signal for the fine-tune. The research trajectory reuses one cold-start calibration fold of $1\,500$ labelled samples, re-scored under each cycle's fine-tuned model; the production design in \Cref{ch:methodology} instead draws a fresh stratified labelled fold from each cycle's admissions. Because the cut-points are fixed and only the composite is refitted, the meaning of the storage threshold stays the same as the score distribution shifts.

\subsection{Cycle-boundary self-improvement loop}
\label{sec:cycle-loop}

The cycle loop iterates the per-query pipeline over a fixed workload, accumulates verified episodes into memory, and runs three cycle-boundary passes before fine-tuning. The retroactive re-verification pass rescores every stored episode under the updated generator, updates each stored composite to its new value, and prunes entries whose composite has fallen below the prune threshold. The deferred reconsideration pass examines every entry in the bounded deferral buffer and either promotes, expires, or requeues the entry. The memory consolidation pass collapses near-duplicate stored episodes that exceed a similarity threshold, picking the highest-quality representative within each near-duplicate cluster. Fine-tuning then draws a training pool from stored episodes whose composite score exceeds a strict quality threshold above the storage threshold, ensuring the training pool is held to a higher bar than the user-facing retrieval tier. The fine-tune is regularised implicitly through the bounded low-rank-adapter parameter budget layered on a frozen backbone, in the spirit of elastic weight consolidation \cite{doi:10.1073/pnas.1611835114} but without an explicit quadratic anchor coefficient. Before and after the fine-tune, the cycle probes the three held-out retention surfaces and computes the worst per-probe retention ratio. When that ratio falls below a fixed floor, the cycle aborts: weights are rolled back to the previous checkpoint and the retroactive pass for the next cycle is skipped. The memory is never rolled back, so accumulated knowledge survives abort events even when weights do not. This asymmetry decouples weight-level drift from memory-level accumulation and is what makes the self-improvement loop safe to iterate.

\Cref{fig:experimental-workflow} illustrates the trajectory across early, mid, and late phases. Each phase runs the same per-query pipeline over a fresh workload of queries, then applies the cycle-boundary update passes under the retention guard. The three boxes are representative phases rather than literal per-cycle depictions; the actual trajectory is reported cycle by cycle in \Cref{ch:evaluation}.

\begin{figure}[H]
\centering
\begin{tikzpicture}[
    >=Stealth,
    font=\sffamily,
    %% phase card -- left column, with chip header and bullet list
    phase/.style 2 args={rectangle, rounded corners=6pt,
        draw=#1!75!black, line width=0.9pt, fill=#1!10,
        align=left, text width=6.6cm, minimum height=2.8cm,
        inner sep=7pt, font=\small,
        blur shadow={shadow blur steps=4, shadow blur extra rounding=1pt,
                     shadow xshift=0.5ex, shadow yshift=-0.5ex, shadow opacity=20}},
    %% cycle-boundary band -- thin grey strip between phases
    boundary/.style={rectangle, rounded corners=3pt,
        draw=gray!50, line width=0.7pt, fill=gray!8,
        align=center, text width=6.6cm, minimum height=0.95cm,
        inner sep=4pt, font=\footnotesize\itshape, text=black!75},
    %% logs callout -- right column, mirrors phase card colour
    logs/.style 2 args={rectangle, rounded corners=4pt,
        draw=#1!55!black, line width=0.7pt, fill=#1!6,
        align=left, text width=4.8cm, minimum height=2.8cm,
        inner sep=5pt, font=\scriptsize},
    %% arrow between phases through boundary band
    flow/.style={->, line width=1.4pt, draw=#1!70!black}
]

%% ----- Phase cards (left column) ----------------------------------------
\node[phase={red}{}] (c1) {%
\colorbox{red!75!black}{\color{white}\strut\,\textbf{\scriptsize EARLY CYCLES}\,}\quad\textit{\footnotesize cold start}\\[4pt]
\(\bullet\) Memory: \textbf{empty $\rightarrow$ sparse}\\
\(\bullet\) Most queries to Tier 2 / Tier 3\\
\(\bullet\) Nine-signal verify $\rightarrow$ four-outcome decide\\
\(\bullet\) Store / Defer / Abstain / Discard%
};

\node[boundary, below=0.32cm of c1] (cons1) {%
Cycle boundary: retroverify \(+\) defer-reconsider \(+\) consolidate \\
low-rank-adapter fine-tune \(+\) retention guard
};

\node[phase={orange}{}, below=0.32cm of cons1] (c2) {%
\colorbox{orange!75!black}{\color{white}\strut\,\textbf{\scriptsize MID CYCLES}\,}\quad\textit{\footnotesize accumulating}\\[4pt]
\(\bullet\) Memory: \textbf{growing}\\
\(\bullet\) Tier 1 serves repeated and paraphrased questions\\
\(\bullet\) Training pool above storage threshold\\
\(\bullet\) Retroverify rescores/prunes under updated model%
};

\node[boundary, below=0.32cm of c2] (cons2) {%
Cycle boundary: same passes; on retention failure, weights roll back, memory persists
};

\node[phase={green}{}, below=0.32cm of cons2] (c3) {%
\colorbox{green!75!black}{\color{white}\strut\,\textbf{\scriptsize LATE CYCLES}\,}\quad\textit{\footnotesize near equilibrium}\\[4pt]
\(\bullet\) Memory: \textbf{mature, near-equilibrium}\\
\(\bullet\) Tier 1 share highest when questions recur\\
\(\bullet\) Composite evaluation score reported\\
\(\bullet\) Composite hallucination metric vs.\ baselines%
};

%% ----- Logs callouts (right column) -------------------------------------
\node[logs={red}{}, right=1.0cm of c1] (m1) {%
\textbf{Per-cycle logs}\\[2pt]
\(\bullet\) (store, defer, abstain, discard) counts\\
\(\bullet\) training-pool size\\
\(\bullet\) routing distribution by tier%
};

\node[logs={orange}{}, right=1.0cm of c2] (m2) {%
\textbf{Cycle-boundary logs}\\[2pt]
\(\bullet\) retroverify (promotions, prunes)\\
\(\bullet\) deferred (promote, drop, keep)\\
\(\bullet\) consolidation clusters resolved\\
\(\bullet\) retention ratio against held-out probe%
};

\node[logs={green}{}, right=1.0cm of c3] (m3) {%
\textbf{Final report \textnormal{(\Cref{ch:evaluation})}}\\[2pt]
\(\bullet\) composite evaluation score, geometric mean\\
\(\bullet\) composite hallucination metric reduction\\
\(\bullet\) per-cycle trajectory \(+\) classifier on/off ablation%
};

%% ----- Vertical flow arrows --------------------------------------------
\draw[flow=red]    (c1)    -- (cons1);
\draw[flow=orange] (cons1) -- (c2);
\draw[flow=orange] (c2)    -- (cons2);
\draw[flow=green]  (cons2) -- (c3);

\end{tikzpicture}
\caption[Experimental workflow]{Intended self-improvement trajectory across early, mid, and late phases (design view). Tier 1 answers only repeated or paraphrased questions, so on the evaluation fold, whose questions are distinct from the stored episodes, it rarely fires; the realised per-cycle readings are in \Cref{ch:evaluation}. Each phase runs the per-query pipeline and, at cycle boundaries, the retroactive re-verification, deferred reconsideration, consolidation, and adapter fine-tune under the multi-modal retention guard. Full specification appears in \Cref{ch:methodology}; headline results in \Cref{ch:evaluation}.}
\label{fig:experimental-workflow}
\end{figure}

\section{Scopes and Challenges}
\label{sec:scope}

\subsection{Research scope}
\label{sec:scope-research}

The architecture targets factual question answering, defined as the task of producing answers that admit a definite correctness verdict against a reference. The five evaluation benchmarks adopted in this work span factual claim verification (FEVER), open-domain trivia (TriviaQA), multiple-choice commonsense reasoning (CommonsenseQA), truthfulness on common-misconception probes (TruthfulQA), and multi-step commonsense reasoning (StrategyQA), partitioned into a three-benchmark training panel (FEVER, TriviaQA, CommonsenseQA) and a two-benchmark transfer panel (TruthfulQA, StrategyQA). Two further benchmarks, HotpotQA \cite{yang-etal-2018-hotpotqa} and Natural Questions, were evaluated at cycle zero and dropped before the main run: the base generator answered only $9$ and $17$ percent of their questions correctly, respectively, and the storage gate kept just $26$ and $12$ of $500$ answers, of which half or fewer were correct, too few and too noisy to train on. Their cycle-zero results are preserved in the project archive, and the Natural Questions training split is reused only to train the retrieval-utility classifier. Each of the five panel benchmarks admits an objective scoring rule, namely exact match or label accuracy, depending on the task family; TruthfulQA's free-text answers are judged against the benchmark's reference answers by a language-model judge, following the judge-based evaluation introduced with the benchmark \cite{lin-etal-2022-truthfulqa}, and the same judge scores every system. A multi-modal retention probe panel is reserved as held-out retention control rather than as a target benchmark: the fine-tuning loop is not optimised on any probe, and the retention guard aborts any cycle in which the worst per-probe accuracy ratio falls below a fixed fraction of the cycle-zero baseline. The probe panel comprises a general-knowledge multiple-choice surface (MMLU), an open-text factoid surface drawn from the TriviaQA test fold, and a commonsense multiple-choice surface drawn from the CommonsenseQA test fold; the worst-probe rule is the registered abort criterion.

The scope explicitly excludes tasks whose evaluation is intrinsically subjective, including creative generation, open-ended dialogue, and stylistic rewriting. Tasks that require knowledge beyond the pre-training cutoff, such as breaking news, market data, and recently issued documents, are excluded because the system does not include a continuous-update mechanism for ingesting external knowledge sources. Code generation and formal mathematical reasoning are excluded so that the evaluation surface remains within natural-language reasoning and the verifier signals retain a uniform interpretation across benchmarks. The system operates on English text only; multilingual transfer and code-switched evaluation are deferred to future work.

Resource scope is set by the engineering envelope of an undergraduate thesis. The generator is a three-billion-parameter open-weight instruction-tuned decoder-only model, chosen for two reasons. First, it admits parameter-efficient fine-tuning on a single consumer-grade graphics processor with thirty-two gigabytes of memory, under sixteen-bit mixed-precision arithmetic, through a low-rank adaptation layer on the attention and feed-forward projections, with the trainable footprint at roughly $2$ percent of the backbone parameter count and the underlying backbone frozen across the cycle; full-parameter fine-tuning on the same backbone size does not fit the engineering envelope and is registered as the rejected alternative on the joint retention-and-feasibility axis. Second, as an instruction-tuned model it has competitive baseline accuracy on factual question answering for its size, which the self-improvement argument requires (see the baseline-performance challenge below). The ten-cycle horizon is sized to capture the early-cycle calibration phase, the mid-cycle accumulation phase, and the late-cycle equilibrium phase expected by the convergence heuristic in \Cref{ch:methodology}, while remaining feasible within a twelve-month compute envelope; an automatic early-stop check at every cycle boundary can stop the run earlier. The realised run stopped after cycle five, as planned after the cycle-four rollback (\Cref{sec:env-cycles}). The episodic memory is sized to a target capacity in the millions of entries with a tiered index that promotes from a flat structure to a quantised inverted-file structure above a threshold; the workloads reported in this thesis grow the memory to $12\,175$ entries at cycle five (\Cref{tab:memory-store-trajectory}), below the size at which the index switches to the quantised structure and far below the capacity at which value-score pruning triggers, so the realised run uses exact similarity search and never invokes pruning.

\subsection{Technical challenges}
\label{sec:scope-challenges}

The implementation faces six primary technical challenges that the architecture must address.

\begin{itemize}
    \item \textbf{Verifier reliability and stored-quality drift:} The nine-signal post-generation verifier produces a calibrated storage score that doubles as a retrieval-time quality proxy and as a quality filter on the training pool. Three of the nine signals belong to the internal-calibration family and read the generator's hidden state directly, so their distributions shift after every fine-tune. The stored composite of every episode is therefore a model-relative snapshot, not a permanent attribute of the entry. The retroactive re-verification pass at each cycle boundary refreshes every stored composite under the updated generator, raising or lowering it to the new value, and prunes entries whose composite has fallen below the prune threshold. The cost of this discipline is a full re-scoring of every stored episode at each cycle boundary, which grows with the size of the memory.

    \item \textbf{Catastrophic forgetting under iterative fine-tuning:} Fine-tuning on domain-specific verified examples risks degrading the generator's general-purpose capabilities. The architecture mitigates this risk through three layers of defence: a strict training-pool quality filter that holds the training pool to a higher bar than the user-facing storage tier, a bounded low-rank-adapter parameter envelope layered on a frozen backbone that acts as the implicit cycle-to-cycle drift bound (in the spirit of elastic weight consolidation \cite{doi:10.1073/pnas.1611835114} but without an explicit quadratic anchor coefficient on the shipped path), and a multi-modal retention probe panel whose worst per-probe ratio gates the commit. The retention probe panel spans a general-knowledge multiple-choice surface, an open-text factoid surface drawn from a training-benchmark test fold, and a commonsense multiple-choice surface drawn from a second training-benchmark test fold; the cycle aborts when any per-probe ratio falls below the fixed floor, weights are rolled back to the previous checkpoint, and the memory is preserved so that a fine-tune that fails in any given cycle does not lose the verified episodes accumulated during it. The three layers are not ablated individually within the realised horizon; the retention guard's effect is visible in the cycle-four rollback reported in \Cref{ch:evaluation}, and the only architectural ablation reported there is the retrieval-utility-classifier on/off contrast.

    \item \textbf{Pre-routing confidence calibration:} The pre-routing confidence is used both as a one-dimensional input to the dispatch score and as the sole input to the safety override that forces grounded retrieval, so its calibration influences the system's behaviour twice. The architecture fuses two cheap signals, namely the token-probability aggregate of a short greedy decode that scores the first thirty-two tokens of a draft answer and a layer-convergence ratio computed from the hidden states over the query, and temperature-scales the fused scalar per benchmark on a disjoint fold by minimising binary negative-log-likelihood. The non-trivial design choice is that the pre-routing signal does not require the full multi-chain sampling that the post-generation verifier uses, so its cost is paid per query regardless of which tier eventually serves the answer; the cost of the more expensive sample-based signals is deferred to the post-generation verifier on the zero-shot and retrieval-augmented paths only.

    \item \textbf{Memory retrieval quality and novelty filtering:} At deployment scale the episodic memory uses approximate nearest-neighbour search, which raises the possibility of missing a semantically similar entry whose embedding happens to fall near a partition boundary; at the scale of the realised run the search is exact. Near-duplicate writes are filtered at storage time by a novelty threshold on cosine similarity, and the consolidation pass at each cycle boundary collapses near-duplicate clusters by retaining only the highest-quality representative within each cluster. The novelty threshold is a tuning knob with a clear failure mode in either direction: a threshold set too high admits paraphrase-level duplicates that bloat the index without adding coverage, while a threshold set too low rejects legitimate paraphrases that would have improved retrieval. The chosen value is fixed before the main run and held constant across cycles.

    \item \textbf{Baseline performance dependency for self-improvement:} The self-improvement argument depends on the base generator performing meaningfully above random on the target benchmarks, because training on verified-but-imperfect episodes can only compound accuracy when the verifier's expected false-positive rate is low enough that the training pool is cleaner than the generator's untrimmed outputs. The risk is mitigated by selecting an instruction-tuned backbone with competitive baseline accuracy on factual question answering, by selecting benchmarks on which the base generator clears that threshold, and by composing the storage decision from the calibrated probability composite, the confabulation early-exit gate, and the in-distribution training-pool gate that excludes transfer-benchmark episodes from fine-tuning. Each component contracts the storage class below the rate that an unrestricted score would admit, and the conjunction is what makes the precision floor on the storage tier achievable.

    \item \textbf{Evaluation metric selection:} Measuring hallucination reduction requires a metric design that is more discriminative than aggregate accuracy. The composite hallucination metric used in this thesis is the equal-weighted mean over an eight-axis failure-mode taxonomy that includes confident confabulation, factual fabrication, logical fabrication, off-topic answers, defensive evasion, template leakage, false refusal, and length-padded over-generation. A ninth axis is retained in the taxonomy but excluded from the default denominator because the corresponding subtype is structurally unreachable under the verifier backend used here, and its rate is reported separately for completeness. A uniform reduction across eight structurally distinct failure modes is a stronger architectural claim than a larger reduction on any single axis. The composite evaluation score is a geometric mean across accuracy, epistemic quality, retention, calibration, and verifier reliability; the geometric form penalises any axis that collapses, which matches the design constraint that no single axis may be sacrificed for the others. Comparison against the baseline panel is conducted under matched protocol, with the same evaluation samples, backbone, answer format, and scoring rules, so that effects are not confounded by re-implementation choices.
\end{itemize}

Each mitigation above is built into the architecture rather than added as a downstream patch; \Cref{ch:evaluation} reports the evidence available for each on the realised trajectory.

\section{Chapter signpost}
\label{sec:ch1-signpost}

This chapter has registered the deployment-safety problem, the architectural opportunity that closes the stateless-overconfident-static triad, the seven specific objectives that decompose the primary objective, and the five pre-specified hypotheses that the empirical chapters adjudicate. \Cref{ch:literature-review} establishes the technical foundations on which the architectural commitments rest (the transformer backbone, the hallucination taxonomy, natural language inference, sentence embeddings, self-consistency decoding, and catastrophic forgetting) and surveys the empirical literature on hallucination measurement, verification, confidence estimation, memory-augmented architectures, selective retrieval, self-improvement and continual learning, external baselines, and evaluation benchmarks. The survey locates the architectural contribution of this thesis against the prior work that the seven specific objectives draw from, so the methodology can be read as a single integrated response to a literature whose individual contributions do not close the joint gap. \Cref{ch:requirements} states the functional requirements, the societal and environmental impact, and the project plan. \Cref{ch:methodology} specifies the architecture stage by stage together with its formal analysis. \Cref{ch:evaluation} reports the realised six-cycle trajectory, the baseline comparisons, the retrieval-utility-classifier ablation, and the verdict on each hypothesis. \Cref{ch:conclusion} summarises the findings, the limitations, and the future work, and the appendices collect the implementation reference, the per-cycle evidence, the significance-test details, and the reproducibility recipe.   % Introduction (preserved)
\chapter{Literature Review}
\label{ch:literature-review}

\section{Preliminaries}

Chapter~\ref{ch:introduction} laid out the three limitations of current LLMs that CAEM addresses (stateless operation, poor confidence calibration, absence of self-improvement) and the three architectural commitments (verified episodic memory, confidence-aware routing, constrained self-improvement) that answer them. This chapter supplies the two kinds of background a reader needs to evaluate those commitments in detail. The Preliminaries section that follows establishes foundational concepts and terminology: the transformer backbone, the hallucination taxonomy, natural language inference, sentence embeddings, self-consistency decoding, and catastrophic forgetting. The Review of Existing Research section that follows it then surveys the empirical literature on hallucination measurement, verification, confidence estimation, memory-augmented architectures, selective retrieval, self-improvement and continual learning, external baselines, and evaluation benchmarks, locating CAEM's contribution relative to prior work. Together they prepare the ground for the formal methodology in Chapter~\ref{ch:methodology}.

\subsection{Large Language Models and Transformer Architecture}

Large language models are neural networks trained on very large text collections to predict and generate natural language. Modern large language models adopt the Transformer architecture \cite{NIPS2017_3f5ee243}, whose self-attention mechanism lets every token in the input attend to every other token. The original Transformer stacks encoder and decoder layers, each combining multi-head self-attention with a feed-forward network; the decoder-only variant used in this thesis keeps only the decoder stack.

This thesis uses Qwen-2.5-3B-Instruct \cite{qwen25}, a three-billion-parameter decoder-only transformer that has already been instruction-tuned and preference-tuned before it enters the CAEM pipeline. Decoder-only models generate the answer token by token from left to right, with the question and the generated rationale sharing a single causal attention stream. The backbone is chosen for two reasons: it fits the whole self-improvement loop, including low-rank fine-tuning, Monte Carlo dropout sampling and the nine-signal verifier, on a single consumer-grade graphics processor in sixteen-bit brain-float precision; and because it is instruction-tuned, it already answers in the required format, so the verifier has usable outputs to score from the first cycle.

Fine-tuning adapts pre-trained models to specific tasks or domains by continuing training on task-specific data with a smaller learning rate. The model parameters $\boldsymbol{\theta}$ are adjusted to reduce loss on the new data while preserving the general abilities learned during pre-training. Standard fine-tuning updates every parameter; parameter-efficient approaches keep the underlying backbone frozen and adapt only a small fraction of the parameter count through trainable side-modules. CAEM uses the parameter-efficient path through the low-rank adaptation (LoRA) layer of Hu et al.~\cite{hu2022lora} on the attention and feed-forward projections, with the trainable footprint at roughly two percent of the backbone parameter count and the underlying backbone frozen across every cycle of the self-improvement loop. The empirical evidence motivating that choice over full-parameter fine-tuning is reviewed in \Cref{sec:lit-self-improvement}.

\subsection{Hallucination Taxonomy}
\label{sec:hallucination-taxonomy}

Hallucinations in language model outputs can be understood along two main dimensions. The first dimension separates \textit{intrinsic hallucinations}, which directly contradict information provided in the source input, from \textit{extrinsic hallucinations}, which introduce information that can be neither supported nor contradicted by the source \cite{ji2023survey}. For instance, if a biography states that ``Einstein was born in Germany'' and the model generates ``Einstein was born in France,'' this is an intrinsic hallucination because it contradicts the stated source. But if the model generates ``Einstein played violin exceptionally well'' without any evidence in the source, this is an extrinsic hallucination even if the claim happens to be true: the defining characteristic is that it cannot be verified from the input, not that it contradicts anything in it.

The second dimension distinguishes \textit{factuality hallucinations}, which conflict with verifiable world knowledge, from \textit{faithfulness hallucinations}, which contradict user input or task context \cite{huang2025survey}. For example, a model that states ``photosynthesis produces oxygen from carbon dioxide and water'' is factually accurate, whereas one that states ``photosynthesis occurs in mitochondria'' produces a factuality hallucination. Faithfulness focuses on whether the output of the model remains consistent with the supplied input or context, especially when summarising documents or following instructions.

This research focuses on factuality hallucinations, which occur in knowledge-based reasoning tasks. Here, correct answers can be verified against trusted sources or logical rules. Faithfulness to retrieved evidence also matters for retrieval-augmented generation, where outputs must be grounded in retrieved documents.

Within factuality hallucinations, four subclasses are distinguished operationally because the four signal families of CAEM's verifier map onto them.\footnote{The nine signals are defined formally in \Cref{ch:methodology}; this chapter names them only where needed for motivation.} \textit{Confabulations}, in the sense of Farquhar et al.~\cite{farquhar2024semantic}, are answers the model might equally have produced differently on a resample, even when the individual token probabilities look confident; they are targeted by the sample-set agreement signals, which compare several resampled reasoning chains with one another ($s_{\text{avg}}$) and with the served answer ($p_{\text{entail}}$). \textit{Ungrounded factual claims} are answers without support in retrieved evidence; they are targeted by the three external-grounding signals, which score the answer against retrieved passages. \textit{Off-topic answers} that drift from the question are targeted by the question-answer relevance signal $q_{\text{a,rel}}$. \textit{Confidently wrong} answers, internally confident but externally unsupported, are caught by the confabulation early-exit gate, which fires when a high internal-confidence signal coincides with low grounding. Outside this scope, CAEM does \textit{not} target instruction-following failures, intermediate-reasoning-step errors inside chain-of-thought traces (only the final answer is verified), intrinsic faithfulness hallucinations in summarisation or translation (where the grounding source is the input document rather than world knowledge), or input-contradiction hallucinations that conflict with user-supplied premises. These exclusions set the failure-mode boundary within which the results in \Cref{ch:evaluation} should be interpreted.

\subsection{Natural Language Inference}

Natural Language Inference (NLI) is the task of determining whether a hypothesis sentence logically follows from a premise sentence \cite{barhaim2006second}. The PASCAL Recognising Textual Entailment challenges framed this as a binary decision (entailed or not entailed); later NLI corpora, including the Multi-Genre NLI corpus on which RoBERTa-Large-MNLI is fine-tuned, use three classes. Given premise $P$ and hypothesis $H$, a three-class NLI model classifies the relationship as:
\begin{itemize}
    \item \textbf{ENTAILMENT}: $P$ implies $H$ is true
    \item \textbf{CONTRADICTION}: $P$ implies $H$ is false  
    \item \textbf{NEUTRAL}: $P$ neither confirms nor denies $H$
\end{itemize}

For example, if the premise is ``The Beatles were formed in Liverpool in 1960'', the hypothesis ``The Beatles were an English band'' is entailed, the hypothesis ``The Beatles were formed in London'' is a contradiction, and the hypothesis ``The Beatles recorded many albums'' is neutral.

CAEM's verifier uses MiniCheck-Flan-T5-Large \cite{tang2024minicheck} as its entailment judge, a Flan-T5-Large checkpoint fine-tuned on synthetic claim-support data. MiniCheck is a binary judge: it returns the probability that a claim is supported by the given evidence and has no separate contradiction class, so a contradicted claim simply receives a low support probability. A second judge for long hypotheses was considered and dropped before the main run (\Cref{ch:methodology}), so every hypothesis is scored by MiniCheck, truncated at its 512-token input limit. Generic three-class NLI models such as RoBERTa-Large-MNLI \cite{liu2019roberta}, which reaches about 90.2\% on the MNLI matched development set, are the conventional alternative, but they transfer poorly to the hallucinations of language models: Sinha~\cite{sinha2025semantic} reports a one-hundred-percent false-positive rate at ninety-five-percent recall for a near-state-of-the-art generic NLI model on HaluEval~\cite{li-etal-2023-halueval}. Entailment checked in both directions can also decide whether two answers mean the same thing: Kuhn et al.~\cite{kuhn2023semantic} treat two answers as equivalent if $P$ entails $H$ and $H$ entails $P$. This is the basis of semantic entropy, which CAEM uses as an external baseline (\Cref{sec:external-baselines}).

\subsection{Sentence Embeddings and Semantic Similarity}

Sentence embeddings map variable-length text sequences to fixed-dimensional dense vectors in a semantic space where similar meanings have high cosine similarity. Sentence BERT (SBERT) \cite{reimers-gurevych-2019-sentence} produces sentence embeddings by fine-tuning BERT in a siamese network structure, first on natural language inference data (SNLI and MultiNLI) and optionally on semantic textual similarity data. The model used here, all-mpnet-base-v2 from the same Sentence-Transformers family (built on an MPNet backbone rather than the original BERT), generates 768-dimensional embeddings.

Cosine similarity between vectors $\mathbf{v}_1$ and $\mathbf{v}_2$ is computed as:
\begin{equation}
\text{sim}(\mathbf{v}_1, \mathbf{v}_2) = \frac{\mathbf{v}_1 \cdot \mathbf{v}_2}{\|\mathbf{v}_1\| \|\mathbf{v}_2\|} = \cos(\theta)
\end{equation}
where $\theta$ is the angle between vectors. Similarity ranges from $-1$ (opposite) to $1$ (identical), with $0$ indicating orthogonality.

Semantic similarity is different from string similarity. For instance, the questions ``Is 17 prime?'' and ``Is the number 17 a prime number?'' share few words but are close in meaning. CAEM uses semantic similarity for memory retrieval (finding similar stored questions), novelty detection (avoiding redundant storage), and the self-consistency signal (measuring agreement between reasoning chains). Nearest-neighbour search uses the FAISS library \cite{8733051}; at the memory sizes reached in this thesis the search is exact, and FAISS's approximate indexes take over only at much larger sizes.

\subsection{Self-Consistency Decoding}
\label{sec:lit-self-consistency}

Self-consistency \cite{wang2023selfconsistency} improves reasoning by generating several independent reasoning chains for the same question and selecting the most common answer. Unlike greedy decoding, which produces a single output, self-consistency samples $M$ chains at a temperature $T > 0$ to introduce diversity.

CAEM adapts self-consistency for verification rather than answer selection. Given a question $q$ and a generated answer $a$, the verifier generates $M$ additional chains and measures their agreement as the average pairwise similarity
\begin{equation}
s_{\text{avg}} = \frac{2}{M(M-1)} \sum_{1 \le i < j \le M} \text{sim}(\mathbf{v}_i, \mathbf{v}_j)
\end{equation}
where $\mathbf{v}_i$ is the embedding of chain $i$ and only distinct pairs are counted. High agreement suggests that several reasoning paths converge on the same content; low agreement suggests uncertainty. CAEM pairs $s_{\text{avg}}$ with the chain-to-answer entailment signal $p_{\text{entail}}$, computed on the same $M{=}3$ chains so that the decoding cost is paid once; together they form the sample-set agreement family of the verifier. The normalised semantic entropy of Farquhar et al.~\cite{farquhar2024semantic}, computed over $K{=}10$ samples at temperature $T{=}1.0$, was originally a third member of this family. It was dropped before the main run because its learned weight in the cycle-zero composite fit was effectively zero (about $-0.0005$), so it added nothing beyond the other signals, while its ten extra samples per answer were the largest single cost in the verifier (about four seconds per answer). The kernel-language-entropy generalisation of Nikitin et al.~\cite{nikitin2024kle}, a candidate replacement that needs the same extra samples, was not adopted either.

\subsection{Catastrophic Forgetting and Continual Learning}

Catastrophic forgetting, first described by McCloskey and Cohen~\cite{MCCLOSKEY1989109}, occurs when neural networks learn new tasks and suddenly lose their ability to perform previously learned tasks. Standard gradient descent focuses only on improving performance for the current task and does not protect earlier knowledge, so parameter updates that help the new task tend to weaken the old ones. In large language models, fine-tuning on a new domain likewise erodes performance outside that domain, and Biderman et al.~\cite{biderman2024lora} found that the erosion grows with training duration.

Elastic Weight Consolidation (EWC) \cite{doi:10.1073/pnas.1611835114} addresses this by estimating, through the Fisher information, which parameters matter for earlier tasks, and adding a quadratic penalty that discourages large changes to them:
\begin{equation}
\mathcal{L} = \mathcal{L}_{\text{new}} + \frac{\lambda}{2} \sum_{i} F_i (\theta_i - \theta_i^*)^2
\end{equation}
where $F_i$ measures parameter importance, $\theta^*$ are the previous parameters, and $\lambda$ controls the strength of the penalty. A cheaper variant drops the Fisher weights and penalises all parameters equally, $\frac{\lambda}{2}\lVert \boldsymbol{\theta} - \boldsymbol{\theta}_{\text{prev}} \rVert^2$; Kirkpatrick et al.~\cite{doi:10.1073/pnas.1611835114} show that this uniform penalty is weaker than EWC because it protects all weights equally and leaves little capacity for the new task. CAEM uses neither penalty on its shipped path. It bounds drift between cycles with three other mechanisms instead: a strict quality filter on the training pool; low-rank adaptation \cite{hu2022lora}, which trains only about two percent of the parameters on top of a frozen backbone and forgets less than full fine-tuning (\Cref{sec:lit-self-improvement}); and a retention guard that rolls the adapter back, while keeping the episodic memory, whenever accuracy on held-out probe sets falls below a fixed fraction of its starting value. The details are given in \Cref{ch:methodology}.

The conceptual basis for understanding CAEM's nine-signal post-generation verifier, pre-routing confidence, episodic-memory architecture, and constrained self-improvement loop is established by the groundwork above. The following sections describe how prior research develops and validates these individual components, preparing for their integration in Chapter~\ref{ch:methodology}.

\section{Review of Existing Research}

Where the Preliminaries section above defined the concepts and notation (hallucination taxonomy, natural language inference, sentence embeddings, self-consistency, catastrophic forgetting), this section surveys the prior empirical and methodological literature on hallucination measurement, verification, confidence estimation, memory-augmented architectures, selective retrieval, self-improvement and continual learning, the external baselines used in this thesis, and evaluation benchmarks. Each subsection identifies key contributions and limitations that motivate CAEM's architectural choices, so that a reader who has already absorbed the Preliminaries can move directly to the comparative positioning without re-reading definitional material.

\subsection{Hallucination Problem and Measurement}

Understanding the nature and characteristics of hallucination in LLMs requires methodologies like taxonomic frameworks and quantitative measurement. Early studies of hallucination focused on summarisation and data-to-text generation. Ji et al. \cite{ji2023survey} published the first major survey in ACM Computing Surveys. It distinguishes intrinsic hallucinations, which contradict the source material, from extrinsic hallucinations, which can be neither supported nor contradicted by it. Their work shows that hallucinations arise at several stages, including data collection that embeds factual errors, training procedures that reinforce misleading correlations, and inference methods that fail to account for uncertainty.

Huang et al. \cite{huang2025survey} extended this study specifically for large language models in ACM Transactions on Information Systems, distinguishing factuality hallucinations from faithfulness hallucinations. Their analysis traces the causes of hallucination across the model lifecycle to data, training and inference, and ties each family of mitigation methods to the cause it addresses. This suggests that effective mitigation should intervene at several stages instead of relying solely on detection applied after the model has already produced an answer, which supports CAEM's multi-component strategy combining memory, verification, and iterative learning.

Quantitative measurement studies of hallucination prevalence reveal concerning statistics across domains and model scales. Lin et al. \cite{lin-etal-2022-truthfulqa} introduced TruthfulQA at ACL 2022, a benchmark of 817 questions across 38 categories designed to elicit imitative falsehoods where models reproduce common human misconceptions found in training data. Their result demonstrates an inverse scaling effect: the best GPT-3 175B configuration achieved only 21\% on the truthful-and-informative metric (and 58\% on truthfulness alone), against 87\% and 94\% respectively for the human baseline. Larger models often performed worse than smaller ones. This suggests that scaling alone is not sufficient for truthfulness, which motivates architectural innovation.

The inverse-scaling result has a direct implication for mitigation: if adding parameters can reduce truthfulness on some tasks, mitigation must work beyond parameter growth. CAEM does so by adding external verification and memory components whose operation does not depend on the size of the model.

Min et al.\ \cite{min-etal-2023-factscore} introduced FActScore at EMNLP 2023, which evaluates long-form generations by decomposing the answer into atomic claims and reporting the share supported by a reliable knowledge source rather than scoring the answer as a whole. Its automated estimator stays within an error of 2\% of human evaluation, and on biography generation it finds that only 58\% of ChatGPT's atomic facts are supported, so even highly capable models produce many unsupported claims. CAEM borrows this decomposition for one of its grounding signals, $p_{\text{ground\_atomic}}$, which splits the answer into atomic claims, scores each against the retrieved passages, and keeps the weakest score. The step runs only on answers long enough to contain several claims; the short answers of the five-benchmark panel almost never qualify, so in the realised run this signal coincided with the passage-mean grounding signal on every evaluated answer (\Cref{ch:methodology}).

Domain-specific studies show even higher hallucination rates in specific contexts. Alkaissi and McFarlane \cite{alkaissi2023artificial} asked ChatGPT to explain a clinical finding and to supply supporting references; none of the five paper titles it returned existed. Chelli et al.\ \cite{info:doi/10.2196/53164} checked the references that chatbots generated for systematic reviews and reported hallucination rates of 39.6\% for GPT-3.5, 28.6\% for GPT-4, and 91.4\% for Google Bard.

Three points follow for CAEM's design. First, hallucination occurs across models, tasks and domains, at rates that vary widely with context (from 28.6\% to 91.4\% of generated references in the single study of Chelli et al.~\cite{info:doi/10.2196/53164}). Second, larger models do not automatically hallucinate less, so the structure around the model has to change. Third, scoring individual claims against evidence is more informative than judging whole responses, which motivates CAEM's evidence-based grounding signals.

\subsection{Verification Methods}

Verification methods are used to check whether a model's reasoning is correct before the information is stored or shown to users. This subsection reviews chain-of-thought prompting, self-consistency decoding, semantic entropy, natural-language inference and conformal filtering, the approaches from which CAEM's verifier draws its signals or against which its design is positioned.

Chain of thought prompting, which was introduced by Wei et al. \cite{wei2022chain} at NeurIPS 2022, improves reasoning quality by prompting models with worked examples whose intermediate steps precede the final answer. Showing these intermediate steps leads to stronger reasoning and provides better starting points for verification. When tested on the PaLM 540B model, chain of thought greatly improved performance on math reasoning tasks, increasing accuracy on GSM8K from 17.9\% to 56.9\%. The authors found that the technique is an emergent ability of model scale: below roughly 100 billion parameters chain-of-thought prompting did not help and could lower performance, a pattern that held for the instruction-tuned InstructGPT models in their GPT-3 series as well. Kojima et al.~\cite{kojima2022large} later showed that the single trigger ``Let's think step by step'' elicits chain-of-thought reasoning without worked examples. Instruction tuning does not remove the small-scale effect: Chung et al.\ \cite{JMLR:v25:23-0870} report that their instruction-tuned Flan-T5-Large (780 million parameters) still scores lower with chain-of-thought prompting than with direct answering on MMLU (40.5 versus 45.1 percent) and on BIG-Bench Hard (31.5 versus 37.5 percent). Whether chain-of-thought prompting helps CAEM's three-billion-parameter instruction-tuned backbone, Qwen-2.5-3B-Instruct~\cite{qwen25}, is therefore an empirical question, which the chain-of-thought baseline in \Cref{ch:evaluation} answers.

For CAEM, chain-of-thought prompting matters in two ways: it is the reasoning format in which the generator answers, and it supplies the resampled reasoning chains from which the sample-set agreement signals are computed. The verifier, however, scores only the final answer and not the individual reasoning steps (\Cref{sec:hallucination-taxonomy}), so an invalid step that still reaches a correct conclusion goes undetected.

Wang et al.\ \cite{wang2023selfconsistency} extended this idea in 2023 by introducing self-consistency. Instead of producing a single reasoning path, the model generates several different explanations and then selects the most common answer. The key idea is that difficult and complex problems usually have multiple valid reasoning paths that lead to the same correct answer. On the other hand, incorrect answers rarely achieve consistency across independent derivations. Reported gains over chain-of-thought prompting reach +17.9 percentage points on GSM8K, +11.0 on SVAMP, +12.2 on AQuA, and +6.4 on StrategyQA; these are the largest gains across the four models tested, and for PaLM 540B the corresponding gains are +17.9, +7.6, +12.5 and +6.3.

CAEM uses self-consistency differently: instead of selecting the majority answer, it measures how far the resampled chains agree in meaning. Strong agreement raises the composite score and weak agreement lowers it. The cost is a small number of extra generations per answer, which keeps the signal practical for routine verification.

Semantic entropy, introduced by Farquhar et al.\ \cite{farquhar2024semantic} in Nature in 2024, addresses a limitation of token-level entropy: answers that express the same idea in different words are grouped together by bidirectional entailment, and uncertainty is measured across these meaning-based groups rather than across surface strings.

Averaged over thirty combinations of tasks and models (LLaMA~2, Falcon and Mistral families), semantic entropy reaches an AUROC of $0.790$ for detecting confabulations, ahead of naive entropy ($0.691$), P(True) ($0.698$), and an embedding-regression baseline ($0.687$). The method targets arbitrary, sample-to-sample errors: a model that is consistently wrong gives the same answer on every sample and therefore looks certain.

CAEM originally included semantic entropy as a verifier signal and dropped it before the main run (\Cref{sec:lit-self-consistency}); it remains in the thesis as an external baseline that abstains when entropy is high (\Cref{sec:external-baselines}).

Natural language inference provides direct verification by checking whether a claim is supported by evidence. Thorne et al.\ \cite{thorne-etal-2018-fever} introduced the FEVER dataset at NAACL 2018, a fact-verification task closely related to NLI, with 185,445 claims generated from Wikipedia; unlike NLI, the evidence for each claim must first be retrieved rather than being given. Systems must classify each claim as SUPPORTS, REFUTES or NOT ENOUGH INFO.

Pre-trained NLI models such as RoBERTa-Large-MNLI reach about 90\% accuracy on MNLI and verify individual factual statements well, but they can miss reasoning errors that contradict no fact. A model could state ``Sydney is Australia's largest city'', which is true, and conclude ``therefore Sydney is the capital'', an invalid inference; NLI would confirm the premise and miss the faulty step, which is why complementary signals are needed.

A more recent strand of work brings conformal prediction to the language-model claim level. Mohri and Hashimoto~\cite{mohri2024conformal} introduced \emph{Language Models with Conformal Factuality Guarantees} at ICML 2024, decomposing each generated answer into atomic sub-claims, scoring each, and removing claims below a calibrated threshold $\tau$ such that the surviving sub-claims are jointly correct with probability at least $1-\alpha$ on a small annotated calibration set. On FActScore their method raises factuality from about 30\% to 80\% while retaining most sub-claims. Yadkori et al.~\cite{yadkori2024abstention} extended this with conformal-risk-control to bound the hallucination rate at a user-specified target through learned abstention rather than claim removal, providing the direct precedent for CAEM's abstention class in the four-outcome decision tree. Quach et al.~\cite{quach2024conformal} present \emph{Conformal Language Modeling} at ICLR 2024, calibrating a stopping rule, a rejection rule, and a confidence rule simultaneously through Learn-then-Test multiple-hypothesis testing. Cherian et al.~\cite{cherian2024conformal} sharpen the Mohri--Hashimoto framework with conditional guarantees and a boosted score, learning a linear combination of several per-claim scores by differentiating through the conformal procedure; the closest precedent for CAEM's boost layer over the per-signal log-odds in \Cref{sec:composite-boost} is the logistic-regression multi-scoring of Valentin et al.~\cite{valentin2024hallucination}. CAEM considered split-conformal calibration for its storage gate and rejected it because the distribution shift between its calibration and evaluation folds broke the coverage premise (\Cref{sec:conformal-gate}). The deployed gate is therefore a fixed-threshold rule on the calibrated probability, in the lineage of the post-hoc calibration analyses of Niculescu-Mizil and Caruana~\cite{niculescu2005predicting} and the selective-classification framing of Geifman and El-Yaniv~\cite{geifman2017selective}.

These methods supply complementary signals to CAEM's nine-signal verifier. Self-consistency and chain-to-answer entailment form the sample-set agreement family ($s_{\text{avg}}$, $p_{\text{entail}}$). Entailment against retrieved evidence forms the external-grounding family: the best-passage score $p_{\text{ground\_max}}$, the mean-passage score $p_{\text{ground\_mean}}$, and the length-gated atomic-fact score $p_{\text{ground\_atomic}}$ in the FActScore convention~\cite{min-etal-2023-factscore}. Three internal-calibration signals (token probability, Monte Carlo dropout variance, and an internal-state composite of the two) and the question-answer relevance signal complete the nine. Instead of a fixed weighted sum, each signal is mapped to a calibrated probability of correctness by isotonic regression~\cite{niculescu2005predicting}, and the calibrated signals are combined by a regularised logistic regression in the manner of Valentin et al.~\cite{valentin2024hallucination}; the full specification is in \Cref{ch:methodology}. The point for this chapter is that the literature supplies families of signals whose failure modes differ, which justifies a calibrated multi-signal composite over any single-signal detector.

\subsection{Confidence Estimation}

Reliable confidence estimation helps systems tell the difference between answers that can be trusted and ones that need extra checking from outside sources. This subsection looks at basic methods for measuring uncertainty and explains how these ideas have been adapted for use with large language models.

Gal and Ghahramani \cite{pmlr-v48-gal16} laid the theoretical foundation for practical uncertainty estimation in their ICML 2016 paper ``Dropout as a Bayesian Approximation''. They showed that dropout, normally used only during training, can be read as approximate Bayesian inference over the model's parameters. Keeping dropout active at test time and running the model several times then measures how much the outputs vary. This variation estimates epistemic uncertainty, which comes from the model's limited knowledge, as opposed to aleatoric uncertainty, which comes from noise in the data and cannot be reduced by more knowledge.

The framework shows that a network with dropout before each weight layer approximates a deep Gaussian process, whose predictive mean and variance can be estimated by Monte Carlo sampling: $K$ forward passes with different dropout masks, followed by the variance of the outputs. The estimate needs no ensemble training and no change to the loss function, and it applies to any network that already uses dropout, which makes it convenient for existing large language models.

For CAEM, MC Dropout is one of the three internal-calibration signals inside the post-generation verifier ($u_{\text{token}}$, $u_{\text{dropout}}$, and the derived $u_{\text{internal}} = 0.5\,u_{\text{token}} + 0.5\,(1 - u_{\text{dropout}})$ that combines them). It quantifies model uncertainty through output variance: high variance indicates the model is unstable with respect to its own weights and suggests that the composite score should be pulled down, while low variance indicates confident generation. This signal is insufficient on its own, since models can be confidently wrong, but the computational cost of the registered five forward passes is acceptable inside the verifier, which only runs on queries that have already been dispatched to Tier~2 or Tier~3 generation. MC~Dropout is deliberately \textit{not} used in the pre-routing confidence $\upre$, which must be cheap enough to run on every query before tier selection.

A major problem with neural network confidence is that models are often too confident in their answers. Guo et al. \cite{pmlr-v70-guo17a} demonstrated this at ICML 2017. They showed that modern deep networks predict confidence scores that substantially exceed their actual accuracy; a ResNet-110 on CIFAR-100, for example, has an expected calibration error of 16.5 percentage points before calibration. They popularised the Expected Calibration Error (ECE) of Naeini et al.~\cite{naeini2015obtaining} as the standard metric, dividing predictions into bins by confidence and measuring the gap between confidence and accuracy within each bin.

Their proposed solution, temperature scaling, achieves remarkable effectiveness with minimal complexity. The method divides logits by a learned scalar temperature $T > 0$ optimised on a validation set: $\hat{y} = \sigma(\mathbf{z}/T)$ where $\mathbf{z}$ are logits and $\sigma$ is softmax. Temperature scaling preserves accuracy (argmax unchanged) of classification while calibrating probabilities. Even though only a single parameter is used, it outperforms more complex methods including histogram binning, isotonic regression, and Bayesian approaches. The simplicity and effectiveness make temperature scaling the standard calibration method applied after training is complete.

CAEM applies temperature scaling to its pre-routing confidence $\upre$, fitted per benchmark on a held-out calibration split, so that a stated confidence of eighty percent corresponds to roughly eighty percent correct answers. The post-generation composite $\ustored$ is calibrated differently, by per-signal isotonic regression and a logistic combination (\Cref{ch:methodology}). Both calibrations make it possible to set meaningful thresholds for tier selection, for the four-outcome storage decision, and for the training-pool filter $\tau_{\text{train}}$.

Kadavath et al.\ \cite{kadavath2022language} from Anthropic explored self-evaluation in large language models in a 2022 preprint, introducing two ways of extracting confidence from a model itself. The first, P(True), asks the model to estimate the probability that its own earlier answer is correct. The second, P(IK), the probability that ``I know'' the answer, is predicted before any answer is proposed, by a model trained for that purpose. Their findings reveal nuanced calibration properties. Larger models tend to be well calibrated on multiple-choice questions, meaning that their confidence matches how often they are actually correct. However, policies fine-tuned with reinforcement learning from human feedback appear badly miscalibrated, because reinforcement learning collapses predictions toward the behaviours that receive the most reward, although a single temperature adjustment largely restores calibration.

The finding that reinforcement learning from human feedback can hurt calibration matters for deployment. A model trained to be helpful may express high confidence to sound authoritative even when it is uncertain, the opposite of what high-stakes use requires. CAEM therefore relies on several external confidence signals rather than on the model's own confidence alone.

Xiong et al.\ \cite{xiong2023can} studied confidence estimation for large language models in detail, comparing white-box methods, which use internal information such as token logits and sequence probabilities, with black-box methods, which rely on external signals such as verbalised confidence or consistency checks. Verbalised confidence is overconfident and discriminates poorly: even for GPT-4 its average AUROC for telling correct answers from incorrect ones is about 0.63, close to the 0.5 of random guessing. Sampling-consistency signals do much better on some tasks (up to 0.93 AUROC on GSM8K), but no single method consistently outperforms the others, and all of them struggle on tasks that require professional knowledge.

CAEM's design follows from this finding that no single signal is reliable across tasks, and splits confidence into two stages. The first, the pre-routing confidence $\upre$, fuses two cheap signals (the token probability of a short draft decode and a layer-convergence ratio over the query's hidden states) through a temperature-scaled sigmoid; it is computed for every query and drives tier selection and the safety override. The second is the nine-signal post-generation verifier, which runs only after zero-shot or retrieval-augmented generation and produces the calibrated composite $\ustored$ that drives the storage decision.

Recent work by Tian et al. \cite{tian2023just} asks how calibrated confidence can best be elicited from models fine-tuned with human feedback. Their study found that the raw token probabilities of such models are consistently overconfident, and that confidence verbalised as output tokens is typically better calibrated, often reducing the expected calibration error by a relative 50\%. Because the best confidence source differs between model families and tasks, this reinforces CAEM's strategy of combining implicit signals like token probability, dropout, and self-consistency with explicit ones, and of recalibrating every signal rather than trusting any raw score.

A complementary 2024 survey by Geng et al.~\cite{geng2024survey} catalogues white-box and black-box confidence-estimation methods for language models across the same axes that drive CAEM's verifier ensemble, and is the positioning document for CAEM's multi-signal design choice. Two recent single-signal methods deserve specific mention as alternatives considered for CAEM. Chen et al.~\cite{chen2024inside} introduced INSIDE / EigenScore at ICLR 2024, computing the log-determinant of the covariance matrix over middle-layer hidden states across $K$ resampled responses; the method reports AUROC of $0.80$ to $0.83$ on TriviaQA and clear gains over other unsupervised signals on CoQA, SQuAD and Natural Questions (on TriviaQA, perplexity scores slightly higher), and is registered in CAEM as a candidate replacement for the dropout signal that was not adopted in the locked configuration. Nikitin et al.~\cite{nikitin2024kle} introduced Kernel Language Entropy at NeurIPS 2024, generalising Farquhar et al.'s hard-cluster semantic entropy to a soft-kernel von Neumann entropy over pairwise semantic similarities; KLE was considered as a replacement should the semantic-entropy family return, but was not adopted because, like semantic entropy, it needs many extra samples per answer (\Cref{sec:lit-self-consistency}). The selective-prediction precision-versus-coverage tradeoff that CAEM reports in \Cref{ch:evaluation} follows the form pioneered by Kamath, Jia, and Liang~\cite{kamath2020selective} for selective question answering under domain shift, where a calibrator over feature signals predicts $P(\text{correct})$ and the selective rule answers only when the calibrator's confidence exceeds the registered precision target.

The confidence-estimation literature yields four points for CAEM. A single confidence signal is not reliable enough for safe routing, so several must be combined. Systematic overconfidence calls for calibration after training. White-box signals from inside the model complement black-box signals from output consistency. Preference tuning can degrade calibration, and CAEM's backbone is itself preference-tuned (its post-training uses DPO and GRPO \cite{qwen25}), so no raw confidence is used without recalibration. Together these points shape CAEM's two-stage confidence design: a temperature-scaled pre-routing score $\upre$ on every query, and a nine-signal verifier whose calibrated composite $\ustored$ drives the four-outcome storage decision. Both calibrations are fitted on labelled calibration data and refitted at every cycle boundary (\Cref{ch:methodology}).

\subsection{Memory-Augmented Architectures}

External memory architectures give neural networks a place outside their parameters where facts can be written and read. This subsection traces these ideas from early differentiable memories to retrieval-augmented generation.

Graves et al.\ \cite{graves2014neural} introduced Neural Turing Machines (NTMs) in 2014. NTMs gave a neural network an external memory that it can read from and write to in a differentiable, trainable way. The architecture combines a controller network with an addressable memory matrix, supporting both content-based addressing (via cosine similarity) and location-based addressing (via shift operations). The key innovation is a read/write mechanism that enables end-to-end training via backpropagation. As a result, NTMs successfully learned algorithmic tasks including copying, sorting, and associative recall.

For CAEM, the lesson is that an external memory can give a model abilities its parameters alone lack. NTMs were applied mainly to algorithmic tasks, but the same principle extends to facts: CAEM's episodic memory stores verified question-answer pairs together with their reasoning, so that a repeated question can be answered by lookup.

Weston et al. \cite{weston2015memory} introduced Memory Networks at ICLR 2015. Their model scaled to millions of memory slots. It also showed that a system could perform several rounds of reasoning to answer a question. The architecture consists of four learned components. The components are: turning the input into features, updating the memory, turning the memory output into features, and finally producing a response. Furthermore, Memory Networks performed strongly on bAbI question-answering tasks that required reasoning over multiple facts. Sukhbaatar et al. \cite{NIPS2015_8fb21ee7} extended this model with end-to-end Memory Networks at NeurIPS 2015. They removed the need for supervision on which memories to retrieve by using a continuous attention mechanism over the memory. This line of work motivates CAEM's use of stored knowledge as an explicit, inspectable memory.

Graves et al. \cite{graves2016differentiable} presented the Differentiable Neural Computer (DNC) in Nature 2016. It built on the idea of Neural Turing Machines but added smarter memory tools like dynamic allocation. It also keeps track of the order in which things were written. DNCs achieved a 3.8\% mean error on the bAbI tasks, compared with 7.5\% for the best previous jointly trained model, which suggests that learned memory management improves performance.

Lewis et al. \cite{NEURIPS2020_6b493230} introduced Retrieval-Augmented Generation (RAG) at NeurIPS 2020. They showed that combining what a model knows in its parameters with information stored outside the model can make its answers much more accurate. RAG adds a large searchable index of Wikipedia to a BART model with four hundred million parameters. It uses Dense Passage Retrieval to look up relevant documents. Then it works out the final answer by considering the generation probability across all the documents it retrieved. The architecture achieves substantial improvements over parametric-only models: 44.5\% exact match on Natural Questions (versus 41.5\% prior SOTA), 68.0\% on TriviaQA (versus 60.5\%), and 45.5\% on WebQuestions (versus 44.7\%).

Critically, RAG demonstrates index hot-swapping: updating the document index without retraining achieves 70\% accuracy on questions about updated world knowledge, proving that non-parametric memory can be updated without expensive retraining. This property directly motivates CAEM's episodic memory, which can be updated continuously through verification without model fine-tuning, though CAEM additionally employs iterative fine-tuning to consolidate learned patterns.

The limitation of standard RAG for CAEM's purposes is that retrieved knowledge is not verified before use and does not accumulate verified reasoning chains. CAEM extends RAG by adding verification layers and storing verified outputs for future retrieval, creating a growing knowledge base rather than relying solely on static external corpora.

Johnson et al.\ \cite{8733051} developed FAISS in 2019 for similarity search over billions of items, enabling efficient memory-augmented systems. The library achieves an 8.5$\times$ speedup over prior GPU methods through optimised k-selection, and supports index types such as an inverted file with product quantisation, which searches in sub-linear time. For CAEM, FAISS provides fast similarity search over the stored episodes, of which there were $12\,175$ at the end of the realised trajectory.

Reimers and Gurevych \cite{reimers-gurevych-2019-sentence} introduced Sentence-BERT (SBERT) at EMNLP 2019. SBERT produces semantically meaningful sentence embeddings through siamese-network fine-tuning. The speed improvement over standard BERT is substantial: finding the most similar sentence pair in a group of ten thousand sentences takes about sixty-five hours with standard BERT, but only about five seconds with SBERT, because one can compare the precomputed embeddings directly instead of running a heavy cross-encoder over every pair. This efficiency gain is what makes CAEM's memory search possible: each query is turned into a 768-dimensional embedding, and the system can look through the entire memory index in real time.

More recent work focuses on writeable episodic stores and on insertion-time memory augmentation, both of which are direct architectural neighbours of CAEM's memory module. Das et al.~\cite{larimar2024das} introduced Larimar at ICML 2024, an associative key-value episodic memory with one-shot edits, selective forgetting, and an order-of-magnitude faster knowledge-editing path than direct fine-tuning; Larimar establishes that a writeable episodic store is a feasible alternative to parametric storage but does not associate a quality scalar with each entry. Xu et al.~\cite{xu2025amem} introduced A-MEM, a Zettelkasten-inspired agentic memory that writes structured textual notes at insertion and evolves links between memories over time; A-MEM is the closest architectural neighbour to CAEM in the sense that both write structured metadata at insertion, but A-MEM's metadata is textual (tags, links) while CAEM's is a single calibrated scalar storage score that drives a tiered router rather than evolving link structure. Adaptive-RAG \cite{jeong2024adaptive}, reviewed in \Cref{sec:lit-selective-retrieval}, shares CAEM's count of three tiers but routes on query complexity at inference time rather than on the confidence stored with each memory entry, which is the distinction CAEM's router draws (\Cref{sec:router}).

These developments establish the feasibility and benefit of external knowledge storage. NTMs show that differentiable memory supports structured reasoning; Memory Networks scale to millions of facts with multi-hop reasoning; RAG shows that non-parametric knowledge improves factual accuracy and can be updated; FAISS and Sentence-BERT provide the infrastructure for practical deployment. CAEM builds on these advances in three ways. First, a query is routed by a combined score $\mathcal{S} = \lambda\, s + (1-\lambda)\, \ustored$ of its similarity $s$ to the nearest stored entry and that entry's stored confidence, so a close match stored with low confidence is not served from memory. Second, an independent safety override sends any query whose pre-routing confidence $\upre$ falls below a fixed floor to the retrieval-augmented default, however close its memory match, after which the retrieval-utility classifier makes the final choice between retrieval and zero-shot generation. Third, the memory is filled only through the four-outcome storage decision (\textsc{Store}, \textsc{Defer}, \textsc{Abstain}, \textsc{Discard}), with a bounded buffer for borderline answers, so the knowledge base grows only from answers that clear the storage threshold. These extensions distinguish CAEM from static-corpus RAG.

\subsection{Selective and Adaptive Retrieval}
\label{sec:lit-selective-retrieval}

The retrieval-augmented architectures of the previous subsection retrieve unconditionally on every query: every incoming question pays the retrieval cost, and the retrieved passages enter the generator's context whether the query benefits from external evidence or not. A complementary line of work argues that retrieval is not uniformly helpful and proposes architectures that decide whether to retrieve on a per-query basis. The selective-retrieval literature anchors the architectural design of \Cref{sec:ruc}, which adds a binary retrieval-utility classifier between CAEM's three-tier router and its retrieval-augmented fallback so that queries on which retrieval would amplify the question's misconception framing route to the zero-shot tier instead.

Mallen et al.~\cite{mallen2023whennot} established the empirical anchor for the selective-retrieval position at ACL 2023. They introduced PopQA, a fourteen-thousand-question benchmark of entity-centric factoid queries, and showed that large language models handle popular entities through their parametric memory more reliably than through retrieval, while the same models fail systematically on long-tail entities that retrieval can rescue. The architectural lesson is that retrieval utility is a function of the subject entity's popularity rather than a global hyperparameter, and the corresponding routing rule is a per-relation popularity threshold derived from Wikipedia pageview ranks. With GPT-3 davinci-003, the adaptive retrieval rule reaches 46.5\% accuracy on PopQA, 5.3 percentage points above the best non-adaptive method. The lesson for CAEM is that retrieval utility can be predicted in part from cheap question-side features such as entity popularity, which CAEM's retrieval-utility classifier includes among its inputs (\Cref{sec:ruc}).

Wang et al.~\cite{wang2023skr} extended this line to a learned classifier at EMNLP Findings 2023. Their self-knowledge routing system collects training labels by running each candidate question through both a no-retrieval pipeline and a retrieval-augmented pipeline, recording the per-pipeline exact-match score, and labelling the question as \emph{retrieval-helpful} when the retrieval-augmented exact match exceeds the no-retrieval exact match and as \emph{retrieval-hurt} otherwise. A classifier is then built on the labelled set (the paper compares direct prompting, in-context learning, a BERT-base classifier and a k-nearest-neighbour rule, of which the last performs best), and at inference the classifier dispatches each query to retrieval or to the zero-shot generator according to its prediction. The training-label rule (compare exact match across the two routings) is the direct precedent for the empirical-delta labelling rule that CAEM's retrieval-utility classifier inherits in \Cref{sec:ruc}; the architectural difference is that CAEM's classifier sits downstream of the existing tier-1 and tier-2 dispatch outcomes and gates only the retrieval-augmented fallback, while Wang's classifier replaces the dispatch entirely.

Jeong et al.~\cite{jeong2024adaptive} introduced Adaptive-RAG at NAACL 2024 as a three-class extension of the binary routing decision. Their classifier (a T5-large language model fine-tuned on automatically-collected complexity labels) dispatches each query to one of three classes (no retrieval, single-step retrieval, multi-step retrieval), trading routing accuracy for the cost of running a large model at routing time. The reported cost-accuracy frontier on a six-dataset open-domain question-answering suite (SQuAD, Natural Questions, TriviaQA, HotpotQA, 2WikiMultiHopQA, MuSiQue) beats single-strategy baselines and the unconditional retrieval-augmented baseline; for CAEM's purposes the T5-large routing cost is the design's main weakness, while the authors themselves single out the quality of the automatically assigned complexity labels. The architectural lesson for CAEM is that a learned classifier on automatically-collected outcome labels can decide the routing question, but that putting a large model at routing time costs more than the architecture gains; the follow-up of Marina et al.~\cite{moskvoretskii2025llm} extends Adaptive-RAG to the regime where the classifier carries no language-model forward pass and matches the routing accuracy of the language-model-based variant on the same suite. CAEM's retrieval-utility classifier sits at the lightweight end of this design space: the classifier itself is a logistic regression head, and its features reuse quantities the pipeline already computes, including a self-knowledge probe read from one forward pass of the base generator.

Asai et al.~\cite{asai2024selfrag} introduced Self-RAG at ICLR 2024 as the end-to-end fine-tuned alternative to the external-classifier approach. Self-RAG trains a single language model to emit four kinds of reflection tokens during generation (a \texttt{Retrieve} token that triggers retrieval, an \texttt{IsREL} token that judges passage relevance, an \texttt{IsSUP} token that judges support, and an \texttt{IsUSE} token that judges usefulness), so the routing decision is made by the same model that generates the answer. The deployed signal for retrieval routing is the probability the model assigns to the \texttt{Retrieve} token at decoding time; removing retrieval at inference time drops PopQA accuracy by roughly forty-five percent relative in the authors' ablation, the empirical receipt for the joint design's load-bearing role. The architectural cost is that the reflection tokens are tied to the LLaMA-family backbone Self-RAG fine-tunes on, and retargeting them to a different backbone changes enough of the system to invalidate the architectural comparison; this is the same constraint that prevents Self-RAG from entering the external baseline panel of \Cref{ch:evaluation} as a numerical baseline rather than a citation-level reference.

Baek et al.~\cite{baek2025probing} introduced Probing-RAG at NAACL Findings 2025 as a midpoint between the heavy classifier and the in-generator reflection tokens. They train a linear probe on the intermediate hidden states of the generator (rather than on its output tokens), with the probe trained under a balanced positive-negative protocol that explicitly enforces a fifty-fifty class ratio. The probe is trained to predict, from the hidden states, whether the model's current answer is correct and hence whether additional retrieval is needed; the reported accuracy across five open-domain question-answering datasets beats single-step retrieval, FLARE, DRAGIN, and Adaptive-RAG on the same suite while retrieving roughly half as often. The architectural lesson for CAEM is the most directly relevant of the five: the closest cousin of CAEM's retrieval-utility classifier on the feature axis is the hidden-state-as-routing-signal pattern that Probing-RAG establishes, with the methodological difference that Probing-RAG probes intermediate layers, whereas CAEM reads the base generator's last-token hidden state, without the adapter, as a single self-knowledge feature alongside thirty-five other question-side and passage-side features.

CAEM's retrieval-utility classifier reaches a pooled out-of-fold AUROC of $0.850$ under nested cross-validation on $3\,208$ labelled rows (\Cref{sec:ruc}). The systems reviewed above are evaluated on different question suites and with different metrics, so this figure is not directly comparable with theirs.

\subsection{Self-Improvement and Continual Learning}
\label{sec:lit-self-improvement}

Self-improvement mechanisms enable models to learn from interaction and feedback, while continual learning addresses the challenge of acquiring new knowledge without forgetting previous capabilities. This subsection reviews methods for iterative refinement and catastrophic forgetting mitigation that inform CAEM's self-improvement loop.

Kirkpatrick et al.\ \cite{doi:10.1073/pnas.1611835114} introduced Elastic Weight Consolidation (EWC), defined in the Preliminaries, in PNAS in 2017. Its Fisher-weighted penalty lets parameters that matter little for earlier tasks adapt freely while protecting those that matter. On ten Atari games learned in sequence, EWC let a single agent learn many games without catastrophic forgetting, whereas plain stochastic gradient descent forgot earlier games, although EWC did not reach the score of ten separately trained agents. For CAEM, EWC supplies the motivation for bounding drift between cycles; as noted in the Preliminaries, CAEM achieves that bound with a low-rank adapter on a frozen backbone, a strict training-pool filter and a retention guard rather than with an explicit penalty. The low-rank choice is reviewed next.

The choice to fine-tune through a low-rank adapter rather than the full backbone is a separate and equally load-bearing decision in the self-improvement loop. Hu et al.~\cite{hu2022lora} introduced low-rank adaptation as a parameter-efficient alternative to full fine-tuning that injects trainable low-rank matrices into the weight matrices of a frozen backbone; the original paper adapts only the attention projections and reports quality on par with full fine-tuning at less than one percent of the trainable parameter count on GPT-3 175B, and later work extends the adapters to the feed-forward projections. The choice between low-rank adaptation and full fine-tuning has empirical consequences for continual learning that the recent literature has measured directly. Biderman et al.~\cite{biderman2024lora} compare low-rank adaptation with full fine-tuning for Llama-2-7B on code and mathematics, sweeping ranks from $16$ to $256$, and report two findings relevant to CAEM's setting: low-rank adaptation learns less of the target domain than full fine-tuning, but it forgets less of the base model's general capability, measured as average HellaSwag, ARC-Challenge, and WinoGrande accuracy. In code instruction tuning on Magicoder data, for example, full fine-tuning falls to 0.414 against 0.509 for low-rank adaptation at rank sixty-four. Forgetting grows with training duration in both cases, and in mathematics low-rank adaptation at rank 256 forgets nearly as much as full fine-tuning, so higher rank trades retention for fit. The QLoRA work of Dettmers et al.~\cite{dettmers2023qlora} establishes the second decisive practical point: applying low-rank adaptation to all linear transformer-block layers, both attention and feed-forward, rather than to the attention projections alone is the most critical configuration choice for matching full fine-tuning, whereas the rank itself has little effect once all layers are adapted. CAEM's adapter target set spans every linear projection inside the transformer block (the four attention projections together with the three feed-forward projections of the gated multi-layer perceptron), which is the all-linear-layer target set that this literature converges on, fitted at rank thirty-two with $\alpha=2r$; with this choice the effective scale $\alpha/r$ stays at two, so the rank-driven gradient collapse that motivates the rank-stabilised alternative for fixed-$\alpha$ adapters does not arise, and rank thirty-two is far below the large ranks at which the theory of Kalajdzievski~\cite{kalajdzievski2023rslora} predicts instability for a constant scale. Low-rank adaptation is also what lets the self-improvement loop fit on a single consumer-grade graphics processor; full fine-tuning at the same backbone size would not fit. The within-adapter analogue of catastrophic forgetting under sequential fine-tuning, documented by Wang et al.~\cite{wang2023olora} and orthogonal-decomposition fixes for it, is registered as future work; the present thesis combines the bounded low-rank-adapter parameter envelope on a frozen backbone, the multi-modal retention guard with adapter rollback, and the strict training threshold above the storage cut-point as the three-mechanism defence against iterative-adapter drift across the realised cycle trajectory.

Natarajan et al. \cite{natarajan2013learning} provided theoretical foundations for learning with noisy labels at NeurIPS 2013. Their unbiased estimator method constructs proxy loss functions recovering optimal classifiers despite label corruption. The key insight is that if the noise process is known or can be estimated (e.g., what fraction of positive labels are actually negative), modified loss functions achieve consistent parameter estimation. On a synthetic non-separable dataset their methods achieve more than 88\% accuracy even when 40\% of the labels are corrupted, and they are competitive with other methods for learning under label noise. 

For CAEM, this framework explains why training on an imperfect pool can still help: no verifier is perfect, so some stored episodes carry wrong answers. Natarajan et al.'s guarantees, however, assume a loss corrected for the noise rates, which CAEM does not apply (\Cref{thm:bayes-purity}). CAEM instead lowers the noise rate directly, by admitting to the training pool only episodes above a stricter threshold $\tau_{\text{train}} > \tau_{\text{store}}$, so that the training pool is purer than the memory as a whole. The realised purity of the training pool is reported in \Cref{ch:evaluation}.

When Zelikman et al. \cite{NEURIPS2022_639a9a17} introduced STaR (Self Taught Reasoner) at NeurIPS 2022, they showed successful bootstrapping using iterative self improvement. STaR loop generates rationales for problems, filters for correct answers, fine tunes on correct rationale-answer pairs and then repeats. A key innovation is rationalisation. For problems the model fails to answer correctly, it is given the correct answer as a hint and generates a rationale that leads to it. This data augmentation adds training signal from problems the model could not yet solve, which the basic loop discards.

On CommonsenseQA, STaR improved over a few-shot baseline by 35.9 percentage points (72.5\% against 36.6\%), showing that bootstrapping from a model's own outputs works when coupled with appropriate filtering. The parallel to CAEM is direct: generate predictions, verify them, fine-tune on verified examples, and repeat. CAEM extends STaR in two ways. First, simple correctness filtering is replaced by the nine-signal verifier, whose calibrated composite $\ustored$ drives the four-outcome storage decision. Second, verified items persist in episodic memory for future retrieval rather than being consumed only as training data.

Recent work on intrinsic self-correction reveals important limitations. Huang et al.\ \cite{huang2024cannot} found at ICLR 2024 that prompting language models to review and correct their own reasoning does not improve accuracy without external feedback, and at times degrades it. Kamoi et al.\ \cite{kamoi-etal-2024-llms} analysed at TACL 2024 when language models can correct their mistakes and found that intrinsic self-correction often fails. The bottleneck is reliable feedback: a model that produces an incorrect first answer usually cannot generate a reliable correction signal without external validation.

This limitation motivates CAEM's design. Instead of asking the model to review itself, CAEM uses external verification: an entailment judge against retrieved evidence, agreement across resampled chains, and recalibrated internal signals. These supply the external feedback that the self-correction literature identifies as a precondition for useful self-improvement, and because their failure modes differ, a calibrated combination of them is expected to be more reliable than any one of them or than the model's own self-evaluation. \Cref{ch:evaluation} reports how well the verifier separates correct from incorrect answers.

The self-improvement literature also reports successes: when external feedback is available, for example code execution results, formal verification or retrieved sources, models can use it to correct themselves. CAEM's retrieval-augmented tier supplies feedback of this kind, since the retrieved evidence both grounds the generation and gives the verifier something to check the answer against.

Domain adaptation research provides additional relevant findings. Gururangan et al. \cite{gururangan-etal-2020-dont} demonstrated at ACL 2020 that continued pre-training on domain-specific unlabelled data (domain-adaptive pre-training) followed by task-specific fine-tuning outperforms direct fine-tuning on limited labelled data. While CAEM does not perform domain-adaptive pre-training due to computational constraints, the finding supports iterative refinement: progressive training on accumulated verified examples enables gradual domain adaptation.

The distinction between CAEM's approach and standard continual learning deserves emphasis. Traditional continual learning addresses learning sequential tasks A, B, C without forgetting previous tasks. CAEM's scenario differs: the task remains constant (question answering), but the model progressively improves at that task through accumulated verified examples. This setting is closer to curriculum learning and iterative refinement than task-sequential continual learning, though catastrophic forgetting concerns still apply since fine-tuning on domain-specific verified examples risks degrading general capabilities.

Several 2024--2025 contributions tighten the theoretical and architectural framing of self-improvement in ways directly relevant to CAEM's theoretical analysis and to its retention guard. Hosseini et al.~\cite{hosseini2024vstar} introduced V-STaR at COLM 2024, training a DPO verifier over correct and incorrect self-generated solutions and using it to rank test-time candidates; CAEM generalises V-STaR by using the verifier for training-data admission rather than only for test-time ranking, and by adding the episodic memory and the retention guard. Song et al.~\cite{song2024gap} formalise the \emph{generation--verification gap}, the improvement obtained by reweighting a model's generations with its own verification, and show empirically that self-improvement requires a positive gap, that small models fail to achieve one, and that iterative self-improvement saturates within two or three rounds; this is the conceptual motivation for CAEM's $\alpha > 1/2$ verifier-purity precondition in \Cref{sec:theory}. Yi et al.~\cite{fu2025collapse} prove theoretically, in linear regression, that a verifier filter over synthetic data prevents model collapse in self-consuming training loops and yields near-term improvement, while in the long run the model converges to the verifier's knowledge centre, so that gains plateau or reverse unless the verifier is unbiased; this is the strongest published safety anchor for CAEM's storage-gate-filtered SIL, and also its main caution. Das et al.~\cite{das2025retraining} prove at ICML 2025 that retraining a model on its own predicted hard labels increases population accuracy in a binary linearly-separable classification setting with randomly flipped labels, and show empirically that retraining only on samples where the prediction agrees with the given label helps further; this is the closest published analogue of CAEM's pool-purity inequality (\Cref{thm:purity}), and CAEM extends the idea from binary classification with self-predicted labels to open-ended generation with an external multi-signal verifier. Huang, Block, Foster et al.~\cite{huang2025sharpening} formalise self-improvement as a sharpening mechanism at ICLR 2025 and prove convergence-style amortisation-via-self-training bounds that are the closest published analogue of CAEM's bounded-band stationarity heuristic on the committed-cycle subsequence (\Cref{thm:convergence}). Lang, Sontag, and Vijayaraghavan~\cite{lang2024w2s} provide a complementary expansion-based bound at NeurIPS 2024 for when a strong student trained on weak-teacher pseudolabels outperforms the teacher; their bound requires expansion assumptions rather than a verifier-accuracy threshold and is therefore complementary to, rather than redundant with, CAEM's data-purity formulation.

These methods inform CAEM's self-improvement loop as follows. EWC motivates bounding drift between cycles, which CAEM achieves with a low-rank adapter on a frozen backbone, an approach the parameter-efficient fine-tuning literature shows forgets less than full fine-tuning. Noisy-label learning shows that learning from imperfect labels remains possible when the noise is bounded, which motivates the stricter training threshold. STaR supplies the generate, filter, fine-tune and repeat cycle. The self-correction literature motivates external verification rather than self-review. Finally, the retention guard, which rolls the adapter back but never the memory, turns forgetting into a detected and rejected event rather than a silent degradation. The trajectory actually achieved is reported in \Cref{ch:evaluation}.

\Cref{tab:positioning} positions CAEM against the closest published systems on the architectural axes that distinguish them: whether the system stores a per-entry confidence at insertion, what signal the router uses, what filter the training-data admission applies, what regulariser preserves general capability under self-improvement, and whether a formal precision or convergence guarantee is supplied. No single prior system shares CAEM on every axis simultaneously; the table makes the differentiation explicit.

\begin{table}[!htbp]
\centering
\footnotesize
\begin{tabular}{@{}p{2.4cm}p{1.7cm}p{2.6cm}p{2.5cm}p{2.0cm}p{2.0cm}@{}}
\toprule
\textbf{System} & \textbf{Stored conf.} & \textbf{Routing signal} & \textbf{Training filter} & \textbf{CL guard} & \textbf{Formal guarantee} \\
\midrule
\textbf{CAEM} (this work) & yes (calibrated $\ustored$) & combined sim $+$ $\ustored$, $\upre$ override & 9-signal cal-prob composite $+$ fixed-threshold gate (per-bench shrinkage; LoRA SIL) & bounded LoRA adapter envelope $+$ multi-modal retention rollback & calibration-consistent precision floor; purification condition; trajectory heuristics \\
A-MEM~\cite{xu2025amem} & no (textual tags) & similarity $+$ link evolution & --- & --- & --- \\
Larimar~\cite{larimar2024das} & no (key-value addr.) & key match & one-shot edit & --- & --- \\
Adaptive-RAG~\cite{jeong2024adaptive} & no & query complexity classifier & --- & --- & --- \\
V-STaR~\cite{hosseini2024vstar} & no & --- & DPO verifier (test-time rank) & --- & --- \\
Mohri-Hashimoto~\cite{mohri2024conformal} & no & --- & split-conformal sub-claim filter & --- & conformal factuality \\
MiniCheck~\cite{tang2024minicheck} & no & --- & NLI judge (fact verification) & --- & --- \\
\bottomrule
\end{tabular}
\caption[Positioning of CAEM against the closest prior systems.]{Positioning of CAEM against the closest prior systems. Insertion-time stored confidence, combined-score routing on a shared memory, the calibrated-probability composite at a fixed-threshold gate, the LoRA self-improvement primitive with an adapter-rollback guard against a multi-modal retention probe panel, and the calibration-consistent precision floor with its purification condition together constitute CAEM's distinguishing combination.}
\label{tab:positioning}
\end{table}

\subsection{External Baselines and Competitive Positioning}
\label{sec:external-baselines}

The related work reviewed above clusters into three families relevant to CAEM's evaluation: prompt-level reasoning methods, retrieval-augmented generation methods, and training-time self-improvement methods, to which a confidence-calibration method is added. \Cref{ch:evaluation} runs representatives of each as external baselines, so that every mechanism CAEM adds to a plain language model is compared with at least one prior system that uses that mechanism alone. This subsection states what each baseline is and what it isolates; the numerical results appear in \Cref{ch:evaluation}.

The prompt-level family contributes three baselines. The first, \textbf{zero-shot Qwen-2.5-3B-Instruct}, uses the same decoder-only backbone CAEM fine-tunes, with no CoT prefix, no retrieval, and no verification. This is the empirical floor: any gain CAEM reports must exceed the zero-shot score on the same backbone, otherwise the gain is attributable to model capacity rather than to the CAEM mechanism. The second, \textbf{Chain-of-Thought} prompting~\cite{wei2022chain}, adds the standard zero-shot ``Let's think step by step'' trigger~\cite{kojima2022large} without any further modification. CoT isolates the contribution of intermediate reasoning on the same backbone and controls for the possibility that CAEM's gains reflect better rationale generation rather than the verifier-memory architecture. The third, \textbf{Five-shot Chain-of-Thought}, prepends five in-context $(\text{question}, \text{reasoning}, \text{answer})$ demonstrations from a benchmark training split before the CoT trigger and is the strongest prompting-only comparison against CAEM's retrieval-free Tier-2.

The retrieval family contributes three baselines. \textbf{DPR-RAG}~\cite{karpukhin-etal-2020-dense, NEURIPS2020_6b493230} retrieves the top $k=5$ Wikipedia passages from a dense DPR index, concatenates them into a 384-token context, and generates from the same Qwen-2.5-3B-Instruct backbone through the Tier-3 prompt builder; this isolates the contribution of external grounding without any verification or memory. \textbf{CoT$+$DPR-RAG} combines the CoT trigger with DPR retrieval and isolates the contribution of reasoning over retrieved evidence. \textbf{FLARE}~\cite{jiang-etal-2023-active} extends DPR-RAG with active retrieval: rather than retrieving once at the start of generation, FLARE generates one sentence at a time and triggers a retrieval whenever the token-level confidence on the upcoming sentence falls below a registered threshold, with retrieved passages prepended to the generation context for the next sentence; this isolates the contribution of selective and step-conditioned retrieval over single-pass retrieval. A fourth retrieval system, \textbf{Self-RAG}~\cite{asai2024selfrag}, is acknowledged as the most ambitious retrieve-and-critique pipeline in the literature but is included only as a citation-level reference rather than as a replicated numerical baseline, because Self-RAG's critic tokens are tied to a LLaMA-family backbone and retargeting them to Qwen-2.5-3B-Instruct would change enough of the system to invalidate the architectural comparison.

The training-time family contributes one baseline. \textbf{Vanilla fine-tuning} fine-tunes all weights of Qwen-2.5-3B-Instruct on raw question-answer pairs from the three training benchmarks for the same number of cycles as CAEM, with no verifier filter, no retention guard and no episodic memory. It isolates the contribution of CAEM's curation and retention machinery: if vanilla fine-tuning degrades on the retention probes or across cycles while CAEM does not, that machinery is responsible. The confidence-calibration family contributes one baseline. \textbf{Semantic entropy}~\cite{farquhar2024semantic} samples several generations at a raised temperature, clusters them into groups of equivalent meaning, and uses the entropy of the cluster distribution as the per-query uncertainty; the baseline runs the same backbone as the zero-shot system and abstains when the entropy is high, which makes it the most direct external comparator to CAEM's calibrated composite. \textbf{STaR}~\cite{NEURIPS2022_639a9a17}, discussed above as the closest prior method in spirit, is not run: its rationalisation step and binary correctness filter differ enough from CAEM's verifier and four-outcome decision that a comparison would mix architectural contribution with rationalisation quality.

Taken together, the eight external baselines (B1 to B8 in \Cref{sec:comp-baselines}) span the prior systems that share at least one mechanism with CAEM: three isolate the prompt-level component (B1 zero-shot, B2 CoT, B5 five-shot CoT), three the retrieval component (B3 DPR-RAG, B4 CoT$+$DPR-RAG, B7 FLARE), one the training-time component (B6 vanilla fine-tuning), and one a sampling-based uncertainty estimate on the same backbone (B8 semantic entropy). Neither Self-RAG nor STaR is run, so neither enters the significance tests. No single baseline combines all four families; CAEM's contribution is their combination, which the comparison in \Cref{ch:evaluation} tests against each family in isolation.

\subsection{Evaluation Benchmarks}

Rigorous evaluation requires benchmarks testing diverse reasoning capabilities with objective correctness criteria. This subsection reviews the five primary factual-QA benchmarks that constitute CAEM's locked evaluation panel (FEVER, TriviaQA, CommonsenseQA training, and TruthfulQA, StrategyQA transfer) and the MMLU general-knowledge probe that sits inside the multi-modal retention guard, explaining how each tests a different aspect of hallucination reduction. Natural Questions is also described for completeness: it and HotpotQA were evaluated at cycle zero and dropped before the main run because the base generator answered too few of their questions correctly for the storage gate to build a usable training pool (\Cref{sec:scope-research}).

Thorne et al.~\cite{thorne-etal-2018-fever} created the FEVER dataset at NAACL 2018 for fact verification, providing 185{,}445 claims generated from Wikipedia requiring classification as SUPPORTS, REFUTES, or NOT ENOUGH INFO relative to Wikipedia evidence. Claims are designed to test factual knowledge; an illustrative pair (not taken from the dataset) is ``Barack Obama was the 44th President of the United States'' (SUPPORTS) and ``The Beatles formed in Manchester'' (REFUTES, since they formed in Liverpool). The NOT ENOUGH INFO category is particularly valuable, testing whether systems recognise insufficient evidence rather than hallucinating conclusions. FEVER directly evaluates the external-grounding family of CAEM's verifier, because the task structure matches the premise-hypothesis entailment check that produces $p_{\text{ground}}$ and $p_{\text{entail}}$.

Joshi et al.~\cite{joshi-etal-2017-triviaqa} introduced TriviaQA at ACL 2017, providing about 95{,}000 trivia question-answer pairs authored by trivia enthusiasts, which with on average six supporting documents each from the web and Wikipedia give over 650{,}000 question-answer-evidence triples. Questions span an unusually wide topical range (history, science, geography, entertainment) and the provided evidence is noisy (only a subset of the retrieved documents actually answers the question), making TriviaQA a realistic test of retrieval-augmented generation under distractor documents. For CAEM, TriviaQA tests whether the Tier-3 retrieval-augmented path can ground answers against noisy evidence and whether the mean-match grounding signal $p_{\text{ground\_mean}}$ remains calibrated when retrieval quality varies across queries.

Kwiatkowski et al.~\cite{kwiatkowski-etal-2019-natural} introduced Natural Questions (NQ) at TACL 2019, providing real anonymised queries to the Google search engine paired with Wikipedia-page long answers and span-level short answers. Unlike synthetic benchmarks, NQ captures the distribution of questions users actually ask, with the accompanying variability in specificity, ambiguity, and presupposition. NQ is the dominant open-domain QA benchmark in recent retrieval-augmented-generation literature~\cite{NEURIPS2020_6b493230, karpukhin-etal-2020-dense}. CAEM evaluated NQ at cycle zero and dropped it before the main run: the base generator answered only $17$ percent of its questions correctly, too few for the storage gate to build a usable training pool. Its training split is reused only to train the retrieval-utility classifier.

Lin et al.~\cite{lin-etal-2022-truthfulqa} designed TruthfulQA specifically to measure how models mimic human falsehoods, providing 817 questions across 38 categories in health, law, finance, politics, and conspiracies. Questions are adversarially constructed to elicit common misconceptions: ``What happens if you crack your knuckles?'' where the false but common answer is ``you get arthritis'' and the true answer is ``nothing harmful happens; no proven link to arthritis.'' The inverse scaling finding (larger models often less truthful) makes TruthfulQA the most challenging benchmark for parametric scaling approaches. Human performance reaches 94\% truthful and 87\% truthful-and-informative, while the best GPT-3 175B configuration achieves 58\% and 21\% respectively. For CAEM, TruthfulQA was registered as the benchmark where the sample-set family of signals was expected to pay off, on the reasoning that a misconception would make the model alternate between the correct answer and the misconception across stochastic decodes, which the sample-set agreement signals $s_{\text{avg}}$ and $p_{\text{entail}}$ would detect. The literature predicts otherwise: semantic entropy targets arbitrary, sample-to-sample errors and explicitly does not cover models that are consistently wrong because of common misconceptions, as Farquhar et al.~\cite{farquhar2024semantic} state, since such a model gives the same wrong answer on every decode and agreement-based signals then read it as confident. The realised results agree with the literature rather than with the registered expectation: entries stored on TruthfulQA at cycle zero are correct only 52.2\% of the time, against 75 to 82\% on the training panel (\Cref{tab:cross-bench-summary}), and the semantic-entropy baseline is among the highest-scoring baselines on TruthfulQA exact match rather than being penalised for heavy abstention (\Cref{ch:evaluation}). TruthfulQA is therefore the benchmark that exposes the blind spot of agreement-based verification.

Geva et al.~\cite{10.1162/tacl_a_00370} introduced StrategyQA at TACL 2021, presenting 2{,}780 questions requiring implicit reasoning where decomposition steps are not mentioned in questions. For example, the dataset question ``Could a llama birth twice during the War in Vietnam?'' requires knowing Vietnam War duration (implicit: 1955--1975) and llama gestation period (implicit: approximately 11 months), then computing whether two births fit the timeframe. Human accuracy reaches 87\% while best models achieve approximately 66\%. The implicit reasoning requirement tests whether models genuinely understand underlying concepts or merely perform shallow pattern matching. For CAEM, StrategyQA evaluates the sample-set agreement signal $s_{\text{avg}}$ because implicit reasoning should converge across multiple chains if the underlying world-knowledge is consistent; divergence across chains is a strong hallucination signal specifically on this benchmark.

Hendrycks et al.~\cite{hendrycks2021measuring} introduced MMLU in 2021, a 57-subject multiple-choice benchmark covering humanities, STEM, social sciences, and professional domains. MMLU is used in CAEM as one of three components of the \textit{multi-modal retention probe panel}, not as a target benchmark. The fine-tune is not optimised for MMLU; instead, MMLU is probed before the first cycle (the pristine baseline) and after every cycle's fine-tune, alongside two further probes drawn from held-out test slices of the training panel: an open-text factoid probe on the held-out TriviaQA test fold, and a commonsense multiple-choice probe on the held-out CommonsenseQA test fold. The cycle aborts and rolls the adapter back to the previous snapshot whenever the worst per-probe ratio against the pristine baseline falls below the fixed floor $\rho_{\min} = 0.93$; the worst-probe rule replaces an earlier single-probe rule because a general-knowledge multiple-choice probe alone missed the failure mode where MCQ accuracy stayed high while open-text generation degraded silently. The episodic memory is preserved across rollbacks. The choice of MMLU as one of the three probes is deliberate: catastrophic forgetting manifests most visibly as erosion of capability on tasks the model was not trained on this cycle, so a broad out-of-distribution probe is a more sensitive trigger than in-distribution QA, and the held-out training-bench test folds add the asymmetric-drift sensitivity that the broad probe alone would miss.

The benchmarks are complementary. FEVER tests entailment-based grounding, including the NOT ENOUGH INFO case where the right behaviour is to recognise missing evidence. TriviaQA tests retrieval-augmented generation under noisy evidence. CommonsenseQA tests commonsense reasoning over short multiple-choice options. TruthfulQA tests agreement-based and relevance signals on imitative falsehoods, as a transfer-only benchmark held out from the SIL training pool. StrategyQA tests sample-set agreement on implicit multi-hop reasoning, also as a transfer-only benchmark. MMLU, the open-text TriviaQA test fold and the multiple-choice CommonsenseQA test fold close the catastrophic-forgetting loop through the retention probe panel. Together they cover a range of factual reasoning tasks, so that results on the panel show whether CAEM's effects carry across reasoning types rather than fitting one task.

\section{Summary of Key Findings}

The literature reviewed in this chapter supports five findings that shape CAEM.

First, hallucination is persistent and is not solved by scale. Rates vary widely by domain, up to 91.4\% of chatbot-generated references in one systematic-review study \cite{info:doi/10.2196/53164}; on TruthfulQA the best GPT-3 configuration was truthful and informative on only 21\% of questions, with larger models often less truthful \cite{lin-etal-2022-truthfulqa}; and FActScore finds only 58\% of ChatGPT's atomic facts supported \cite{min-etal-2023-factscore}. This motivates changing the structure around the model rather than relying on scale.

Second, no single verification or confidence signal is reliable across tasks. Verbalised confidence discriminates weakly \cite{xiong2023can}, modern networks are systematically overconfident \cite{pmlr-v70-guo17a}, preference-tuned models are miscalibrated until rescaled \cite{kadavath2022language, tian2023just}, and uncertainty methods that do well on open generation do poorly on constrained outputs and vice versa \cite{vashurin2024polygraph}. Combining calibrated signals has published precedent in the logistic multi-scoring of \cite{valentin2024hallucination}, and generic NLI models transfer poorly to language-model hallucinations \cite{sinha2025semantic}, which favours a judge trained on claim-support data \cite{tang2024minicheck}. CAEM therefore combines nine recalibrated signals into one composite.

Third, external memory and retrieval improve factual accuracy and can be updated without retraining \cite{graves2016differentiable, NEURIPS2020_6b493230}, but standard retrieval-augmented generation neither verifies what it retrieves nor accumulates verified answers, and retrieval does not help every query \cite{mallen2023whennot, wang2023skr}. CAEM therefore stores only verified answers, routes on similarity combined with stored confidence, and decides per query whether retrieval is worth using.

Fourth, self-improvement works when it receives reliable external feedback and fails when a model only reviews itself \cite{NEURIPS2022_639a9a17, huang2024cannot, kamoi-etal-2024-llms}; verifier-filtered retraining can help but saturates within a few rounds \cite{song2024gap, fu2025collapse}. CAEM therefore filters its training data through an external verifier, with a stricter threshold than its memory uses.

Fifth, fine-tuning causes forgetting, which low-rank adaptation reduces but does not remove \cite{hu2022lora, biderman2024lora}, and weight-regularisation penalties add little once the backbone is pre-trained \cite{mehta2023empirical, wu2022pretrained}. CAEM therefore relies on a low-rank adapter on a frozen backbone and on a retention guard that rolls back the adapter, but never the memory, when held-out accuracy falls.

The five-benchmark panel and the three-probe retention panel exercise these mechanisms on different failure modes, and the baselines of \Cref{sec:external-baselines} each isolate one family of prior methods. No prior system combines all of these elements (\Cref{tab:positioning}), which is the gap CAEM addresses.

\section{Chapter signpost}
\label{sec:ch2-signpost}

The preliminaries section and the literature review have established the technical foundations on which the architecture rests and have located CAEM's joint intervention against the verification, retrieval, memory-augmented, and continual-learning prior work whose individual contributions do not close the deployment-safety gap registered in \Cref{ch:introduction}. \Cref{ch:requirements} translates the architectural opportunity surfaced by the literature into a registered specification: seven functional requirements paired with named architectural components, five non-functional requirements with empirical anchors, six production-parity design commitments, and the societal, environmental, ethical, standards, project-management, risk, and economic readings that bound the engineering envelope. The specification is then discharged end-to-end by the methodology chapter that follows.   % Literature Review (preserved, frozen)
% ============================================================================
% Chapter 3 — Requirements, Impact, and Project Management
% Rubric-mandated section structure preserved at the top level; depth lives
% in subsections.
%
% Course outcomes covered: CO2 (specifications and standards), CO3
% (societal impact), CO4 (sustainability and environment), CO10
% (resource management), CO11 (economic analysis), CO12 (ethics).
% Section titles 3.1 through 3.8 are fixed by the rubric.
% ============================================================================

\chapter{Requirements, Impact, and Project Management}
\label{ch:requirements}

% Chapter scope: states the functional and non-functional requirements
% the architecture must meet, surveys the societal, environmental, and
% ethical considerations, declares the engineering standards followed,
% and lays out the project management plan, risk register, and
% economic analysis. This chapter operationalises the Course Outcomes
% CO2, CO3, CO4, CO10, CO11, and CO12.

\section{Final Specifications and Requirements}
\label{sec:specifications}

\subsection{Functional requirements}
\label{sec:func-reqs}

The architecture must perform seven functions at deployment, each derived from the problem statement in \Cref{ch:introduction}; \Cref{tab:func-reqs} lists them with the component that delivers each. Three of them need a word of explanation. Abstention (FR2) means that, when the four-outcome decision is to abstain, the system returns an explicit insufficient-evidence message instead of a hedged or speculative answer. Recall (FR4) answers from memory only when a stored episode is close enough to the new question; otherwise the router sends the query to zero-shot generation or to retrieval-augmented generation, and the retrieval-utility classifier decides between those two for the queries that the router cannot settle on confidence alone. Rendering (FR7) shows the user the calibrated probability of correctness as a percentage, together with one of four reliability tags, on every answer; under abstention the answer text is hidden but the probability is still shown.

\begin{table}[!htbp]
\centering
\small
\begin{tabular}{@{}p{0.8cm}p{4.4cm}p{8cm}@{}}
\toprule
\textbf{ID} & \textbf{Requirement} & \textbf{Architectural component (specified in \Cref{ch:methodology})} \\
\midrule
FR1 & Calibrated answer per query & Nine-signal verifier feeding the per-benchmark calibrated probability composite (with shrinkage prior) \\
FR2 & Explicit abstention under uncertainty & Fixed-threshold storage gate $+$ confabulation early-exit rule \\
FR3 & Commit verified episodes to memory as one record (question, answer, supporting context, storage score) & Four-outcome decision tree (store branch) \\
FR4 & Recall on subsequent queries & Three-tier router with combined-score dispatch and safety override; the retrieval-utility classifier re-decides, at threshold $\tau_{\text{RUC}} = 0.375$, the queries that reach the safety override or fall through the score formula \\
FR5 & Cycle-boundary fine-tune on verified episodes only & Self-improvement loop with low-rank adapter, pool-reweighting layer, and strict training threshold \\
FR6 & Retention guard with adapter rollback that keeps the memory & Multi-modal retention probe panel (worst-probe rule) and asymmetric memory-versus-adapter rollback \\
FR7 & Calibrated probability rendered to the user with every answer & Production response envelope: four-class decision tag $+$ continuous $u_{\text{stored}}$ percentage; insufficient-evidence message replaces the answer text under abstain \\
\bottomrule
\end{tabular}
\caption[Functional requirements and the architectural components that deliver them.]{Functional requirements and the architectural components that deliver them.}
\label{tab:func-reqs}
\end{table}

\subsection{Non-functional requirements}
\label{sec:nonfunc-reqs}

Five non-functional requirements bound the operating envelope of the architecture (\Cref{tab:nonfunc-reqs}).

\textbf{Latency:} The cost of a query depends on its tier: the memory-direct tier costs an embedding lookup and a similarity search, the zero-shot tier adds one decode, and the retrieval-augmented tier adds passage retrieval and a grounded decode. Per-tier latencies are reported in \Cref{ch:evaluation}.

\textbf{Storage precision:} Among the answers that the gate stores, the fraction that are correct should be high. The requirement is stated on the calibration fold at the locked cut-point $\tau_{\text{store}} = 0.60$, where the storage precision $\pi_{\text{store}}$ is
\begin{equation}
\Pr\!\big(\mathrm{EM} = 1 \,\big|\, u_{\text{stored}} \ge \tau_{\text{store}}\big)_{\text{cal}} \;=\; \pi_{\text{store}}.
\label{eq:nfr-precision-floor}
\end{equation}
The cut-point was chosen on the cycle-zero calibration fold before the main run, and the realised precision at every cycle is reported in \Cref{ch:evaluation}.

\textbf{Retention:} After every fine-tune, accuracy on each probe of the multi-modal probe panel must stay at or above $93$ percent of its value before any fine-tuning; otherwise the cycle is aborted and the adapter is rolled back:
\begin{equation}
\min_{p \in \mathcal{P}}\; \rho_{p}(M_{t}) \;:=\; \min_{p \in \mathcal{P}}\; \frac{\mathrm{EM}(M_{t},\,p)}{\mathrm{EM}(M_{0},\,p)} \;\ge\; \rho_{\min} = 0.93,
\label{eq:nfr-retention-floor}
\end{equation}
where $M_{t}$ is the generator after the fine-tune at cycle $t$ and $\mathcal{P}$ is the probe panel: a general-knowledge multiple-choice probe (MMLU), an open-text factoid probe drawn from the TriviaQA test fold, and a commonsense multiple-choice probe drawn from the CommonsenseQA test fold.

\textbf{Reproducibility:} Every released artefact, including the locked cut-points, the calibrated composite, the seed-pool composition and the per-cycle snapshots, must be regenerable from one recorded random seed and a published code commit.

\textbf{Integrity:} No sample in the evaluation fold may appear in the calibration fold, the seed pool or the per-cycle stream chunks, so that calibration-fold and evaluation-fold results are measured on disjoint samples.

\begin{table}[!htbp]
\centering
\small
\begin{tabular}{@{}p{0.8cm}p{3.5cm}p{4.5cm}p{4.5cm}@{}}
\toprule
\textbf{ID} & \textbf{Requirement} & \textbf{Target} & \textbf{Where it is checked} \\
\midrule
NFR1 & Per-query latency & Tier-dependent: embedding lookup (Tier~1), single decode (Tier~2), retrieval $+$ grounded decode (Tier~3) & Per-tier latency in \Cref{ch:evaluation} \\
NFR2 & Storage precision at $\tau_{\text{store}} = 0.60$ & Calibration-fold precision of the stored class, with the cut-point fixed before the main run & Per-cycle precision in \Cref{ch:evaluation} \\
NFR3 & Retention floor (multi-modal probe panel) & $\ge 0.93$ of the pre-fine-tune accuracy on every probe; the worst probe decides & Per-cycle probe ratios in \Cref{ch:evaluation}; one breach (cycle four) rolled back \\
NFR4 & Reproducibility & Every artefact regenerable from one seed and a published commit & Locked configuration files and archived snapshots (\Cref{app:code-data-availability}) \\
NFR5 & Integrity & Evaluation, calibration, seed and stream-chunk folds disjoint & Disjointness verified per benchmark \\
\bottomrule
\end{tabular}
\caption[Non-functional requirements, their targets, and where each is checked.]{Non-functional requirements, their targets, and where each is checked.}
\label{tab:nonfunc-reqs}
\end{table}

\subsection{Constraints and design commitments}
\label{sec:constraints}

The requirements imply six design commitments, all fixed before the main run.
\begin{itemize}
    \item \textbf{Single backbone:} The generator backbone is the same in every cycle and every tier; only the low-rank adapter is trained, so reported gains come from memory, calibration and adapter fine-tuning rather than from a change of backbone.
    \item \textbf{Frozen cut-points:} The storage and deferral cut-points, $\tau_{\text{store}} = 0.60$ and $\tau_{\text{defer}} = 0.45$ on the calibrated probability scale, are fixed once before the main run and never changed, so the operating point of the storage decision does not drift. Only the calibrated composite is refitted at each cycle boundary, on the held-out calibration fold re-scored with the current adapter, with each benchmark's curve shrunk toward the pooled fit so that benchmarks with little data do not overfit.
    \item \textbf{Local execution:} No external model service sits in the inference path: the generator, the verifier ensemble, the embedding encoder and the passage index all run from open weights on one rented graphics processor. The language-model judge that scores TruthfulQA belongs to the evaluation, not to the system.
    \item \textbf{Bounded cycles:} The self-improvement loop is capped at ten cycles and may stop earlier on an automatic early-stop check (\Cref{ch:methodology}).
    \item \textbf{Matched protocol:} Every external baseline is evaluated on the same samples, prompts and scoring rules as the architecture, so that the comparison in \Cref{ch:evaluation} is not confounded by re-implementation choices.
    \item \textbf{Single seed:} All randomness derives from one recorded seed, and the per-cycle artefacts are archived so that a third party can regenerate the trajectory.
\end{itemize}

\section{Societal Impact}
\label{sec:societal-impact}

\subsection{Deployment benefits}
\label{sec:soc-benefits}

The main societal benefit of the architecture is a reduction in confident misinformation from deployed language models on factual question answering. Every answer carries a calibrated probability and a reliability tag, and the answer is withheld when the evidence is insufficient. This explicit refusal channel is itself a benefit, because it turns otherwise hidden uncertainty into a signal that users such as educators, journalists and information specialists can read and act on. Fluent but ungrounded answers have caused documented harm in health information: Alkaissi and McFarlane~\cite{alkaissi2023artificial} report fabricated clinical references, and Chelli et al.~\cite{info:doi/10.2196/53164} report high rates of hallucinated references in systematic-review searches. In such settings, and in others such as legal triage and educational tutoring, a system that abstains is preferable to one that fabricates. The episodic memory can also share verification cost across users: once an answer has been verified, it can be served directly when the same or a paraphrased question is asked again, so the verification done for an early user benefits later users who ask the same thing.

\subsection{Dual-use considerations}
\label{sec:soc-dual-use}

A verified-memory system that learns from its own interactions could memorise harmful content if the verifier is misconfigured or if adversarial inputs reach its training pool. Three mechanisms limit this risk.

\textbf{Retroactive re-verification:} At each cycle boundary every stored episode is re-scored with the updated generator and the refitted composite; its score can go down as well as up, and entries whose score falls below $0.50$ are removed, so a bad entry does not persist unchallenged.

\textbf{Strict training threshold:} The training pool admits only episodes whose storage score is at least $0.70$, above the storage cut-point of $0.60$, so marginal episodes stay out of the fine-tune even when they are admissible for retrieval.

\textbf{Retention guard:} If a fine-tune degrades any probe in the multi-modal panel (general-knowledge multiple choice plus held-out test slices of two training benchmarks) below the floor, the cycle is aborted and the adapter is rolled back, so one damaging cycle cannot silently carry forward.

These mechanisms reduce the risk but do not remove it; how far they reach depends on the verifier's coverage, which is reported empirically and discussed as a threat to validity in \Cref{ch:evaluation}.

\subsection{Marginalised-user impact}
\label{sec:soc-marginalised}

The harm caused by confident hallucination is not evenly distributed. Users who cannot independently check a confident-sounding answer, such as students using language-model tutors, patients using language-model health agents, and applicants using language-model administrative agents, are the most exposed, because they have no domain expert to fall back on. The calibrated abstention class is a partial response: by returning an explicit refusal rather than a hedged guess when confidence is low and little supporting evidence is found, the system reduces the most dangerous failure mode for exactly these users. The abstention is governed by a fixed cut-point on a calibrated probability rather than by a soft heuristic, and the abstention rate is reported in the four-outcome decision distribution in \Cref{ch:evaluation}. The architecture does not solve unequal access to verification, but it narrows one channel through which fluency is mistaken for reliability.

\section{Environmental Impact}
\label{sec:environment}

\subsection{Compute footprint of the experiment}
\label{sec:env-experiment}

The compute of the empirical programme falls into three phases: the cycle-zero calibration phase, the main run, and the external baseline panel. The retrieval-utility-classifier ablation needed no run of its own: it compares the main run, at the cycle-three close, with an earlier trajectory that ran without the classifier (the main run itself continued to cycle five, and all its other results cover cycles zero to five). Everything ran on a single rented graphics processor of the consumer $32$-gigabyte class (an RTX 5090). The main run was the largest: according to its run log, cycles zero to five took $6.0$, $29.0$, $35.4$, $43.2$, $24.3$ and $54.6$ hours, about $193$ hours (eight days) in total. Part of the growth comes from the retroactive re-verification pass, which grows with the memory (about $38$ minutes for $431$ entries at cycle one, about $16$ hours for $10{,}214$ entries at cycle five). The main run was not the first. Three earlier trajectories were rejected or superseded: the first benchmark panel, a run stopped in May 2026, and the trajectory without the classifier. The last took about $134$ hours over cycles zero to five according to its run log, and its first three cycles supply the classifier-off arm of the ablation; the run logs of the first two were not archived, so their compute cannot be reported exactly. Counting these runs, the programme used several times the compute of the main run alone. The carbon footprint of a phase can be estimated with the method of \textcite{strubell2019energy},
\begin{equation}
\mathrm{kgCO_{2}e}_{\,\phi} \;=\; T_{\phi} \cdot P_{\mathrm{TDP}} \cdot \varepsilon_{\mathrm{grid}},
\label{eq:carbon-per-phase}
\end{equation}
where $T_{\phi}$ is the wall-clock duration of phase $\phi$ in hours, $P_{\mathrm{TDP}}$ is the rated power of the graphics processor in kilowatts, and $\varepsilon_{\mathrm{grid}}$ is the local grid emissions factor in kilograms of carbon-dioxide equivalent per kilowatt-hour. At the card's rated board power of $0.575$ kilowatts, the main run used at most about $111$ kilowatt-hours of graphics-processor energy. The grid factor of the rented host is not known, so no carbon figure is reported.

\subsection{Cycle budget rationale}
\label{sec:env-cycles}

The ten-cycle cap bounds the compute that the loop may spend. The convergence heuristic of \Cref{ch:methodology} (\Cref{thm:convergence}; a heuristic, not a proven result) suggests that gains flatten within a few cycles, and Song et al.~\cite{song2024gap} find that iterative self-improvement saturates within a few rounds; a cap of ten leaves room above that while bounding cost. Each cycle adds its own fine-tuning and re-verification energy, so stopping once gains have flattened is an environmental choice as well as an experimental one. The realised run stopped after cycle five, as planned once cycle four had been aborted by the retention guard and its weights rolled back. The automatic early-stop check fired at the same point, but two of its three inputs were wired to the wrong quantities: its storage-saturation signal read the storage rate of the evaluation fold, where nothing is ever stored, instead of the stream admission rate (about half of each chunk), and its retention signal compared raw MMLU accuracy (about $0.61$) with the $0.93$ ratio floor instead of the retention ratio (between $0.94$ and $1.00$). Both signals were therefore certain to fire as soon as the check first ran, at cycle five, so its firing is not evidence that the trajectory had levelled off. The wiring was corrected after the run; replaying the check on the realised cycle values with the corrected inputs, none of its three signals fires at cycle five. The trajectory is reported in \Cref{ch:evaluation}.

\subsection{Retention versus retraining}
\label{sec:env-retention}

The retention guard with rollback also saves compute. Without it, a loop that regressed would either keep training from damaged weights, making every later cycle pay to recover the lost capability, or need a restart from a clean checkpoint, discarding all the fine-tuning done since. With it, a failed cycle costs only its own fine-tuning: the adapter returns to the previous checkpoint, while the memory, including the episodes verified during the failed cycle, is kept. The argument is qualitative; the realised run had one such rollback, at cycle four.

\section{Ethical Issues}
\label{sec:ethics}

\subsection{Alignment and refusal posture}
\label{sec:ethics-alignment}

The architecture treats explicit abstention as a first-class action rather than a fallback, and this choice has an ethical justification beyond its effect on accuracy. A model that hedges or speculates under uncertainty exposes the user to harm without telling them that uncertainty is present, whereas a model that abstains hands the uncertainty to the user as a visible signal, on which they can act by seeking a second opinion or a human expert. This connects to the question, studied by Kadavath et al.~\cite{kadavath2022language}, of whether language models know what they know. The storage and deferral cut-points are fixed values on the calibrated probability scale, chosen on the cycle-zero calibration fold before the main run rather than tuned by hand during it. The ethical commitment behind this posture is that a model which cannot be confident should not present its answer as if it were.

\subsection{Bias propagation through episodic memory}
\label{sec:ethics-bias}

A system that commits its own outputs to memory and then trains on that memory can amplify biases in the verifier or in the underlying data. Three mechanisms limit this. The retroactive re-verification pass re-scores every stored episode at every cycle boundary and removes entries whose score falls below $0.50$, so a bias that enters the stored pool does not persist unchallenged. The per-cycle refit of the calibrated composite, with each benchmark's curve shrunk toward the pooled fit, follows score-distribution drift by re-scoring the held-out calibration fold with each cycle's adapter, while the storage cut-point stays fixed so that the gate's operating point does not move. The strict training threshold holds the training pool to a higher bar than the storage tier (a storage score of at least $0.70$ against $0.60$ for storage), so marginal episodes cannot enter the next fine-tune even when they pass the retrieval gate. None of these mechanisms removes bias inherited from the pre-training corpus, which remains a limitation of the approach; what they do is limit the rate at which new bias can enter through the self-improvement loop itself.

\subsection{Data provenance and licensing}
\label{sec:ethics-provenance}

The empirical programme uses publicly released benchmarks and a public retrieval corpus throughout. The five evaluation benchmarks are FEVER \cite{thorne-etal-2018-fever}, TriviaQA \cite{joshi-etal-2017-triviaqa}, CommonsenseQA \cite{talmor-etal-2019-commonsenseqa}, TruthfulQA \cite{lin-etal-2022-truthfulqa}, and StrategyQA \cite{10.1162/tacl_a_00370}, all distributed under licences that permit research use with attribution. Two further benchmarks, HotpotQA and Natural Questions, were evaluated at cycle zero and dropped before the main run (\Cref{sec:scope-research}); their cycle-zero results are kept in the project archive. The retrieval corpus is a Wikipedia-derived passage collection distributed under a Creative Commons licence. No private user data, scraped social-media content or proprietary corpus enters the calibration fold, the seed pool or the cycle pool. The cold-start memory that seeds the self-improvement loop is built from the benchmarks' training splits, with disjointness from the evaluation fold enforced at the sample level. The code is public, and the configuration files, the per-cycle artefacts and the pre-self-improvement snapshot are archived in a shared folder, so that the data path can be audited end to end (\Cref{app:code-data-availability}).

\section{Standards and Codes of Practice}
\label{sec:standards}

\subsection{Software-engineering standards}
\label{sec:std-software}

The implementation is written in Python and broadly follows the PEP~8 conventions for layout and naming. The project does not ship a pinned requirements file; the versions of the main libraries in the realised environment are recorded in \Cref{app:repro-versions}. All randomness is seeded from one value at process start (the Python, NumPy and PyTorch generators), and deterministic cuDNN kernels are selected where available; the attention kernels on the graphics processor are not bitwise deterministic, so repeated runs agree up to floating-point noise rather than bit for bit. The source is kept under Git version control in a public repository, with one branch per project phase, and each cycle's artefacts are archived at the cycle boundary, so that the empirical record can be traced to the code that produced it.

\subsection{Reproducibility standards}
\label{sec:std-reproducibility}

Reproducibility is treated as a deliverable rather than a courtesy. Four kinds of artefact support it.
\begin{itemize}
    \item \textbf{Recipe:} \Cref{app:reproducibility} gives the run sequence from data preparation to the final tables, together with a checklist for comparing reproduced numbers against the thesis.
    \item \textbf{Pre-self-improvement snapshot:} The cold-start memory, the cycle-zero calibrated composite and the configuration at the start of the main run are archived, so that a run can resume from that point instead of from scratch.
    \item \textbf{Per-cycle artefacts:} Each cycle's run summary, memory snapshot, retention-probe readings, precision audit and decision log are kept together.
    \item \textbf{Locked configuration:} One configuration file records the storage and deferral cut-points, the shrinkage prior, the boost-layer regulariser and the low-rank adapter settings used at every downstream scoring and training step.
\end{itemize}
Together these let a reader check a specific claim in \Cref{ch:evaluation} by re-running the stage that produced it rather than by trusting the reported number.

\subsection{Evaluation-protocol standards}
\label{sec:std-evaluation}

The empirical comparisons follow four protocol standards.

\textbf{Matched protocol:} Every external baseline is evaluated on the same samples, the same prompt template family and the same scoring rules as the architecture.

\textbf{Pre-specification:} The calibration sweep grid, the cut-point targets and the expected effect bands of the ablation were recorded in the runbook before the corresponding results were read, and the locked configuration was applied unchanged in every cycle of the main run; outcomes outside the expected bands are reported rather than tuned away. Throughout the thesis, a quantity described as pre-specified or registered was fixed in the project's configuration and runbook before the main evaluation run; the term does not refer to an entry in an external pre-registration registry.

\textbf{Selective prediction:} Abstentions are reported per benchmark alongside accuracy, following the selective-prediction framing of Kamath et al.~\cite{kamath2020selective}, in which a system that may abstain is judged both on how often it answers and on how accurate its answers are.

\textbf{Multiple comparisons:} Within each benchmark, the baseline contrasts are corrected with the Holm--Bonferroni procedure \cite{holm1979simple}, so that a significance verdict cannot be inflated by the number of comparisons. Each contrast uses the paired McNemar test with the continuity correction of Edwards~\cite{edwards1948note}, and a bootstrap $95$ percent confidence interval on the exact-match difference is reported alongside the corrected $p$-value, so that effect size is visible as well as significance.

\section{Project Management Plan}
\label{sec:project-management}

\subsection{Phase breakdown}
\label{sec:pm-phases}

The project has five phases, each producing what the next consumes.
\begin{enumerate}
    \item \textbf{Pre-thesis:} Literature review, problem statement, methodology design, and the per-stage pipeline that the empirical phases depend on.
    \item \textbf{Calibration:} Fit the calibrated composite on the disjoint cycle-zero calibration fold, fix the storage and deferral cut-points, and lock the configuration for the main run (only the composite is refitted in later cycles).
    \item \textbf{Main run:} Iterate the per-query pipeline and the cycle-boundary passes for up to ten cycles, producing the per-cycle artefacts.
    \item \textbf{Comparison:} Run the external baseline panel and the retrieval-utility-classifier on/off ablation under matched protocol, followed by the statistical analysis.
    \item \textbf{Write-up:} The thesis report, the defence presentation, and a journal manuscript planned for after the defence.
\end{enumerate}

\subsection{Milestone schedule}
\label{sec:pm-milestones}

The pre-thesis milestone, a locked methodology design, was met on schedule. The calibration milestone, the cycle-zero validation gate, took two iterations: the first failed because of a bug in the step that writes out the verifier signals, and the second, with the bug fixed and a wider parameter grid, produced the locked configuration that the main run uses. The slip is reported here rather than smoothed over, and \Cref{ch:evaluation} gives the diagnostics of both iterations. The main run completed six cycles (cycles zero to five) on 27 May 2026, and the comparison phase completed the eight-baseline panel and the retrieval-utility-classifier ablation. The write-up milestone is the submission of this report and its companion artefacts.

\subsection{Deliverables}
\label{sec:pm-deliverables}

The project releases five deliverables.
\begin{itemize}
    \item \textbf{Source code:} public at \url{https://github.com/aksaN000/caem-thesis}: the implementation of the architecture, the cycle-zero calibration scripts, the cycle-boundary passes, and the baseline and ablation drivers.
    \item \textbf{Artefact bundle:} stored in the shared Google Drive folder \url{https://drive.google.com/drive/folders/1xErhsRDYe_qO0ZgUBn0WXK81eHc5jZCa}: the cycle-zero calibrated composite, the locked cut-points, the cold-start memory snapshot, and the per-cycle artefacts of the realised trajectory. The $64$-gigabyte Wikipedia passage index is mirrored on a private Hugging Face repository, available on request (\Cref{app:code-data-availability}).
    \item \textbf{Thesis report:} this document.
    \item \textbf{Defence presentation:} a talk summarising the headline result, the method and the threats to validity.
    \item \textbf{Journal manuscript:} planned for after the defence and prepared from the evaluation chapter.
\end{itemize}

\section{Risk Management}
\label{sec:risk}

\subsection{Technical risks}
\label{sec:risk-technical}

Six technical risks were identified before the main run. Each is paired with its mitigation below, and \Cref{tab:risk-register} records whether it materialised.
\begin{itemize}
    \item \textbf{Calibration drift:} The composite, refitted every cycle, could lose its precision at the frozen cut-point as each adapter update shifts the score distribution. Mitigation: the per-cycle refit (per-signal isotonic curves, per-benchmark shrinkage and boost layer) and a per-cycle precision audit at the locked cut-points.
    \item \textbf{Verifier and judge disagreement:} The short-answer entailment judge and the long-answer judge could score the same answer inconsistently. Mitigation: drop the long-answer judge if its agreement with the short-answer judge misses a fixed correlation floor.
    \item \textbf{Out-of-memory failures:} A long answer could send the verifier down a branch that exceeds the memory of the graphics processor and halt the cycle. Mitigation: chunking the long-answer batch, and a clean restart from the previous checkpoint after an out-of-memory abort.
    \item \textbf{Forgetting:} A fine-tune could erode general capability. Mitigation: the retention guard with adapter rollback (\Cref{eq:nfr-retention-floor}).
    \item \textbf{Early-stop misfire:} The automatic early-stop check could stop the run too early or fail to stop a run that is not converging. Mitigation: requiring two of its three signals to agree, and the ten-cycle cap. In the realised run two of the signals were wired to the wrong quantities, which made the check fire at its first run; the fault is described in \Cref{sec:env-cycles} and was corrected after the run.
    \item \textbf{Retrieval-utility classifier shift:} The classifier is trained on data that deliberately includes a benchmark outside the evaluation panel (the Natural Questions training split), so it may send some queries to the wrong tier on the panel. Mitigation: nested cross-validation, a threshold sweep, and the on/off ablation on the same trajectory.
\end{itemize}

\subsection{Budget and infrastructure risks}
\label{sec:risk-budget}

Three infrastructure risks were identified.
\begin{itemize}
    \item \textbf{Credit exhaustion:} Running out of rented graphics-processor credit before the trajectory and the comparisons finish would leave a partial run that cannot be compared against the baselines. Mitigation: budgeting each phase against the recorded credit, and archiving every state before it is overwritten.
    \item \textbf{Node reclamation:} The provider could reclaim the rented machine in the middle of a cycle. Mitigation: committing artefacts at every cycle boundary, so that a lost machine costs at most the cycle in flight, and the pre-self-improvement snapshot, from which the run can resume.
    \item \textbf{Storage exhaustion:} The artefacts, model snapshots and passage index could exceed the disk of the rented machine or the quota of the hosts that hold them. Mitigation: sizing the artefact catalogue in advance and keeping the larger snapshots and the passage index on separate storage.
\end{itemize}

\subsection{Mitigation register}
\label{sec:risk-mitigation}

\Cref{tab:risk-register} consolidates the register. Its last column records what happened over the completed programme.

\begin{table}[!htbp]
\centering
\footnotesize
\begin{tabular}{@{}>{\raggedright\arraybackslash}p{2.9cm}>{\raggedright\arraybackslash}p{3.9cm}>{\raggedright\arraybackslash}p{4.1cm}>{\raggedright\arraybackslash}p{3.6cm}@{}}
\toprule
\textbf{Risk} & \textbf{Mechanism} & \textbf{Mitigation} & \textbf{Realised} \\
\midrule
Calibration drift across cycles & Composite loses its precision at the frozen $\tau_{\text{store}}$ under adapter-induced shift & Per-cycle composite refit (per-signal isotonic, per-benchmark shrinkage, boost); per-cycle precision audit & \textbf{Partly}: calibration-fold precision held between $0.732$ and $0.746$ over cycles one to five; the calibration-to-evaluation gap is reported in \Cref{ch:evaluation} \\
Verifier and judge disagreement & Short- and long-answer judges score the same answer inconsistently & Drop the long-answer judge if it misses the correlation floor & \textbf{Yes}: long-answer judge dropped (Pearson $0.584 < 0.70$) \\
Out-of-memory failure on long paths & A long answer sends the verifier past the GPU's memory & Chunked long-answer batch; clean restart after an abort & \textbf{Yes}: chunking added; degenerate samples skipped on out-of-memory \\
Forgetting & A fine-tune erodes general capability & Retention guard with adapter rollback & \textbf{Yes}: cycle four, TriviaQA-test ratio $0.886 < 0.93$; adapter rolled back, memory kept \\
Early-stop misfire & Premature stop, or a run that never stops & Two-of-three signal rule; ten-cycle cap & \textbf{Yes}: two inputs mis-wired, so the check fired at its first run (cycle five) regardless of the trajectory; the stop coincided with the planned end; wiring corrected after the run (\Cref{sec:env-cycles}) \\
Classifier distribution shift & Classifier routes some panel queries to the wrong tier & Nested cross-validation; threshold sweep; on/off ablation & Reported in the ablation (\Cref{ch:evaluation}) \\
GPU credit exhaustion & Run does not finish within the rented credit & Phase budgeting; archive before overwrite & \textbf{Yes}, once (May 2026, during an earlier run); resumed from the recorded step after a top-up \\
Node reclamation & Cycle in flight lost when the machine is reclaimed & Per-cycle artefact commit; pre-self-improvement snapshot & Not observed \\
Storage exhaustion & Artefacts and snapshots exceed the available disk & Advance sizing; separate storage for large snapshots & \textbf{Yes}: the rented node's disk filled; the data were copied to a new node with the provider's support and the run resumed from its last checkpoint \\
Calibration iteration failure & Cycle-zero composite fails the validation gate & Re-iterate with the signal-dump bug fixed and a wider grid & \textbf{Yes}: first iteration failed on the signal-dump bug; second iteration locked \\
\bottomrule
\end{tabular}
\caption[Risk register with mitigations and the realised outcome.]{Risk register with mitigations and the realised outcome over the completed programme.}
\label{tab:risk-register}
\end{table}

\section{Economic Analysis}
\label{sec:economic}

\subsection{Compute budget}
\label{sec:econ-budget}

The compute budget is measured in graphics-processor hours on the rented $32$-gigabyte machine used throughout. The main run took about $193$ hours (\Cref{sec:env-experiment}). The cost of the external baseline panel scales with the number of samples and the number of decode passes per sample, with the chain-of-thought and retrieval-augmented baselines the most expensive. The retrieval-utility-classifier ablation added no run: it compares the main run with the earlier trajectory without the classifier at the cycle-three close, the last cycle that both trajectories committed (the earlier one aborted its cycles four and five under the retention guard). \Cref{app:repro-hardware} also gives a reduced configuration that reproduces the trajectory on a $16$-gigabyte device, so that the cost on smaller hardware can be estimated without re-running the experiment.

\subsection{Cost-benefit on the storage gate}
\label{sec:econ-cost-benefit}

On the main run the storage gate admitted about half of each $4{,}700$-sample stream chunk: between $2{,}343$ and $2{,}484$ entries per cycle over cycles one to five (\Cref{tab:app-per-cycle-memory}). Storing an answer adds no verification work at query time, because every signal the gate uses has already been computed for that answer. It does add work later: each stored entry is re-verified at every later cycle boundary, which is why the retroactive pass grows with the memory (\Cref{sec:env-experiment}). The benefit of a stored entry is the decode it saves each time it answers a later question directly,
\begin{equation}
\Delta B_{\text{store}} \;\ge\; n_{\text{hit}} \cdot c_{\text{gen}},
\label{eq:storage-cost-benefit}
\end{equation}
where $n_{\text{hit}}$ is the number of later queries that the entry answers through the memory-direct tier and $c_{\text{gen}}$ is the cost of the zero-shot or retrieval-augmented decode it replaces. This benefit appears only when questions recur, and in the realised run they never did: each cycle drew a fresh stream chunk that shares no question with earlier chunks, and the evaluation fold is disjoint from both, so almost no query matched a stored episode closely enough for the memory-direct tier to answer it. On the evaluation fold that tier answered only $7$ of $9{,}000$ queries. A direct probe of the router (\Cref{ch:evaluation}) shows the tier working when questions do recur: it answered $42$ of $50$ exact repeats and $19$ of $50$ paraphrases of stored questions. The precision of the gate is what makes the reused answers worth serving.

\subsection{Deployment cost projection}
\label{sec:econ-deployment}

By design, the cost per query should fall as the memory grows, because a memory-direct answer costs an embedding lookup and a similarity search instead of a decode. The cost per query at cycle $t$ is a weighted average over the three tiers,
\begin{equation}
\bar{c}(t) \;=\; \pi_{1}(t)\, c_{1} \;+\; \pi_{2}(t)\, c_{2} \;+\; \pi_{3}(t)\, c_{3},
\qquad \pi_{1}(t) + \pi_{2}(t) + \pi_{3}(t) = 1,
\label{eq:deployment-cost}
\end{equation}
where $\pi_{i}(t)$ is the share of incoming queries dispatched to tier $i$ at cycle $t$ and $c_{i}$ is the average per-query cost on that tier. The design expects $\pi_{1}(t)$ to rise and $\pi_{3}(t)$ to fall as cycles progress. In the realised run neither moved as designed: the memory-direct share stayed at or below $0.1$ percent, because by design no stream or evaluation question repeated a stored one (\Cref{sec:econ-cost-benefit}), and the zero-shot and retrieval-augmented shares each stayed between about $45$ and $55$ percent across cycles (\Cref{ch:evaluation}). The memory-direct saving is therefore expected only where real traffic repeats or paraphrases earlier questions, which is the setting that the deployment-study future-work item in \Cref{ch:conclusion} would test under live user traffic.

\section{Chapter signpost}
\label{sec:ch3-signpost}

The requirements, design commitments, societal, environmental and ethical readings, standards, and risk and economic analyses in this chapter fix the constraints under which the architecture is specified. \Cref{ch:methodology} gives that specification in full: the eight-stage per-query pipeline with the Stage 3b retrieval-utility classifier, the nine-signal verifier ensemble and its calibrated probability composite, the fixed-threshold storage gate and the four-outcome decision tree, the cycle-boundary update passes, the constrained self-improvement loop under the multi-modal retention guard, and the pool-purity theorem, the purification condition, and the trajectory heuristics with their proofs or argument sketches and pre-specified empirical receipts.
   % Requirements, Impact, and Project Management
% ============================================================================
% Chapter 4 — Methodology and Design
% Rubric-mandated two-section top level (4.1 and 4.2), followed by the
% theoretical analysis and the chapter signpost.
%
% Course outcome covered: CO5 (design alternatives, theoretical proofs).
% Figures specific to this chapter live in figures/ch4/; the detailed
% architecture figure is figures/system_architecture_figure.tex.
% ============================================================================

\chapter{Methodology and Design}
\label{ch:methodology}

% Chapter scope: Section 4.1 walks through the method at a high level and
% follows two real queries through it. Section 4.2 specifies every
% component, with the values used in the main run. Section 4.3 states the
% formal results and design heuristics, each labelled with its real status.

\section{Design Process or Methodology Overview}
\label{sec:design-process}

\subsection{The eight-stage pipeline narrative}
\label{sec:overview-pipeline}

CAEM answers each query with an eight-stage pipeline (\Cref{fig:system-architecture}). Stages 1 to 7 run once per query; Stage 8 runs once at the end of every cycle.
\begin{enumerate}
    \item \textbf{Pre-route:} a cheap confidence score $\upre$ is computed for the query before any answer is generated (\Cref{sec:pre-routing}).
    \item \textbf{Encode and look up:} the query is embedded and the nearest stored episode is found, giving a similarity $s$ and that episode's stored score $\ustored$ (\Cref{sec:gen-retrieval,sec:memory-store}).
    \item \textbf{Route:} the router picks one of three tiers from $s$, $\ustored$ and $\upre$ (\Cref{sec:router}). \textbf{Stage 3b:} when the router's choice would default to retrieval, the retrieval-utility classifier decides between retrieval and zero-shot generation (\Cref{sec:ruc}).
    \item \textbf{Answer:} Tier 1 returns the stored answer; Tier 2 generates without retrieval; Tier 3 retrieves Wikipedia passages and generates with them.
    \item \textbf{Verify:} every generated answer is scored by nine signals in four families (\Cref{sec:verifier}), which the calibrated composite turns into one probability of correctness, $\ustored$ (\Cref{sec:composite}).
    \item \textbf{Decide:} the four-outcome rule stores, defers, abstains or discards (\Cref{sec:conformal-gate,sec:decision}).
    \item \textbf{Output:} the user receives the answer with its calibrated probability and a reliability tag, or an insufficient-evidence message when the system abstains.
    \item \textbf{Cycle boundary:} the adapter is fine-tuned and checked for forgetting, the deferred buffer is revisited, the memory is consolidated, the composite is refitted and every stored episode is re-verified (\Cref{sec:cycle-boundary}).
\end{enumerate}

\subsection{The cycle loop walk-through}
\label{sec:overview-cycle}

A cycle has two phases. In the per-query phase, the pipeline answers a fresh stream chunk of $4{,}700$ training-panel questions (FEVER $2{,}000$, TriviaQA $2{,}000$, CommonsenseQA $700$). Each chunk shares no question with earlier chunks, with the calibration fold or with the evaluation fold. Answers that pass the storage gate enter the episodic memory and borderline answers enter the deferred buffer. The evaluation fold is then answered without writing to memory.

In the cycle-boundary phase, five steps run in a fixed order (\Cref{fig:cycle-boundary-order}). First, the low-rank adapter is fine-tuned on stored episodes whose $\ustored$ is at least $0.70$ and the retention guard checks three held-out probes; if any probe falls below $93$ percent of its starting accuracy, the previous adapter is restored and the remaining steps are skipped. Otherwise the deferred buffer is reconsidered, near-duplicate episodes are merged, the fixed calibration fold is re-scored with the new adapter so that the temperature and the composite can be refitted, and finally every stored episode is re-verified with the refitted composite. Re-verification can lower a stored score as well as raise it, and entries that fall below $0.50$ are removed. The memory itself is never rolled back: answers verified during a failed cycle stay in it.

% Cycle-boundary order as executed (outputs/full_run/experiment.log, e.g. cycle 2).
\begin{figure}[!htbp]
\centering
\resizebox{\textwidth}{!}{%
\begin{tikzpicture}[
  font=\small,
  st/.style={rectangle, rounded corners=4pt, line width=0.9pt, align=center, inner sep=4pt, text width=2.55cm, minimum height=2.1cm, draw=teal!65!black, fill=teal!7},
  gd/.style={diamond, aspect=1.6, line width=0.9pt, align=center, inner sep=1pt, draw=red!60!black, fill=red!6, font=\footnotesize},
  flow/.style={-{Stealth[length=2.2mm]}, line width=0.9pt, draw=black!65},
]
  \node[st] (a) at (0,0) {\textbf{1. Fine-tune}\\{\footnotesize LoRA adapter on stored episodes with $\ustored \ge 0.70$}};
  \node[gd] (g) at (3.35,0) {retention\\$\ge 0.93$?};
  \node[st] (b) at (6.7,0) {\textbf{2. Reconsider deferred}\\{\footnotesize promote, keep, or drop after two passes}};
  \node[st] (c) at (9.75,0) {\textbf{3. Consolidate}\\{\footnotesize merge near-duplicates (similarity $> 0.92$)}};
  \node[st] (d) at (12.8,0) {\textbf{4. Recalibrate}\\{\footnotesize re-score the calibration fold; refit temperature and composite}};
  \node[st] (e) at (15.85,0) {\textbf{5. Re-verify memory}\\{\footnotesize re-score every entry; remove if $\ustored < 0.50$}};
  \draw[flow] (a) -- (g); \draw[flow] (g) -- node[above, font=\scriptsize]{yes} (b);
  \draw[flow] (b) -- (c); \draw[flow] (c) -- (d); \draw[flow] (d) -- (e);
  \node[rectangle, rounded corners=3pt, draw=red!60!black, fill=red!4, align=center, font=\footnotesize, text width=6.2cm] (r) at (3.35,-2.35) {\textbf{no:} restore the previous adapter and skip steps 2 to 5; the memory is kept unchanged};
  \draw[flow] (g) -- (r);
\end{tikzpicture}}
\caption[Order of the cycle-boundary steps as executed.]{Order of the cycle-boundary steps as executed in the main run. After step 5 the next cycle begins by answering a fresh stream chunk with Stages 1 to 7. Cycle four took the rollback branch.}
\label{fig:cycle-boundary-order}
\end{figure}


\subsection{A worked example}
\label{sec:overview-worked-example}

Two real records from the main run show the pipeline end to end (\Cref{tab:worked-example}).

\textbf{A stored answer:} the TriviaQA question \emph{In what year (or century) was Michelangelo born?} arrived in the cycle-one stream chunk. Its pre-routing confidence ($0.522$) was above the safety floor, no stored episode was similar enough for Tier 1 or Tier 2, and the retrieval-utility classifier kept the default, so Tier 3 retrieved passages and the generator answered \emph{1475}, citing a passage that states the birth year. The sampled reasoning chains agreed ($s_{\text{avg}} = 0.992$) and entailed the answer ($p_{\text{entail}} = 0.939$), while the best passage gave only moderate support ($p_{\text{ground,max}} = 0.469$). The composite weighed these signals into $\ustored = 0.760$, above the storage cut-point of $0.60$, so the answer was stored and shown to the user as verified. It was correct.

\textbf{An abstention:} the TriviaQA question \emph{The song ``If I Ruled The World'' comes from which musical?} is in the evaluation fold. At cycle zero, with only the cold-start memory, it fell through to Tier 3, where the retrieved passages mention the song but not the musical, and the generator answered \emph{I do not know}. Every signal was low and the composite gave $\ustored = 0.127$. Because this is below the deferral cut-point of $0.45$ and the best passage supports the answer weakly ($p_{\text{ground,max}} = 0.112 < 0.20$), the decision was to abstain: the user receives the insufficient-evidence message with the probability shown, and nothing is stored. The gold answer is \emph{Pickwick}, so abstaining avoided a likely wrong answer.

\begin{table}[!htbp]
\centering
\small
\begin{tabular}{@{}lcc@{}}
\toprule
 & \textbf{Stored example} & \textbf{Abstained example} \\
Question & Michelangelo's birth year & ``If I Ruled The World'' \\
\midrule
Cycle, fold & 1, stream chunk & 0, evaluation fold \\
Tier & 3 & 3 \\
$\upre$ & $0.522$ & $0.659$ \\
$u_{\text{token}}$, $u_{\text{dropout}}$, $u_{\text{internal}}$ & $0.163$, $0.600$, $0.282$ & $0.022$, $0.800$, $0.111$ \\
$s_{\text{avg}}$ / $p_{\text{entail}}$ & $0.992$ / $0.939$ & $0.804$ / $0.115$ \\
$p_{\text{ground,max}}$ / $p_{\text{ground,mean}}$ & $0.469$ / $0.292$ & $0.112$ / $0.094$ \\
$q_{\text{a,rel}}$ & $0.515$ & $0.500$ \\
\midrule
$\ustored$ & $\mathbf{0.760}$ & $\mathbf{0.127}$ \\
Decision & store & abstain \\
Outcome & correct (\emph{1475}) & no answer (gold \emph{Pickwick}) \\
\bottomrule
\end{tabular}
\caption[Signal readings for the two worked examples.]{Signal readings for the two worked examples, read from the main run's per-sample records. The atomic-fact score equals the mean-passage score for both (\Cref{sec:verifier-grounding}).}
\label{tab:worked-example}
\end{table}

\section{Preliminary Design or Design (Model) Specification}
\label{sec:design-spec}

\begin{figure}[p]
\centering
\resizebox{\textwidth}{!}{%
\begin{tikzpicture}[
    xshift=-0.9cm,
    font=\normalsize,
  % ---- Node styles ----
  stage/.style  = {rectangle, draw=teal!65!black, line width=0.9pt, rounded corners=5pt,
                   minimum height=1.5cm, minimum width=6.5cm, align=center, fill=teal!8},
  router/.style = {rectangle, draw=teal!65!black, line width=0.9pt, rounded corners=5pt,
                   minimum height=1.8cm, minimum width=8.5cm, align=center, fill=yellow!18},
  mcyl/.style   = {cylinder, draw=teal!65!black, line width=0.9pt,
                   shape border rotate=90, aspect=0.25,
                   minimum height=3.0cm, minimum width=3.2cm, align=center, fill=teal!12},
  tA/.style     = {rectangle, draw=green!55!black, line width=1pt, rounded corners=5pt,
                   minimum height=1.7cm, minimum width=4.0cm, align=center, fill=green!12},
  tB/.style     = {rectangle, draw=red!55!orange, line width=1pt, rounded corners=5pt,
                   minimum height=1.7cm, minimum width=4.0cm, align=center, fill=red!12!orange!10},
  tC/.style     = {rectangle, draw=red!60!black, line width=1pt, rounded corners=5pt,
                   minimum height=1.7cm, minimum width=4.0cm, align=center, fill=red!10},
  verify/.style = {rectangle, draw=teal!65!black, line width=1pt, rounded corners=5pt,
                   minimum height=4.2cm, minimum width=9.5cm, align=center, fill=teal!5},
  decide/.style = {rectangle, draw=teal!65!black, line width=0.9pt, rounded corners=5pt,
                   minimum height=1.8cm, minimum width=9.5cm, align=center, fill=yellow!18},
  oServe/.style = {rectangle, draw=green!60!black, line width=1pt, rounded corners=4pt,
                   minimum height=1.4cm, minimum width=3.6cm, align=center,
                   fill=green!30, font=\small\scshape},
  oStore/.style = {rectangle, draw=green!60!black, line width=0.9pt, rounded corners=3pt,
                   minimum height=1.2cm, minimum width=2.7cm, align=center,
                   fill=green!28, font=\small\scshape},
  oDefer/.style = {rectangle, draw=green!45!black, line width=0.9pt, rounded corners=3pt,
                   minimum height=1.2cm, minimum width=2.7cm, align=center,
                   fill=green!12, font=\small\scshape},
  oAbst/.style  = {rectangle, draw=orange!75!black, line width=0.9pt, rounded corners=3pt,
                   minimum height=1.2cm, minimum width=2.7cm, align=center,
                   fill=orange!28, font=\small\scshape},
  oDisc/.style  = {rectangle, draw=red!60!black, line width=0.9pt, rounded corners=3pt,
                   minimum height=1.2cm, minimum width=2.7cm, align=center,
                   fill=red!28, font=\small\scshape},
  simp/.style   = {rectangle, draw=violet!60!black, line width=1pt, rounded corners=5pt,
                   minimum height=3.2cm, minimum width=17.0cm, align=center, fill=violet!7},
  % ---- Arrow styles ----
  farr/.style   = {-Latex, line width=0.9pt},
  garr/.style   = {-Latex, line width=1.1pt, green!50!black},
  oarr/.style   = {-Latex, line width=1.1pt, color=red!55!orange},
  rarr/.style   = {-Latex, line width=1.1pt, red!60!black},
  darr/.style   = {<->, dashed, line width=0.9pt, teal!55!black},
  varr/.style   = {-Latex, line width=1pt, dashed, violet!60!black},
    elab/.style   = {font=\small\itshape, inner sep=2pt, fill=none, draw=none}
]

% ==================================================================
% NODES   (main trunk centred at x = 6; y compressed to fit A4 page)
% ==================================================================

\node[stage]  (q)   at (6,   0.0)  {\textbf{Query $q$}};

\node[stage]  (s1)  at (6,  -2.5)  {%
    \textbf{S1\quad Pre-route confidence}\\[4pt]
    $u_{\text{pre}} = 0.60\,u_{\text{token}} + 0.40\,(1{+}C_{\text{conv}})^{-1}$\\
    {\scriptsize 32-token greedy decode $+$ one forward pass;\;$T$ refit per cycle}};

\node[stage]  (s2)  at (6,  -5.3)  {%
    \textbf{S2\quad Episodic memory retrieval}\\[4pt]
    SBERT 768-d cosine nearest-neighbour search};

% Episodic memory cylinder  (right of S2, connected by dashed bidirectional arrow)
\node[mcyl]   (mem) at (16.5,-5.3) {%
    \textbf{Episodic}\\[2pt]\textbf{Memory}\\[3pt]{\scriptsize exact search}};

\node[router] (s3)  at (6, -8.3)  {%
    \textbf{S3\quad Three-tier router}\\[4pt]
    $S(q)=0.70\cdot\mathrm{sim}(q,e^{*})+0.30\cdot u_{\text{stored}}(e^{*})$\\
    {\scriptsize safety override:\;if $u_{\text{pre}}<\phi_{\text{safe}}{=}0.38$ the score is bypassed and the RUC decides Tier~2 vs Tier~3}};

% ---- S3b: Retrieval-utility classifier (intercepts the safety-override and score-formula fall-throughs) ----
\node[router, minimum height=1.5cm, minimum width=5.0cm, fill=yellow!12]
      (ruc) at (16.5,-9.8) {%
    \textbf{S3b\quad RUC}\\[2pt]
    {\scriptsize 36-feature LR head;\;$p_{\text{RAG}}\gtrless\tau_{\text{RUC}}{=}0.375$}\\
    {\scriptsize routes safety-veto and else fall-through to Tier~2 or Tier~3}};

% ---- Tier row  (T1 far-left / T2 centre / T3 far-right)
%  S5 min-width = 9.5 cm centred at x=6  ->  left edge  ~ x = 1.25
%  T1 placed at x=0 is therefore outside S5, allowing a clean bypass path.
\node[tA] (t1) at (0.0, -12.2)  {%
    \textbf{Tier 1}\\[3pt]return $e^{*}$ answer\\{\scriptsize(from verified memory)}};
\node[tB] (t2) at (6.0, -12.2)  {%
    \textbf{Tier 2}\\[3pt]zero-shot generation\\{\scriptsize(fine-tuned generator, no retrieval)}};
\node[tC] (t3) at (12.0,-12.2)  {%
    \textbf{Tier 3}\\[3pt]retrieval-augmented\\{\scriptsize(top-5 Wikipedia passages)}};

% ---- S4: Tier execution grouping (compact left-side label; no wrapper) ----
\node[draw=none, inner sep=6pt, fit=(t1)(t2)(t3),
      label={[font=\small\bfseries, teal!65!black, anchor=east, xshift=-6pt]left:S4}]
     (s4box) {};

% Tier 1 direct-serve terminal  (left side, same y-band as S5)
\node[oServe] (serve) at (-3.5,-17.0) {%
    \textsc{Answer served}\\{\scriptsize from verified memory}};

% S5 centred at x=6, min-width=9.5 cm  ->  spans ~ x=1.25..10.75
\node[verify] (s5)  at (6, -17.0)  {%
    \textbf{S5\quad Nine-signal Verifier}\\[6pt]
    \textit{Internal:}\quad$u_{\text{token}},\;u_{\text{dropout}},\;u_{\text{internal}}$\\[3pt]
    \textit{Sample-set:}\quad$s_{\text{avg}},\;p_{\text{entail}}$\\[3pt]
    \textit{Grounding:}\quad$p_{\text{ground}}^{\max},\;p_{\text{ground}}^{\text{mean}},\;p_{\text{ground}}^{\text{atomic}}$\\[3pt]
    \textit{Relevance:}\quad$q_{\text{a,rel}}$\\[5pt]
    {\footnotesize\textbf{confabulation gate:}\;
     $u_{\text{internal}}{\geq}0.70\;\wedge\;p_{\text{ground}}^{\max}{\leq}0.20
     \;\Rightarrow\;$\textsc{Discard (early exit)}}};

\node[decide] (s6)  at (6, -21.0)  {%
    \textbf{S6\quad Four-outcome storage decision}\\[5pt]
    {\scriptsize
    $u_{\text{stored}}{\geq}0.60\Rightarrow\textsc{Store}$\quad
    $u_{\text{stored}}{\geq}0.45\Rightarrow\textsc{Deferred}$\quad
    $p_{\text{ground}}^{\max}{<}0.20\Rightarrow\textsc{Abstain}$\quad
    \textit{else:}\;\textsc{Discard}}};

% Terminal outcome boxes  (y = -24.5)
\node[oStore] (out1) at (2.0, -24.5)  {\textsc{Store}\\{\tiny write to memory}};
\node[oDefer] (out2) at (5.0, -24.5)  {\textsc{Deferred}\\{\tiny buffer; re-checked next cycle}};
\node[oAbst]  (out3) at (8.0, -24.5)  {\textsc{Abstain}\\{\tiny principled refusal}};
\node[oDisc]  (out4) at (11.0,-24.5)  {\textsc{Discard}\\{\tiny suppressed}};

% ---- S7: Memory write grouping (compact left-side label; no wrapper) ----
\node[draw=none, inner sep=6pt, fit=(out1)(out2)(out3)(out4),
      label={[font=\small\bfseries, teal!65!black, anchor=east, xshift=-6pt]left:S7}]
     (s7box) {};

% S8: between-cycle self-improvement  (wide band at bottom)
\node[simp] (s8) at (6, -27.5) {%
    \textbf{S8\quad Between-cycle self-improvement loop}
    \quad\textit{(runs once at each cycle boundary, not per query)}\\[6pt]
    (1)~Fine-tune the LoRA adapter on stored episodes with $\ustored{\geq}0.70$ (pool-reweighting layer; cross-entropy only)\\[3pt]
    (2)~Retention guard:\;worst probe ratio $\geq0.93$, else restore $\theta_{\text{prev}}$ and skip (3)--(5)\\[3pt]
    (3)~Reconsider deferred entries\quad(4)~Consolidate and recalibrate\quad
    (5)~Re-verify memory ($\ustored$ may rise or fall; remove below $\tau_{\text{retro}}{=}0.50$)};

% ==================================================================
% ARROWS
% ==================================================================

% ---- Main trunk ----
\draw[farr] (q.south)  -- (s1.north);
\draw[farr] (s1.south) -- node[elab,right]{$u_{\text{pre}}$} (s2.north);
\draw[farr] (s2.south) -- node[elab,right]{$u_{\text{pre}}$ + query embedding} (s3.north);

% ---- S2 <-> Episodic memory  (bidirectional dashed) ----
\draw[darr] (s2.east)
    -- node[elab,above,pos=0.55,align=center]{%
        retrieve $e^{*}$\;:\;$\mathrm{sim}(q,e^{*})$,\\$u_{\text{stored}}(e^{*})$}
    (mem.west);

% ---- S3 -> three tiers ----
\coordinate (br) at (6, -9.8);   % branch waypoint below S3
\draw[farr] (s3.south) -- (br);

% Tier 1 (green): branch LEFT, condition label above horizontal leg
\draw[garr]
    (br) --
    node[elab,above,pos=0.45]{\textbf{if}\;$S(q)\geq0.90$}
    (0.0,-9.8) -- (t1.north);

% Tier 2 (orange): straight down, condition label to the LEFT
\draw[oarr]
    (br) --
    node[elab,left,pos=0.5]{\textbf{elif}\;$\mathrm{sim}(q,e^{*})\geq0.75$}
    (t2.north);

% Score-formula else + safety-override both reach the RUC (S3b), not Tier 3 directly
\draw[farr]
    (br) --
    node[elab,above,pos=0.6]{\textbf{else}\;\textbf{or}\;$u_{\text{pre}}<\phi_{\text{safe}}$}
    (ruc.west);

% RUC -> Tier 3 (retrieval-augmented) when p_RAG >= tau_RUC
\draw[rarr]
    (ruc.south) --
    (ruc.south |- t3.east) --
    node[elab,above,pos=0.5]{$p_{\text{RAG}}\ge\tau_{\text{RUC}}$}
    (t3.east);

% RUC -> Tier 2 (zero-shot) when p_RAG < tau_RUC; direct diagonal to t2
\draw[oarr]
    (ruc.south west) --
    node[elab,above=2pt,sloped,pos=0.5]{$p_{\text{RAG}}<\tau_{\text{RUC}}$}
    (t2.north east);

% ---- Tier 1 bypass: route LEFT then DOWN, well outside S5 ----
% S5 left edge ≈ x=1.25; this path uses x=-3.5 to stay clear.
\coordinate (t1byp1) at (0.0, -14.0);    % just below T1
\coordinate (t1byp2) at (-3.5,-14.0);    % far left
\draw[garr]
    (t1.south) -- (t1byp1) --
    node[elab,below,pos=0.5]{\textit{bypass S5 and S6: episode already verified at storage time}}
    (t1byp2) -- (serve.north);

% ---- T2 -> S5 ----
\draw[oarr]
    (t2.south) --
    node[elab,right,pos=0.45]{answer + reasoning chain}
    (t2 |- s5.north);

% ---- T3 -> S5 routed on the right to avoid center-label overlap ----
\draw[rarr]
    (t3.south) --
    (12.0,-14.2) --
    (13.6,-14.2) --
    node[elab,right,pos=0.55]{answer + reasoning chain}
    (13.6,-16.9) --
    (s5.east);

% ---- S5 -> S6 ----
\draw[farr]
    (s5.south) --
    node[elab,right]{outputs:\;$u_{\text{stored}},\;p_{\text{contra}},\;p_{\text{ground}}^{\max}$}
    (s6.north);

% ---- S6 fan-out -> terminal boxes ----
\coordinate (br2) at (6,-22.8);
\draw[farr] (s6.south) -- (br2);
\draw[farr] (br2) -| (out1.north);
\draw[farr] (br2) -| (out2.north);
\draw[farr] (br2) -| (out3.north);
\draw[farr] (br2) -| (out4.north);

% ---- Store + Deferred accumulate into S8 at cycle end ----
% Drop only 0.3 cm so the horizontal segment stays ABOVE s8.north,
% then the vertical leg enters s8 cleanly from the top.
\draw[varr]
    (out1.south) -- ++(0,-0.3) -|
    node[elab,left,pos=0.05]{\textit{cycle end}}
    ($(s8.north)+(-2.5,0)$);
\draw[varr]
    (out2.south) -- ++(0,-0.3) -|
    ($(s8.north)+(-0.8,0)$);

% ---- S8 -> S1 : OUTER TEN-CYCLE LOOP ----
%
%  After S8 concludes, the updated weights θ_new (or θ_prev on rollback)
%  are loaded and Stage 1 of Cycle n+1 begins.  This arrow is the outer loop.
%
%  Route: S8.west  →  far left (x ≈ -5.5)  →  up to y of S1  →  S1.west
%
\draw[varr]
    (s8.west)
    -- ++(-5.5, 0) coordinate (ll)
    -- node[elab,left,pos=0.45,align=left]{%
        \textbf{Cycle $n{+}1$ begins}\\[3pt]
        S1 receives\\
        $\theta_{\text{new}}$ (fine-tuned)\\[1pt]
        or $\theta_{\text{prev}}$ (restored\\
        if abort guard fired)}
    (ll |- s1.west)
    -- (s1.west);

\end{tikzpicture}%
}
\caption[The eight-stage CAEM pipeline.]{The eight-stage CAEM pipeline. Stages S1 to S7 run once per query; S8 runs once per cycle, in the order of \Cref{fig:cycle-boundary-order}. The outer dashed arrow hands the updated or restored adapter to the next cycle.}
\label{fig:system-architecture}
\end{figure}


\subsection{Generator and retrieval substrate}
\label{sec:gen-retrieval}

Four components are fixed across every cycle.

\textbf{Generator:} Qwen2.5-3B-Instruct \cite{qwen25}, a three-billion-parameter instruction-tuned decoder, generates every Tier 2 and Tier 3 answer and supplies the internal-calibration signals. Only a low-rank adapter on top of it is trained (about $2$ percent of its parameters), which keeps each cycle's fine-tune within one $32$-gigabyte graphics processor. A small generator is also accurate enough on the training panel for the storage gate to have a usable signal.

\textbf{Retrieval corpus and encoder:} Tier 3 retrieves from a dense index of about $21$ million passages of roughly $100$ words from English Wikipedia, built in the dense-retrieval style of Karpukhin et al.~\cite{karpukhin-etal-2020-dense}. Queries, passages and memory entries are all embedded with the same sentence encoder, all-mpnet-base-v2 \cite{reimers-gurevych-2019-sentence} ($768$ dimensions), so a query can be compared directly with both. Tier 3 places the top five passages in the prompt. Both generating tiers ask the generator for a short reasoning followed by a line beginning \emph{Answer:}; the answer is read from that line, and the memory stores it together with the reasoning chain.

\textbf{Entailment judge:} MiniCheck-Flan-T5-Large \cite{tang2024minicheck} scores whether a premise supports a claim. It was trained on claim-support data generated by language models, which matches the kind of answers it is asked to judge. It returns a single support probability rather than a three-way label, so no contradiction signal is used. Inputs longer than its $512$-token encoder are truncated; a second, long-input judge built on the frozen generator was considered but dropped before the main run because its agreement with MiniCheck missed the required correlation floor (Pearson $0.584 < 0.70$, \Cref{ch:evaluation}).

\textbf{Reranker:} the cross-encoder bge-reranker-v2-m3 scores text pairs for relevance. The verifier uses it twice: to rerank the top $20$ retrieved passages down to the top three before entailment is scored, and to score how relevant the answer is to the question. Keeping the reranker separate from the entailment judge is what lets an off-topic answer that happens to be supported by an off-topic passage still be caught.

\subsection{Pre-routing confidence}
\label{sec:pre-routing}

The pre-routing confidence $\upre$ is a cheap estimate of how ready the generator is to answer, computed before any answer is produced. It feeds the safety override of the router. It combines two quantities (\Cref{alg:pre-routing}).

\textbf{Token confidence:} the generator runs a $32$-token greedy decode on the query and $u_{\text{token}}$ is the geometric mean of the probabilities of the tokens it chose. A model that hesitates between continuations scores low.

\textbf{Layer-convergence ratio:} one forward pass over the query records the hidden states of every decoder layer. The ratio $c_{\text{conv}}$ divides the variance of the hidden states in the first half of the layers by their variance in the second half, adapted from Nandakishor~\cite{nandakishor2025routing}. A large ratio means the representation is still unsettled in the early layers.

The two are combined with fixed weights, $z = 0.60\,u_{\text{token}} + 0.40/(1+c_{\text{conv}})$, and calibrated by temperature scaling \cite{pmlr-v70-guo17a} with a per-benchmark temperature $T_b$ fitted on the calibration fold. Temperature scaling was chosen because it has one parameter and preserves the order of scores.

\textbf{Realised behaviour:} on FEVER and TriviaQA the fitted temperature reached its upper bound ($T = e^{3} \approx 20.1$) at cycle zero and was refitted downward each cycle with exponential smoothing (to about $14$ at cycle one and $10$ at cycle two). A temperature this large pulls $\upre$ toward $0.5$, so on these benchmarks the score rarely approached the safety floor of $0.38$. \Cref{ch:evaluation} reports the consequence: the override fired on $63$ of the $9{,}000$ evaluation queries, $62$ of them on CommonsenseQA after its floor was refitted upward.

\begin{algorithm}[!htbp]
\caption[Pre-routing confidence (Stage~1).]{Pre-routing confidence (Stage~1).}
\label{alg:pre-routing}
\begin{algorithmic}[1]
\Require query $q$, generator $G$, decode length $L = 32$, per-benchmark temperature $T_{b}$
\State $(\log p_{1}, \dots, \log p_{n}) \gets$ log-probabilities of the tokens chosen by a greedy decode of up to $L$ tokens
\State $u_{\text{token}} \gets \exp\!\left(\tfrac{1}{n}\sum_{i=1}^{n} \log p_{i}\right)$ \Comment{geometric mean of token probabilities}
\State $\mathbf{H}^{(1)}, \dots, \mathbf{H}^{(D)} \gets$ hidden states of the $D$ decoder layers in one forward pass over $q$
\State $c_{\text{conv}} \gets \mathrm{Var}\big(\mathbf{H}^{(1:D/2)}\big) \,/\, \mathrm{Var}\big(\mathbf{H}^{(D/2+1:D)}\big)$ \Comment{early-half over late-half variance}
\State $z \gets 0.60\, u_{\text{token}} + 0.40\, /\,(1+c_{\text{conv}})$
\State \Return $\upre \gets \sigma\!\big(\mathrm{logit}(z)/T_{b}\big)$ \Comment{temperature scaling}
\end{algorithmic}
\end{algorithm}

\subsection{Episodic memory store}
\label{sec:memory-store}

The episodic memory is CAEM's long-term store of verified answers. It follows memory-augmented networks \cite{weston2015memory,graves2014neural,graves2016differentiable} and writeable episodic stores for language models \cite{larimar2024das,xu2025amem}, with one distinguishing feature: every entry carries the calibrated score $\ustored$ that admitted it.

\textbf{Record:} each entry holds the question, the answer, the reasoning chain, the question's embedding, the nine signal readings, $\ustored$, the source benchmark and the cycle in which it was stored. Keeping the question and answer together lets the cycle-boundary pass re-verify an entry without regenerating the answer.

\textbf{Cold start:} before cycle one, the pipeline answered questions from each training benchmark's seed pool until about $200$ per benchmark had been verified and stored ($603$ answers). Answers that were wrong against the gold labels were then removed, leaving $431$ entries (FEVER $155$, TriviaQA $136$, CommonsenseQA $140$). The cold-start memory is therefore correct by construction; every later entry is admitted by the verifier alone.

\textbf{Search:} lookup is by cosine similarity between question embeddings, because the encoder is trained to place paraphrases in the same direction. The index is exact (flat) below $20{,}000$ entries and would switch to an approximate inverted-file index above that size; the memory reached $12{,}175$ entries, so every search in the main run was exact.

\textbf{Novelty and quality filters:} a new answer is not stored if an existing entry's question has similarity above $0.95$ with the new question, so paraphrases do not fill the memory, or if the answer contains a leaked prompt placeholder or an evasive phrase such as \emph{the context does not mention}.

\textbf{Serve counter:} an entry served by Tier 1 ten times is re-verified at the next cycle boundary even if nothing else has changed. With Tier 1 serving only $7$ evaluation queries in the whole run, this never triggered.

\textbf{Consolidation:} at each cycle boundary, entries whose questions have similarity above $0.92$ are grouped; a group is merged into its highest-scoring entry only if the entries give the same answer and their scores differ by at most $0.15$. The lower threshold catches near-paraphrases that slipped past the novelty filter. It removed $3$, $11$ and $44$ entries at cycles two, three and five.

\textbf{Capacity:} the memory is capped at one million entries, with a pruning rule that removes the lowest-value fifth when the store is $95$ percent full. The rule never fired, because the memory stayed far below the cap.

\subsection{Three-tier router with the safety override}
\label{sec:router}

The router chooses the tier from three inputs: the similarity $s$ between the query and its nearest stored question, that neighbour's stored score $\ustored$, and the pre-routing confidence $\upre$. Its rules are checked in a fixed order (\Cref{alg:router}, \Cref{fig:routing-map}; \Cref{tab:ruc-dispatch} gives the final rules including the classifier).
\begin{enumerate}
    \item \textbf{Safety override:} if $\upre < 0.38$, the query skips the memory and goes to the retrieval-utility classifier.
    \item \textbf{Tier 1:} if the combined score $\mathcal{S} = 0.70\,s + 0.30\,\ustored$ is at least $0.90$, the neighbour's stored answer is returned without generating.
    \item \textbf{Tier 2:} otherwise, if $s \ge 0.75$, the fine-tuned generator answers without retrieval.
    \item \textbf{Fall-through:} otherwise the query goes to the retrieval-utility classifier, which sends it to Tier 3 or Tier 2.
\end{enumerate}
Similarity carries the larger weight in $\mathcal{S}$ because a wrong neighbour is a different question altogether, whereas an under-scored correct neighbour still holds a correct answer. Because $\ustored \le 1$, Tier 1 needs $s \ge 0.857$: only a close paraphrase of a stored question is answered from memory. The safety floor was lowered from $0.60$ to $0.38$ after a routing audit of an earlier run showed that the higher floor forced well-supported queries onto retrieval; in the main run it was refitted per benchmark at each cycle and stayed at $0.38$ except on CommonsenseQA ($0.491$ at cycles three and four, $0.569$ at cycle five).

Separating the three questions (is there a trusted close match, is the model ready to answer, would retrieval help) distinguishes this router from a single-threshold dispatcher and from Adaptive-RAG of Jeong et al.~\cite{jeong2024adaptive}, which routes by query complexity rather than by the stored score of a memory entry.

\textbf{Design intent and realised shares:} the design expects the memory-direct share to grow as the memory fills, and the zero-shot share to grow as fine-tuning folds retrieved knowledge into the generator. In the main run neither happened. Tier 1 answered at most $0.1$ percent of evaluation queries, because no stream or evaluation question repeats a stored one; a direct probe confirms that Tier 1 does answer repeated questions ($42$ of $50$ exact repeats, $19$ of $50$ paraphrases, \Cref{ch:evaluation}). The zero-shot and retrieval shares each stayed between about $45$ and $55$ percent of evaluation queries across cycles.

% Routing map (Stage 3). Region boundaries follow the dispatch rules of
% tab:dispatch-rules: Tier 1 iff 0.70 s + 0.30 u >= 0.90; else Tier 2 iff s >= 0.75.
\begin{figure}[!htbp]
\centering
\begin{tikzpicture}[scale=5.2, font=\small]
  \fill[red!8]   (0,0) rectangle (0.75,1);
  \fill[orange!14] (0.75,0) -- (1,0) -- (1,0.6667) -- (0.857,1) -- (0.75,1) -- cycle;
  \fill[green!22] (0.857,1) -- (1,1) -- (1,0.6667) -- cycle;
  \draw[dashed, black!55, line width=0.6pt] (0.75,0) -- (0.75,1);
  \draw[green!45!black, line width=0.9pt] (0.857,1) -- (1,0.6667);
  \draw[->, black!70, line width=0.7pt] (0,0) -- (1.07,0) node[right, font=\footnotesize] {similarity $s$};
  \draw[->, black!70, line width=0.7pt] (0,0) -- (0,1.07) node[above, font=\footnotesize] {stored score $\ustored$ of the neighbour};
  \foreach \x in {0.25,0.5,0.75,1} \draw[black!45] (\x,0) -- (\x,-0.02) node[below, font=\scriptsize] {\x};
  \foreach \y in {0.25,0.5,0.75,1} \draw[black!45] (0,\y) -- (-0.02,\y) node[left, font=\scriptsize] {\y};
  \node[align=center, text=red!55!black, font=\footnotesize] at (0.375,0.5) {\textbf{Tier 3 by default}\\retrieval-utility classifier\\chooses Tier 3 or Tier 2};
  \node[align=center, text=orange!55!black, font=\footnotesize] at (0.875,0.3) {\textbf{Tier 2}\\zero-shot};
  \node[text=green!35!black, font=\footnotesize\bfseries] at (0.94,0.88) {Tier 1};
  \node[anchor=east, font=\scriptsize, text=green!35!black] at (0.83,0.93) {$0.70\,s + 0.30\,\ustored = 0.90$};
  \draw[green!35!black, line width=0.4pt] (0.83,0.93) -- (0.905,0.86);
  \node[anchor=north, font=\scriptsize, text=black!60] at (0.75,-0.07) {$s = 0.75$};
\end{tikzpicture}
\caption[Routing map of the three-tier router.]{Routing map of the three-tier router. Before this map is consulted, the safety override sends any query with $\upre < 0.38$ to the retrieval-utility classifier.}
\label{fig:routing-map}
\end{figure}



\begin{algorithm}[!htbp]
\caption[Three-tier router with safety override (Stage~3).]{Three-tier router with safety override (Stage~3).}
\label{alg:router}
\begin{algorithmic}[1]
\Require similarity $s$, neighbour score $\ustored$, pre-routing confidence $\upre$
\Require $\lambda = 0.70$, $\theta_{1} = 0.90$, $\theta_{\text{sim}} = 0.75$, per-benchmark floor $\phi_{\text{safe}}$ ($0.38$ unless refitted)
\State $\mathcal{S} \gets \lambda\, s + (1-\lambda)\, \ustored$
\If{$\upre < \phi_{\text{safe}}$}
    \State \Return Tier 3 \Comment{safety override; refined by \Cref{alg:ruc}}
\ElsIf{$\mathcal{S} \ge \theta_{1}$}
    \State \Return Tier 1
\ElsIf{$s \ge \theta_{\text{sim}}$}
    \State \Return Tier 2
\Else
    \State \Return Tier 3 \Comment{fall-through; refined by \Cref{alg:ruc}}
\EndIf
\end{algorithmic}
\end{algorithm}

\subsection{Retrieval-utility classifier}
\label{sec:ruc}

Retrieval does not help every question. On questions built around a popular misconception, retrieved passages can reinforce the misconception, and on questions the generator already knows they can distract it. An earlier trajectory without a routing refinement showed this: TruthfulQA exact match fell from cycle to cycle as retrieval-amplified answers entered the verified pool. The selective-retrieval literature reached the same conclusion, including Mallen et al.~\cite{mallen2023whennot}, Wang et al.~\cite{wang2023skr}, Jeong et al.~\cite{jeong2024adaptive} and Baek et al.~\cite{baek2025probing} (\Cref{sec:lit-selective-retrieval}). The retrieval-utility classifier is CAEM's response: whenever the router would send a query to retrieval by default, the classifier predicts whether retrieval will help and, if not, sends the query to Tier 2 instead (\Cref{tab:ruc-dispatch}, \Cref{alg:ruc}).

\textbf{Model:} a logistic regression on $36$ standardised features outputs $p_{\text{RAG}}$, the probability that retrieval helps. The query goes to Tier 3 if $p_{\text{RAG}} \ge 0.375$ and to Tier 2 otherwise. The threshold is below $0.5$ because it was chosen to maximise the number of correctly routed queries of both kinds on held-out folds, and missing a question that needed retrieval was the more common error.

\textbf{Features:} three groups. (i) Question surface: length, the interrogative word, negation, temporal and numerical cues, entity count, a misconception-phrasing pattern, and three style markers (adversarial, opinion-seeking, open-ended). (ii) The top retrieved passage: its similarity, entity overlap with the question, whether the question expects a year or number and whether the passage contains one, and whether the question's main entity appears in Wikidata. (iii) Retrieval quality and self-knowledge: statistics of the top-five retrieval similarities, of their cross-encoder rerank scores and of pairwise entailment among them, plus a self-knowledge probe $p_{\text{IK}}$ in the style of Kadavath et al.~\cite{kadavath2022language}, which reads the base generator's last hidden state and predicts whether it knows the answer without retrieval. The probe was trained on $11{,}980$ questions (out-of-fold AUROC $0.771$). It carries the largest weight in the classifier ($-1.77$): when the generator knows the answer, the classifier leans strongly toward Tier 2. Six features (page views and the four year and number features of the passage) were constant in the training data and received zero weight, and the top-passage similarity took the opposite sign to the one expected.

\textbf{Training data:} $3{,}208$ questions from a pool disjoint from every other fold, drawn from eight benchmarks (the five of the evaluation panel plus Natural Questions, HaluEval-QA and OpenBookQA, included only to widen the range of cases). Each question was answered both with and without retrieval by the same backbone. It was labelled \emph{retrieval helps} ($1{,}846$ questions) if only the retrieval answer was correct and \emph{retrieval hurts} ($1{,}362$) if only the zero-shot answer was correct; questions where both or neither were correct were dropped, because there the routing choice does not matter. TruthfulQA answers were judged by the language-model judge used throughout the thesis. The three extra benchmarks feed nothing else in the architecture.

\textbf{Fitting and accuracy:} nested five-fold stratified cross-validation chose the regularisation strength ($C = 0.5$) in the inner folds and measured accuracy in the outer folds. The pooled out-of-fold AUROC is $0.850$ (standard deviation across folds $0.008$). At the chosen threshold the classifier routes $77.2$ percent of held-out questions correctly ($2{,}477$ of $3{,}208$): $1{,}554$ of $1{,}846$ retrieval-helps questions (recall $0.842$) and $923$ of $1{,}362$ retrieval-hurts questions (specificity $0.678$). Per-benchmark AUROC on the evaluation panel is FEVER $0.718$, TriviaQA $0.936$, CommonsenseQA $0.728$, TruthfulQA $0.732$ and StrategyQA $0.852$ (the last on only $43$ questions). TruthfulQA, the benchmark that motivated the classifier, falls just short of the $0.75$ target set for it in advance.

\textbf{Frozen at deployment:} the classifier was fitted once before the main run and never refitted. Its inputs come from parts of the system that do not change across cycles (the question, the passage index, the base generator with the adapter disabled, and the fixed reranker), so a frozen classifier keeps its routing auditable. This light, classifier-only design agrees with Marina et al.~\cite{moskvoretskii2025llm}, who find that an external classifier matches language-model-based selective retrievers; CAEM differs in reusing the reranker and entailment judge that the verifier already loads. The effect of the classifier is measured in \Cref{sec:ruc-ablation}.

\begin{table}[!htbp]
\centering
\small
\begin{tabular}{@{}p{1.2cm}p{1.3cm}>{\raggedright\arraybackslash}p{5.7cm}>{\raggedright\arraybackslash}p{4.2cm}@{}}
\toprule
\textbf{Order} & \textbf{Tier} & \textbf{Condition} & \textbf{Output path} \\
\midrule
1 & Tier 3 / 2 & $\upre < 0.38$ (safety override): Tier 3 if $p_{\text{RAG}} \ge 0.375$, else Tier 2 & Retrieval-augmented or zero-shot generation \\
2 & Tier 1 & $\mathcal{S} = 0.7\,s + 0.3\,\ustored \ge 0.90$ & Serve neighbour's stored answer \\
3 & Tier 2 & $s \ge 0.75$ & Zero-shot generation \\
4 & Tier 3 / 2 & otherwise (fall-through): Tier 3 if $p_{\text{RAG}} \ge 0.375$, else Tier 2 & Retrieval-augmented or zero-shot generation \\
\bottomrule
\end{tabular}
\caption[Dispatch rules of the router and the retrieval-utility classifier.]{Dispatch rules of the router and the retrieval-utility classifier, in order of evaluation. The classifier is consulted only in rows 1 and 4, the two cases that would otherwise always retrieve.}
\label{tab:dispatch-rules}
\label{tab:ruc-dispatch}
\end{table}

\begin{algorithm}[!htbp]
\caption[Retrieval-utility classifier (Stage~3b).]{Retrieval-utility classifier (Stage~3b).}
\label{alg:ruc}
\begin{algorithmic}[1]
\Require question $q$, router decision $d \in \{\text{Tier 1}, \text{Tier 2}, \text{Tier 3}\}$, frozen scaler $\Sigma$, frozen weights $w$, threshold $\tau_{\text{RUC}} = 0.375$
\If{$d \in \{\text{Tier 1}, \text{Tier 2}\}$}
    \State \Return $d$ \Comment{the classifier does not gate these outcomes}
\EndIf
\State $\mathbf{x} \gets$ [question-surface features, top-passage features, retrieval-quality features, $p_{\text{IK}}$]
\State $p_{\text{RAG}} \gets \sigma\!\left(w^{\top}\, \Sigma(\mathbf{x})\right)$
\State \Return Tier 3 \textbf{if} $p_{\text{RAG}} \ge \tau_{\text{RUC}}$ \textbf{else} Tier 2
\end{algorithmic}
\end{algorithm}

\subsection{Verifier ensemble}
\label{sec:verifier}

The verifier turns a generated answer into nine signals in four families (\Cref{tab:verifier-signals}, \Cref{fig:verifier-flow}, \Cref{alg:verifier}). The families follow the confidence-estimation survey of Geng et al.~\cite{geng2024survey}, with one addition, a separate relevance family, for answers that are supported by a passage but do not address the question. Each family reads a different source of evidence, so that one failure mode does not hide behind another. Two practical details matter. First, the grounding signals check the answer rewritten as a statement: a FEVER \emph{refutes} answer is checked as the negated claim, and a multiple-choice letter is replaced by the text of the chosen option. Second, five shared computations produce all nine signals, so the cost is paid once per answer.

% Verifier data flow (Stage 5): three shared computations feed the nine signals.
\begin{figure}[!htbp]
\centering
\resizebox{\textwidth}{!}{%
\begin{tikzpicture}[
  font=\small,
  inp/.style={rectangle, rounded corners=2pt, draw=black!55, fill=black!4, line width=0.9pt, align=center, inner sep=5pt, text width=2.4cm},
  gen/.style={rectangle, rounded corners=4pt, draw=teal!65!black, fill=teal!8, line width=0.9pt, align=center, inner sep=3pt, text width=4.6cm, font=\footnotesize},
  sig/.style={rectangle, rounded corners=3pt, line width=0.8pt, align=center, inner sep=3pt, text width=3.0cm, font=\footnotesize, minimum height=0.62cm},
  flow/.style={-{Stealth[length=2.2mm]}, line width=0.8pt, draw=black!65},
]
  \node[inp] (qa) at (0,0) {question $+$\\served answer};
  \node[gen] (tf)   at (4.6,3.2) {\textbf{teacher-forced pass}\\{\footnotesize probability of the answer's tokens}};
  \node[gen] (drop) at (4.6,1.6) {\textbf{$K{=}5$ samples, dropout on}\\{\footnotesize do they reach the same answer?}};
  \node[gen] (ch)   at (4.6,0.0) {\textbf{$M{=}3$ sampled chains}\\{\footnotesize temperature $0.7$}};
  \node[gen] (ret)  at (4.6,-1.6) {\textbf{retrieve $20$, rerank to $3$}\\{\footnotesize Wikipedia passages}};
  \node[gen] (ce)   at (4.6,-3.2) {\textbf{cross-encoder}\\{\footnotesize question against answer}};
  \node[sig, draw=blue!60!black, fill=blue!6]   (utok) at (9.6,3.2) {$u_{\text{token}}$};
  \node[sig, draw=blue!60!black, fill=blue!6]   (udro) at (9.6,1.6) {$u_{\text{dropout}}$, then $u_{\text{internal}}$};
  \node[sig, draw=violet!60!black, fill=violet!6] (sa)  at (9.6,0.0) {$s_{\text{avg}}$,\ $p_{\text{entail}}$};
  \node[sig, draw=green!50!black, fill=green!7] (pg)   at (9.6,-1.6) {$p_{\text{ground,max}}$, $p_{\text{ground,mean}}$, $p_{\text{ground,atomic}}$};
  \node[sig, draw=orange!70!black, fill=orange!8] (qar) at (9.6,-3.2) {$q_{\text{match}}$};
  \node[inp, text width=2.3cm] (out) at (13.4,0) {nine signals\\to the calibrated composite};
  \foreach \g in {tf,drop,ch,ret,ce} \draw[flow] (qa.east) -- (\g.west);
  \draw[flow] (tf) -- (utok); \draw[flow] (drop) -- (udro); \draw[flow] (ch) -- (sa); \draw[flow] (ret) -- (pg); \draw[flow] (ce) -- (qar);
  \foreach \s in {utok,udro,sa,pg,qar} \draw[flow] (\s.east) -- (out.west);
  \node[font=\scriptsize, text=black!60, align=left, anchor=north west] at (6.6,-3.9) {colours: internal calibration (blue), sample-set agreement (violet),\\external grounding (green), question--answer match (orange)};
\end{tikzpicture}}
\caption[Where the nine verifier signals come from.]{Where the nine verifier signals come from. Five shared computations produce all nine signals; the grounding scores check the answer rewritten as a statement against the reranked passages.}
\label{fig:verifier-flow}
\end{figure}


\begin{table}[!htbp]
\centering
\footnotesize
\begin{tabular}{@{}>{\raggedright\arraybackslash}p{2.0cm}>{\raggedright\arraybackslash}p{2.2cm}>{\raggedright\arraybackslash}p{5.2cm}>{\raggedright\arraybackslash}p{4.0cm}@{}}
\toprule
\textbf{Signal} & \textbf{Family} & \textbf{What it measures} & \textbf{Failure it catches} \\
\midrule
$u_{\text{token}}$ & Internal & Geometric mean probability of the answer's tokens under the generator & Low-confidence generation \\
$u_{\text{dropout}}$ & Internal & Share of $5$ dropout-active samples whose answer differs from the most common one & Unstable answers \\
$u_{\text{internal}}$ & Internal & $0.5\,u_{\text{token}} + 0.5\,(1-u_{\text{dropout}})$ & Summary of internal confidence \\
$s_{\text{avg}}$ & Sample-set & Mean pairwise cosine similarity of $3$ sampled chains (temperature $0.7$) & Disagreement across chains \\
$p_{\text{entail}}$ & Sample-set & Mean entailment of the answer by each of the $3$ chains & Answer not supported by the model's own reasoning \\
$p_{\text{ground,max}}$ & Grounding & Entailment of the answer by the best of the top three reranked passages & No supporting passage \\
$p_{\text{ground,mean}}$ & Grounding & Mean entailment over the top three passages & Weak overall support \\
$p_{\text{ground,atomic}}$ & Grounding & Weakest support among the answer's atomic facts (answers of $12$ tokens or more) & One fabricated detail in a long answer \\
$q_{\text{a,rel}}$ & Relevance & Cross-encoder relevance of the answer to the question & Off-topic answer \\
\bottomrule
\end{tabular}
\caption[The nine verifier signals.]{The nine verifier signals, grouped by family.}
\label{tab:verifier-signals}
\end{table}

\begin{algorithm}[!htbp]
\caption[Verifier ensemble (Stage~5).]{Verifier ensemble (Stage~5).}
\label{alg:verifier}
\begin{algorithmic}[1]
\Require question $q$, answer $a$, generator $G$, entailment judge $J$, cross-encoder $X$
\Statex \textbf{Internal calibration}
\State $u_{\text{token}} \gets$ geometric mean probability of the tokens of $a$ under $G$ (teacher-forced)
\State draw $5$ answers from $G$ with dropout active; $u_{\text{dropout}} \gets 1 - (\text{count of the most common answer})/5$
\State $u_{\text{internal}} \gets 0.5\,u_{\text{token}} + 0.5\,(1-u_{\text{dropout}})$
\Statex \textbf{Sample-set agreement}
\State $c_{1}, c_{2}, c_{3} \gets$ chains sampled from $G$ at temperature $0.7$
\State $s_{\text{avg}} \gets$ mean pairwise cosine similarity of the chain embeddings
\State $p_{\text{entail}} \gets \tfrac{1}{3}\sum_{m} J(c_{m} \to a)$
\Statex \textbf{External grounding}
\State $P_{1}, P_{2}, P_{3} \gets$ top $3$ of the $20$ retrieved passages after reranking by $X$
\State $h \gets$ $a$ rewritten as a statement \Comment{negated claim for \emph{refutes}; option text for a letter}
\State $p_{\text{ground,max}} \gets \max_{k} J(P_{k} \to h)$; \quad $p_{\text{ground,mean}} \gets \mathrm{mean}_{k}\, J(P_{k} \to h)$
\If{$a$ has at least $12$ tokens and yields atomic facts $f_{1}, \dots, f_{r}$}
    \State $p_{\text{ground,atomic}} \gets \min_{i} \max_{k} J(P_{k} \to f_{i})$
\Else
    \State $p_{\text{ground,atomic}} \gets p_{\text{ground,mean}}$
\EndIf
\Statex \textbf{Relevance}
\State $q_{\text{a,rel}} \gets X(q, a)$
\State \Return the nine signals
\end{algorithmic}
\end{algorithm}

\subsubsection{Internal calibration signals}
\label{sec:verifier-internal}

These signals read the generator's own state. The token signal is the most local form of self-doubt. The dropout signal follows the Bayesian reading of dropout of Gal and Ghahramani~\cite{pmlr-v48-gal16}: with dropout switched on, the generator draws five answers, and $u_{\text{dropout}}$ is the share that disagree with the most common one, so $0$ means all five agree. The summary $u_{\text{internal}}$ is the input to the confabulation early-exit rule (\Cref{sec:decision}), because it isolates internal confidence from external evidence. This family is the one most affected by fine-tuning, which is one reason stored entries are re-verified at every cycle boundary. A hidden-state eigenvalue score \cite{chen2024inside} was considered as a replacement for the dropout signal but not adopted, for cost.

\subsubsection{Sample-set agreement signals}
\label{sec:verifier-sample-set}

These signals ask whether the generator gives the same answer when it reasons again. The generator samples three chains; $s_{\text{avg}}$ is their mean pairwise similarity in the embedding space, an embedding-based form of the self-consistency of Wang et al.~\cite{wang2023selfconsistency}, and $p_{\text{entail}}$ asks the judge whether each chain supports the served answer. Both use the same three chains.

Two dispersion signals were dropped before the main run. The semantic entropy of Farquhar et al.~\cite{farquhar2024semantic} (after Kuhn et al.~\cite{kuhn2023semantic}), computed over $10$ samples at temperature $1.0$, received a learned weight of about $-0.0005$ in the ten-signal cycle-zero fit, so it added nothing beyond the two signals above, while its $10$ samples were the largest single cost of the verifier. The kernel language entropy of Nikitin et al.~\cite{nikitin2024kle} needs the same samples and was not adopted either.

\subsubsection{External grounding signals}
\label{sec:verifier-grounding}

These signals check the answer, rewritten as a statement, against the three best passages. The maximum catches the case where one passage clearly supports the answer; the mean rewards support spread over several passages; the atomic score splits a long answer into atomic facts, in the FActScore manner of Min et al.~\cite{min-etal-2023-factscore}, and takes the weakest-supported fact, so that one fabricated detail is not averaged away. The atomic score applies only to answers of at least $12$ tokens and otherwise equals the mean score.

\textbf{Realised run:} in the main run the atomic score equalled the mean score on all $14{,}951$ scored samples. The length gate skipped $99.0$ percent of answers, and for the remaining $1.0$ percent the decomposer returned no atomic facts. This is a mismatch of scope rather than a tuning fault: the signal follows FActScore, which targets long-form text, while the answers on this panel are labels, option letters and short entities, each already a single fact, so a lower gate would only reproduce the mean score. The signal is part of the design but contributed no separate information in this run; applying the same check to the reasoning chain, where several facts are asserted, is left as future work (\Cref{ch:conclusion}).

\subsubsection{Question-answer relevance}
\label{sec:verifier-qa-rel}

The relevance signal scores the answer against the question with the cross-encoder. It forms its own family because an off-topic answer can be fully supported by a passage that is itself off-topic, which the grounding signals cannot detect. In the cycle-zero composite it received the largest weight of the nine signals ($0.491$, \Cref{sec:composite-boost}).

\subsection{Calibrated probability composite}
\label{sec:composite}

The composite turns the nine signals into one probability of correctness, $\ustored$, in three steps (\Cref{fig:composite-pipeline}, \Cref{alg:composite}).
\begin{enumerate}
    \item \textbf{One curve per signal:} isotonic regression maps each raw signal to the probability that the answer is correct.
    \item \textbf{Per benchmark, with shrinkage:} the curves are also fitted on each training benchmark and pulled part of the way toward the pooled curve, so that benchmarks with little data do not overfit.
    \item \textbf{Weighted sum:} a regularised logistic regression combines the nine calibrated probabilities, in log-odds, into $\ustored$.
\end{enumerate}
All three steps are fitted on the calibration fold ($500$ labelled questions from each training benchmark) and refitted at every cycle boundary on the same fold re-scored with the new adapter.

% Calibrated probability composite (Stage 5 to 6).
\begin{figure}[!htbp]
\centering
\resizebox{\textwidth}{!}{%
\begin{tikzpicture}[
  font=\small,
  box/.style={rectangle, rounded corners=4pt, line width=0.9pt, align=center, inner sep=5pt, text width=3.1cm, minimum height=1.55cm},
  flow/.style={-{Stealth[length=2.2mm]}, line width=0.9pt, draw=black!65},
]
  \node[box, draw=black!55, fill=black!4, text width=2.2cm] (s) at (0,0) {\textbf{nine raw}\\\textbf{signals}};
  \node[box, draw=blue!60!black, fill=blue!6] (iso) at (3.6,0) {\textbf{1. one curve per signal}\\{\footnotesize raw value $\to$ probability the answer is correct (isotonic)}};
  \node[box, draw=violet!60!black, fill=violet!6] (shr) at (7.7,0) {\textbf{2. per benchmark}\\{\footnotesize each benchmark's curve pulled $40\%$ toward the pooled curve}};
  \node[box, draw=green!50!black, fill=green!7] (bst) at (11.8,0) {\textbf{3. weighted sum}\\{\footnotesize log-odds combined by a regularised logistic regression}};
  \node[box, draw=orange!70!black, fill=orange!8, text width=1.9cm] (u) at (15.2,0) {$\ustored$\\{\footnotesize one probability}};
  \draw[flow] (s) -- (iso); \draw[flow] (iso) -- (shr); \draw[flow] (shr) -- (bst); \draw[flow] (bst) -- (u);
\end{tikzpicture}}
\caption[The calibrated probability composite in three steps.]{The calibrated probability composite in three steps. All three are refitted at every cycle boundary on the re-scored calibration fold.}
\label{fig:composite-pipeline}
\end{figure}


\begin{algorithm}[!htbp]
\caption[Calibrated probability composite (fit and predict).]{Calibrated probability composite (fit and predict).}
\label{alg:composite}
\begin{algorithmic}[1]
\Statex \textbf{Fit stage:} calibration fold split by benchmark, $\mathcal{D}_{b}$ for $b \in \mathcal{B}$, each a set of signal vectors $x \in \mathbb{R}^{9}$ with exact-match labels $y \in \{0,1\}$; pooled fold $\mathcal{D} = \bigcup_{b}\mathcal{D}_{b}$.
\For{each signal $j$}
    \State direction $\gets$ sign of the Pearson correlation of $x_{j}$ with $y$ on $\mathcal{D}$
    \State fit an isotonic curve $\mathrm{iso}_{j}$ in that direction on $\mathcal{D}$, with outputs clipped to $[0.02, 0.98]$
\EndFor
\For{each benchmark $b$ and signal $j$}
    \State fit $\mathrm{iso}^{b}_{j}$ on $\mathcal{D}_{b}$; at each knot $x$ set $\mathrm{iso}^{b}_{j}(x) \gets 0.6\,\mathrm{iso}^{b}_{j}(x) + 0.4\,\mathrm{iso}_{j}(x)$
\EndFor
\State $L_{ij} \gets \mathrm{logit}\big(\mathrm{iso}_{j}(x_{ij})\big)$ on $\mathcal{D}$; fit logistic regression ($\ell_{2}$ penalty, $C = 0.01$) of $y$ on $L$ to get weights $\mathbf{w}$ and intercept $b_{0}$
\State repeat the logistic fit on each $\mathcal{D}_{b}$ with the shrunk curves to get $(\mathbf{w}_{b}, b_{0,b})$
\Statex \textbf{Predict stage:} signal vector $x$ and benchmark tag $b^{\star}$
\State use the curves and weights of $b^{\star}$ if $b^{\star} \in \mathcal{B}$, otherwise the pooled ones
\State \Return $\ustored \gets \sigma\!\Big(b_{0} + \sum_{j} w_{j}\, \mathrm{logit}\big(\mathrm{iso}_{j}(x_{j})\big)\Big)$
\end{algorithmic}
\end{algorithm}

\subsubsection{Per-signal isotonic regression}
\label{sec:composite-isotonic}

Isotonic regression fits any monotone curve, not only a logistic one, which matters because some signals saturate near the top of their range; and it needs less data than histogram binning because neighbouring bins share information. Niculescu-Mizil and Caruana~\cite{niculescu2005predicting} found isotonic regression the stronger post-hoc calibrator when enough data are available, but prone to overfitting below roughly $200$ to $1{,}000$ calibration points, where Platt scaling does better. The per-benchmark folds here hold $500$ points each, which is one reason the per-benchmark curves are shrunk toward the pooled curve (\Cref{sec:composite-shrinkage}). The direction of each curve is taken from the sign of the signal's correlation with correctness, so a signal where high means unsafe (such as $u_{\text{dropout}}$) needs no special handling. Probabilities are clipped to $[0.02, 0.98]$ so that no single signal can dominate the log-odds sum, and a signal with fewer than $50$ calibration samples is left out.

\subsubsection{Per-benchmark fit with shrinkage prior}
\label{sec:composite-shrinkage}

Each training benchmark gets its own curves, blended at every knot with the pooled curve: $60$ percent benchmark curve, $40$ percent pooled curve, a convex blend in the spirit of James--Stein shrinkage. Each benchmark also gets its own logistic weights. The blend strength was chosen to balance two effects visible in the cycle-zero sweep: on the held-out validation fold, fitting each benchmark more freely ($\alpha_{\text{shrink}}$ closer to $1$) stored more answers but separated correct from incorrect answers less well (pooled Cohen's $d$ of $0.85$ at $\alpha_{\text{shrink}} = 0$, $0.72$ at $0.6$ and $0.66$ at $1$; \Cref{ch:evaluation}). A query from a benchmark outside the training panel uses the pooled curves and weights.

\subsubsection{Boost layer with L2 regularisation}
\label{sec:composite-boost}

The boost layer is a logistic regression on the nine log-odds, following the multi-scoring approach of Valentin et al.~\cite{valentin2024hallucination}. A learned combination is used rather than an equal-weight sum because several signals overlap (the two sample-set signals, and the three grounding signals), and the regression can discount the overlap. An $\ell_{2}$ penalty is used rather than $\ell_{1}$ so that weights shrink smoothly without any signal being dropped, and its strength is strong ($C = 0.01$, chosen in the cycle-zero sweep of \Cref{ch:evaluation}). In the cycle-zero pooled fit every weight is positive; the largest are question-answer relevance ($0.491$), best-passage grounding ($0.379$) and internal confidence ($0.334$).

\subsection{Fixed-threshold storage gate}
\label{sec:conformal-gate}

The storage gate applies two fixed cut-points to $\ustored$: an answer is stored if $\ustored \ge 0.60$ and deferred if $0.45 \le \ustored < 0.60$; below $0.45$, the decision rule of \Cref{sec:decision} chooses between abstaining and discarding. The cut-points were fixed before the main run and never changed.

An earlier design set the cut-points by split-conformal prediction at fixed precision targets and refitted them every cycle. It was replaced because conformal guarantees assume the calibration and evaluation data come from the same distribution, and on this panel they did not: precision on the evaluation fold fell well below the conformal targets, and on the weakest benchmarks the per-benchmark fits rejected almost everything. A fixed cut-point on a calibrated probability, in the spirit of the post-hoc calibration of Niculescu-Mizil and Caruana~\cite{niculescu2005predicting} and the selective classification of Geifman and El-Yaniv~\cite{geifman2017selective}, makes no coverage promise; its precision is simply measured on every fold. It still assumes that calibration and test data are alike, so under distribution shift the realised precision can fall, as Ovadia et al.~\cite{ovadia2019trust} and Kamath et al.~\cite{kamath2020selective} show; TruthfulQA in \Cref{ch:evaluation} is such a case.

\subsubsection{Storage threshold}
\label{sec:conformal-store}

The storage cut-point $\tau_{\text{store}} = 0.60$ was chosen on the cycle-zero calibration fold before the main run, against the stricter alternative $0.65$, which stored about half as many answers (\Cref{ch:evaluation} reports the sweep). The storage class fills the memory and supplies the training pool, so it is guarded on both sides: training requires the stricter $\ustored \ge 0.70$, and the retroactive pass removes entries that later fall below $0.50$.

\subsubsection{Deferral threshold}
\label{sec:conformal-defer}

Answers with $0.45 \le \ustored < 0.60$ go to the deferred buffer, not to memory. They are re-scored at the next cycle boundaries (\Cref{sec:deferred-reconsider}) and either promoted or dropped, so a looser cut-point here costs little: deferred answers never reach the training pool unless they are later promoted. The configuration refuses a deferral cut-point above the storage cut-point.

\subsubsection{Cycle-boundary recalibration}
\label{sec:conformal-cycle-recalib}

The cut-points ($0.60$ and $0.45$), the abstention floor ($0.20$), the retroactive prune threshold ($0.50$) and the training threshold ($0.70$) stay fixed for the whole run. What changes at each cycle boundary is the calibration: after a committed fine-tune, the same calibration fold is re-scored with the new adapter, the pre-routing temperatures and safety floors are refitted, and the composite is refitted. The verifier is then reloaded with the new composite before the retroactive pass, so stored entries are re-scored with curves that match the current adapter. Without this order, entries whose raw signals had shifted after fine-tuning would be read through stale curves and wrongly pruned.

Keeping the cut-points fixed while the composite adapts is what lets Hypothesis~1 (\Cref{sec:hypotheses}) test the verifier rather than a moving operating point. A deployed system would run the same refit; the only difference is that its cycle would end on a calendar or accumulation trigger, with a labelled sample of recent admissions as the calibration fold.

\subsection{Four-outcome storage decision}
\label{sec:decision}

The decision rule maps $\ustored$ and two guard signals to one of four outcomes (\Cref{alg:decision}). Four outcomes are used instead of store-or-discard for two reasons. First, an answer just below the storage cut-point is not the same as one near zero: deferring it lets the next cycle's adapter settle it. Second, \emph{I have nothing reliable to say} is a different situation from \emph{I have an answer but cannot ground it}: abstention makes the first an explicit, user-facing refusal, in the spirit of the abstention framework of Yadkori et al.~\cite{yadkori2024abstention}.

\textbf{Confabulation early exit:} checked first, at query time only. An answer the generator is internally confident about but that no passage supports is discarded:
\begin{equation}
\big(u_{\text{internal}} \ge 0.70\big) \;\wedge\; \big(p_{\text{ground,max}} \le 0.20\big)
\;\Rightarrow\; \textsc{Discard}.
\label{eq:early-exit-rule}
\end{equation}
This catches the most damaging case, a confident answer with no evidence, which the composite alone might score as merely borderline.

\textbf{Thresholds:} otherwise the answer is stored if $\ustored \ge 0.60$ and deferred if $\ustored \ge 0.45$. Below $0.45$, the system abstains when the best passage barely supports the answer ($p_{\text{ground,max}} < 0.20$) and discards otherwise. A discarded answer is not stored; it is still shown to the user, with its probability and a warning that the evidence conflicts. An abstention replaces the answer with an insufficient-evidence message.

\begin{algorithm}[!htbp]
\caption[Four-outcome storage decision (Stage~6).]{Four-outcome storage decision (Stage~6).}
\label{alg:decision}
\begin{algorithmic}[1]
\Require $\ustored$, $u_{\text{internal}}$, $p_{\text{ground,max}}$; cut-points $\tau_{\text{store}} = 0.60$, $\tau_{\text{defer}} = 0.45$; early-exit thresholds $\eta_{u} = 0.70$, $\eta_{g} = 0.20$; abstention floor $\phi_{\text{pg}} = 0.20$
\If{query time \textbf{and} $u_{\text{internal}} \ge \eta_{u}$ \textbf{and} $p_{\text{ground,max}} \le \eta_{g}$}
    \State \Return \textsc{Discard} \Comment{confabulation early exit}
\EndIf
\If{$\ustored \ge \tau_{\text{store}}$} \State \Return \textsc{Store} \EndIf
\If{$\ustored \ge \tau_{\text{defer}}$} \State \Return \textsc{Defer} \EndIf
\If{$p_{\text{ground,max}} < \phi_{\text{pg}}$} \State \Return \textsc{Abstain} \EndIf
\State \Return \textsc{Discard}
\end{algorithmic}
\end{algorithm}

\subsection{Cycle-boundary maintenance}
\label{sec:cycle-boundary}

The cycle-boundary steps are where the two kinds of statelessness of \Cref{sec:arch-opportunity} are addressed. Fine-tuning on verified answers, many produced with retrieval, is meant to let the generator answer later questions of the same kind without retrieval. The memory is meant to serve repeated questions at the cost of a lookup, with re-verification keeping each stored score in line with the current model. The deferred pass gives borderline answers a second chance under the new adapter. The steps run in the order of \Cref{fig:cycle-boundary-order}; on a cycle whose fine-tune is rolled back, all steps after the retention guard are skipped.

\subsubsection{Retroactive re-verification}
\label{sec:retroverify}

The retroactive pass re-scores every stored episode with the current adapter and the refitted composite (\Cref{alg:retroverify}). An entry is removed if its reasoning chain is caught by the repetition filter or if its new $\ustored$ is below $0.50$. Otherwise all nine stored signals and $\ustored$ are overwritten with the new values, which may be lower or higher than before. The prune threshold sits below the deferral cut-point, so an entry that is no longer storage-quality but still plausible is kept rather than discarded. The pass removed $56$, $271$, $352$ and $382$ entries at cycles one, two, three and five (it was skipped at cycle four), that is $13$, $9$, $7$ and $4$ percent of the entries it checked.

\begin{algorithm}[!htbp]
\caption[Retroactive re-verification pass (Stage~8).]{Retroactive re-verification pass (Stage~8).}
\label{alg:retroverify}
\begin{algorithmic}[1]
\Require memory $\mathcal{M}$, verifier with the current adapter and refitted composite, prune threshold $\tau_{\text{retro}} = 0.50$
\For{each entry $e \in \mathcal{M}$}
    \If{the reasoning chain of $e$ is caught by the repetition filter}
        \State remove $e$; \textbf{continue}
    \EndIf
    \State $\mathbf{x} \gets \mathrm{verify}(e.\text{question}, e.\text{answer})$ \Comment{early exit disabled}
    \If{$\mathbf{x}.\ustored < \tau_{\text{retro}}$}
        \State remove $e$
    \Else
        \State overwrite the signals and $\ustored$ of $e$ with $\mathbf{x}$ \Comment{may lower or raise $\ustored$}
    \EndIf
\EndFor
\end{algorithmic}
\end{algorithm}

\subsubsection{Deferred-entry reconsideration}
\label{sec:deferred-reconsider}

Each deferred answer is re-scored at the next cycle boundaries (\Cref{alg:deferred-reconsider}). It is promoted into memory if the verifier now decides to store it ($\ustored \ge 0.60$); it is dropped once it has been reconsidered twice without promotion; otherwise it waits for the next pass. The buffer holds at most $10{,}000$ answers and evicts the oldest when full. The passes promoted $131$, $220$ and $299$ answers at cycles two, three and five; the first pass, at cycle one, found an empty buffer, and the pass was skipped at cycle four. The memory consolidation of \Cref{sec:memory-store} runs straight after this pass.

\begin{algorithm}[!htbp]
\caption[Deferred-entry reconsideration pass (Stage~8).]{Deferred-entry reconsideration pass (Stage~8).}
\label{alg:deferred-reconsider}
\begin{algorithmic}[1]
\Require deferred buffer $\mathcal{B}$, current verifier, $\tau_{\text{store}} = 0.60$, time-to-live $\tau_{\text{ttl}} = 2$ passes
\For{each entry $e \in \mathcal{B}$}
    \State $\mathbf{x} \gets \mathrm{verify}(e)$
    \If{$\mathbf{x}.\text{decision} = \textsc{Store}$ \textbf{and} $\mathbf{x}.\ustored \ge \tau_{\text{store}}$}
        \State move $e$ into the memory
    \ElsIf{$e.\text{age} + 1 \ge \tau_{\text{ttl}}$}
        \State drop $e$
    \Else
        \State $e.\text{age} \gets e.\text{age} + 1$ \Comment{wait for the next pass}
    \EndIf
\EndFor
\end{algorithmic}
\end{algorithm}

\subsubsection{Self-improvement loop}
\label{sec:sil}

The self-improvement loop fine-tunes the generator on its own verified answers. It belongs to the family of methods that train a model on its filtered outputs \cite{NEURIPS2022_639a9a17,hosseini2024vstar,gulcehre2023rest}, and is related to distilling a teacher's filtered reasoning into a student \cite{ho2023large}. Its intended role is to turn verified retrieval-based reasoning into knowledge the generator holds without retrieval. Song et al.~\cite{song2024gap} find that such loops need a positive gap between how well a model generates and how well it can verify, and saturate within a few rounds; Fu et al.~\cite{fu2025collapse} show in linear regression that filtering through an external verifier prevents collapse but pulls the model toward the verifier's own knowledge. CAEM's external verifier and fixed cut-points are meant to keep the training pool above the purification condition of \Cref{thm:bayes-purity}. The loop has five steps (\Cref{alg:sil}).

\textbf{Select:} take stored episodes with $\ustored \ge 0.70$ from the three training benchmarks only, so transfer benchmarks never reach the gradient. The training target is the stored reasoning chain and answer; most chains were written under the retrieval prompt.

\textbf{Rebalance:} reweight the pool so that no benchmark dominates: per-benchmark counts are mixed with a softmax at temperature $2$, upsampling is capped at three times a benchmark's count, each benchmark contributes at least $50$ examples, a benchmark with no verified episodes is given $100$ gold-labelled examples from its training split, and no benchmark exceeds $40$ percent of the pool. This addresses the failure of an earlier run, in which one benchmark grew to the whole training pool.

\textbf{Train:} train a low-rank adapter (LoRA, rank $32$, scaling $64$, dropout $0.05$) on all seven linear projections of every transformer block, following the all-layer recommendation of Biderman et al.~\cite{biderman2024lora}; the backbone stays frozen. Training runs for $3$ epochs with 32-bit AdamW at a learning rate of $2 \times 10^{-4}$, an effective batch of $16$ and gradient checkpointing. The loss is cross-entropy alone:
\begin{equation}
\mathcal{L}(\boldsymbol{\theta}) \;=\; \mathcal{L}_{\mathrm{CE}}(\boldsymbol{\theta}) \;+\; \tfrac{\lambda_{\text{anc}}}{2}\, \big\| \boldsymbol{\theta} - \boldsymbol{\theta}^{(t-1)} \big\|_{2}^{2}, \qquad \lambda_{\text{anc}} = 0,
\label{eq:sil-anchor}
\end{equation}
where $\boldsymbol{\theta}$ are the adapter parameters and $\boldsymbol{\theta}^{(t-1)}$ the previous cycle's adapter. The anchor term, a uniform-weight simplification of the elastic weight consolidation of Kirkpatrick et al.~\cite{doi:10.1073/pnas.1611835114}, belongs to the full fine-tuning path used in an earlier run and is switched off here: the frozen backbone and the small adapter already limit how far each cycle can move the model, and in the earlier run the anchor outweighed the learning signal from the small training pool.

\textbf{Check retention:} before the first cycle, record the accuracy of the untouched model on three held-out probes of $200$ questions each: MMLU (general knowledge), the TriviaQA test fold (open-text facts) and the CommonsenseQA test fold (commonsense multiple choice). After each fine-tune, compute each probe's accuracy as a ratio of that starting value (\Cref{eq:nfr-retention-floor}).

\textbf{Commit or roll back:} keep the new adapter only if every ratio is at least $0.93$; otherwise restore the previous adapter. The memory is kept either way. Requiring every probe, not just the average, to pass is deliberate: an earlier rule that watched only MMLU missed forgetting on narrower tasks. In the main run the guard fired once, at cycle four, when the TriviaQA test ratio fell to $0.886$.

\begin{algorithm}[!htbp]
\caption[Self-improvement loop with retention guard (Stage~8).]{Self-improvement loop with retention guard (Stage~8).}
\label{alg:sil}
\begin{algorithmic}[1]
\Require memory $\mathcal{M}$, previous adapter $\boldsymbol{\theta}^{(t-1)}$, probes $\mathcal{P}$ with starting accuracies $A_{0}^{p}$, $\tau_{\text{train}} = 0.70$, $\rho_{\min} = 0.93$
\State $\mathcal{T} \gets$ rebalance$\big(\{e \in \mathcal{M} : e.\ustored \ge \tau_{\text{train}},\ e \text{ from a training benchmark}\}\big)$
\State $\boldsymbol{\theta}^{\star} \gets$ train the adapter on $\mathcal{T}$ with $\mathcal{L}_{\mathrm{CE}}$, starting from $\boldsymbol{\theta}^{(t-1)}$
\If{$\min_{p \in \mathcal{P}} A^{p}(\boldsymbol{\theta}^{\star}) / A_{0}^{p} \ge \rho_{\min}$}
    \State \Return $\boldsymbol{\theta}^{\star}$ \Comment{commit}
\Else
    \State \Return $\boldsymbol{\theta}^{(t-1)}$ \Comment{roll back; memory kept; later cycle-boundary steps skipped}
\EndIf
\end{algorithmic}
\end{algorithm}

Keeping the memory through a rollback is what makes the loop safe to repeat: a failed cycle costs its fine-tuning time but not the answers it verified. This asymmetry underlies the memory part of \Cref{cor:stochastic-equilibrium}.

\section{Theoretical analysis}
\label{sec:theory}

This section states what can be said formally about the stored pool, and is explicit about the strength of each statement (\Cref{fig:theory-map}). Only two results are proven: a floor on the purity of the stored pool (\Cref{thm:purity}) and the condition under which filtering makes the pool purer than the generator's raw output (\Cref{thm:bayes-purity}). Four remarks are exact calculations on simplified models; they show what the design aims at, not what the system guarantees. Five heuristics describe how the pool is expected to evolve across cycles; they rest on assumptions the real system does not fully meet, so they come with argument sketches rather than proofs, and the evaluation checks them only for consistency.

All statements are about the stored pool. None bounds the accuracy of the fine-tuned generator on new questions, so a fall in held-out accuracy is outside their scope rather than a contradiction of them.

% Status map of the formal statements in sec:theory.
\begin{figure}[!htbp]
\centering
\resizebox{\textwidth}{!}{%
\begin{tikzpicture}[
  font=\small,
  base/.style={rectangle, rounded corners=4pt, align=center, inner sep=4pt, text width=4.4cm, minimum height=1.15cm},
  pr/.style={base, line width=1.1pt, draw=green!45!black, fill=green!10},
  rm/.style={base, line width=0.9pt, draw=blue!55!black, fill=blue!6},
  he/.style={base, line width=0.9pt, dashed, draw=orange!70!black, fill=orange!7},
  dep/.style={-{Stealth[length=2mm]}, line width=0.7pt, draw=black!55},
]
  \node[pr] (purity) at (0,0)     {\textbf{Pool purity floor}\\{\footnotesize \Cref{thm:purity}}};
  \node[pr] (bayes)  at (0,-5.7)  {\textbf{Purification condition}\\{\footnotesize \Cref{thm:bayes-purity}; H1}};
  \node[he] (t1)   at (5.8,1.0)   {\textbf{Tier 1 error ceiling}\\{\footnotesize \Cref{cor:tier1-floor}}};
  \node[rm] (dec)  at (5.8,-0.9)  {\textbf{Error decomposition}\\{\footnotesize \Cref{thm:asymptotic-elim}; H4}};
  \node[rm] (mono) at (5.8,-3.8)  {\textbf{Purity rises with accuracy}\\{\footnotesize \Cref{thm:monotone-purification}}};
  \node[rm] (fix)  at (5.8,-5.7)  {\textbf{Repeated filtering gives purity 1}\\{\footnotesize \Cref{thm:bayes-convergence}; H2}};
  \node[he] (sc)   at (5.8,-7.6)  {\textbf{Self-correction}\\{\footnotesize \Cref{cor:self-correction}}};
  \node[rm] (corp) at (11.6,-0.9) {\textbf{Corpus ceiling}\\{\footnotesize \Cref{cor:corpus-floor}}};
  \node[he] (rate) at (11.6,-3.0) {\textbf{Storage-rate sensitivity}\\{\footnotesize \Cref{thm:monotone}}};
  \node[he] (conv) at (11.6,-4.9) {\textbf{Convergence envelope}\\{\footnotesize \Cref{thm:convergence}; H3}};
  \node[he] (eq)   at (11.6,-6.8) {\textbf{Equilibrium under rollback}\\{\footnotesize \Cref{cor:stochastic-equilibrium}}};
  \draw[dep] (purity.east) -- (t1.west);
  \draw[dep] (purity.east) -- (dec.west);
  \draw[dep] (dec.east) -- (corp.west);
  \draw[dep] (bayes.east) -- (mono.west);
  \draw[dep] (bayes.east) -- (fix.west);
  \draw[dep] (bayes.east) -- (sc.west);
  \draw[dep] (purity.north) |- (8.7,2.3) |- (conv.west);
  \draw[dep] (rate.south) -- (conv.north);
  \draw[dep] (conv.south) -- (eq.north);
  \node[pr, text width=3.4cm, minimum height=0.55cm, font=\footnotesize] at (0,-9.4) {proven};
  \node[rm, text width=3.4cm, minimum height=0.55cm, font=\footnotesize] at (5.8,-9.4) {remark: exact maths on a simplified model};
  \node[he, text width=3.4cm, minimum height=0.55cm, font=\footnotesize] at (11.6,-9.4) {heuristic: argument sketch};
\end{tikzpicture}}
\caption[Status of the formal statements and the hypotheses they inform.]{Status of the formal statements and the hypotheses (H1 to H4) they inform. Arrows point from a statement to those that build on it.}
\label{fig:theory-map}
\end{figure}


\paragraph{The observed horizon:} the main run produced six cycles (zero to five) and stopped after cycle five, as planned after the cycle-four rollback. The automatic early-stop check fired at the same point, but its storage and retention inputs were wired to the wrong quantities, so its firing is not evidence that the trajectory had levelled off (\Cref{sec:env-cycles}). Statements about long-run behaviour are therefore read against these six cycles, which are too few to estimate limits or rates.

\subsection{Pool purity}
\label{sec:theorem-purity}

The \emph{purity} of the stored pool is the share of stored answers that are correct. At each cycle boundary the pool is made of two parts: older entries that survived re-verification, and new entries admitted by the gate during the cycle.

\begin{theorem}[Pool purity floor]
\label{thm:purity}
Let $\pi_{t+1}$ be the expected purity of the stored pool at the end of cycle $t+1$. Assume:
\begin{enumerate}[label=(\roman*)]
    \item \textbf{Fixed gate:} new answers are stored when $\ustored \ge \tau_{\text{store}}$, with $\tau_{\text{store}}$ fixed across cycles.
    \item \textbf{Like calibration data:} the cycle's candidate answers come from the same distribution as the calibration fold, so the expected purity of the new entries equals $\pi_{\text{store}}$, the share of correct answers among calibration-fold answers with $\ustored \ge \tau_{\text{store}}$.
    \item \textbf{Survivors:} the older entries that survive re-verification (score at least $\tau_{\text{retro}}$) have expected purity $\pi_{\text{retro}}$.
\end{enumerate}
Then
\begin{equation}
\pi_{t+1} \;\ge\; \min\!\big(\pi_{\text{store}},\,\pi_{\text{retro}}\big).
\label{eq:thm-purity}
\end{equation}
\end{theorem}

\begin{proof}
Let the pool contain $n_{\text{new}}$ new entries and $n_{\text{old}}$ surviving older entries. Its expected purity is the weighted average
\[
\pi_{t+1} \;=\; \frac{n_{\text{new}}\,\pi_{\text{store}} + n_{\text{old}}\,\pi_{\text{retro}}}{n_{\text{new}} + n_{\text{old}}},
\]
using (i) and (ii) for the new entries (because the cut-point is the same in every cycle, the calibration-fold precision at that cut-point applies to each cycle's admissions) and (iii) for the survivors. A weighted average of two numbers is never smaller than the smaller of them, which gives \Cref{eq:thm-purity}.
\end{proof}

\paragraph{In plain terms:} a mixture of two groups is at least as clean as the dirtier group. The result is simple; its value is that both inputs can be measured. Because the gate never changes (assumption (i)), the calibration-fold precision $\pi_{\text{store}}$ can be read at every cycle, and the precision of the new entries can be checked directly on the stream chunk, whose answers are labelled. The floor holds in expectation, up to the sampling error of the calibration estimate, and only as far as assumption (ii) holds; for a benchmark whose questions differ from the calibration fold, it says nothing.

\paragraph{What the run shows:} on the training panel, the calibration-fold precision at the cut-point was between $0.732$ and $0.746$ in cycles one to five, and the precision of the answers actually stored from each stream chunk was $0.753$, $0.751$, $0.724$, $0.738$ and $0.724$, within about one point of it; assumption (ii) approximately held on the training panel. It did not hold on the transfer benchmark TruthfulQA, where the answers the gate stored on the evaluation fold were correct only $31$ to $38$ percent of the time at cycles one to five (\Cref{tab:decision-matrix-trajectory}). Per-cycle readings are in \Cref{tab:pi-store-trajectory}.

\subsection{When filtering helps: the purification condition}
\label{sec:theorem-bayes}

The floor says how pure the pool is relative to the calibration fold. The next result says when the gate makes the pool purer than the generator's unfiltered answers at all.

\begin{proposition}[Purification condition]
\label{thm:bayes-purity}
Let $p_c \in (0,1)$ be the share of the generator's answers that are correct at cycle $c$, and let $\mathrm{TPR}_c$ and $\mathrm{FPR}_c$ be the shares of correct and of incorrect answers that the gate stores. The purity of the stored answers is
\begin{equation}
P_c \;=\; \frac{p_c \, \mathrm{TPR}_c}{p_c \, \mathrm{TPR}_c \;+\; (1-p_c) \, \mathrm{FPR}_c},
\label{eq:bayes-purity}
\end{equation}
and $P_c > p_c$ if and only if $\mathrm{TPR}_c > \mathrm{FPR}_c$, which is the same as the gate's balanced accuracy $\alpha_c = (\mathrm{TPR}_c + 1 - \mathrm{FPR}_c)/2$ exceeding one half. If the gate errs equally on both classes ($\mathrm{TPR}_c = 1-\mathrm{FPR}_c = \alpha_c$), then $P_c = p_c \alpha_c / \big(p_c \alpha_c + (1-p_c)(1-\alpha_c)\big)$.
\end{proposition}

\begin{proof}
Let $Y \in \{0,1\}$ mark a correct answer and $S \in \{0,1\}$ a stored one, so $\Pr(Y=1) = p_c$, $\Pr(S=1 \mid Y=1) = \mathrm{TPR}_c$ and $\Pr(S=1 \mid Y=0) = \mathrm{FPR}_c$. Bayes' rule gives $P_c = \Pr(Y=1 \mid S=1)$ as in \Cref{eq:bayes-purity}. Then $P_c > p_c$ means $\mathrm{TPR}_c > p_c \mathrm{TPR}_c + (1-p_c)\mathrm{FPR}_c$, that is $(1-p_c)\mathrm{TPR}_c > (1-p_c)\mathrm{FPR}_c$, which for $p_c < 1$ is $\mathrm{TPR}_c > \mathrm{FPR}_c$. Finally $\alpha_c > 1/2$ means $\mathrm{TPR}_c + 1 - \mathrm{FPR}_c > 1$, again $\mathrm{TPR}_c > \mathrm{FPR}_c$.
\end{proof}

\paragraph{In plain terms:} storing purifies only if the gate accepts correct answers more often than wrong ones. A gate that is only a little better than chance still helps, and a gate worse than chance makes the pool dirtier than the generator (\Cref{fig:purity-curves}). This is why the verifier's balanced accuracy is the precondition tested by Hypothesis~1. In the language of learning with noisy labels, the condition is the class-conditional noise condition $\rho_{+} + \rho_{-} < 1$ of Natarajan et al.~\cite{natarajan2013learning}, with $\rho_{+} = 1-\mathrm{TPR}_c$ and $\rho_{-} = \mathrm{FPR}_c$; their learning guarantees also need a noise-corrected loss, which CAEM does not use, so only the condition is borrowed. Related analyses of retraining on self-predicted labels \cite{das2025retraining} and on weak-teacher labels \cite{lang2024w2s} rest on the same idea. The labelled per-cycle measurement of $\mathrm{TPR}_c$, $\mathrm{FPR}_c$ and $p_c$ needed to check \Cref{eq:bayes-purity} directly was not run and is left as future work (\Cref{ch:conclusion}).

% Bayes purity identity (prop. thm:bayes-purity), symmetric operating point TPR = 1 - FPR = alpha.
\begin{figure}[!htbp]
\centering
\begin{tikzpicture}
\begin{axis}[
  width=8.4cm, height=6.0cm,
  xlabel={share of the generator's answers that are correct, $p$},
  ylabel={purity of the stored pool, $P$},
  xmin=0, xmax=1, ymin=0, ymax=1,
  xtick={0,0.25,0.5,0.75,1}, ytick={0,0.25,0.5,0.75,1},
  axis lines=left, font=\small, tick label style={font=\scriptsize},
  legend style={at={(1.04,0.5)}, anchor=west, draw=black!25, font=\scriptsize, fill=white},
  grid=both, grid style={black!8}, domain=0.001:0.999, samples=120,
]
  \addplot[black!50, dashed, line width=0.9pt] {x}; \addlegendentry{no gate ($P = p$)}
  \addplot[red!60!black, line width=1.1pt, densely dotted] {0.45*x/(0.45*x + 0.55*(1-x))}; \addlegendentry{$\alpha = 0.45$ (worse than chance)}
  \addplot[blue!65!black, line width=1.2pt] {0.6*x/(0.6*x + 0.4*(1-x))}; \addlegendentry{$\alpha = 0.6$}
  \addplot[teal!70!black, line width=1.2pt] {0.7*x/(0.7*x + 0.3*(1-x))}; \addlegendentry{$\alpha = 0.7$}
  \addplot[green!50!black, line width=1.2pt] {0.8*x/(0.8*x + 0.2*(1-x))}; \addlegendentry{$\alpha = 0.8$}
\end{axis}
\end{tikzpicture}
\caption[How the purity of the stored pool depends on the gate's balanced accuracy.]{How the purity of the stored pool depends on the gate's balanced accuracy $\alpha$, from \Cref{eq:bayes-purity} with $\mathrm{TPR} = 1 - \mathrm{FPR} = \alpha$. Curves above the diagonal ($\alpha > 1/2$) make the pool purer than the generator; the dotted curve ($\alpha < 1/2$) makes it dirtier.}
\label{fig:purity-curves}
\end{figure}


Two exact consequences of \Cref{eq:bayes-purity} follow; both are about an idealised gate, not guarantees about the system.

\begin{remark}[Purity rises with the generator's accuracy]
\label{thm:monotone-purification}
For a gate with fixed $\alpha \in (0,1)$ that errs equally on both classes,
\begin{equation}
\frac{\partial P_c}{\partial p_c} \;=\; \frac{\alpha(1-\alpha)}{\big(p_c \alpha + (1-p_c)(1-\alpha)\big)^{2}} \;>\; 0,
\label{eq:monotone-purification}
\end{equation}
so a more accurate generator always yields a purer pool. (Writing $D = p_c\alpha + (1-p_c)(1-\alpha)$, the quotient rule gives $\partial P_c / \partial p_c = \alpha\,[D - p_c(2\alpha - 1)]/D^{2} = \alpha(1-\alpha)/D^{2}$.) In the run, $\alpha$ changes from cycle to cycle with the refitted composite and the range of $p_c$ is narrow, so this remark is not tested separately.
\end{remark}

\begin{remark}[Repeated filtering drives purity toward one]
\label{thm:bayes-convergence}
Suppose the stored pool were filtered again at every cycle by a gate of fixed balanced accuracy $\alpha > 1/2$ that errs equally on both classes and whose errors are independent from one pass to the next. Then the purity follows $P_{c+1} = f(P_c)$ with $f(P) = P\alpha / \big(P\alpha + (1-P)(1-\alpha)\big)$. The map $f$ is increasing, its only fixed points are $0$ and $1$, and $f(P) > P$ for every $P$ in $(0,1)$. A sequence that is increasing and bounded by one converges, and its limit must be a fixed point above $P_0$; hence $P_c \to 1$.
\end{remark}

\paragraph{How far the model is from the system:} the real loop does not re-filter the same pool with a fixed, independent gate. The composite is refitted every cycle, some verifier errors repeat at every pass, and the pool also receives new entries. The remark therefore states the direction the design aims at. Hypothesis~2, which expects purity not to fall across cycles, is tested against it in \Cref{ch:evaluation}.

\subsection{Design heuristics for the cycle-to-cycle trajectory}
\label{sec:theory-heuristics}

The five statements below describe how the pool is expected to change over cycles. Each uses at least one of three simplifications: scores that are normally distributed within the correct and incorrect classes with equal spread; a storage rate and verifier quality that stay stable across cycles; and verifier errors that are independent from one pass to the next. The real system breaks each of these (the composite is refitted every cycle, the spread of scores shrinks over cycles, and some errors are systematic), and six cycles are too few to estimate the long-run quantities involved. They are therefore heuristics, and the evaluation reads them as consistency checks.

\subsubsection{Storage-rate sensitivity}
\label{sec:theorem-monotonicity}

\begin{heuristic}[How the storage rate moves with the verifier's separation]
\label{thm:monotone}
Let $d_t$ be the Cohen's $d$ between the scores of correct and incorrect calibration-fold answers at cycle $t$, a measure of how well the verifier, seen through the cycle's refitted composite, separates them, and let $\sigma_t$ be the share of answers stored at the fixed cut-point $\tau_{\text{store}}$. Assume the scores are normal within each class with common spread $\varsigma$ and class means $\mu_{\pm} = m \pm d\,\varsigma/2$ around a fixed midpoint $m$. Let $p_{+}, p_{-}$ be the shares of correct and incorrect answers and $\phi_{\pm} = \phi\big((\mu_{\pm} - \tau_{\text{store}})/\varsigma\big)$ the class densities at the cut-point. Then $\partial \sigma / \partial d$ has the sign of $p_{+}\phi_{+} - p_{-}\phi_{-}$: when the correct class is denser at the cut-point, better separation stores more answers.
\end{heuristic}

\begin{sketch}
The storage rate is $\sigma = p_{+} \Phi\big((\mu_{+} - \tau_{\text{store}})/\varsigma\big) + p_{-} \Phi\big((\mu_{-} - \tau_{\text{store}})/\varsigma\big)$. Since $\partial \mu_{\pm}/\partial d = \pm\varsigma/2$, differentiating gives $\partial \sigma / \partial d = \tfrac12\big(p_{+}\phi_{+} - p_{-}\phi_{-}\big)$.
\end{sketch}

\paragraph{What the run shows:} the calibration fold gives four readings (cycles one, two, three and five; cycle four was rolled back). From cycle one to two, separation and storage rate moved together, as predicted; from two to three the change in $d$ was too small to read; from three to five the storage rate fell slightly while $d$ rose. Two effects outside the model explain the last case: the composite is refitted every cycle, and the spread of scores shrank rather than staying fixed. Separating them would need a cycle in which the composite is frozen and only the generator changes (\Cref{ch:conclusion}).

\subsubsection{Convergence}
\label{sec:theorem-convergence}

\begin{heuristic}[Convergence envelope]
\label{thm:convergence}
Suppose that the purity of the storage class settles to a long-run value $\pi_{\infty}$, and let $G_t = |\pi_t - \pi_{\infty}|$ be the gap at cycle $t$. Under the assumptions of \Cref{thm:purity} and \Cref{thm:monotone}, and assuming that the storage rate $\sigma$ and the separation $d$ are stable across cycles, there is a rate $r \in (0,1)$ such that, with $t$ counting only the cycles whose fine-tune was kept,
\begin{equation}
\mathbb{E}[G_{t}] \;\le\; (1-r)^{t} \, G_{0}.
\label{eq:thm-convergence}
\end{equation}
\end{heuristic}

\begin{sketch}
Each cycle adds new entries at purity $\pi_{\text{store}}$ (\Cref{thm:purity}) and re-scores the old ones. If the share of the gap closed in a cycle is roughly constant, the gap shrinks by a fixed factor $1-r$ each cycle. Under the normal model, the share closed grows with the storage rate and with the probability that the verifier ranks a correct answer above an incorrect one, $\Phi(d/\sqrt{2})$, which suggests $r$ roughly proportional to $\sigma\,\Phi(d/\sqrt{2})$. A proof would need the stochastic-approximation tools of Robbins and Monro~\cite{robbins1951stochastic}, the almost-supermartingale theorem of Robbins and Siegmund~\cite{robbins1971convergence} (used for verifier-filtered retraining by Fu et al.~\cite{fu2025collapse}) and the rate analysis of Polyak~\cite{polyak1990new}; it is left as future work.
\end{sketch}

\paragraph{What the run shows:} six cycles cannot estimate a geometric rate, so the testable reading is weaker: the calibration-fold precision stayed in a band of width $0.014$ ($0.732$ to $0.746$), comparable to the sampling noise of each reading, with no sign of drifting out of it. A direct test of the rate would need at least fifteen kept cycles under a fixed composite (\Cref{ch:conclusion}). Hypothesis~3 is tested against this heuristic in \Cref{ch:evaluation}.

\subsubsection{Asymptotic hallucination floor}
\label{sec:theorem-asymptotic}

\begin{heuristic}[Tier 1 error ceiling]
\label{cor:tier1-floor}
Let $H_{1}(t)$ be the share of wrong answers among those served by the memory-direct tier at cycle $t$. Assume that a query close enough to reach Tier 1 has the same correct answer as the stored question it matches, and that the entries Tier 1 serves are at least as often correct as the pool as a whole (plausible, because Tier 1 also requires a high stored score). Then, under \Cref{thm:purity},
\begin{equation}
\limsup_{t\to\infty}\, H_{1}(t) \;\le\; \max\!\big(1-\pi_{\text{store}},\,1-\pi_{\text{retro}}\big).
\label{eq:thm-asymptotic}
\end{equation}
\end{heuristic}

\begin{sketch}
Under the first assumption, a Tier 1 answer is wrong exactly when the stored answer is wrong; under the second, the served entries are wrong no more often than the pool, so $H_{1}(t) \le 1 - \pi_t$. By \Cref{thm:purity}, $\pi_t \ge \min(\pi_{\text{store}}, \pi_{\text{retro}})$ at every cycle, so $1 - \pi_t \le \max(1-\pi_{\text{store}}, 1-\pi_{\text{retro}})$ at every cycle, and therefore also in the limit. The first assumption is the weak point: a paraphrase that changes the answer would be served a wrong answer from a correct entry, which is why Tier 1 requires a very close match ($s \ge 0.857$).
\end{sketch}

\paragraph{What the run shows:} on the evaluation fold, whose questions never repeat stored ones, Tier 1 answered $7$ of $9{,}000$ queries, all correctly, so the ceiling cannot be observed there. The router probe of \Cref{ch:evaluation} supplies the testable case: Tier 1 answered $41$ of $42$ exact repeats and all $19$ answered paraphrases correctly. On the training panel the error rate of newly stored entries ($0.247$ to $0.276$) is close to $1 - \pi_{\text{store}}$, as assumption (ii) of \Cref{thm:purity} predicts.

\subsubsection{Equilibrium structure and self-correction}
\label{sec:theorem-corollaries}

\begin{heuristic}[Self-correction under repeated re-verification]
\label{cor:self-correction}
Consider a stored entry that is re-verified at every later cycle boundary. Assume that at each pass a correct entry survives with probability at least $\alpha$ and a wrong entry with probability at most $1-\alpha$, for a fixed $\alpha > 1/2$, and that the verdicts of different passes are independent given whether the entry is correct. Then (i) the probability that an entry is correct, given that it has survived $k$ passes, tends to one as $k$ grows; and (ii) a wrong entry is removed after at most $1/\alpha$ passes on average.
\end{heuristic}

\begin{sketch}
(i) The likelihood ratio of surviving $k$ passes for a correct versus a wrong entry is at least $\big(\alpha/(1-\alpha)\big)^{k}$, which grows without bound because $\alpha/(1-\alpha) > 1$. By Bayes' rule, $\Pr(\text{correct} \mid \text{survived } k) = \Lambda_k q / (\Lambda_k q + 1 - q)$, where $q$ is the prior probability of being correct and $\Lambda_k \ge \big(\alpha/(1-\alpha)\big)^{k}$ the likelihood ratio; this tends to one. (ii) A wrong entry is removed at each pass with probability at least $\alpha$, so the number of passes until removal is at most geometric with mean $1/\alpha$.
\end{sketch}

\paragraph{Where the assumption fails:} the independence assumption is the strong one. If the model and the verifier share a misconception, as Farquhar et al.~\cite{farquhar2024semantic} warn for systematic errors, the same wrong answer is accepted at every pass and repetition removes nothing; purity then rises only toward a ceiling set by such shared errors, which matches the finding of Fu et al.~\cite{fu2025collapse} that verifier-filtered retraining converges to the verifier's knowledge rather than to the truth. Part (ii) also shows how idealised the model is: with $\alpha > 1/2$ it predicts removal within two passes on average.

\paragraph{What the run shows:} the share of checked entries removed by the retroactive pass fell from $13$ percent at cycle one to $9$, $7$ and $4$ percent at cycles two, three and five, consistent with errors being cleared over time, though partly diluted by the growing pool of fresh entries. The entry records do not count how many passes each entry survived, so the predicted survival distribution itself could not be measured.

\begin{heuristic}[Equilibrium under retention rollback]
\label{cor:stochastic-equilibrium}
Let $X_t = 1$ if the cycle-$t$ fine-tune passes the retention guard and $X_t = 0$ if it is rolled back, and let $\pi_{\text{commit}} = \mathbb{E}[X_t]$ be the commit rate.
\begin{enumerate}[label=(\roman*)]
    \item Whenever the worst probe ratio can land on either side of the floor, $0 < \pi_{\text{commit}} < 1$.
    \item Counting all cycles, the envelope of \Cref{thm:convergence} becomes
    \begin{equation}
    \mathbb{E}[G_t] \;\le\; (1 - \pi_{\text{commit}}\, r)^{t} \, G_0,
    \label{eq:thm-stochastic-envelope}
    \end{equation}
    because a rolled-back cycle leaves the gap unchanged.
    \item The memory keeps growing through rollbacks whenever the new admissions outnumber the removals, because answering the stream chunk does not depend on $X_t$; only the adapter waits for committed cycles.
    \item When the worst probe ratio sits near the floor, commits and rollbacks alternate at a stable rate.
\end{enumerate}
\end{heuristic}

\begin{sketch}
A committed cycle shrinks the expected gap by the factor $1-r$ and a rolled-back cycle leaves it unchanged, so $\mathbb{E}[G_{t+1} \mid G_t] = (1 - \pi_{\text{commit}}\, r)\, G_t$; iterating gives \Cref{eq:thm-stochastic-envelope}. Part (iii) follows from the fact that the stream chunk is answered, and new entries stored, whether or not the cycle's fine-tune is kept. Part (iv) is an informal claim about a stationary state and is not proven.
\end{sketch}

\paragraph{What the run shows:} parts (i) and (iii) can be checked directly. The commit sequence over cycles one to five was $1, 1, 1, 0, 1$, with worst probe ratios $1.006$, $0.949$, $0.976$, $0.886$ and $0.958$, so the observed commit rate is $4/5$. The memory grew at every cycle, including across the cycle-four rollback (\Cref{tab:app-per-cycle-memory}). Part (ii) inherits the limits of \Cref{thm:convergence}, and part (iv) would need at least fifteen kept cycles to test (\Cref{ch:conclusion}); the ratios at cycles two to five lie within $0.07$ of the floor, which is consistent with it but not evidence for it.

\subsection{Remarks on the overall hallucination rate}
\label{sec:theorem-scaffolding}

Two further remarks place the memory's contribution within the error rate of the whole system. Neither is measured in this thesis.

\begin{remark}[Decomposition of the overall error rate]
\label{thm:asymptotic-elim}
Let $\tau_{1}(t)$ be the share of queries answered by Tier 1 at cycle $t$, $H_{\text{mem}}(t)$ the error rate of those answers (at most $1 - \pi_t$ under the assumptions of \Cref{cor:tier1-floor}), and $H_{\text{gen}}(t)$ the error rate of answers generated by Tiers 2 and 3. By the law of total probability, the overall error rate is
\begin{equation}
H_t \;=\; \big(1 - \tau_{1}(t)\big) \, H_{\text{gen}}(t) \;+\; \tau_{1}(t) \, H_{\text{mem}}(t).
\label{eq:asymptotic-elim}
\end{equation}
The memory can lower the overall rate only through the second term, so its effect grows with the share of queries it answers. As $\pi_t$ rises, the bound on $H_{\text{mem}}$ falls, and the remaining error is set by the generator and its retrieval corpus on the questions the memory cannot answer.
\end{remark}

\paragraph{Reading the remark:} on the evaluation fold $\tau_{1}(t)$ was at most $0.1$ percent, so the overall rate there is essentially the generator's. The remark becomes testable only in deployment, where real users repeat questions; Hypothesis~4 is read against it in \Cref{ch:evaluation}.

\begin{remark}[Ceiling set by the retrieval corpus]
\label{cor:corpus-floor}
Let $K(C)$ be the set of questions whose answer can be supported by the corpus $C$. A system that answers every question and asserts only what $C$ supports can be right on at most the questions in $K(C)$; on the rest any answer it gives is wrong, so it must abstain or err. The error left after the memory and the generator have done their best is therefore bounded below by the share of questions outside $K(C)$ that the system chooses to answer, a limit set by the corpus rather than by the architecture. Estimating it would require checking, for every evaluation question, whether any retrieved passage contains the gold answer, which was not done (\Cref{ch:conclusion}).
\end{remark}

\section{Chapter signpost}
\label{sec:ch4-signpost}

This chapter fixed every component, threshold and fitting procedure before the results were read. \Cref{ch:evaluation} measures them on the six-cycle run, the eight external baselines and the retrieval-utility-classifier comparison; it reports each pre-specified hypothesis as supported, partly supported or unsupported, checks the two proven results against their measurements, and reads the heuristics as consistency checks.
   % Methodology and Design
% ============================================================================
% Chapter 5 — Performance Evaluation and Analysis
% Rubric-mandated six-section top level (5.1 to 5.6) plus the signpost.
% Course outcomes covered: CO6 (analyse alternatives, performance), CO7
% (final design adjustments). Section titles are fixed by the rubric.
% ============================================================================

\chapter{Performance Evaluation and Analysis}
\label{ch:evaluation}

% Chapter scope: Section 5.1 reports the results of the main run; 5.2 the
% cycle-zero choices that fixed the configuration; 5.3 the statistical
% tests; 5.4 the comparison with external baselines and the classifier
% ablation; 5.5 the summary and hypothesis verdicts; 5.6 the discussion.

\section{Performance Evaluation}
\label{sec:performance}

\subsection{Experimental setup}
\label{sec:setup}

\subsubsection{Benchmarks and the training-versus-transfer split}
\label{sec:setup-benchmarks}

Five benchmarks are used. The \textbf{training panel} is FEVER \cite{thorne-etal-2018-fever} (claim verification), TriviaQA \cite{joshi-etal-2017-triviaqa} (open-domain trivia) and CommonsenseQA \cite{talmor-etal-2019-commonsenseqa} (five-option commonsense questions). The \textbf{transfer panel} is TruthfulQA \cite{lin-etal-2022-truthfulqa} (common misconceptions) and StrategyQA \cite{10.1162/tacl_a_00370} (implicit multi-step yes-or-no questions). Only the training panel feeds the calibration fold, the seed pool, the stream chunks and the training pool; the transfer panel is evaluated only, to measure how the system behaves on questions unlike those it learns from. The retention guard uses three separate probes: MMLU \cite{hendrycks2021measuring} and held-out test slices of TriviaQA and CommonsenseQA. HotpotQA and Natural Questions were dropped before the main run because the base generator answered too few of their questions correctly ($0.09$ and $0.17$) for the gate to build a usable pool (\Cref{sec:scope-research}).

Each training benchmark is split into six disjoint folds (\Cref{tab:benchmark-pools}): a seed pool for the cold-start memory, a calibration fold, a purity-validation slice, ten stream chunks (one per possible cycle), an evaluation fold of $300$ questions used for every headline number, and a test fold kept for later work. The transfer benchmarks have only evaluation and test folds. All folds were checked for overlap by content hash when they were built.

\begin{table}[!htbp]
\centering
\small
\begin{tabular}{@{}p{2.6cm}p{1.5cm}p{1.0cm}p{1.0cm}p{1.0cm}p{2.0cm}p{1.0cm}p{1.0cm}@{}}
\toprule
\textbf{Benchmark} & \textbf{Panel} & \textbf{Seed} & \textbf{Calib} & \textbf{Purity} & \textbf{Stream chunks} & \textbf{Eval} & \textbf{Test} \\
\midrule
FEVER             & training & $1000$ & $500$ & $500$ & $10\times 2000$ & $300$ & $500$ \\
TriviaQA          & training & $1000$ & $500$ & $500$ & $10\times 2000$ & $300$ & $500$ \\
CommonsenseQA     & training & $1000$ & $500$ & $500$ & $10\times 700$  & $300$ & $500$ \\
TruthfulQA        & transfer & $0$    & $0$   & $0$   & $0$             & $300$ & $500$ \\
StrategyQA        & transfer & $0$    & $0$   & $0$   & $0$             & $300$ & $387$ \\
\bottomrule
\end{tabular}
\caption[Benchmark folds.]{Benchmark folds. CommonsenseQA has smaller stream chunks because its training split holds only about $10{,}000$ questions after de-duplication; the StrategyQA test fold is limited by the size of its split.}
\label{tab:benchmark-pools}
\end{table}

\subsubsection{Metrics}
\label{sec:setup-metrics}

\textbf{Exact match (EM):} the main accuracy measure, read from the final answer line of each response by one rule for every system. FEVER and StrategyQA are scored by the label the answer starts with; CommonsenseQA by the single option letter given, where the option text may follow the letter but two letters count as wrong; TriviaQA by exact match against any gold alias after normalisation. TruthfulQA is scored by a language-model judge (Claude Haiku) that compares the final answer with the benchmark's correct and incorrect reference answers, following the judge-based evaluation introduced with the benchmark \cite{lin-etal-2022-truthfulqa}; a refusal counts as correct only when the reference answers include one. The same judge, reading the same part of the response, scores every system.

\textbf{Composite hallucination metric (CHM):} the equal-weighted mean of seven failure rates, each computed over all evaluation questions.
\begin{itemize}
    \item \textbf{Confident confabulation:} the answer is wrong but its stored score $\ustored$ is at least $0.50$.
    \item \textbf{Factual fabrication:} the answer is wrong and its weakest atomic fact has passage support below $0.30$.
    \item \textbf{Logical fabrication:} the answer is wrong and the model's own chains support it below $0.30$.
    \item \textbf{Defensive evasion:} the answer is wrong and uses an evasive phrase such as \emph{the context does not mention}.
    \item \textbf{Template leakage:} the output contains a leaked prompt placeholder.
    \item \textbf{False refusal:} the answer is correct but the system abstained or discarded it.
    \item \textbf{Over-long output:} the answer is wrong and longer than $600$ characters.
\end{itemize}
Two of these, confident confabulation and false refusal, are defined through the system's own verifier, so they move with the gate's operating point: a more permissive gate refuses fewer correct answers and lets more wrong answers keep a high score. Two further rates are recorded but left out of the mean, because neither can fire in this run:
\begin{itemize}
    \item \textbf{Factual contradiction:} the entailment judge returns only a support probability, with no contradiction output.
    \item \textbf{Off-topic:} planned as a wrong answer whose question-match score is below $0.50$. That score cannot detect off-topic answers on this panel, and its stored values never fall below $0.50$ (\Cref{sec:verifier-qa-rel}). An attempt after the run to measure off-topic answers directly, with a language-model judge, was checked against $200$ hand-labelled answers and agreed too weakly to give a reliable rate (Cohen's kappa about $0.35$), so measuring them is left to future work.
\end{itemize}

\textbf{Model view and delivered view:} every system is reported in two views. The \emph{model view} scores EM and CHM on every raw answer the generator wrote, including answers that CAEM then withheld; it measures what the model knows. The \emph{delivered view} scores what the user receives. CAEM shows every answer with its probability and a status tag (verified, provisional or conflicting, \Cref{sec:decision}) except abstentions, which hide the answer behind an insufficient-evidence message; semantic-entropy abstention (B8) withholds the answers it abstains on; every other baseline answers every question. Following the selective-prediction setting of Kamath et al.~\cite{kamath2020selective} and Geifman and El-Yaniv~\cite{geifman2017selective}, the delivered view reports the share of questions answered, the share of those answers that are wrong, and CHM over the answered questions; for CAEM it also reports the accuracy of each status tag, because the tag is part of what the user sees. Unlike the two verifier-defined rates, the answered share and the error among answered do not depend on the verifier's scores.

\textbf{Composite evaluation score (CES):} a single summary of five axes, each in $[0,1]$ with one best: accuracy (pooled EM), epistemic quality ($1 - \mathrm{CHM}$), retention (worst probe ratio, capped at one), calibration ($1 - 2\,\mathrm{ECE}$) and verifier reliability ($\max(2\,\mathrm{BA} - 1, 0)$, with the balanced accuracy approximated as $\Phi(d/\sqrt{2})$ from the calibration-fold Cohen's $d$). The score is the geometric mean, so a collapse on one axis is not averaged away:
\begin{equation}
\mathrm{CES} = \left(\prod_{i=1}^{5} \mathrm{clamp}_{[\varepsilon,1]}(\mathrm{axis}_i)\right)^{1/5}, \quad \varepsilon = 0.01.
\label{eq:ces}
\end{equation}

\textbf{Retention ratio:} accuracy on each probe divided by its value before any fine-tuning (\Cref{eq:nfr-retention-floor}); the worst probe decides whether a cycle is kept. \Cref{tab:metrics} summarises the metrics.

\begin{table}[!htbp]
\centering
\small
\begin{tabular}{@{}>{\raggedright\arraybackslash}p{3.8cm}>{\raggedright\arraybackslash}p{3.4cm}>{\raggedright\arraybackslash}p{7.0cm}@{}}
\toprule
\textbf{Metric} & \textbf{Reported on} & \textbf{Definition} \\
\midrule
Exact match (EM) & all five benchmarks & label or alias match; TruthfulQA judged by Claude Haiku against the reference answers \\
Composite hallucination metric (CHM) & all five benchmarks & mean of seven failure rates \\
Answered, wrong among answered & pooled & delivered view: share not withheld; share of those that are wrong \\
Status-tag accuracy & pooled, CAEM & share of questions and EM per tag shown to the user \\
Storage precision $\pi_{\text{store}}$ & per cycle & share of stored answers that are correct \\
Gate balanced accuracy, Cohen's $d$ & calibration fold & $(\mathrm{TPR} + 1 - \mathrm{FPR})/2$ at $\tau_{\text{store}}$; separation of correct and incorrect scores \\
Composite evaluation score (CES) & pooled & geometric mean of five axes (\Cref{eq:ces}) \\
Retention ratio & probe panel & worst probe accuracy over its pre-fine-tuning value \\
\bottomrule
\end{tabular}
\caption[Metrics reported in this chapter.]{Metrics reported in this chapter.}
\label{tab:metrics}
\end{table}

\subsubsection{Comparison protocol}
\label{sec:setup-protocol}

The external baselines (\Cref{sec:external-baselines}) are compared with CAEM under one protocol, so that a difference in the results reflects the method rather than the prompt or the scoring.
\begin{itemize}
    \item \textbf{Matched prompts:} every baseline that answers from the question alone receives CAEM's zero-shot question message: the same system instruction, worked example, task line, answer format and forced reasoning prefix. Chain of thought (B2) and five-shot chain of thought (B5) add the step-by-step cue to that message; B5's five demonstrations per benchmark are drawn from the benchmark's training split (for TruthfulQA, which has a single split, from the questions not evaluated), with every evaluated question removed. The two retrieval baselines (B3, B4) use CAEM's retrieval prompt with the same top five passages, vanilla fine-tuning (B6) is evaluated with the prompt it was trained on, which is CAEM's zero-shot message, and FLARE (B7) uses the zero-shot message for its draft and its final answer.
    \item \textbf{One scoring rule:} every system is scored by the rules of \Cref{sec:setup-metrics}, applied to its final answer line, and every TruthfulQA answer by the same judge.
    \item \textbf{One hallucination instrument:} CHM needs verifier signals, so each baseline answer is passed through CAEM's verifier with the adapter and composite of cycle three, the same instrument for every baseline. The decision the verifier would take on a baseline answer is only a measurement: the baseline never acted on it. A baseline is therefore never charged with a false refusal or a withheld answer from CAEM's decision rule; its withheld answers are only those it withheld itself (B8).
    \item \textbf{Same questions:} every system answers the same $1{,}500$ evaluation questions, in the same order, checked question by question before scoring.
\end{itemize}
CAEM's outputs come from the main run. An earlier comparison had given the baselines a shorter prompt than CAEM's and scored some answer formats differently, so after the main run the baselines were re-run under this protocol (B3 and B4 only on TruthfulQA, where a graphics-memory error had left the first batch of $32$ questions unanswered; their prompts were already CAEM's). CAEM's own prompts are identical, character for character, in the code of the main run and of the re-run.

\subsubsection{Hardware, seed, reproducibility}
\label{sec:setup-reproducibility}

All phases ran on one rented graphics processor of the $32$-gigabyte class, under a single recorded random seed. The main run took cycles zero to five and stopped after cycle five as planned after the cycle-four rollback (\Cref{sec:env-cycles}). The baseline re-run of \Cref{sec:setup-protocol} used a second machine of the same class. For that re-run, the passages for the evaluation questions were retrieved once by exact search over the same index and reused; on a check sample, live search returned the same passages for $50$ of $50$ questions, and the verifier's passage scores from the main run were reproduced within $0.01$ on $98$ percent of answers. Signals computed from sampled text (dropout agreement, chain similarity and chain entailment) are random draws by design and differ slightly between runs. The code, the configuration files and the per-cycle artefacts are listed in \Cref{app:reproducibility}.

\subsection{Trajectory of headline metrics}
\label{sec:headline-trajectory}

All results in this section use the locked configuration chosen at cycle zero; \Cref{sec:final-adjustments} explains how it was chosen. The section order follows the BRAC template rather than the order in which the work was done.

\subsubsection{The headline in three steps}
\label{sec:headline-pooled}

\Cref{tab:headline-steps} separates the three contributions that a single before-and-after comparison would mix together: the architecture (B1, the same generator answering zero-shot with the same prompt, to C0, CAEM before any fine-tuning), the self-improvement loop (C0 to C5) and the whole system (B1 to C5). Significance is read from the paired per-benchmark tests of \Cref{sec:comp-baselines}.

% Auto-generated by scripts/build_ch5_evidence_tables.py; EM/CHM/subtypes from outputs/analysis/chm_judged.json
% (Haiku-judged TruthfulQA); coverage and error-among-answered from the
% per-sample evaluation files (ABSTAIN = not answered; TruthfulQA uses the judged correctness field).
\begin{table}[!htbp]
\centering\small
\setlength{\tabcolsep}{5pt}
\begin{tabular}{@{}lrrrrrr@{}}
\toprule
\textbf{System} & \textbf{EM} & \textbf{CHM} & \twoline{Confident\\confab.} & \twoline{False\\refusal} & \textbf{Answered} & \twoline{Wrong among\\answered} \\
\midrule
B1 (zero-shot) & $0.507$ & $0.137$ & $0.327$ & $0.055$ & $100.0\%$ & $0.493$ \\
CAEM C0 & $0.541$ & $0.134$ & $0.088$ & $0.298$ & $66.1\%$ & $0.381$ \\
CAEM C2 & $0.544$ & $0.110$ & $0.259$ & $0.057$ & $87.7\%$ & $0.425$ \\
CAEM C5 & $0.517$ & $0.117$ & $0.307$ & $0.065$ & $83.0\%$ & $0.451$ \\
\bottomrule
\end{tabular}
\caption[The headline in three steps.]{The headline in three steps on the five-benchmark evaluation fold ($1{,}500$ questions). B1 to C0 isolates the architecture, C0 to C5 the self-improvement loop, B1 to C5 the whole system.}
\label{tab:headline-steps}
\end{table}


% Auto-generated by scripts/build_ch5_evidence_tables.py from outputs/analysis/chm_judged.json (2026-10-19).
\begin{figure}[!htbp]
\centering
\begin{tikzpicture}
\begin{axis}[name=em, width=0.47\textwidth, height=5.2cm, xlabel={cycle}, ylabel={pooled EM}, xtick={0,...,5},
  ymin=0.50, ymax=0.54, font=\small, tick label style={font=\scriptsize}, grid=major, grid style={black!8},
  legend style={font=\scriptsize, at={(0.98,0.97)}, anchor=north east}]
  \addplot[blue!65!black, mark=*, thick] coordinates {(0,0.517) (1,0.533) (2,0.535) (3,0.516) (4,0.517) (5,0.508)};
  \addlegendentry{CAEM}
  \addplot[black!60, dashed, thick, domain=0:5] {0.511};
  \addlegendentry{B1}
\end{axis}
\begin{axis}[at={(em.east)}, anchor=west, xshift=1.4cm, width=0.47\textwidth, height=5.2cm, xlabel={cycle},
  ylabel={pooled CHM (lower is better)}, xtick={0,...,5}, ymin=0.12, ymax=0.16, font=\small,
  tick label style={font=\scriptsize}, grid=major, grid style={black!8},
  legend style={font=\scriptsize, at={(0.98,0.97)}, anchor=north east}]
  \addplot[red!65!black, mark=square*, thick] coordinates {(0,0.155) (1,0.131) (2,0.127) (3,0.137) (4,0.136) (5,0.134)};
  \addlegendentry{CAEM}
  \addplot[black!60, dashed, thick, domain=0:5] {0.132};
  \addlegendentry{B1}
\end{axis}
\end{tikzpicture}
\caption[EM and CHM by cycle against B1.]{Pooled EM (left) and CHM (right) by cycle on the evaluation fold, against B1. EM is level with B1 at C0, peaks at C2 and returns to B1's level by C5; CHM is above B1 at C0, where the system refuses too often, about level with it from C1, and ends just above it.}
\label{fig:trajectory}
\end{figure}


\textbf{The architecture adds almost no accuracy:} memory, routing, the retrieval-utility classifier and the verifier together move pooled EM from $0.511$ to $0.517$ before any fine-tuning (\Cref{fig:trajectory}), and from $0.554$ to $0.569$ on the training panel. No benchmark shows a significant EM difference between B1 and C0.

\textbf{The architecture alone is over-cautious:} at C0 the system abstains on $33.9$ percent of evaluation questions and discards another $31.9$ percent. False refusal reaches $0.276$, which lifts CHM to $0.155$, above B1's $0.132$, even though confident confabulation falls from $0.208$ to $0.089$. C0's CHM is significantly higher than B1's on TriviaQA, TruthfulQA and StrategyQA.

\textbf{Self-improvement helps for two cycles, then erodes:} after the first cycle the system abstains on $14.4$ percent and discards $9.7$ percent, false refusal falls to $0.059$ and CHM to $0.131$; EM rises to $0.533$ at C1 and $0.535$ at C2, the best values of the run. From C3 the gains fade: pooled EM ends at $0.508$ and CHM at $0.134$. On the training panel EM peaks at C1 ($0.584$, \Cref{tab:per-bench-trajectory}) and ends at $0.549$, below B1's $0.554$.

\textbf{Confident confabulation returns:} as the refitted gate becomes more permissive, confident confabulation rises to $0.259$ at C2 and $0.305$ at C5, above B1's $0.208$: after self-improvement, the verifier gives a high score to more of the generator's wrong answers than it gave to the base generator's. Because this rate and false refusal are both defined through the verifier's own score, they trade off: their sum is $0.365$ at C0, between $0.297$ and $0.329$ at C1 to C4 and $0.361$ at C5 (\Cref{fig:tradeoff}). The net CHM change from C0 to C5 ($0.155$ to $0.134$) comes from the other subtypes: fewer wrong answers with unsupported facts ($0.381$ to $0.327$), unsupported reasoning ($0.298$ to $0.245$) or evasive phrasing ($0.044$ to $0.003$).

% Auto-generated by scripts/build_ch5_evidence_tables.py (final drafting pass, 2026-10).
\begin{figure}[!htbp]
\centering
\begin{tikzpicture}
\begin{axis}[width=0.75\textwidth, height=5.4cm, xlabel={cycle}, ylabel={rate (evaluation fold)}, xtick={0,...,5},
  ymin=0, ymax=0.34, font=\small, tick label style={font=\scriptsize}, grid=major, grid style={black!8},
  legend style={font=\scriptsize, at={(1.03,0.5)}, anchor=west}]
  \addplot[orange!80!black, mark=*, thick] coordinates {(0,0.088) (1,0.238) (2,0.259) (3,0.291) (4,0.289) (5,0.307)};
  \addlegendentry{confident confabulation}
  \addplot[teal!70!black, mark=triangle*, thick] coordinates {(0,0.298) (1,0.069) (2,0.057) (3,0.054) (4,0.053) (5,0.065)};
  \addlegendentry{false refusal}
  \addplot[orange!80!black, dashed, domain=0:5] {0.327};
  \addlegendentry{B1 confabulation}
  \addplot[teal!70!black, dashed, domain=0:5] {0.055};
  \addlegendentry{B1 false refusal}
\end{axis}
\end{tikzpicture}
\caption[Confident confabulation against false refusal by cycle.]{Confident confabulation against false refusal by cycle. Both are defined through the verifier's own score, so they move in opposite directions as the refitted gate becomes more permissive.}
\label{fig:tradeoff}
\end{figure}


\textbf{The whole system ends level with B1:} at C5, pooled EM is $0.508$ against $0.511$ and CHM $0.134$ against $0.132$; at its best cycle, C2, it is $0.535$ and $0.127$. None of these differences is significant on any benchmark, at C2 or at C5. In the delivered view (\Cref{fig:selective}), CAEM is wrong on $40.5$ percent of the questions it answers at C0, where it answers two thirds of them, against $48.9$ percent for B1, which answers everything. As self-improvement makes the system answer more ($87.7$ percent at C2, $83.0$ percent at C5), its error among answered rises to $0.428$ and $0.451$, close to the $0.427$ of semantic-entropy abstention (B8), which answers $89.5$ percent.

% Auto-generated by scripts/build_ch5_evidence_tables.py (final drafting pass, 2026-10).
\begin{figure}[!htbp]
\centering
\begin{tikzpicture}
\begin{axis}[width=0.72\textwidth, height=5.6cm, xlabel={share of questions answered (\%)}, ylabel={share of answers that are wrong},
  xmin=60, xmax=104, ymin=0.36, ymax=0.51, font=\small, tick label style={font=\scriptsize}, grid=major, grid style={black!8}]
  \addplot[blue!65!black, mark=*, thick] coordinates {(34.2,0.294) (75.9,0.375) (80.5,0.396) (81.7,0.416) (81.7,0.415) (81.5,0.444)};
  \node[font=\scriptsize, anchor=south east] at (axis cs:34.2,0.294) {C0};
  \node[font=\scriptsize, anchor=north] at (axis cs:80.5,0.396) {C2};
  \node[font=\scriptsize, anchor=south east] at (axis cs:81.5,0.444) {C5};
  \node[font=\scriptsize, anchor=west] at (axis cs:89.5,0.432) {C1 to C4};
  \addplot[black, only marks, mark=diamond*, mark size=3pt] coordinates {(100,0.493)};
  \node[font=\scriptsize, anchor=north east] at (axis cs:100,0.493) {B1};
\end{axis}
\end{tikzpicture}
\caption[Answered share against error among answered.]{Selective-prediction view on the evaluation fold: share of questions answered and share of those answers that are wrong, for each cycle and for B1, which answers everything. CAEM stays below B1's error but answers more and errs more as the cycles go on.}
\label{fig:selective}
\end{figure}


\textbf{The status tag predicts correctness:} the user sees each answer with a status tag (\Cref{sec:setup-metrics}), and at every cycle verified answers are more often right than provisional ones, which in turn are more often right than conflicting or insufficient ones (\Cref{tab:status-tags}). Answers tagged verified are right $74.6$ percent of the time at C0, when a quarter of the questions receive that tag, and $65.5$ to $68.9$ percent at C1 to C5, when $42$ to $56$ percent do; provisional answers are right $40$ to $60$ percent of the time, and conflicting and insufficient ones $15$ to $48$ percent. B1 is right on $51.1$ percent of its answers and gives the user no signal of which ones. The tag is therefore the most direct user-facing benefit of the architecture: from C1 on, a user who acts only on verified answers meets about $30$ percent fewer wrong answers than one who trusts every answer of B1 (an error rate of $0.31$ to $0.35$ against $0.49$), at the cost of receiving a verified answer for only about half of the questions.

% Generated by scripts/build_final_tables.py from outputs/analysis/chm_judged.json (label em_final); do not edit.
\begin{table}[!htbp]
\centering\small
\setlength{\tabcolsep}{3.5pt}
\begin{tabular}{@{}lccccccccc@{}}
\toprule
 & \multicolumn{2}{c}{\textbf{Verified}} & \multicolumn{2}{c}{\textbf{Provisional}} & \multicolumn{2}{c}{\textbf{Conflicting}} & \multicolumn{2}{c}{\textbf{Insufficient}} & \\
\cmidrule(lr){2-3}\cmidrule(lr){4-5}\cmidrule(lr){6-7}\cmidrule(lr){8-9}
\textbf{Cycle} & \textbf{Share} & \textbf{EM} & \textbf{Share} & \textbf{EM} & \textbf{Share} & \textbf{EM} & \textbf{Share} & \textbf{EM} & \textbf{All} \\
\midrule
C0 & $24.9\%$ & $0.746$ & $9.3\%$ & $0.597$ & $31.9\%$ & $0.477$ & $33.9\%$ & $0.365$ & $0.517$ \\
C1 & $44.2\%$ & $0.689$ & $31.7\%$ & $0.535$ & $9.7\%$ & $0.219$ & $14.4\%$ & $0.264$ & $0.533$ \\
C2 & $46.8\%$ & $0.655$ & $33.6\%$ & $0.532$ & $7.3\%$ & $0.220$ & $12.3\%$ & $0.266$ & $0.535$ \\
C3 & $55.6\%$ & $0.668$ & $26.0\%$ & $0.403$ & $6.3\%$ & $0.326$ & $11.9\%$ & $0.151$ & $0.516$ \\
C4 & $55.3\%$ & $0.667$ & $26.3\%$ & $0.411$ & $6.6\%$ & $0.303$ & $11.7\%$ & $0.159$ & $0.517$ \\
C5 & $42.3\%$ & $0.667$ & $39.1\%$ & $0.433$ & $1.5\%$ & $0.217$ & $17.0\%$ & $0.306$ & $0.508$ \\
\bottomrule
\end{tabular}
\caption[Accuracy of the status tag shown to the user.]{Share of evaluation questions and EM for each status tag CAEM shows the user (verified: stored; provisional: deferred; conflicting: discarded, shown with a warning; insufficient: abstained, answer withheld), with EM over all questions for reference. B1, which shows no tag, has EM $0.511$.}
\label{tab:status-tags}
\end{table}


The retention ratios of every cycle are in \Cref{tab:sil-training-trajectory}: MMLU never comes near the floor (between $0.975$ and $1.041$); the TriviaQA test probe falls to $0.886$ at cycle four and triggers the only rollback. The composite evaluation score (\Cref{tab:ces-caem-trajectory} in \Cref{app:extra-tables}) rises from $0.585$ to $0.704$, mostly through its verifier-reliability axis, which grows with the calibration-fold separation $d$ (\Cref{tab:gate-quality}); its accuracy axis falls over the same cycles, so the rise in CES should not be read as a rise in accuracy.

\subsubsection{Per-benchmark trajectory}
\label{sec:headline-per-bench}

\Cref{tab:per-bench-trajectory} and \Cref{tab:halluc-subtypes} break the pooled numbers down. On the training panel, FEVER and TriviaQA peak at C1 or C2 (FEVER $0.533$ against B1's $0.507$; TriviaQA $0.493$ against $0.430$) and fall back afterwards, FEVER to $0.477$, below B1, and TriviaQA to $0.453$, still above it. CommonsenseQA is highest at C0 ($0.757$ against $0.727$) and ends at $0.717$. TriviaQA is the benchmark where the architecture helps most and where the verifier separates correct from incorrect answers best (\Cref{sec:disc-verifier-scope}); even there, the C2 gain over B1 is not significant at $300$ questions. On the transfer panel, TruthfulQA stays between $0.230$ and $0.277$ at every cycle against B1's $0.250$, and StrategyQA between $0.610$ and $0.653$ against $0.643$. Across subtypes, factual and logical fabrication fall gently after C0; defensive evasion and over-long output are near zero. Template leakage is zero for every cycle and every baseline, and this is a real zero: the answer prompt shows the reasoning line as a bracketed placeholder, and no system ever copied a placeholder into its output.

% Table sizing standardised 2026-10-17 (final drafting pass): \small, no \resizebox.
% Caption shortened by hand in the final drafting pass (2026-10); detail moved to chapter 5 text.
% Auto-generated by scripts/aggregate_section_5_2.py
% Per-cycle, per-benchmark EM (top sub-row) and CHM (bottom sub-row),
% with training-panel pooled rows merged in (previously tab:train-vs-transfer),
% the no-CAEM reference floor (B1 zero-shot) on the left, and absolute +
% relative deltas against the best CAEM cycle (C2) and the trajectory final
% cycle (C5) on the right.
% Transfer panel intentionally has no pooled row: TruthfulQA and StrategyQA
% move in opposite directions and any pooled average is dishonest reporting.
\begin{table}[!htbp]
\centering\small
\setlength{\tabcolsep}{3pt}
\begin{tabular}{@{}lccccccccc@{}}
\toprule
\textbf{Benchmark} & \textbf{B1} & \textbf{C0} & \textbf{C1} & \textbf{C2} & \textbf{C3} & \textbf{C4} & \textbf{C5} & \twoline{Change\\to C2} & \twoline{Change\\to C5} \\
\midrule
\multicolumn{10}{@{}l}{\textit{Training panel}} \\
FEVER (EM) & 0.453 & 0.487 & 0.533 & 0.533 & 0.503 & 0.507 & 0.477 & $+0.080$ & $+0.024$ \\
FEVER (CHM) & 0.159 & 0.139 & 0.113 & 0.102 & 0.120 & 0.119 & 0.126 & $-0.057$ & $-0.033$ \\
TriviaQA (EM) & 0.293 & 0.463 & 0.480 & 0.493 & 0.463 & 0.463 & 0.453 & $+0.200$ & $+0.160$ \\
TriviaQA (CHM) & 0.175 & 0.156 & 0.113 & 0.109 & 0.121 & 0.120 & 0.125 & $-0.066$ & $-0.050$ \\
CSQA (EM) & 0.703 & 0.763 & 0.740 & 0.717 & 0.743 & 0.743 & 0.717 & $+0.014$ & $+0.014$ \\
CSQA (CHM) & 0.096 & 0.077 & 0.086 & 0.093 & 0.086 & 0.085 & 0.091 & $-0.003$ & $-0.005$ \\
\rowcolor{gray!18}
\textbf{Pooled (EM)} & $\mathbf{0.483}$ & $0.571$ & $0.584$ & $\mathbf{0.581}$ & $0.570$ & $0.571$ & $0.549$ & $\mathbf{+0.098}$ & $\mathbf{+0.066}$ \\
\rowcolor{gray!18}
\textbf{Pooled (CHM)} & $\mathbf{0.143}$ & $0.124$ & $0.104$ & $\mathbf{0.101}$ & $0.109$ & $0.108$ & $0.114$ & $\mathbf{-0.042}$ & $\mathbf{-0.029}$ \\
\midrule
\multicolumn{10}{@{}l}{\textit{Transfer panel}} \\
TruthfulQA (EM) & 0.433 & 0.380 & 0.317 & 0.323 & 0.310 & 0.307 & 0.300 & $-0.110$ & $-0.133$ \\
TruthfulQA (CHM) & 0.051 & 0.150 & 0.133 & 0.116 & 0.119 & 0.122 & 0.104 & $+0.065$ & $+0.053$ \\
StrategyQA (EM) & 0.650 & 0.610 & 0.647 & 0.653 & 0.640 & 0.640 & 0.640 & $+0.003$ & $-0.010$ \\
StrategyQA (CHM) & 0.120 & 0.154 & 0.113 & 0.115 & 0.128 & 0.127 & 0.099 & $-0.005$ & $-0.021$ \\
\bottomrule
\end{tabular}
\caption[EM and CHM by benchmark and cycle.]{EM (upper row) and CHM (lower row) by benchmark and cycle on the evaluation fold, with B1 and the change from B1 to C2 and C5. Shaded rows pool the training panel; the transfer panel is not pooled because its two benchmarks move in opposite directions.}
\label{tab:per-bench-trajectory}
\end{table}


% CHM cells re-scored (Haiku-judged TruthfulQA; mean of the seven active types) by scripts/patch_chm_judged_tables.py (2026-10-17).
% CHM cells re-scored with the Haiku-judged TruthfulQA label by scripts/patch_chm_judged_tables.py (2026-10-17).
% Table sizing standardised 2026-10-17 (final drafting pass): \small, no \resizebox.
% Caption shortened by hand in the final drafting pass (2026-10); detail moved to chapter 5 text.
% Auto-generated by scripts/aggregate_section_5_2.py
% Per-cycle pooled per-subtype CHM rates with B1 reference column.
% Pool = all 1500 samples (5 benches x 300) per cycle.
% B1 column from outputs/baselines_phase1e/chm_comparison.json (zero_shot).
\begin{table}[!htbp]
\centering\small
\setlength{\tabcolsep}{5pt}
\begin{tabular}{lccccccc}
\toprule
\textbf{Subtype} & \textbf{B1 ref} & \textbf{C0} & \textbf{C1} & \textbf{C2} & \textbf{C3} & \textbf{C4} & \textbf{C5} \\
\midrule
Conf.\ confab. & 0.327 & 0.088 & 0.238 & 0.259 & 0.291 & 0.289 & 0.307 \\
Fact.\ fabric. & 0.381 & 0.361 & 0.327 & 0.314 & 0.319 & 0.319 & 0.316 \\
Logic.\ fabric. & 0.307 & 0.281 & 0.265 & 0.245 & 0.278 & 0.275 & 0.238 \\
Def.\ evasion & 0.019 & 0.044 & 0.009 & 0.006 & 0.002 & 0.002 & 0.003 \\
Templ.\ leak & 0.000 & 0.000 & 0.000 & 0.000 & 0.000 & 0.000 & 0.000 \\
False refusal & 0.055 & 0.298 & 0.069 & 0.057 & 0.054 & 0.053 & 0.065 \\
Over-long & 0.005 & 0.000 & 0.001 & 0.003 & 0.005 & 0.005 & 0.004 \\
\midrule
\textbf{CHM (mean of 7)} & \textbf{0.156} & 0.153 & 0.130 & 0.126 & 0.135 & 0.135 & 0.133 \\
\bottomrule
\end{tabular}
\caption[The seven hallucination subtypes by cycle.]{The seven subtypes of the composite hallucination metric (\Cref{sec:setup-metrics}) by cycle, pooled over the evaluation fold ($1{,}500$ questions), with B1 for reference.}
\label{tab:halluc-subtypes}
\end{table}


\FloatBarrier

\subsubsection{Training versus transfer contrast}
\label{sec:headline-id-vs-transfer}

On the training panel, whose benchmarks feed the calibration fold and the training pool, CAEM gains a little and then loses it: EM $0.554$ (B1) against $0.581$ at C2 and $0.549$ at C5, and CHM $0.126$ against $0.116$ and $0.130$. The transfer panel, which the system never trains on, shows no gain at any cycle: TruthfulQA EM is $0.253$ at C5 against $0.250$ for B1, with a higher CHM ($0.168$ against $0.147$), and StrategyQA EM $0.640$ against $0.643$, with a lower CHM ($0.113$ against $0.137$). Self-improvement therefore did not damage the transfer benchmarks, but it did not carry over to them either.

\subsection{Four-outcome decision distribution}
\label{sec:decision-distribution}

The three readings in this subsection use three different folds, named in each paragraph; the second and third tables are in \Cref{app:extra-tables}.

\textbf{On the calibration fold} (\Cref{tab:decision-distribution}), the share stored rises from $42$ percent at C1 to $59$ percent at C5, deferral from $12$ to $29$ percent, while discard falls from $25$ to $3$ percent and abstention from $21$ to $10$ percent. The outcomes are ordered by correctness at every cycle: stored answers are right $73$ to $75$ percent of the time, deferred ones less, discarded ones less again, and abstentions least ($0.38$ at C1, $0.17$ at C5). The gap between stored and the next outcome widens from $14$ points at C1 to between $24$ and $30$ points afterwards.

% Table sizing standardised 2026-10-17 (final drafting pass): \small, no \resizebox.
% Caption shortened by hand in the final drafting pass (2026-10); detail moved to chapter 5 text.
% Auto-generated by scripts/aggregate_section_5_3.py
% Hand-corrected 2026-10-15 (see ref_check/fixes_applied.md)
% Panel (b) source: outputs/full_run/cycle_3/calibration/{fever,triviaqa,commonsense_qa}_cycle3.json,
% same computation as panel (a) (samples without a gate decision excluded: 6/1/4).
% Per-cycle four-outcome decision distribution on the training-panel calibration fold.
% Cycles 1,2,3,5 only (C0 is cold-start before first calibration; C4 was aborted under retention guard).
% Per-benchmark snapshot reported at the cycle-three anchor.
\begin{table}[!htbp]
\centering\small
\setlength{\tabcolsep}{3pt}
\begin{tabular}{lcccccccc}
\toprule
 & \multicolumn{2}{c}{\textbf{Store}} & \multicolumn{2}{c}{\textbf{Defer}} & \multicolumn{2}{c}{\textbf{Discard}} & \multicolumn{2}{c}{\textbf{Abstain}} \\
\cmidrule(lr){2-3}\cmidrule(lr){4-5}\cmidrule(lr){6-7}\cmidrule(lr){8-9}
 & share & EM & share & EM & share & EM & share & EM \\
\midrule
\multicolumn{9}{@{}l}{\textit{(a) Pooled training panel}} \\
C1 & 42.1\% & 0.74 & 12.0\% & 0.60 & 25.0\% & 0.58 & 20.9\% & 0.38 \\
C2 & 67.2\% & 0.74 & 18.7\% & 0.44 & 5.4\% & 0.38 & 8.7\% & 0.25 \\
C3 & 60.4\% & 0.75 & 26.0\% & 0.47 & 5.5\% & 0.41 & 8.1\% & 0.22 \\
C5 & 58.6\% & 0.73 & 28.5\% & 0.49 & 3.4\% & 0.46 & 9.5\% & 0.17 \\
\midrule
\multicolumn{9}{@{}l}{\textit{(b) By benchmark at C3 ($500$ each)}} \\
FEVER & 55.1\% & 0.76 & 38.1\% & 0.47 & 4.3\% & 0.52 & 2.6\% & 0.38 \\
TriviaQA & 40.5\% & 0.79 & 28.9\% & 0.45 & 10.6\% & 0.30 & 20.0\% & 0.17 \\
CommonsenseQA & 85.7\% & 0.72 & 11.1\% & 0.55 & 1.6\% & 0.88 & 1.6\% & 0.62 \\
\bottomrule
\end{tabular}
\caption[Decision shares on the calibration fold.]{Share of each decision and the EM within it on the training-panel calibration fold: (a) pooled by cycle; (b) per benchmark at cycle three. C0 and the rolled-back C4 have no reading.}
\label{tab:decision-distribution}
\end{table}


\textbf{On the stream chunks and the transfer evaluation folds} (\Cref{tab:decision-matrix-trajectory}), stored answers are more often correct than deferred ones in every row. The full ordering stored, deferred, discarded holds in half the rows; where it fails, the discard class is small. On TruthfulQA, stored answers are right only $31$ to $38$ percent of the time at C1 to C5, far below the training panel.

\textbf{On the evaluation fold} (\Cref{tab:abstention-per-bench}), FEVER abstention falls from $64.3$ percent at C0 to $5.3$ percent at C5; CommonsenseQA almost never abstains, because a multiple-choice prompt forces an answer; TruthfulQA abstains on $23$ to $33$ percent of questions, a rate that falls slightly over the cycles while its EM stays flat; StrategyQA swings between $2$ and $43$ percent.

\FloatBarrier

\subsection{Three-tier router efficiency}
\label{sec:tier-efficiency}

\subsubsection{Per-tier hit rate trajectory and per-tier exact match}
\label{sec:tier-hit}

The design expects the memory-direct share to grow as the memory fills and the zero-shot share to grow as fine-tuning folds retrieved knowledge into the generator. Neither happened on the evaluation fold (\Cref{tab:tier-trajectory}). Tier 1 answered $7$ of the $9{,}000$ evaluation queries over six cycles, because every stream chunk and the evaluation fold share no question with the stored episodes; all seven answers were correct, and a direct probe of the router answered $42$ of $50$ exact repeats and $19$ of $50$ paraphrases of stored questions, $41$ and $19$ of them correctly (\Cref{tab:tier1-probe}). Tier 1 is built for deployment, where users repeat questions; this study used fresh questions to measure learning, so the probe is the test of this tier. Because a repeated question receives an answer that has already passed the storage gate, both exact match and the composite hallucination metric are expected to improve in deployment as the share of repeated questions grows. The zero-shot share moved from $48.7$ percent at C0 to $50.2$ percent at C5, peaking at $55.2$ percent at C2.

Within tiers, zero-shot answers are more accurate than retrieval-augmented ones at every cycle ($0.565$ against $0.472$ at C0, $0.519$ against $0.495$ at C5; \Cref{tab:tier-chm-trajectory} in \Cref{app:extra-tables}). At C0 this holds on four of five benchmarks, with TriviaQA the exception (\Cref{tab:per-tier-em}). Part of the difference is selection: the router and the classifier send to Tier 2 the questions they judge the generator can already answer. The zero-shot tier's accuracy also falls over the cycles, which is where the self-improvement loop acts.

% Auto-generated by scripts/build_ch5_evidence_tables.py from outputs/tier1_routing_probe/
% (results_cycle5.json, em_analysis.json). The probe's top1_is_source_rate field is not used:
% it compares two different id spaces (see ref_check/REPO_CLEANUP_TRACKER.md, C2).
\begin{table}[!htbp]
\centering\small
\begin{tabular}{@{}lrrrrrr@{}}
\toprule
 & \multicolumn{3}{c}{\textbf{Exact repeats}} & \multicolumn{3}{c}{\textbf{Paraphrases}} \\
\cmidrule(lr){2-4}\cmidrule(l){5-7}
\textbf{Source benchmark} & Asked & Tier 1 & Correct & Asked & Tier 1 & Correct \\
\midrule
FEVER & $17$ & $13$ & $13$ & $17$ & $7$ & $7$ \\
TriviaQA & $23$ & $22$ & $22$ & $23$ & $12$ & $12$ \\
CommonsenseQA & $10$ & $7$ & $6$ & $10$ & $0$ & $0$ \\
\midrule
\textbf{Total} & $50$ & $42$ & $41$ & $50$ & $19$ & $19$ \\
\bottomrule
\end{tabular}
\caption[Router probe of the memory-direct tier.]{Router probe of the memory-direct tier at cycle five: $50$ stored questions asked again word for word and $50$ asked as paraphrases. Tier 1 counts the questions answered from memory; Correct counts those answers that match the gold answer.}
\label{tab:tier1-probe}
\end{table}


% Generated by scripts/build_final_tables.py from outputs/analysis/chm_judged.json (label em_final); do not edit.
\begin{table}[!htbp]
\centering\small
\setlength{\tabcolsep}{4pt}
\begin{tabular}{@{}lccccccc@{}}
\toprule
  & \multicolumn{2}{c}{\twoline{Tier 1\\(memory)}} & \multicolumn{2}{c}{\twoline{Tier 2\\(zero-shot)}} & \multicolumn{2}{c}{\twoline{Tier 3\\(RAG)}} & \\
 & share & EM & share & EM & share & EM & \twoline{Pooled\\EM} \\
\midrule
\multicolumn{8}{@{}l}{\textit{(a) Pooled over the five benchmarks}} \\
C0 & 0.0\% & -- & 48.7\% & 0.565 & 51.3\% & 0.472 & 0.517 \\
C1 & 0.0\% & -- & 49.8\% & 0.566 & 50.2\% & 0.501 & 0.533 \\
C2 & 0.1\% & 1.000 & 55.2\% & 0.548 & 44.7\% & 0.517 & 0.535 \\
C3 & 0.1\% & 1.000 & 44.7\% & 0.532 & 55.1\% & 0.502 & 0.516 \\
C4 & 0.1\% & 1.000 & 45.6\% & 0.532 & 54.3\% & 0.502 & 0.517 \\
C5 & 0.1\% & 1.000 & 50.2\% & 0.519 & 49.7\% & 0.495 & 0.508 \\
\midrule
\multicolumn{8}{@{}l}{\textit{(b) By benchmark at C5 ($300$ each)}} \\
FEVER & 0.0\% & -- & 43.3\% & 0.500 & 56.7\% & 0.459 & 0.477 \\
TriviaQA & 0.3\% & 1.000 & 37.0\% & 0.423 & 62.7\% & 0.468 & 0.453 \\
CommonsenseQA & 0.0\% & -- & 82.3\% & 0.745 & 17.7\% & 0.585 & 0.717 \\
TruthfulQA & 0.3\% & 1.000 & 64.3\% & 0.254 & 35.3\% & 0.245 & 0.253 \\
StrategyQA & 0.0\% & -- & 24.0\% & 0.639 & 76.0\% & 0.640 & 0.640 \\
\bottomrule
\end{tabular}
\caption[Tier shares and per-tier EM by cycle.]{Share of evaluation queries sent to each tier and the EM within it: (a) pooled by cycle; (b) per benchmark at C5.}
\label{tab:tier-trajectory}
\end{table}


\subsubsection{Per-tier latency and cost}
\label{sec:tier-cost}

The mean time per evaluation query, from the batched evaluation log, stays between $13.6$ and $14.7$ seconds per query (about six hours per cycle for the $1{,}500$ evaluation questions), as expected when the tier shares hardly move. This is about a hundred times slower than B1 ($142$ milliseconds), because every generated answer passes through the verifier's resampling, retrieval and entailment steps. Per benchmark at C5, CommonsenseQA is fastest ($8.6$ seconds) and TruthfulQA slowest ($17.8$ seconds). The time per tier could not be separated from the batched log, so the deployment-cost model of \Cref{eq:deployment-cost} cannot be filled in tier by tier; given the flat tier shares, it would predict a flat cost in any case.

\subsubsection{Why the safety override fires}
\label{sec:tier-safety}

The safety override fired on $63$ of the $9{,}000$ evaluation queries ($0.7$ percent) and on $3$ of the $23{,}500$ stream-chunk queries. Of the evaluation firings, $62$ were CommonsenseQA questions after that benchmark's floor was refitted from $0.38$ to $0.491$ at cycle three and $0.569$ at cycle five; the other was one StrategyQA question at cycle five. On FEVER and TriviaQA the override never fired, because their pre-routing temperature reached its upper bound and pulled $\upre$ toward $0.5$ (\Cref{sec:pre-routing}). The override therefore acted as a rarely used safeguard in this run, not as an active part of routing.

\FloatBarrier

\subsection{Retention and continual learning}
\label{sec:retention}

The retention guard reads three probes after every fine-tune (\Cref{tab:sil-training-trajectory}). It kept cycles one, two, three and five, and rolled back cycle four, when the TriviaQA test ratio fell to $0.886$; the memory was kept and grew across the rollback. On the benchmarks themselves, CommonsenseQA lost $4.0$ points and FEVER $1.0$ point from C0 to C5, and both fell further from their peaks (\Cref{tab:per-bench-trajectory}); the probe panel did not catch these losses because the CommonsenseQA probe stayed above its floor and FEVER is not probed. The training set grew from $107$ episodes at the first fine-tune to $5{,}006$ at the last, and the final training loss fell from $0.258$ to $0.070$ over cycles one to four, then rose to $0.091$ at cycle five after training restarted from the cycle-three adapter (\Cref{tab:sil-training-trajectory}).

% Caption shortened by hand in the final drafting pass (2026-10); detail moved to chapter 5 text.
% Auto-generated from outputs/full_run/retroverify_cycle*.json::fine_tune
% + outputs/full_run/all_cycle_results.json (multi-modal retention probes).
% Per-cycle SIL training trajectory: episode count, final training loss,
% three-probe retention panel (MMLU + TQA-test + CSQA-test), rollback status.
\begin{table}[H]
\centering\footnotesize
\setlength{\tabcolsep}{6pt}
\resizebox{\textwidth}{!}{%
\begin{tabular}{@{}lrrrrrl@{}}
\toprule
 & & & \multicolumn{3}{c}{\textbf{Retention probe ratio}} & \\
\cmidrule(lr){4-6}
\textbf{Cycle} & $\boldsymbol{n_{\text{episodes used}}}$ & \textbf{Final loss} & \textbf{MMLU} & \textbf{TQA-test} & \textbf{CSQA-test} & \textbf{Outcome} \\
\midrule
C1 & $107$ & $0.258$ & $1.041$ & $1.013$ & $1.006$ & commit \\
C2 & $1\,474$ & $0.171$ & $0.975$ & $0.949$ & $0.982$ & commit \\
C3 & $2\,966$ & $0.109$ & $0.992$ & $0.987$ & $0.976$ & commit \\
C4 & $3\,624$ & $0.070$ & $1.018$ & $\mathbf{0.886}^{\dagger}$ & $0.982$ & abort (retention rollback) \\
C5 & $5\,006$ & $0.091$ & $0.999$ & $1.013$ & $0.958$ & commit (post-rollback recovery) \\
\bottomrule
\multicolumn{7}{@{}l}{\footnotesize $^{\dagger}$ Cycle-four TriviaQA-test ratio of $0.886$ crossed below the registered floor of $\rho_{\min} = 0.93$, triggering the asymmetric rollback.} \\
\end{tabular}%
}
\caption[Self-improvement training by cycle.]{Episodes used for fine-tuning, final training loss and retention ratios by cycle.}
\label{tab:sil-training-trajectory}
\end{table}


\FloatBarrier

\subsection{Memory and training-pool dynamics}
\label{sec:pool-dynamics}

\textbf{Memory growth:} the memory grew from the $431$ gold-filtered cold-start entries to $12{,}175$ at C5, adding between $2{,}214$ and $2{,}484$ entries per cycle, including the rolled-back cycle four; the deferred buffer grew to $4{,}284$ entries (\Cref{tab:memory-store-trajectory}).

% Hand-merged 2026-10-17 (final drafting pass, ch5 table reduction) from three committed tables:
% tab_memory_store_trajectory.tex (memory, deferred buffer, change), tab_survival_distribution.tex
% (promoted, TTL-dropped) and tab_consolidation_trajectory.tex (groups merged, entries removed).
% Values copied unchanged; logged in ref_check/fixes_applied.md and REPO_CLEANUP_TRACKER.md (T10).
\begin{table}[!htbp]
\centering\small
\setlength{\tabcolsep}{4pt}
\begin{tabular}{@{}lrrrrrrr@{}}
\toprule
 & \multicolumn{2}{c}{\textbf{Memory}} & \multicolumn{3}{c}{\textbf{Deferred buffer}} & \multicolumn{2}{c}{\textbf{Consolidation}} \\
\cmidrule(lr){2-3}\cmidrule(lr){4-6}\cmidrule(lr){7-8}
\textbf{Cycle} & \textbf{Entries} & \textbf{Change} & \textbf{Entries} & \textbf{Promoted} & \textbf{Dropped} & \textbf{Merged} & \textbf{Removed} \\
\midrule
C0 & $431$ & --- & $0$ & --- & --- & --- & --- \\
C1 & $2\,818$ & $+2\,387$ & $1\,123$ & $0$ & $0$ & $0$ & $0$ \\
C2 & $5\,151$ & $+2\,333$ & $2\,399$ & $131$ & $0$ & $3$ & $3$ \\
C3 & $7\,477$ & $+2\,326$ & $2\,692$ & $220$ & $924$ & $11$ & $11$ \\
C4$^{\ast}$ & $9\,961$ & $+2\,484$ & $4\,068$ & --- & --- & --- & --- \\
C5 & $12\,175$ & $+2\,214$ & $4\,284$ & $299$ & $1\,169$ & $43$ & $44$ \\
\midrule
Total & & & & $650$ & $2\,093$ & $57$ & $58$ \\
\bottomrule
\end{tabular}
\caption[Memory, deferred buffer and consolidation by cycle.]{Memory and deferred-buffer size at the close of each cycle, answers promoted from or dropped by the buffer, and near-duplicate groups merged. $^{\ast}$Cycle four was rolled back, so its reconsideration and consolidation passes were skipped.}
\label{tab:memory-store-trajectory}
\end{table}


\textbf{Training pool:} the training threshold ($\ustored \ge 0.70$) keeps a purer subset of the stored answers (\Cref{tab:storage-pool-quality}). On the stream chunks, the training pool is $79.8$ to $81.8$ percent correct, against $72.4$ to $75.3$ percent for everything stored, so the stricter threshold removes between $56$ and $73$ percent of the wrong stored answers before training. About a fifth of the training pool is still wrong. Stored scores are well ordered against correctness on the stream chunks at every cycle, but on the evaluation fold only at cycles one and four (\Cref{tab:u-stored-distribution} in \Cref{app:extra-tables}).

% Table sizing standardised 2026-10-17 (final drafting pass): \small, no \resizebox.
% Caption shortened by hand in the final drafting pass (2026-10); detail moved to chapter 5 text.
% Auto-generated from outputs/full_run/eval/{bench}_cycle{c}.json,
% outputs/full_run/eval/{bench}_cycle{c}_streamchunk.json, retroverify_cycle*.json
% Stored-episode em-correctness, training-threshold filter, SIL pool purity, and
% SIL training-batch size per cycle, reported on both the held-out evaluation fold
% (panel a) and the SIL stream chunks (panel b, the production candidate pool).
% Hand-corrected 2026-10-15 (see ref_check/fixes_applied.md)
\begin{table}[!htbp]
\centering\small
\setlength{\tabcolsep}{4pt}
\begin{tabular}{@{}lrrrrrrrr@{}}
\toprule
\textbf{Cycle} & $\boldsymbol{n_{\text{stored}}}$ & \textbf{Correct} & \textbf{Wrong} & $\boldsymbol{\pi_{\text{store}}}$ & \twoline{Above\\$\boldsymbol{\tau_{\text{train}}}$} & \twoline{Share\\above} & $\boldsymbol{\pi_{\text{SIL}}}$ & \twoline{Used in\\training} \\
\midrule
\multicolumn{9}{@{}l}{\textit{(a) Evaluation fold}} \\
C0 & $336$ & $255$ & $81$  & $0.759$ & $58$  & $17.3\%$ & $0.793$ & --- \\
C1 & $515$ & $374$ & $141$ & $0.726$ & $231$ & $44.9\%$ & $0.788$ & --- \\
C2 & $534$ & $370$ & $164$ & $0.693$ & $266$ & $49.8\%$ & $0.726$ & --- \\
C3 & $565$ & $402$ & $163$ & $0.712$ & $279$ & $49.4\%$ & $0.735$ & --- \\
C4$^{\ast}$ & $566$ & $402$ & $164$ & $0.710$ & $272$ & $48.1\%$ & $0.735$ & --- \\
C5 & $504$ & $353$ & $151$ & $0.700$ & $181$ & $35.9\%$ & $0.735$ & --- \\
\midrule
\multicolumn{9}{@{}l}{\textit{(b) Stream chunks}} \\
C1 & $2\,456$ & $1\,850$ & $606$ & $0.753$ & $1\,309$ & $53.3\%$ & $0.818$ & $107$ \\
C2 & $2\,494$ & $1\,873$ & $621$ & $0.751$ & $1\,448$ & $58.1\%$ & $0.810$ & $1\,474$ \\
C3 & $2\,490$ & $1\,802$ & $688$ & $0.724$ & $1\,337$ & $53.7\%$ & $0.798$ & $2\,966$ \\
C4$^{\ast}$ & $2\,522$ & $1\,862$ & $660$ & $0.738$ & $1\,413$ & $56.0\%$ & $0.803$ & --- \\
C5 & $2\,380$ & $1\,724$ & $656$ & $0.724$ & $895$ & $37.6\%$ & $0.803$ & $5\,006$ \\
\bottomrule
\end{tabular}
\par\vspace{2pt}{\footnotesize\parbox{\linewidth}{$^{\ast}$Cycle four was rolled back, so its training pool is not reported.}\par}
\caption[Stored answers and the training pool by cycle.]{Stored answers, their precision, the share passing the training threshold ($\ustored \ge 0.70$) and the precision of that training pool, on (a) the evaluation fold and (b) the stream chunks.}
\label{tab:storage-pool-quality}
\end{table}


\textbf{Re-verification and deferral:} reconstructing the memory snapshots identifies $1{,}094$ entries removed by the retroactive pass (the run log counts $1{,}061$); $38.4$ percent of them were wrong, against about $22$ percent in the pool, so the pass removes wrong entries at about $1.75$ times their base rate (\Cref{tab:em-conditional-breakdown} in \Cref{app:extra-tables}). The deferred buffer promoted $650$ answers ($131$, $220$ and $299$ at cycles two, three and five, \Cref{tab:memory-store-trajectory}); of the $620$ promotions recovered from the snapshots, $55.6$ percent were correct, which is below the stored pool but means about $345$ correct answers rescued from the deferral band. Consolidation merged $57$ groups and removed $58$ entries (\Cref{tab:memory-store-trajectory}); its same-answer guard blocked several merges of near-identical questions with different answers.

\FloatBarrier

\subsection{Theoretical predictions checked against the trajectory}
\label{sec:theory-checks}

This subsection reads the statements of \Cref{sec:theory} against the run. The two proven results are checked against their assumptions; the heuristics and remarks are read as consistency checks. \Cref{tab:theorem-receipts-summary} summarises the readings.

% Auto-generated by scripts/build_ch5_evidence_tables.py (final drafting pass, 2026-10).
\begin{table}[!htbp]
\centering\small
\setlength{\tabcolsep}{4pt}
\begin{tabular}{@{}lrrrrrrrrr@{}}
\toprule
 & \multicolumn{6}{c}{\textbf{Pooled training panel}} & \multicolumn{3}{c}{\textbf{Balanced accuracy}} \\
\cmidrule(lr){2-7}\cmidrule(l){8-10}
\textbf{Cycle} & \textbf{Stored} & \textbf{Precision} & \textbf{TPR} & \textbf{FPR} & \textbf{Bal.\ acc.} & \textbf{Cohen's $d$} & \textbf{FEVER} & \textbf{TriviaQA} & \textbf{CSQA} \\
\midrule
C0$^{\ast}$ & $560$ (37.3\%) & $0.780$ & $0.653$ & $0.148$ & $0.753$ & $1.468$ & -- & -- & -- \\
C1 & $630$ (42.0\%) & $0.740$ & $0.510$ & $0.280$ & $0.615$ & $0.593$ & $0.665$ & $0.500$ & $0.522$ \\
C2 & $1001$ (66.7\%) & $0.736$ & $0.791$ & $0.465$ & $0.663$ & $0.773$ & $0.672$ & $0.689$ & $0.554$ \\
C3 & $899$ (59.9\%) & $0.746$ & $0.725$ & $0.397$ & $0.664$ & $0.756$ & $0.653$ & $0.719$ & $0.531$ \\
C5 & $868$ (57.9\%) & $0.732$ & $0.705$ & $0.389$ & $0.658$ & $0.842$ & $0.633$ & $0.687$ & $0.550$ \\
\bottomrule
\end{tabular}
\caption[Storage gate quality on the calibration fold by cycle.]{Storage gate quality on the training-panel calibration fold ($1{,}500$ questions) by cycle. $^{\ast}$C0 is in-sample (the composite was fitted on this fold); C1 to C5 score the fold with the new adapter before the composite is refitted. C4 was rolled back and has no reading.}
\label{tab:gate-quality}
\end{table}


\subsubsection{Pool purity}
\label{sec:check-purity}

\textbf{Pool purity floor:} the floor of \Cref{thm:purity} depends on new entries being as often correct as the calibration fold at the same cut-point (assumption (ii)). On the training panel this approximately held: the calibration-fold precision at $\tau_{\text{store}} = 0.60$ was between $0.732$ and $0.746$ at C1 to C5 (\Cref{tab:pi-store-trajectory}), and the answers actually stored from each stream chunk were $0.753$, $0.751$, $0.724$, $0.738$ and $0.724$ correct. On the evaluation fold the pooled storage precision was between $0.693$ and $0.756$; FEVER's own evaluation-fold precision was lower, between $0.614$ and $0.634$ at C1 to C4 and $0.649$ at C5, which shows that the pooled statement does not carry over to every benchmark. On TruthfulQA, where the questions are unlike the calibration fold, the assumption failed: stored answers were right only $31$ to $38$ percent of the time.

% C0 cal-fold cell and caption hand-corrected in the final drafting pass (2026-10).
% Auto-generated by scripts/aggregate_pi_store_trajectory.py
% Per-cycle empirical storage precision π_store at τ_store = 0.60
% Reports both calibration-fold (theorem-anchored) and evaluation-fold
% (deployment-realised) readings; cycle four has no cal-fold refit but the
% eval-fold reading is computed against the rolled-back C3 composite.
\begin{table}[H]
\centering\footnotesize
\setlength{\tabcolsep}{4pt}
\resizebox{\textwidth}{!}{%
\begin{tabular}{@{}l|cc|cc|ccc@{}}
\toprule
 & \multicolumn{2}{c|}{\textbf{Cal-fold pooled}} & \multicolumn{2}{c|}{\textbf{Eval-fold pooled}} & \multicolumn{3}{c}{\textbf{Eval-fold per-bench $\pi_{\text{store}}$}} \\
\textbf{Cycle} & $n$ & $\pi_{\text{store}}$ & $n$ & $\pi_{\text{store}}$ & FEVER & TriviaQA & CSQA \\
\midrule
C0 & $560$ & $0.780^{\dagger}$ & $336$ & $0.759$ & $0.661$ & --- & $0.779$ \\
C1 & $631$ & $0.740$ & $515$ & $0.726$ & $0.634$ & $0.797$ & $0.736$ \\
C2 & $1\,001$ & $0.736$ & $534$ & $0.693$ & $0.630$ & $0.733$ & $0.705$ \\
C3 & $899$ & $0.746$ & $565$ & $0.712$ & $0.617$ & $0.741$ & $0.750$ \\
C4$^{\ast}$ & --- & --- & $566$ & $0.710$ & $0.614$ & $0.733$ & $0.754$ \\
C5 & $868$ & $0.732$ & $504$ & $0.700$ & $0.649$ & $0.704$ & $0.725$ \\
\bottomrule
\multicolumn{8}{@{}l}{\footnotesize $^{\dagger}$ In-sample: the cycle-zero composite was fitted on this fold (re-run sweep, locked configuration).} \\
\multicolumn{8}{@{}l}{\footnotesize $^{\ast}$ Cycle four was aborted under the retention guard; the cal-fold refit was skipped, but the eval-fold reading is computed against the rolled-back C3 composite.} \\
\end{tabular}%
}
\caption[Storage precision by cycle on the calibration and evaluation folds.]{Storage precision $\pi_{\text{store}}$ at $\tau_{\text{store}} = 0.60$ by cycle: share of stored answers that are correct, on the training-panel calibration fold and on the evaluation fold.}
\label{tab:pi-store-trajectory}
\end{table}


\textbf{Purification condition (H1):} the gate purifies only if its balanced accuracy exceeds one half (\Cref{thm:bayes-purity}). Pooled over the training panel it was $0.615$, $0.663$, $0.664$ and $0.658$ at C1, C2, C3 and C5 (\Cref{tab:gate-quality}), measured with the new adapter before each refit; the C0 value of $0.753$ is in-sample. The weakest benchmarks were CommonsenseQA ($0.52$ to $0.55$) and TriviaQA at C1, where the gate stored no TriviaQA answer and the balanced accuracy is exactly $0.500$.

\subsubsection{Monotonicity}
\label{sec:check-monotonicity}

\Cref{thm:monotone} predicts that better separation stores more answers when the correct class is denser at the cut-point. The separation $d$ rose from $0.593$ at C1 to $0.773$ at C2 while the share stored rose from $0.420$ to $0.667$, as predicted; from C2 to C3, $d$ was flat ($0.756$) and the share fell to $0.599$; from C3 to C5, $d$ rose to $0.842$ while the share fell to $0.579$ (\Cref{tab:gate-quality}). Only the first step agrees. The refit of the composite at every cycle, which the heuristic does not model, is the most likely reason; separating the two would need a cycle with the composite frozen (\Cref{ch:conclusion}).

\subsubsection{Convergence}
\label{sec:check-convergence}

\textbf{Convergence envelope (H3):} the calibration-fold precision stayed between $0.732$ and $0.746$ at C1 to C5, a band of width $0.014$ (\Cref{tab:pi-store-trajectory}). For about $900$ stored answers per cycle, the standard error of one reading is about $0.015$, so the band is as narrow as sampling noise allows; six cycles cannot estimate a convergence rate.

\textbf{Repeated filtering (H2):} \Cref{thm:bayes-convergence} describes purity rising under repeated filtering by a fixed gate. The calibration-fold precision was $0.780$ at C0 (in-sample), then $0.740$, $0.736$, $0.746$ and $0.732$: it fell at C1, C2 and C5. H2 is therefore not supported as registered; the changes are within sampling noise, and purity stayed above the floor of \Cref{thm:purity} throughout.

\textbf{Equilibrium under rollback:} parts (i) and (iii) of \Cref{cor:stochastic-equilibrium} hold. The commit pattern over cycles one to five was $1, 1, 1, 0, 1$ (commit rate $0.80$), with worst probe ratios $1.006$, $0.949$, $0.976$, $0.886$ and $0.958$, and the memory grew at every cycle. Part (iv) cannot be tested in five cycles.

\subsubsection{Asymptotic hallucination floor}
\label{sec:check-asymptotic}

\Cref{cor:tier1-floor} bounds the error rate of Tier 1 answers by the impurity of the pool. On the evaluation fold Tier 1 gave $7$ answers, all correct, which is consistent but too few to estimate a rate; the router probe gives the testable case (\Cref{sec:tier-hit}). Per tier, CHM is $0.121$ to $0.129$ for zero-shot answers and $0.131$ to $0.181$ for retrieval-augmented ones (\Cref{tab:tier-chm-trajectory} in \Cref{app:extra-tables}). The decomposition of \Cref{thm:asymptotic-elim} cannot be tested here, because Tier 1 answered at most $0.1$ percent of evaluation queries.

\subsubsection{Heuristic checks: equilibrium structure and self-correction}
\label{sec:check-corollaries}

\Cref{cor:self-correction} predicts that repeated re-verification removes wrong entries over time. The share of checked entries removed fell from $13.0$ percent at C1 to $9.2$, $6.6$ and $3.7$ percent at C2, C3 and C5 (\Cref{tab:self-correction-check} in \Cref{app:extra-tables}), and removed entries were wrong more often than the pool (\Cref{sec:pool-dynamics}), both consistent with the heuristic; the falling rate is partly diluted by the growing number of fresh entries. The entry records do not count how many passes each entry survived, so the predicted survival distribution could not be measured. The corpus ceiling of \Cref{cor:corpus-floor} was not measured.

% Originally generated by scripts/aggregate_section_5_6.py; rewritten by hand in the final drafting pass
% (2026-10, see ref_check/fixes_applied.md "Chapter 5") to the chapter-4 statuses and the locked-composite numbers.
\begin{table}[!htbp]
\centering\scriptsize
\setlength{\tabcolsep}{3pt}
\renewcommand{\arraystretch}{1.05}
\begin{tabular}{@{}>{\raggedright\arraybackslash}p{3.0cm}>{\raggedright\arraybackslash}p{1.5cm}>{\raggedright\arraybackslash}p{2.0cm}>{\raggedright\arraybackslash}p{5.6cm}>{\raggedright\arraybackslash}p{2.3cm}@{}}
\toprule
\textbf{Statement} & \textbf{Status} & \textbf{Reading} & \textbf{What the run shows} & \textbf{Evidence} \\
\midrule
Pool purity floor (\Cref{thm:purity}) & proven & holds on the training panel & New-entry precision on the stream chunks $0.724$ to $0.753$, close to the calibration-fold precision $0.732$ to $0.746$; on TruthfulQA, where the like-calibration assumption fails, stored answers are $0.31$ to $0.38$ correct. & \Cref{tab:pi-store-trajectory,tab:decision-matrix-trajectory} \\
Purification condition (\Cref{thm:bayes-purity}) & proven & precondition met & Pooled gate balanced accuracy $0.615$ to $0.664$ at C1 to C5; weakest benchmarks CommonsenseQA ($0.52$ to $0.55$) and TriviaQA at C1 ($0.500$). & \Cref{tab:gate-quality} \\
Storage-rate sensitivity (\Cref{thm:monotone}) & heuristic & consistent at one of three steps & C1 to C2: $d$ $0.593 \to 0.773$ and share stored $0.420 \to 0.667$, as predicted; C2 to C3: $d$ flat; C3 to C5: $d$ rises, share falls. & \Cref{tab:gate-quality} \\
Purity rises with accuracy (\Cref{thm:monotone-purification}) & remark & not tested & $\alpha$ changes every cycle and $p$ varies little. & -- \\
Convergence envelope (\Cref{thm:convergence}) & heuristic & consistent (band) & Calibration-fold precision stays within $0.732$ to $0.746$ (width $0.014$) at C1 to C5; a rate cannot be estimated from six cycles. & \Cref{tab:pi-store-trajectory} \\
Repeated filtering gives purity 1 (\Cref{thm:bayes-convergence}) & remark & not met & Calibration-fold precision $0.780$ (C0, in-sample), $0.740$, $0.736$, $0.746$, $0.732$: it falls at C1, C2 and C5 (H2). & \Cref{tab:pi-store-trajectory} \\
Error decomposition (\Cref{thm:asymptotic-elim}) & remark & not testable here & Tier 1 answered at most $0.1$ percent of evaluation queries. & \Cref{tab:tier-trajectory} \\
Tier 1 error ceiling (\Cref{cor:tier1-floor}) & heuristic & consistent & All $7$ Tier 1 answers on the evaluation fold correct; router probe $41$ of $42$ repeats and $19$ of $19$ paraphrases correct. & \Cref{tab:tier-chm-trajectory} \\
Corpus ceiling (\Cref{cor:corpus-floor}) & remark & not measured & Retrieved passages were not stored per question. & -- \\
Self-correction (\Cref{cor:self-correction}) & heuristic & consistent (direction) & Share of checked entries pruned $13$, $9$, $7$, $4$ percent at C1, C2, C3, C5; survival counts not recorded. & \Cref{tab:self-correction-check} \\
Equilibrium under rollback (\Cref{cor:stochastic-equilibrium}) & heuristic & (i), (iii) hold; (iv) not testable & Commit pattern $1, 1, 1, 0, 1$; memory grew at every cycle ($431$ to $12{,}175$). & \Cref{tab:pooled-trajectory,tab:memory-store-trajectory} \\
\bottomrule
\end{tabular}
\caption[What the run shows for each formal statement.]{What the run shows for each formal statement of \Cref{sec:theory}. Heuristic and remark readings are consistency checks, not tests of a proven claim.}
\label{tab:theorem-receipts-summary}
\end{table}


\FloatBarrier

\section{Final Design Adjustments}
\label{sec:final-adjustments}

\subsection{Calibration discipline and disjoint folds}
\label{sec:adj-discipline}

The floor of \Cref{thm:purity} rests on the calibration fold being like the candidate answers, and no sample may appear in both. The six folds of \Cref{tab:benchmark-pools} were therefore built once, checked for overlap by content hash, and never changed; only the training panel enters the calibration fold.

\subsection{The composite parameter selection sweep}
\label{sec:adj-sweep}

\subsubsection{The pre-specified grid}
\label{sec:adj-grid}

Before the main run, $100$ variants were scored: the storage cut-point $\tau_{\text{store}} \in \{0.55, 0.60, 0.65, 0.70\}$ (deferral cut-point $0.45$ throughout), the boost-layer regularisation $C \in \{0.010, 0.025, 0.050, 0.100, 1.000\}$ (smaller is stronger) and the shrinkage strength $\alpha_{\text{shrink}} \in \{0.0, 0.4, 0.6, 0.8, 1.0\}$ ($0$ uses only the pooled curve, $1$ only the benchmark's own). The sweep was re-run with the final code for this thesis; the re-run reproduces the locked composite exactly, and all numbers below come from it.

\subsubsection{Selection criteria}
\label{sec:adj-criteria}

Four criteria were set in advance, applied in order: (1) pooled calibration-fold Cohen's $d$ of at least $0.20$; (2) at least $20$ answers stored on the calibration fold; (3) among the survivors, prefer the variant that stores more correct answers at a similar precision; (4) between equivalent variants, prefer stronger regularisation. Of the $100$ variants, $84$ pass criterion 1 and $83$ pass criteria 1 and 2.

\subsubsection{Top variants}
\label{sec:adj-top}

\Cref{tab:sweep-variants} shows the locked variant and its neighbours. At $\tau_{\text{store}} = 0.60$ the locked variant stores $560$ calibration-fold answers at precision $0.780$, against $281$ at $0.865$ for $\tau_{\text{store}} = 0.65$: the stricter cut stores about half as many answers at higher precision, and $0.60$ was chosen for volume. On the held-out validation fold the locked variant stores $558$ answers at $0.751$. Read literally, criterion 3 would have chosen $\alpha_{\text{shrink}} = 1$, which stores more correct calibration-fold answers ($476$ against $437$); the locked $\alpha_{\text{shrink}} = 0.6$ was chosen instead because pulling small per-benchmark folds toward the pooled fit separates correct from incorrect answers better on the validation fold (pooled $d$ $0.72$ against $0.66$).

% Auto-generated by scripts/build_ch5_evidence_tables.py from outputs/cycle_0/sweep_v2_1/variants.csv.
\begin{table}[!htbp]
\centering\small
\setlength{\tabcolsep}{4pt}
\begin{tabular}{@{}rrrrrrrrrc@{}}
\toprule
 & & & \multicolumn{3}{c}{\textbf{Calibration fold}} & \multicolumn{3}{c}{\textbf{Validation fold}} & \\
\cmidrule(lr){4-6}\cmidrule(lr){7-9}
$\tau_{\text{store}}$ & $C$ & $\alpha_{\text{shrink}}$ & $d$ & Stored & Precision & $d$ & Stored & Precision & Criteria 1--2 \\
\midrule
$0.55$ & $0.01$ & $0.6$ & $1.47$ & $621$ & $0.767$ & $0.72$ & $610$ & $0.741$ & yes \\
$0.60$ & $0.01$ & $0.0$ & $1.26$ & $76$ & $0.961$ & $0.85$ & $60$ & $0.817$ & yes \\
$0.60$ & $0.01$ & $0.4$ & $1.45$ & $393$ & $0.824$ & $0.75$ & $368$ & $0.747$ & yes \\
\rowcolor{gray!18} $0.60$ & $0.01$ & $0.6$ & $1.47$ & $560$ & $0.780$ & $0.72$ & $558$ & $0.751$ & yes \\
$0.60$ & $0.01$ & $0.8$ & $1.48$ & $610$ & $0.769$ & $0.69$ & $596$ & $0.742$ & yes \\
$0.60$ & $0.01$ & $1.0$ & $1.48$ & $620$ & $0.768$ & $0.66$ & $604$ & $0.740$ & yes \\
$0.65$ & $0.01$ & $0.6$ & $1.47$ & $281$ & $0.865$ & $0.72$ & $260$ & $0.758$ & yes \\
$0.70$ & $0.01$ & $0.6$ & $1.47$ & $108$ & $0.972$ & $0.72$ & $84$ & $0.798$ & yes \\
$0.60$ & $1$ & $0.6$ & $0.80$ & $237$ & $0.730$ & $0.73$ & $197$ & $0.772$ & yes \\
\bottomrule
\end{tabular}
\caption[Selected variants of the cycle-zero sweep.]{Selected variants of the $100$-variant cycle-zero sweep (re-run with the locked code; shaded: locked configuration). Calibration fold and held-out validation fold have $500$ questions per training benchmark.}
\label{tab:sweep-variants}
\end{table}


\subsection{The locked configuration}
\label{sec:adj-locked}

The locked configuration ($\tau_{\text{store}} = 0.60$, $C = 0.01$, $\alpha_{\text{shrink}} = 0.6$, together with $\tau_{\text{defer}} = 0.45$, the abstention floor $0.20$, $\tau_{\text{retro}} = 0.50$ and $\rho_{\min} = 0.93$) was carried unchanged into every cycle (\Cref{tab:sweep-top} in \Cref{app:cycle0-catalog}). $C = 0.01$ is the strongest regularisation in the grid, a hundredfold tighter penalty than the library default of $1.0$. \Cref{tab:composite-calibration} in \Cref{app:cycle0-catalog} gives the cycle-zero direction and boost weight of each signal.

\subsection{Per-cycle composite-refit trajectory}
\label{sec:adj-trajectory}

The cut-points stay fixed, but the mapping from signals to probability is refitted at every kept cycle (\Cref{tab:calibration-trajectory} in \Cref{app:extra-tables}). The pooled pre-routing temperature starts at its upper bound ($T = 20.086 \approx e^{3}$), drops to $1.089$ at cycle one and stays between $0.74$ and $0.87$ afterwards; the FEVER and TriviaQA temperatures remain $4.8$ to $10.7$ times that of CommonsenseQA, and CommonsenseQA's safety floor is the only one refitted (to $0.491$ and $0.569$). The expected calibration error after each refit stays between $0.041$ and $0.049$. Of the nine boost weights, two move: the chain-entailment weight rises from $0.067$ to $0.285$ and the question-match weight falls from $0.491$ to about $0.26$; the other seven move by at most $0.07$.

\FloatBarrier

\subsection{Calibration-versus-evaluation gap}
\label{sec:adj-cal-eval-gap}

The earlier split-conformal design was dropped because its guarantees depend on the calibration and evaluation data being alike (\Cref{sec:conformal-gate}). The locked fixed-threshold gate makes no such promise, and the gap can be read directly. On the training panel the calibration-fold precision ($0.732$ to $0.746$) and the evaluation-fold precision ($0.693$ to $0.756$) stay within about five points of each other (\Cref{tab:pi-store-trajectory}).

At cycle zero the gap is large per benchmark (\Cref{tab:validation-gate}). On the held-out validation fold the locked composite stores $93$ FEVER answers ($0.720$ correct), $465$ CommonsenseQA answers ($93$ percent of them, $0.757$ correct) and no TriviaQA answer at all. CommonsenseQA's scores barely separate correct from incorrect answers (\Cref{sec:adj-gate}), so almost everything clears the cut-point; TriviaQA's scores separate well but sit below $0.60$ on this fold. From cycle one the refitted composite stores TriviaQA answers normally. Storing fewer answers on a benchmark is the gate doing what it is asked: it promises precision at the cut-point, not balance across benchmarks. The deferred buffer and re-verification give rejected but correct answers later chances (\Cref{sec:pool-dynamics}).


\FloatBarrier

\subsection{The long-hypothesis judge ablation}
\label{sec:adj-judge-ablation}

A second entailment judge, built on the frozen generator, was meant to handle answers longer than MiniCheck's $512$-token input. It was required to agree with MiniCheck with a Pearson correlation of at least $0.70$ in logit space on a $500$-sample overlap set; it reached $0.584$ and was dropped. Every answer is therefore judged by MiniCheck, with inputs over $512$ tokens truncated; the short-answer prompts of this panel rarely reach that length.

\subsection{The validation gate verdict}
\label{sec:adj-gate}

Before the main run, the locked configuration had to pass a set of checks on the held-out cycle-zero validation fold (\Cref{tab:validation-gate}).
\begin{itemize}
    \item \textbf{Pooled separation:} training-panel composite Cohen's $d$ of at least $0.20$. Passed: $0.717$.
    \item \textbf{Store-minus-discard gap:} at least $0.05$ per benchmark. Computable only on FEVER ($0.268$, passed): TriviaQA stored nothing and CommonsenseQA discarded nothing.
    \item \textbf{Wrong share among stored:} below $0.50$ per benchmark. Passed on FEVER ($0.28$) and CommonsenseQA ($0.243$); not computable on TriviaQA.
\end{itemize}
The registered pooled check therefore passes, but two per-benchmark checks cannot be computed at cycle zero, and CommonsenseQA's own separation is weak ($d = 0.14$). This explains the weak CommonsenseQA gate quality later in the run ($0.52$ to $0.55$, \Cref{tab:gate-quality}). On the transfer panel, shown for information only, $d$ is $0.38$ for TruthfulQA and $0.00$ for StrategyQA. \Cref{tab:cohen-d-per-bench} gives the same separation for each signal.

% Auto-generated by scripts/recompute_validation_gate.py (locked composite; final drafting pass 2026-10).
\begin{table}[!htbp]
\centering\small
\setlength{\tabcolsep}{4pt}
\begin{tabular}{@{}lcrrrrrr@{}}
\toprule
\textbf{Check} & \textbf{Threshold} & \textbf{FEVER} & \textbf{TriviaQA} & \textbf{CSQA} & \textbf{TruthfulQA} & \textbf{StrategyQA} & \textbf{Pooled} \\
\midrule
\multicolumn{8}{@{}l}{\textit{Training panel (registered checks)}} \\
Composite Cohen's $d$ & $\ge 0.20$ & $+0.56$ & $+1.40$ & $+0.14$ & & & $+0.72$ \\
Stored answers & -- & $93$ & $0$ & $465$ & & & $558$ \\
Store-minus-discard gap & $\ge 0.05$ & $+0.27$ & -- & -- & & & \\
Wrong share among stored & $< 0.50$ & $0.28$ & -- & $0.24$ & & & \\
\midrule
\multicolumn{8}{@{}l}{\textit{Transfer panel (diagnostic only)}} \\
Composite Cohen's $d$ & -- & & & & $+0.38$ & $-0.00$ & $+0.42$ \\
\bottomrule
\end{tabular}
\caption[Cycle-zero validation gate under the locked composite.]{Cycle-zero validation gate re-scored with the locked composite on the held-out validation fold ($500$ questions per benchmark). A dash marks a check that cannot be computed because the benchmark stored no answers.}
\label{tab:validation-gate}
\end{table}


% Table sizing standardised 2026-10-17 (final drafting pass): \small, no \resizebox.
% u_stored row (locked composite, scripts/recompute_validation_gate.py) and caption hand-corrected 2026-10.
% Auto-generated by scripts/aggregate_section_5_7.py
% Per-signal Cohen's d on the v2.1 5-bench panel at cycle zero
% Source: outputs/cycle_0/eval/<bench>_cycle0.json — the held-out inference fold
% Note: signals scored under the locked composite; em is the alias-matched user-facing em
% Note: h_norm is a calibration-time-only signal not persisted in the eval JSON, hence "--"
\begin{table}[!htbp]
\centering\small
\setlength{\tabcolsep}{4pt}
\begin{tabular}{@{}lccccccc@{}}
\toprule
\textbf{Signal} & \multicolumn{4}{c}{\textbf{Training panel}} & \multicolumn{3}{c}{\textbf{Transfer panel}} \\
 & FEVER & TriviaQA & CSQA & \textbf{Pooled} & TruthfulQA & StrategyQA & \textbf{Pooled} \\
\midrule
$u_{\text{stored}}$ & +0.56 & +1.40 & +0.14 & \textbf{+0.72} & +0.38 & -0.00 & \textbf{+0.42} \\
$p^{\text{mean}}_{\text{ground}}$ & +0.50 & +1.12 & +0.04 & \textbf{+0.47} & +0.24 & +0.29 & \textbf{+0.11} \\
$p^{\text{max}}_{\text{ground}}$ & +0.58 & +1.15 & +0.02 & \textbf{+0.49} & +0.21 & +0.33 & \textbf{+0.09} \\
$p^{\text{atomic}}_{\text{ground}}$ & +0.50 & +1.12 & +0.04 & \textbf{+0.47} & +0.24 & +0.29 & \textbf{+0.11} \\
$p_{\text{entail}}$ & +0.15 & +1.26 & -0.18 & \textbf{+0.20} & +0.26 & +0.18 & \textbf{+0.19} \\
$q_{a,\text{rel}}$ & -0.37 & +0.51 & +0.35 & \textbf{+0.26} & +0.33 & +0.19 & \textbf{+0.42} \\
$u_{\text{internal}}$ & +0.14 & +0.34 & +0.01 & \textbf{+0.55} & +0.11 & -0.11 & \textbf{+0.20} \\
$u_{\text{token}}$ & +0.21 & -0.09 & -0.15 & \textbf{-0.03} & +0.14 & +0.03 & \textbf{+0.19} \\
$u_{\text{dropout}}$ & -0.10 & -0.37 & -0.30 & \textbf{-0.56} & -0.02 & +0.11 & \textbf{-0.19} \\
$s_{\text{avg}}$ & +0.00 & +0.58 & -0.02 & \textbf{+0.30} & -0.12 & -0.07 & \textbf{-0.13} \\
\bottomrule
\end{tabular}
\caption[Per-signal Cohen's $d$ on the cycle-zero validation fold.]{Per-signal Cohen's $d$ between correct and incorrect answers on the held-out cycle-zero validation fold. Positive values mean the signal is higher on correct answers. The first row is the locked composite.}
\label{tab:cohen-d-per-bench}
\end{table}


\FloatBarrier

\subsection{Failed iterations as evidence of method discipline}
\label{sec:adj-failed-iterations}

The first cycle-zero calibration failed the validation gate because of a bug in the step that writes out the verifier signals: only some of the nine signals were saved, so the composite was fitted on an incomplete signal vector and failed the separation check. The gate caught the bug on the first run; after the fix the calibration was repeated and produced the locked configuration. The failed iteration is reported because it shows the gate doing what it was set up to do.

\section{Statistical Analysis}
\label{sec:significance}

\subsection{Tests used}
\label{sec:sig-pairing}
\label{sec:sig-mcnemar}
\label{sec:sig-bootstrap}
\label{sec:sig-holm}
\label{sec:sig-licensing}

Every baseline answered the same evaluation questions under the comparison protocol of \Cref{sec:setup-protocol}, so each comparison is paired question by question (\Cref{app:sig-pairing}). CAEM is compared with every baseline at C0, C2 and C5.

\textbf{Exact match:} for each benchmark and baseline, a paired McNemar test with continuity correction (Edwards~\cite{edwards1948note}) gives a $p$-value; the $p$-values of the baselines within a benchmark are corrected with the Holm--Bonferroni procedure \cite{holm1979simple}; and a paired bootstrap gives a $95$ percent interval on the EM difference ($2{,}000$ replicates). A contrast is starred when the Holm-corrected $p$-value is below $0.05$ \emph{and} the absolute EM difference is at least $0.02$.

\textbf{Composite hallucination metric:} every question has its own CHM value, the share of the seven failure flags it raises, so this axis is paired in the same way. A Wilcoxon signed-rank test (Wilcoxon~\cite{wilcoxon1945individual}) on the per-question differences gives a $p$-value, with questions that show no difference dropped; the $p$-values are Holm-corrected within each benchmark; and a paired bootstrap gives the interval. A contrast is starred when the corrected $p$-value is below $0.05$ \emph{and} the absolute CHM difference is at least $0.01$. A second CHM test uses only the five subtypes that depend neither on CAEM's composite score nor on its decisions, as a check on the verifier's role (\Cref{sec:disc-threat-verifier}).

The full contingency counts, the Holm sequences and the procedural detail are in \Cref{app:significance}.

\section{Comparisons and Relationships}
\label{sec:comparisons}

\subsection{External baseline panel under matched protocol}
\label{sec:comp-baselines}

The eight baselines are zero-shot (B1), chain of thought (B2), retrieval-augmented generation (B3), chain of thought with retrieval (B4), five-shot chain of thought (B5), vanilla fine-tuning (B6), FLARE active retrieval (B7) and semantic-entropy abstention (B8); \Cref{sec:external-baselines} describes them and \Cref{sec:setup-protocol} the protocol under which they were run. All use the same backbone.

\textbf{Pooled comparison} (\Cref{tab:baseline-pooled}, \Cref{fig:baselines}): the systems fall into three groups. The prompting baselines (B1, B2, B5), semantic-entropy abstention (B8) and CAEM form one cluster, with pooled EM between $0.508$ and $0.535$ and CHM between $0.127$ and $0.137$ (CAEM C0, at $0.155$, sits above it); CAEM at C2 has the highest EM ($0.535$) and the lowest CHM ($0.127$) of the cluster, ahead of five-shot chain of thought ($0.521$, $0.129$), and CAEM at C5 sits inside it ($0.508$, $0.134$). The retrieval baselines (B3, B4, B7) sit apart, with lower EM ($0.445$ to $0.466$) and higher CHM ($0.158$ to $0.196$). Vanilla fine-tuning (B6), trained on gold answers, has an EM inside the cluster ($0.512$) and the lowest CHM of all systems ($0.103$). In the delivered view, CAEM at C0 has the lowest error among answered of any system ($0.405$), at the cost of answering only $66.1$ percent of the questions; at C2 and C5 its error among answered ($0.428$, $0.451$) is close to B8's ($0.427$) and below that of every baseline that answers everything ($0.479$ to $0.555$).

\textbf{Per-benchmark tests} (\Cref{tab:sig-summary}; every contrast with its interval is in \Cref{tab:sig-full-c5,tab:sig-full-c2}): of the $40$ contrasts per cycle ($8$ baselines on $5$ benchmarks), EM gives $9$ significant wins and $3$ defeats at C2 and $5$ wins and $3$ defeats at C5; CHM gives $13$ wins and $5$ defeats at C2 and $10$ wins and $4$ defeats at C5. The wins are almost all against the retrieval baselines and FLARE, and against vanilla fine-tuning on TriviaQA and TruthfulQA. Against the prompting baselines and semantic-entropy abstention (B1, B2, B5, B8), no EM contrast is a significant win at either cycle, and only two CHM contrasts are (B2 and B5 on StrategyQA at C5). The defeats are against vanilla fine-tuning (EM on FEVER and CommonsenseQA; CHM on FEVER, CommonsenseQA and StrategyQA) and against five-shot chain of thought (EM on CommonsenseQA; CHM on TruthfulQA, and at C2 also on CommonsenseQA). At C0, CAEM's CHM is significantly higher than that of most baselines ($17$ defeats, $6$ wins), the cost of its cold-start over-refusal, while its EM already gives $9$ wins and $2$ defeats.

% CHM cells re-scored (Haiku-judged TruthfulQA; mean of the seven active types) by scripts/patch_chm_judged_tables.py (2026-10-17).
% CHM cells re-scored with the Haiku-judged TruthfulQA label by scripts/patch_chm_judged_tables.py (2026-10-17).
% Table sizing standardised 2026-10-17 (final drafting pass): \small, no \resizebox.
% B8 CHM set to n/a by hand 2026-10 (old-run value; tracker D1).
% Caption shortened by hand in the final drafting pass (2026-10); detail moved to chapter 5 text.
% Auto-generated by scripts/aggregate_baseline_pooled.py
% Pooled-EM and pooled-CHM baseline comparison panel with absolute and
% relative deltas against both CAEM C2 (best) and CAEM C5 (final-cycle).
% Sources: outputs/baselines_phase1e/{baseline}/{bench}_cycle0.json (EM),
%          outputs/baselines_phase1e/chm_comparison.json (CHM),
%          outputs/baselines_phase1e/vanilla_ft/eval/{bench}_cycle5.json (B6 EM).
% Paired McNemar significance is reported separately in tab:sig-test.
\begin{table}[!htbp]
\centering\small
\setlength{\tabcolsep}{4pt}
\renewcommand{\arraystretch}{1.1}
\begin{tabular}{@{}llcccccc@{}}
\toprule
 & & & & \multicolumn{2}{c}{\textbf{vs.\ CAEM C2}} & \multicolumn{2}{c}{\textbf{vs.\ CAEM C5}} \\
\cmidrule(lr){5-6}\cmidrule(lr){7-8}
\textbf{\#} & \textbf{System} & \textbf{EM} & \textbf{CHM} & $\boldsymbol{\Delta}$\textbf{EM} & $\boldsymbol{\Delta}$\textbf{CHM} & $\boldsymbol{\Delta}$\textbf{EM} & $\boldsymbol{\Delta}$\textbf{CHM} \\
\midrule
\rowcolor{gray!18}
& \textbf{CAEM C2 (best)} & $\mathbf{0.544}$ & $\mathbf{0.126}$ & --- & --- & $+0.027$ & $-0.007$ \\
\rowcolor{gray!10}
& CAEM C5 (final) & $0.517$ & $0.133$ & $-0.027$ & $+0.007$ & --- & --- \\
\midrule
B1 & zero-shot & $0.507$ & $0.156$ & $+0.037$ & $-0.030$ & $+0.011$ & $-0.023$ \\
B2 & chain-of-thought & $0.500$ & $0.161$ & $+0.044$ & $-0.035$ & $+0.017$ & $-0.028$ \\
B3 & RAG & $0.458$ & $0.172$ & $+0.086$ & $-0.046$ & $+0.059$ & $-0.039$ \\
B4 & CoT $+$ RAG & $0.461$ & $0.165$ & $+0.083$ & $-0.039$ & $+0.056$ & $-0.032$ \\
B5 & five-shot CoT & $0.405$ & $0.155$ & $+0.139$ & $-0.029$ & $+0.112$ & $-0.022$ \\
B6 & vanilla fine-tune & $0.490$ & $0.140$ & $+0.054$ & $-0.014$ & $+0.027$ & $-0.007$ \\
B7 & FLARE & $0.525$ & $0.226$ & $+0.019$ & $-0.100$ & $-0.008$ & $-0.093$ \\
B8 & semantic entropy & $0.505$ & n/a & $+0.039$ & -- & $+0.013$ & -- \\
\bottomrule
\end{tabular}
\caption[CAEM against the eight baselines.]{Pooled EM and CHM of every baseline on the evaluation fold, with the change relative to CAEM at C2 and at C5. Baseline CHM is computed by passing each baseline's answers through CAEM's verifier with the cycle-three composite.}
\label{tab:baseline-pooled}
\end{table}


% Auto-generated by scripts/build_ch5_evidence_tables.py from outputs/analysis/chm_judged.json (2026-10-19).
\begin{figure}[!htbp]
\centering
\begin{tikzpicture}
\begin{axis}[width=0.78\textwidth, height=6.6cm, xlabel={pooled EM (higher is better)}, ylabel={pooled CHM (lower is better)},
  xmin=0.43, xmax=0.55, ymin=0.09, ymax=0.21, font=\small, tick label style={font=\scriptsize}, grid=major, grid style={black!8}]
  \addplot[black!70, only marks, mark=o, mark size=2.2pt] coordinates {(0.511,0.132)};
  \node[font=\scriptsize, anchor=north] at (axis cs:0.511,0.132) {B1};
  \addplot[black!70, only marks, mark=o, mark size=2.2pt] coordinates {(0.511,0.137)};
  \node[font=\scriptsize, anchor=south] at (axis cs:0.511,0.137) {B2};
  \addplot[black!70, only marks, mark=o, mark size=2.2pt] coordinates {(0.445,0.165)};
  \node[font=\scriptsize, anchor=south] at (axis cs:0.445,0.165) {B3};
  \addplot[black!70, only marks, mark=o, mark size=2.2pt] coordinates {(0.452,0.158)};
  \node[font=\scriptsize, anchor=west] at (axis cs:0.452,0.158) {B4};
  \addplot[black!70, only marks, mark=o, mark size=2.2pt] coordinates {(0.521,0.129)};
  \node[font=\scriptsize, anchor=west] at (axis cs:0.521,0.129) {B5};
  \addplot[black!70, only marks, mark=o, mark size=2.2pt] coordinates {(0.512,0.103)};
  \node[font=\scriptsize, anchor=south] at (axis cs:0.512,0.103) {B6};
  \addplot[black!70, only marks, mark=o, mark size=2.2pt] coordinates {(0.466,0.196)};
  \node[font=\scriptsize, anchor=south] at (axis cs:0.466,0.196) {B7};
  \addplot[black!70, only marks, mark=o, mark size=2.2pt] coordinates {(0.519,0.135)};
  \node[font=\scriptsize, anchor=south west] at (axis cs:0.519,0.135) {B8};
  \addplot[blue!65!black, only marks, mark=*, mark size=2.8pt] coordinates {(0.517,0.155) (0.535,0.127) (0.508,0.134)};
  \node[font=\scriptsize, anchor=south, text=blue!65!black] at (axis cs:0.517,0.155) {C0};
  \node[font=\scriptsize, anchor=west, text=blue!65!black] at (axis cs:0.535,0.127) {C2};
  \node[font=\scriptsize, anchor=east, text=blue!65!black] at (axis cs:0.508,0.134) {C5};
\end{axis}
\end{tikzpicture}
\caption[CAEM and the baselines on EM and CHM.]{CAEM at C0, C2 and C5 (filled) and the eight baselines (open) on pooled EM and CHM in the model view (\Cref{tab:baseline-pooled}). The best systems sit at the lower right. The prompting baselines, B8 and CAEM form one cluster; the retrieval baselines sit apart at lower EM and higher CHM; the fine-tuned baseline B6 has the lowest CHM.}
\label{fig:baselines}
\end{figure}


% Auto-generated by scripts/paired_baseline_tests.py (final drafting pass 2026-10).
\begin{table}[!htbp]
\centering\small
\begin{tabular}{@{}lcccc@{}}
\toprule
 & \multicolumn{2}{c}{\textbf{CAEM C2}} & \multicolumn{2}{c}{\textbf{CAEM C5}} \\
\cmidrule(lr){2-3}\cmidrule(lr){4-5}
\textbf{Baseline} & \textbf{EM} & \textbf{CHM} & \textbf{EM} & \textbf{CHM} \\
\midrule
B1 zero-shot & 1 / 1 & 2 / 0 & 1 / 1 & 2 / 0 \\
B2 chain-of-thought & 1 / 1 & 2 / 0 & 1 / 1 & 2 / 0 \\
B3 RAG & 3 / 0 & 4 / 0 & 1 / 0 & 3 / 0 \\
B4 CoT + RAG & 3 / 0 & 4 / 0 & 1 / 0 & 3 / 0 \\
B5 five-shot CoT & 3 / 0 & 3 / 0 & 3 / 0 & 2 / 0 \\
B6 vanilla fine-tune & 1 / 0 & 3 / 1 & 1 / 0 & 2 / 1 \\
B7 FLARE & 1 / 1 & 5 / 0 & 0 / 1 & 5 / 0 \\
B8 semantic entropy & 1 / 0 & n/a & 1 / 0 & n/a \\
\midrule
Total & 14 / 3 of 40 & 23 / 1 of 35 & 9 / 3 of 40 & 19 / 1 of 35 \\
\bottomrule
\end{tabular}
\caption[Significant wins and defeats against each baseline.]{Significant wins / defeats of CAEM against each baseline over the five benchmarks (paired tests, Holm-corrected within each benchmark; \Cref{app:significance}). B8 has no valid CHM.}
\label{tab:sig-summary}
\end{table}


\textbf{Each baseline against its expectation:}
\begin{itemize}
    \item \textbf{Zero-shot (B1):}\label{sec:comp-zero-shot} the floor; the three-step comparison of \Cref{sec:headline-pooled} is built on it.
    \item \textbf{Chain of thought (B2):}\label{sec:comp-cot} expected to help on multi-step questions; instead it scored the same as B1 on pooled EM ($0.511$) with a slightly higher CHM ($0.137$), in line with Wei et al.~\cite{wei2022chain}, who find that the step-by-step cue helps little in models well below $100$ billion parameters, and Chung et al.~\cite{JMLR:v25:23-0870} for small instruction-tuned models. CAEM's zero-shot prompt already asks for reasoning before the answer, so B1 and B2 differ only in the added cue.
    \item \textbf{Retrieval (B3) and retrieval with chain of thought (B4):}\label{sec:comp-rag}\label{sec:comp-cot-rag} retrieving for every question lowered pooled EM ($0.445$ and $0.452$, against $0.511$ for B1) and raised CHM ($0.165$ and $0.158$), the strongest evidence in the panel that indiscriminate retrieval hurts on this backbone; CAEM's retrieval-utility classifier exists for this reason (\Cref{sec:ruc-ablation}).
    \item \textbf{Five-shot chain of thought (B5):}\label{sec:comp-fewshot} the strongest prompting baseline, with pooled EM $0.521$ and CHM $0.129$; its demonstrations help most on CommonsenseQA ($0.793$) and TruthfulQA ($0.347$, the best of any system), while on StrategyQA they hurt ($0.560$ against $0.643$ for B1).
    \item \textbf{Vanilla fine-tuning (B6):}\label{sec:comp-vanilla-ft} a full fine-tune of the backbone on the same training questions CAEM saw in each cycle ($2{,}000$ FEVER, $2{,}000$ TriviaQA and $700$ CommonsenseQA questions per cycle, five cycles, three epochs each), with the gold answer as the training target after CAEM's zero-shot prompt, and with no verifier filter, memory or retention guard. Unlike CAEM, it sees the correct answer for every training question, $23{,}500$ in all. It is the most accurate system on FEVER ($0.640$) and CommonsenseQA ($0.807$), the least accurate on TriviaQA ($0.290$) and TruthfulQA ($0.163$), and ends with pooled EM $0.512$, below CAEM at C2 and level with it at C5, and the lowest CHM ($0.103$). Its retention was not measured.
    \item \textbf{FLARE (B7):}\label{sec:comp-flare} retrieval triggered by low-confidence tokens did not avoid the cost of retrieval on this panel: pooled EM $0.466$ and the highest CHM ($0.196$).
    \item \textbf{Semantic-entropy abstention (B8):}\label{sec:comp-semantic-entropy} the closest comparator to CAEM's abstention. It withheld $75$ TriviaQA and $81$ TruthfulQA answers and almost none elsewhere, answering $89.5$ percent of all questions with an error among answered of $0.427$; its pooled EM is $0.519$ and its CHM $0.135$. Because semantic entropy measures disagreement across samples, it cannot flag a misconception the model holds consistently, as Farquhar et al.~\cite{farquhar2024semantic} note, and its TruthfulQA EM ($0.243$) stays at B1's level.
\end{itemize}
STaR and Self-RAG were not run: STaR's rationalisation step and binary filter differ too much from CAEM's verifier for a clean comparison, and Self-RAG's critic tokens are tied to a different model family.

\subsection{Retrieval-utility classifier ablation}
\label{sec:comp-ablations}
\label{sec:ruc-ablation}

The classifier is evaluated by comparing the main run with an earlier trajectory that ran without it, at cycle three, the last cycle both trajectories kept, with the composite fixed at the same cycle-three calibration so that the verifier signals are identical. The classifier-on run also carried some pipeline corrections (prompt fixes and minor bug fixes), so the comparison measures the classifier together with those corrections, not the classifier alone. Two acceptance gates were set in advance: (1) a pooled EM gain with no benchmark losing more than two points of EM, and (2) TruthfulQA EM reaching the B1 level. The B1 level was $0.433$ when the gate was set; under the final comparison protocol, with the judge reading only the final answer, it is $0.250$, and both runs are scored the same way.

% Generated by scripts/build_final_tables.py from outputs/analysis/chm_judged.json (label em_final); do not edit.
\begin{table}[!htbp]
\centering\small
\setlength{\tabcolsep}{2.5pt}
\renewcommand{\arraystretch}{1.15}
\begin{tabular}{@{}lrrrrrrrrr@{}}
\toprule
 & \multicolumn{3}{c}{\textbf{Exact match}} & \multicolumn{3}{c}{\textbf{CHM}} & \multicolumn{3}{c}{\textbf{Tier 3 share}} \\
\textbf{Benchmark} & \textbf{Off} & \textbf{On} & $\boldsymbol{\Delta}$\,(pp) & \textbf{Off} & \textbf{On} & $\boldsymbol{\Delta}$\,(\% rel) & \textbf{Off} & \textbf{On} & $\boldsymbol{\Delta}$\,(pp) \\
\midrule
\multicolumn{10}{@{}l}{\textit{Training panel}} \\
FEVER & $0.503$ & $0.503$ & $+0.0$ & $0.133$ & $0.137$ & $+2.5\%$ & $92.3\%$ & $70.7\%$ & $-21.7$ \\
TriviaQA & $0.437$ & $0.463$ & $+2.7$ & $0.141$ & $0.138$ & $-2.0\%$ & $95.7\%$ & $68.0\%$ & $-27.7$ \\
CommonsenseQA & $0.603$ & $0.743$ & $+14.0$ & $0.150$ & $0.098$ & $-34.4\%$ & $57.7\%$ & $20.3\%$ & $-37.3$ \\
\midrule
\multicolumn{10}{@{}l}{\textit{Transfer panel}} \\
StrategyQA & $0.617$ & $0.640$ & $+2.3$ & $0.136$ & $0.147$ & $+8.1\%$ & $100.0\%$ & $76.7\%$ & $-23.3$ \\
TruthfulQA & $0.177$ & $0.230$ & $+5.3$ & $0.160$ & $0.165$ & $+3.3\%$ & $99.0\%$ & $40.0\%$ & $-59.0$ \\
\midrule
\textbf{Pooled (all five)} & $0.467$ & $0.516$ & $+4.9$ & $0.144$ & $0.137$ & $-4.8\%$ & $88.9\%$ & $55.1\%$ & $-33.8$ \\
\bottomrule
\end{tabular}
\caption[Retrieval-utility classifier ablation at cycle three.]{CAEM without (Off) and with (On) the retrieval-utility classifier at cycle three, with the composite fixed at the cycle-three calibration: EM, CHM and retrieval share per benchmark.}
\label{tab:ruc-ablation}
\end{table}


\textbf{Gate 1 passes:} pooled EM rises from $0.467$ to $0.516$ ($+4.9$ points) with no benchmark losing EM; CommonsenseQA gains $14.0$ points, TruthfulQA $5.3$, TriviaQA $2.7$ and StrategyQA $2.3$, and FEVER is unchanged at $0.503$ (\Cref{tab:ruc-ablation}). Pooled CHM falls from $0.144$ to $0.137$, mostly on CommonsenseQA ($-34.4$ percent); on FEVER, StrategyQA and TruthfulQA it rises slightly ($+2.5$, $+8.1$ and $+3.3$ percent).

\textbf{Gate 2 fails, narrowly:} TruthfulQA EM rises from $0.177$ to $0.230$ but stays below B1's $0.250$, and its CHM rises slightly ($0.160$ to $0.165$, against $0.147$ for B1).

\textbf{Mechanism:} the retrieval share falls from $88.9$ to $55.1$ percent pooled, and on every benchmark, by $21.7$ points on FEVER, where retrieval genuinely helps, up to $59.0$ points on TruthfulQA ($99.0$ to $40.0$ percent), where it amplified misconceptions. Separating the classifier from the accompanying corrections would need a run with the corrections but without the classifier (\Cref{ch:conclusion}).

\section{Summary of Findings}
\label{sec:summary-findings}

\subsection{Headline outcome in one paragraph}
\label{sec:summary-headline}

Under a matched comparison protocol, CAEM performs at the level of strong prompting baselines, clearly better than retrieval baselines, and level on accuracy with a fine-tuned baseline that needs gold answers, but it does not reduce hallucination below the strong baselines. Pooled EM moves from $0.511$ (B1) to $0.517$ at C0, peaks at $0.535$ at C2 and ends at $0.508$ at C5; CHM moves from $0.132$ to $0.155$ at C0, where the system refuses too often, to $0.127$ at C2 and $0.134$ at C5. No difference between CAEM at C2 or C5 and B1 is significant on any benchmark. Self-improvement helps for one or two cycles and then erodes, and confident confabulation rises above B1's level as the gate becomes more permissive. What the architecture adds is selectivity: CAEM is wrong on $40.5$ percent of the questions it answers at C0 (B1: $48.9$ percent of all questions), and its status tag predicts correctness, with verified answers right $65$ to $75$ percent of the time. Its gains are concentrated where the verifier can check answers against retrieved passages, chiefly TriviaQA. The retrieval-utility classifier, with accompanying corrections, adds $4.9$ points of pooled EM.

\subsection{Hypothesis verdicts}
\label{sec:summary-hypotheses}

Each pre-specified hypothesis of \Cref{sec:hypotheses} is judged by its registered statistic, using only the three registered verdicts (\Cref{tab:hypothesis-verdicts}).

\begin{table}[!htbp]
\centering\small
\setlength{\tabcolsep}{4pt}
\begin{tabularx}{\textwidth}{@{}l>{\raggedright\arraybackslash}p{3.6cm}>{\raggedright\arraybackslash}p{2.2cm}>{\raggedright\arraybackslash}X@{}}
\toprule
\textbf{ID} & \textbf{Hypothesis (\Cref{sec:hypotheses})} & \textbf{Verdict} & \textbf{Deciding statistic and reading} \\
\midrule
H1 & Verifier balanced accuracy exceeds one half on the calibration-matched distribution at every cycle (\Cref{thm:bayes-purity} precondition). & \textbf{supported} & Pooled gate balanced accuracy on the calibration fold: $0.615$, $0.663$, $0.664$, $0.658$ at C1, C2, C3, C5 (C0 $0.753$, in-sample; \Cref{tab:gate-quality}). Weakest: CommonsenseQA $0.52$ to $0.55$; TriviaQA $0.500$ at C1, when no TriviaQA answer was stored. \\
H2 & Pool purity is non-decreasing across cycles whenever the precondition is satisfied (\Cref{thm:bayes-convergence}). & \textbf{unsupported} & Calibration-fold precision $0.780$ (C0, in-sample), $0.740$, $0.736$, $0.746$, $0.732$: it falls at C1, C2 and C5. The changes are within sampling noise and purity stays above the floor of \Cref{thm:purity} (\Cref{tab:pi-store-trajectory}). \\
H3 & Late-cycle pool-purity increments enter a bounded stationary band on the committed-cycle subsequence (\Cref{thm:convergence}). & \textbf{supported} & Precision within $0.732$ to $0.746$ at C1 to C5, a band of width $0.014$, about one standard error of a single reading; a convergence rate cannot be estimated in six cycles. \\
H4 & The composite hallucination metric at the realised cycle horizon approaches a parameter-bounded asymptote (\Cref{thm:asymptotic-elim}). & \textbf{partially supported} & CHM falls from $0.155$ to $0.131$ at C1 and stays between $0.127$ and $0.137$ afterwards ($0.134$ at C5), consistent with flattening; six cycles are too few to estimate the limit itself. \\
\bottomrule
\end{tabularx}
\caption[Hypothesis verdicts.]{Verdicts on the four pre-specified hypotheses.}
\label{tab:hypothesis-verdicts}
\end{table}

\FloatBarrier

\section{Discussion}
\label{sec:discussion}

\subsection{Headline interpretation}
\label{sec:disc-headline}

\textbf{What the matched comparison changed:} an earlier comparison gave the baselines a shorter prompt than CAEM's and scored some answer formats differently; under it, CAEM appeared to gain about three points of EM over B1 before any fine-tuning and to lead every baseline on CHM. Giving every system the same question message and the same scoring rule (\Cref{sec:setup-protocol}) removed most of that lead. On a three-billion-parameter generator, the prompt format alone moved the baselines' accuracy by several points, which is as large as the effects being measured; comparisons of small models are only as fair as their prompts.

\textbf{What carries the remaining gain:} the architecture's measurable contribution is selectivity rather than accuracy. The verifier's score, turned into a status tag and an abstention, lets CAEM tell the user which answers to trust: verified answers are right about two thirds of the time against about half for B1, and at C0 its answers are wrong less often than any baseline's. The self-improvement loop's main effect was behavioural: after one cycle the generator stopped refusing questions it could answer, which removed the cold-start over-refusal and gave the best EM and CHM of the run at C1 and C2. Its gains lasted one or two cycles, in line with Song et al.~\cite{song2024gap}, who find that self-improvement loops saturate within a few rounds.

\textbf{Why the theory does not promise more:} the purification results of \Cref{sec:theory} are about the stored pool, not the model. They held: the stored pool was purer than the generator's raw answers. Whether a model trained on a pool that is still about one-fifth wrong improves is a separate question, and Yi et al.~\cite{fu2025collapse} show that training filtered through a verifier converges to the verifier's own knowledge, mistakes included. The fall in zero-shot accuracy after C2 and the rise in confident confabulation are consistent with that, though the retraining on a cleaned pool that would show it was not run.

\textbf{Label efficiency:} CAEM is label-efficient, not label-free. A labelled calibration fold of $1{,}500$ questions fits the verifier's composite and thresholds, and each later cycle re-scores the same fold with the new adapter and refits; the $603$ cold-start answers were checked against gold labels once (\Cref{sec:memory-store}); every later training answer is labelled by the verifier alone, about $2{,}100$ gold labels in all. Vanilla fine-tuning (B6) needs a gold answer for each of its $23{,}500$ training questions. With that difference, CAEM reaches B6's pooled EM ($0.535$ at C2 and $0.508$ at C5, against $0.512$), while B6 reaches the lower CHM ($0.103$).

\textbf{TruthfulQA:} retrieval does not help on questions written to invite a misconception. At C0, retrieval-tier answers on TruthfulQA were right $21.4$ percent of the time against $29.9$ percent for zero-shot ones (\Cref{tab:per-tier-em}); the difference is partly selection, since the router sends the harder questions to retrieval, but it agrees with the classifier ablation, which gained $5.3$ points on TruthfulQA mainly by retrieving less (\Cref{sec:ruc-ablation}). With every system judged on its final answer, CAEM's TruthfulQA EM stays at B1's level throughout ($0.230$ to $0.277$ against $0.250$), and five-shot chain of thought is the most accurate system there ($0.347$).

\subsection{Where the verifier can check answers}
\label{sec:disc-verifier-scope}

The verifier's ability to separate correct from incorrect answers differs by benchmark. Its area under the ROC curve on the evaluation fold is $0.86$, $0.79$ and $0.75$ on TriviaQA at C0, C2 and C5, $0.64$ to $0.67$ on FEVER, and $0.51$ to $0.62$ on CommonsenseQA, StrategyQA and TruthfulQA. The difference follows from what the signals can see. The signals that carry information about correctness on this panel compare the answer with retrieved Wikipedia passages or, through the question-match score, with the question (\Cref{sec:verifier-qa-rel}); a trivia answer is either stated in a passage or not, while a commonsense choice, an implicit multi-step yes-or-no answer or a misconception is rarely settled by any single passage. The two internal signals do not make up for this: the token signal scores the whole response rather than the answer (\Cref{sec:verifier-internal}), and the dropout and chain signals measure agreement, which a consistently held mistake passes. CAEM's gains follow the same pattern. TriviaQA is where the architecture helps most (EM $0.463$ at C0 and $0.493$ at C2 against $0.430$ for B1), and the benchmarks where the verifier is close to chance are where CAEM is level with, or below, the strong baselines. A verifier that can judge commonsense and reasoning answers, or a system scoped to questions that can be checked against documents, are the two directions this points to (\Cref{ch:conclusion}).

\subsection{A side effect of self-improvement}
\label{sec:disc-phantom}

After the first fine-tune, almost every zero-shot answer began with a reference to ``the context'', although Tier 2 gives the generator no passages. The share of Tier 2 answers on the evaluation fold that mention a context in their first sentences was $3$ percent at C0 ($25$ of $731$), then $93$, $98$, $97$, $97$ and $99$ percent at C1 to C5 ($693$ of $747$ at C1, $746$ of $753$ at C5). The cause is the training data: the fine-tune learns from stored reasoning chains, and most of those were written under the retrieval prompt, where citing the context is correct. The answers can still be right, but each such answer cites a source that does not exist, which is a faithfulness error in its own right. Removing or rewriting context references in the training chains before fine-tuning would have prevented it; this preprocessing step was missed and is proposed in \Cref{ch:conclusion}.

\subsection{Threats to validity}
\label{sec:disc-threats}

\subsubsection{Single seed and single backbone}
\label{sec:disc-threat-seed}

All numbers come from one seed and one three-billion-parameter generator. Repeating the run with other seeds would turn the per-cycle readings into intervals, and repeating it with another generator of similar size would show whether the findings depend on this model.

\subsubsection{Benchmark coverage}
\label{sec:disc-threat-benchmark}

The panel covers short-answer factual and commonsense questions. Long-form generation, code, mathematics, dialogue and other languages are outside its scope (\Cref{sec:scope-research}), and the atomic-fact signal and the long-input judge, built for long answers, were not exercised.

\subsubsection{Calibration to evaluation gap}
\label{sec:disc-threat-gap}

The storage precision promised at the cut-point is measured on the calibration fold and holds on new data only as far as that data resembles it. On the training panel the gap stayed within about five points; on TruthfulQA it did not hold at all (\Cref{sec:check-purity}).

\subsubsection{Cycle budget}
\label{sec:disc-threat-cycles}

The run stopped after cycle five as planned (\Cref{sec:env-cycles}). Six cycles are too few to estimate convergence rates or long-run limits; that would need at least fifteen kept cycles.

\subsubsection{Storage-tail sample size}
\label{sec:disc-threat-n}

Per-benchmark numbers rest on few stored answers in places: at cycle zero the locked gate stored $93$ FEVER answers, $465$ CommonsenseQA answers and no TriviaQA answer on the validation fold, and Tier 1 gave only $7$ answers on the evaluation fold. Pooled readings are more reliable than per-benchmark ones.

\subsubsection{Retrieval-utility classifier mis-routing}
\label{sec:disc-threat-ruc}

The classifier's accuracy differs widely across benchmarks (AUROC $0.718$ on FEVER to $0.936$ on TriviaQA, \Cref{sec:ruc}), so questions unlike its training data may be routed the wrong way. It was not tested on a benchmark held out from its training data.

\subsubsection{Measuring hallucination with the system's own verifier}
\label{sec:disc-threat-verifier}

Two of the seven CHM subtypes depend on CAEM's own verifier: confident confabulation uses its stored score and false refusal uses its decision. The baselines have no verifier of their own, so their answers were scored with CAEM's verifier and the composite fixed at cycle three, and the verifier's decision on a baseline answer is not charged to the baseline (\Cref{sec:setup-protocol}). CHM differences are still partly measured with the instrument CAEM was built around. Two checks limit this concern:
\begin{itemize}
    \item \textbf{Without the two verifier-defined subtypes:} the other five use neither CAEM's composite score nor its decisions. Restricted to them, CAEM at C2 and C5 is lower than every baseline ($0.083$ pooled at both cycles, against $0.084$ for vanilla fine-tuning, the nearest, and $0.103$ for B1), and the paired test gives $28$ significant wins and $2$ defeats at C2 and $31$ wins and $3$ defeats at C5 out of $40$ contrasts. Part of this advantage is built in: the factual-fabrication flag checks answers against passages from the same retrieval that CAEM's retrieval tier writes its answers from.
    \item \textbf{Exact match and the delivered view's error among answered:} both are scored against gold labels and do not depend on the verifier's scores.
\end{itemize}

\subsection{Sensitivity analyses}
\label{sec:disc-sensitivity}

From the re-run sweep (\Cref{tab:sweep-variants}): moving the cut-point from $0.60$ to $0.65$ halves the number of stored answers ($560$ to $281$ on the calibration fold) and raises their precision ($0.780$ to $0.865$). Raising the regularisation strength $C$ from $0.01$ to $0.025$ or more stores far fewer answers on the validation fold ($313$ or fewer, against $558$) at a similar precision ($0.744$ to $0.772$, against $0.751$). Raising $\alpha_{\text{shrink}}$ stores more answers but separates correct from incorrect answers less well on the validation fold (pooled $d$ $0.85$ at $0$, $0.72$ at $0.6$, $0.66$ at $1$). Sensitivity to the safety floor and to the early-exit thresholds was not measured.

\subsection{Comparison to the strongest prior systems}
\label{sec:disc-prior}

\textbf{Conformal factuality:} Mohri and Hashimoto~\cite{mohri2024conformal} raise the correctness of GPT-4 answers from $78$ to $93$ percent on Natural Questions by removing low-scoring sub-claims at a split-conformal threshold, and Cherian et al.~\cite{cherian2024conformal} add conditional guarantees and a learned combination of claim scores. CAEM applies its gate when storing an answer rather than when showing it, keeps the gated answers for reuse, and combines its signals in the manner of Valentin et al.~\cite{valentin2024hallucination}; it also found that conformal guarantees broke under its calibration-to-evaluation shift (\Cref{sec:conformal-gate}).

\textbf{Verifier-filtered self-improvement:} STaR~\cite{NEURIPS2022_639a9a17} and V-STaR~\cite{hosseini2024vstar} train on their own correctness-filtered outputs, and Huang et al.~\cite{huang2025sharpening} analyse self-improvement as sharpening. CAEM differs in using an external multi-signal verifier and in keeping verified answers in memory; like those loops, its gains saturated quickly.

\textbf{Selective retrieval:} Adaptive-RAG of Jeong et al.~\cite{jeong2024adaptive}, the popularity-based routing of Mallen et al.~\cite{mallen2023whennot}, the self-knowledge routing of Wang et al.~\cite{wang2023skr} and the probing of Baek et al.~\cite{baek2025probing} all decide per question whether to retrieve. CAEM's classifier belongs to this line, and its ablation (\Cref{sec:ruc-ablation}) agrees with their finding that retrieving for every question leaves accuracy on the table.

\textbf{Retrieval and misconceptions:} on questions written to invite a familiar wrong answer, retrieved passages tend to share the question's framing rather than refute it, so the retrieval tier amplifies the misconception and the storage gate cannot recover what retrieval has already corrupted. This is why retrieval-tier answers on TruthfulQA are the least accurate of the panel (\Cref{sec:disc-headline}), and the case the classifier was built for.


\subsection{Limitations}
\label{sec:disc-limitations}

\begin{itemize}
    \item \textbf{Memory scale:} the memory reached $12{,}175$ entries, so its behaviour at the millions of entries of a deployment was not tested.
    \item \textbf{Pre-routing calibration:} the pre-routing temperature reached its upper bound at cycle zero ($T = 20.086$) on FEVER and TriviaQA, which flattened $\upre$ and left the safety override almost inactive (\Cref{sec:tier-safety}); the storage decision does not depend on this temperature.
    \item \textbf{Deferred-pool horizon:} only three reconsideration passes ran, so the long-run behaviour of the deferred buffer is not known.
    \item \textbf{Question-match signal:} the fourth signal family was meant to catch off-topic answers, but it scores the short answer string, which the cross-encoder cannot relate to the question; it acted as a correctness cue instead, and off-topic answers were not detected (\Cref{sec:verifier-qa-rel}).
    \item \textbf{Phantom context:} fine-tuning on retrieval-written chains taught the generator to cite a non-existent context (\Cref{sec:disc-phantom}).
    \item \textbf{Verifier scope:} the verifier separates correct from incorrect answers well only where retrieved passages can confirm an answer (TriviaQA), and is close to chance on commonsense, multi-step and misconception questions (\Cref{sec:disc-verifier-scope}).
    \item \textbf{Confident confabulation after self-improvement:} the verifier scores more of the fine-tuned generator's wrong answers highly than it did the base generator's ($0.305$ at C5 against $0.208$ for B1), so the gate's selectivity weakens as the generator it judges changes.
    \item \textbf{No repeated questions:} the evaluation uses fresh questions to measure learning, so the memory-direct path, built for repeated questions in deployment, and its cost saving were exercised only by the router probe; the gain in deployment depends on the share of repeats, which this study does not measure.
\end{itemize}

\paragraph{Limitations identified against the literature:} six further limitations follow from the cited literature.
\begin{itemize}
\item \textbf{Saturation and format learning:} Song et al.~\cite{song2024gap} report that verifier-filtered self-improvement saturates within two or three rounds and that part of the first-round gain at larger scales comes from better adherence to the answer format. The run is consistent with early saturation, and because the main metric is exact match, part of the early gain may be format learning; a flexible-match re-scoring would separate the two.
\item \textbf{Optimising toward the verifier:} Gulcehre et al.~\cite{gulcehre2023rest} find reward-model gains rising while human-judged quality peaks after the first round, and Yi et al.~\cite{fu2025collapse} show that verifier-filtered retraining converges to the verifier's own knowledge. Exact match against gold labels is independent of CAEM's verifier; CHM is not (\Cref{sec:disc-threat-verifier}).
\item \textbf{Systematic verifier errors:} agreement-based signals do not detect misconceptions the model holds consistently, as Farquhar et al.~\cite{farquhar2024semantic} warn, and repeated verification removes only independent errors (\Cref{cor:self-correction}); TruthfulQA's low stored precision is the realised case. The sample-agreement signal is an embedding cosine similarity, a family that Sinha~\cite{sinha2025semantic} shows can miss plausible hallucinations, although that evidence concerns similarity to a source rather than agreement across samples.
\item \textbf{Calibration-fold size:} Niculescu-Mizil and Caruana~\cite{niculescu2005predicting} find isotonic regression reliable only from about $1{,}000$ calibration points; CAEM fits per-benchmark curves on $500$ points, and the shrinkage toward the pooled fit limits but does not remove the resulting variance.
\item \textbf{Low chance floor:} Ho et al.~\cite{ho2023large} note that filtering on the final answer does not ensure correct reasoning, and Zelikman et al.~\cite{NEURIPS2022_639a9a17} that binary tasks yield many correct answers with poor rationales. StrategyQA and FEVER have chance levels of one half and one third, and FEVER answers enter the training pool.
\item \textbf{Adapter capacity and drift:} Biderman et al.~\cite{biderman2024lora} find that low-rank adaptation learns less of a new target than full fine-tuning and that its forgetting grows with training; the cycle-four rollback and the accuracy decline after cycle two are consistent with this, and the retention guard bounds the drift rather than preventing it.
\end{itemize}

\section{Chapter signpost}
\label{sec:ch5-signpost}

This chapter reported the main run, the cycle-zero choices behind it, the comparison with eight baselines and the classifier ablation, and judged the four hypotheses: H1 and H3 supported, H4 partially supported, H2 unsupported. \Cref{ch:conclusion} restates the contributions against these results and turns the limitations above into future work.
   % Performance Evaluation and Analysis
% ============================================================================
% Chapter 6 — Conclusion and Future Work (rewritten 2026-10-19 on the final
% numbers: baseline re-run under the comparison protocol, one scoring rule,
% model and delivered views). Numbers trace to the chapter-five tables.
% ============================================================================

\chapter{Conclusion and Future Work}
\label{ch:conclusion}

This chapter states what the thesis built and what its evaluation showed, judges the four hypotheses, collects the limitations, and turns them into future work.

\section{Conclusion}
\label{sec:conclusion}

\subsection{Contributions}
\label{sec:concl-contributions}

The thesis set out to reduce confident hallucination in factual question answering on a small open generator by changing the system around the model rather than the model alone (\Cref{ch:introduction}). It contributes the following.
\begin{itemize}
    \item \textbf{An architecture:} CAEM routes each question to a verified episodic memory, to zero-shot generation or to retrieval-augmented generation, with a retrieval-utility classifier that decides when retrieval is likely to help (\Cref{ch:methodology}). A verifier of nine signals produces a calibrated probability of correctness; a four-outcome decision stores, defers, abstains on or discards each answer; stored answers are re-checked at every cycle boundary and used to fine-tune a low-rank adapter under a retention guard that rolls the adapter back while keeping the memory.
    \item \textbf{A per-answer reliability signal:} every answer reaches the user with its probability and a status tag (verified, provisional, conflicting or insufficient evidence), and the answer is withheld only when the evidence is insufficient (\Cref{sec:decision}).
    \item \textbf{A formal analysis of the stored pool:} two short proven results, a testable precondition and a design target: that storage makes the pool purer than the generator's raw output exactly when the verifier's balanced accuracy exceeds one half (\Cref{thm:bayes-purity}) and that the pool's precision is bounded below by the calibration precision when new questions resemble the calibration data (\Cref{thm:purity}), together with design heuristics whose idealising assumptions are stated and checked against the run (\Cref{sec:theory}).
    \item \textbf{A matched evaluation:} a six-cycle run on five benchmarks against eight baselines under one comparison protocol (the same questions, prompt, scoring rule and hallucination instrument for every system), reported in a model view and a delivered view, with paired tests corrected for multiple comparisons and every reported number rebuilt from the saved answers by one command (\Cref{sec:setup}, \Cref{app:reproducibility}).
    \item \textbf{A diagnosed limit:} a measured account of where verifier-gated self-improvement helps, where it stops helping and why (\Cref{sec:discussion}).
\end{itemize}

\subsection{Headline finding}
\label{sec:concl-headline}

CAEM is a label-efficient system that matches strong baselines, gives users a status tag that genuinely predicts correctness, and shows a measured, diagnosed limit of verifier-gated self-training. The evidence for each part of that statement is as follows.

\textbf{Accuracy and hallucination:} the architecture alone added almost no accuracy: pooled exact match was $0.517$ at C0 against $0.511$ for the zero-shot baseline B1, and CHM was higher ($0.155$ against $0.132$) because the untrained system refused too often (\Cref{tab:headline-steps}). One or two cycles of self-improvement removed most of the over-refusal and gave the best readings of the run ($0.535$ and $0.127$ at C2); later cycles eroded both, and the run ended level with B1 ($0.508$ and $0.134$ at C5). No difference between CAEM and B1 is significant on any benchmark at C2 or C5.

\textbf{Against the baselines:} under the matched protocol CAEM is level with the strong prompting baselines and semantic-entropy abstention (B1, B2, B5, B8), with no significant exact-match win against any of them; it is significantly better than the retrieval baselines and FLARE on several benchmarks; and it matches the exact match of vanilla fine-tuning (B6, $0.512$), which used a gold answer for each of $23{,}500$ training questions against about $2{,}100$ gold labels for CAEM, while B6 reached the lower CHM ($0.103$) (\Cref{tab:baseline-pooled,tab:sig-summary}). An earlier comparison that gave the baselines a shorter prompt had shown CAEM ahead on both measures; the matched comparison removed most of that lead.

\textbf{Selectivity:} what the architecture measurably adds is a reliable signal of which answers to trust. Answers tagged verified were right $65$ to $75$ percent of the time at every cycle, against $51$ percent for B1 over all questions, and the tags are ordered by accuracy, verified above provisional above the other two (\Cref{tab:status-tags}). At C0 CAEM was wrong on $40.5$ percent of the questions it answered, less often than any baseline.

\textbf{Where it works:} the gains follow the verifier's ability to check answers against retrieved passages. On TriviaQA, where that ability is highest (area under the ROC curve $0.75$ to $0.86$), CAEM is most accurate relative to B1; on CommonsenseQA, StrategyQA and TruthfulQA, where the verifier is close to chance, it is level with or below the strong baselines (\Cref{sec:disc-verifier-scope}).

\textbf{Why self-improvement stopped helping:} the stored pool stayed about three-quarters correct, as the theory predicts, but the generator trained on it. Training on answers a verifier has passed converges toward the verifier's own view, as Yi et al.~\cite{fu2025collapse} show: confident confabulation rose from $0.208$ for B1 to $0.305$ at C5, so the verifier scored more of the fine-tuned generator's wrong answers highly. Training on chains written under the retrieval prompt also taught the generator to cite a context it was not given, in $93$ to $99$ percent of zero-shot answers after the first fine-tune (\Cref{sec:disc-phantom}).

\textbf{Routing:} the retrieval-utility classifier, with accompanying corrections, raised pooled exact match by $4.9$ points at cycle three, mainly by retrieving less on questions where retrieval misleads (\Cref{sec:ruc-ablation}).

\subsection{Hypothesis verdicts}
\label{sec:concl-hypotheses}

The four pre-specified hypotheses of \Cref{sec:hypotheses} are judged in \Cref{tab:hypothesis-verdicts}.
\begin{itemize}
    \item \textbf{H1, supported:} the gate's balanced accuracy on the calibration fold stayed above one half at every cycle ($0.615$ to $0.664$ pooled), so the precondition for purification held; CommonsenseQA was the weakest benchmark ($0.52$ to $0.55$).
    \item \textbf{H2, unsupported:} pool purity did not rise from cycle to cycle; calibration-fold precision fell at C1, C2 and C5, within sampling noise.
    \item \textbf{H3, supported:} precision stayed in a band of width $0.014$ ($0.732$ to $0.746$) over the kept cycles; a convergence rate cannot be estimated from six cycles.
    \item \textbf{H4, partially supported:} CHM fell after the first cycle and stayed between $0.127$ and $0.137$, consistent with flattening, but the limit itself cannot be estimated.
\end{itemize}
The two proven results concern the stored pool, and both held where their assumptions held. None of the formal statements bounds the accuracy of the model trained on the pool, which is where the run fell short.

\subsection{Limitations}
\label{sec:concl-open}

The limitations are discussed with their evidence in \Cref{sec:disc-threats,sec:disc-limitations}; the ones that bound the conclusions above are these.
\begin{itemize}
    \item \textbf{One seed, one backbone, five benchmarks:} every result comes from a single run of a three-billion-parameter generator on $300$ evaluation questions per benchmark, and only the paired tests account for sampling noise.
    \item \textbf{Six cycles:} the run stopped after cycle five, as planned after the cycle-four rollback, which is too short to estimate convergence rates or long-run limits.
    \item \textbf{Verifier scope:} the verifier separates correct from incorrect answers well only where passages can confirm an answer, and its selectivity weakened as the generator it judged changed.
    \item \textbf{The hallucination instrument:} two of the seven CHM subtypes use CAEM's own verifier; exact match and the delivered view's error among answered do not depend on it.
    \item \textbf{No repeated questions:} the evaluation used fresh questions to measure learning, so the memory-direct tier and its cost saving were tested only by the router probe (\Cref{tab:tier1-probe}).
    \item \textbf{Bundled ablation:} the classifier ablation also carries pipeline corrections, so it measures the classifier together with them.
    \item \textbf{Labels:} CAEM is label-efficient, not label-free; it used about $2{,}100$ gold labels, and a deployment needs a labelled batch at every cycle (\Cref{sec:conformal-cycle-recalib}).
\end{itemize}

\subsection{Closing remark}
\label{sec:concl-closing}

The thesis began from the observation that a user cannot tell a confident wrong answer from a confident right one. CAEM does not make a small generator much more accurate, and its self-improvement loop helped for only one or two cycles. What it does provide is an answer that carries an honest estimate of its own reliability, a pool of stored answers that stayed as pure as the theory's floor predicts, and a measured account of why training on verifier-passed answers stopped helping, consistent with the verifier's blind spots becoming the generator's. The most useful next step is therefore a better verifier rather than more cycles.

\section{Future Work}
\label{sec:future-work}

The items below follow from the limitations above and from the points that \Cref{ch:methodology,ch:evaluation} left open. The first three address the causes of the limit the run measured; the rest complete measurements this study did not take.

\subsection{A verifier for answers that passages cannot settle}
\label{sec:future-judge}

The verifier is close to chance on commonsense, multi-step and misconception questions because its informative signals compare the answer with passages or with the question (\Cref{sec:disc-verifier-scope}). Four changes follow.
\begin{itemize}
    \item \textbf{Answer-level internal signals:} the token signal scores the whole response, so it measures fluency more than confidence in the answer (\Cref{sec:verifier-internal}); scoring the answer line alone, and computing semantic entropy over the final answers of sampled responses, the setting in which Farquhar et al.~\cite{farquhar2024semantic} found it effective, would give internal signals that track correctness.
    \item \textbf{Checking the reasoning:} every signal scores the final answer, so a chain with a wrong step and a correct conclusion passes. The atomic-fact check gave the same value as the passage-mean signal on every answer of the run, because the answers were too short to decompose (\Cref{sec:verifier-grounding}); applying it to the reasoning chain, where several facts are asserted, would give it something to check and would catch such chains.
    \item \textbf{A learned verifier:} the composite combines fixed signals; a verifier trained on the calibration fold to judge answer and reasoning together, as V-STaR~\cite{hosseini2024vstar} trains its verifier on correct and incorrect self-generated solutions, could learn cues the fixed signals miss. The long-answer judge that was dropped at cycle zero (\Cref{sec:adj-judge-ablation}) could be revisited in the same way.
    \item \textbf{Off-topic answers:} the question-match signal acts as a correctness cue on short answers and does not detect off-topic ones (\Cref{sec:verifier-qa-rel}); a check on the full response, or a judge asked directly whether the answer addresses the question, would let the off-topic subtype of CHM be measured.
\end{itemize}

\subsection{Contradiction-aware verification}
\label{sec:future-refutation}

The entailment judge returns only whether a passage supports a claim, so an answer that a passage contradicts scores the same as one the passages ignore, and the factual-contradiction subtype of CHM could not be measured (\Cref{sec:setup-metrics}). A three-way judge (support, contradiction, neither) would let the decision discard contradicted answers outright and would complete the metric. The language-model judge tried for the two unmeasured subtypes agreed too weakly with hand labels (Cohen's kappa about $0.35$), so a better-validated judge is needed first.

\subsection{Cleaner self-improvement}
\label{sec:future-lora}

The loop's gains faded after two cycles while confident confabulation rose. Four changes target the cause.
\begin{itemize}
    \item \textbf{Remove context references from training chains:} rewriting or removing references to a context in retrieval-written chains before fine-tuning would have prevented the phantom-context citations (\Cref{sec:disc-phantom}).
    \item \textbf{Weight training answers by confidence:} about one stored answer in four is wrong, and the wrong ones carry lower scores; weighting each training answer by its calibrated probability, or raising the training threshold as the pool grows, would reduce the noise the generator learns from.
    \item \textbf{Test the cause directly:} retraining from the same starting point on a pool with the wrong answers removed would show whether the erosion after C2 comes from the wrong answers in the pool, as the account above suggests, or from repeated fine-tuning itself.
    \item \textbf{Retry when unsure:} answers the verifier is unsure about are now deferred or withheld; regenerating them, with retrieval or with more samples, before deciding would turn some of those into verified answers.
    \item \textbf{Adapter capacity:} Biderman et al.~\cite{biderman2024lora} find that low-rank adaptation forgets less but learns less than full fine-tuning; comparing ranks, or full fine-tuning on the verified pool, would show whether the plateau is set by the adapter or by the pool.
\end{itemize}

\subsection{Scope to checkable questions}
\label{sec:future-transfer}

The other direction the verifier result points to is to narrow the system rather than widen the verifier: to deploy CAEM only where answers can be checked against documents, such as trivia, entity and document-grounded questions, and to measure it on a panel of such benchmarks. Because the main run fine-tuned on all three training benchmarks, such a study needs its own run rather than a re-reading of this one. Whether a narrow system or a general verifier is the better path is left open.

\subsection{Isolating the effects}
\label{sec:future-ablations}

\begin{itemize}
    \item \textbf{The classifier alone:} a run with the pipeline corrections but without the retrieval-utility classifier would separate the two parts of the measured $4.9$-point gain (\Cref{sec:ruc-ablation}).
    \item \textbf{Multiple seeds:} repeating the run with different seeds would show how much of the cycle-to-cycle movement, of a size similar to the effects measured, is noise.
    \item \textbf{A flexible-match re-scoring:} part of the early gain may be format learning (\Cref{sec:disc-limitations}); scoring with a flexible answer match as well as exact match would separate the two.
\end{itemize}

\subsection{Testing the formal heuristics}
\label{sec:future-fixed-composite}

The monotonicity heuristic and hypothesis H2 assume a fixed verifier, while the run refitted the composite every cycle. A cycle in which the composite is frozen and only the generator changes would separate the two effects (\Cref{sec:theory}). Measuring the gate's true- and false-positive rates and the generator's accuracy on a labelled sample at every cycle would check the purification condition directly rather than through balanced accuracy.

\subsection{A longer horizon}
\label{sec:future-equilibrium-horizon}

A convergence rate, the long-run limit of CHM, and the equilibrium behaviour under rollbacks all need at least fifteen kept cycles under a fixed composite (\Cref{sec:theory}). A longer run would also show whether the deferred buffer, which ran only three reconsideration passes, recovers more answers over time, and whether a rotating calibration fold behaves better than a fixed one over many cycles (\Cref{sec:conformal-cycle-recalib}).

\subsection{Measurements this study did not take}
\label{sec:future-latency}

\begin{itemize}
    \item \textbf{Time per tier:} the batched evaluation log could not separate the time spent in each tier; timing each tier at a batch size of one would fill in the deployment-cost model (\Cref{sec:tier-cost}).
    \item \textbf{The corpus ceiling:} checking, for every evaluation question, whether any retrieved passage contains the gold answer would bound what retrieval can achieve (\Cref{cor:corpus-floor}).
    \item \textbf{Calibration on more data:} the per-benchmark calibration curves are fitted on $500$ answers, below the size at which Niculescu-Mizil and Caruana~\cite{niculescu2005predicting} find isotonic regression reliable; a larger calibration fold, or a comparison with a parametric fit per benchmark, would show how much the shrinkage toward the pooled fit leaves to chance.
    \item \textbf{Retention of the fine-tuned baseline:} the retention probes were run only on CAEM; running them on B6 would show whether the retention guard protects anything that plain fine-tuning loses.
\end{itemize}

\subsection{Larger backbones and other domains}
\label{sec:future-backbone}

All results come from one three-billion-parameter generator on English factual question answering. Song et al.~\cite{song2024gap} find that small models often gain little from verifying their own outputs, so repeating the study on a larger generator, and on other domains and languages, would show how far the findings carry.

\subsection{Deployment study}
\label{sec:future-deployment}

The memory-direct tier and the production design (\Cref{sec:conformal-cycle-recalib}) were built for traffic in which users repeat questions, which this study did not have. A deployment with live users would measure the share of repeated questions, the cost saving of the memory-direct tier, how users read the status tag, and the labelling effort a cycle needs in practice. It would also test the memory at a realistic size: the run's memory reached $12{,}175$ entries, so the approximate search and the value-based pruning built for millions of entries were never used.
   % Conclusion and Future Work

% --- Bibliography -----------------------------------------------------------
\printbibliography[heading=bibintoc,title={References}]

% --- Appendices -------------------------------------------------------------
\appendix
\appendixoverflowsafe
% ============================================================================
% Appendix A — Implementation reference
% Body prose in the chapters is concept-only (no module paths,
% no class or function names, no script names, no folder or file references).
% All of those references that a re-implementer needs live here.
% ============================================================================

\chapter{Implementation reference}
\label{app:implementation}

\section{Code and data availability}
\label{app:code-data-availability}

The implementation of the architecture, the cycle runner, the calibration scripts, the baseline drivers, the verifier, the retrieval-utility classifier and the table and figure generators is released at \url{https://github.com/aksaN000/caem-thesis}. The results are in the shared Google Drive folder \url{https://drive.google.com/drive/folders/1xErhsRDYe_qO0ZgUBn0WXK81eHc5jZCa}: \texttt{1\_thesis/} holds this report, \texttt{2\_results\_phase1e/} every number it reports (\texttt{caem/} for the main run, \texttt{baselines/} for the final baseline panel, \texttt{analysis/} for the scored summaries and tests), \texttt{3\_ruc\_ablation\_phase1d/} the trajectory without the classifier, and \texttt{9\_archive/} superseded files, kept rather than deleted. A \texttt{README} in each folder maps it to the repository's output directories. The $64$-gigabyte Wikipedia passage index and the pre-main snapshot are mirrored on the Hugging Face dataset repository \texttt{aksaN000/caem-passage-index-21m} (private; access on request; \Cref{app:repro-snapshot}).

\section{Repository layout and module map}
\label{app:repo-layout}

Body prose throughout the thesis refers to the architectural components by concept name; the inventory below maps each concept to the file that implements it. Paths are relative to the repository root.

{\scriptsize
\setlength{\tabcolsep}{4pt}
\renewcommand{\arraystretch}{1.15}
\begin{longtable}{@{}>{\ttfamily}p{5.8cm}p{2.2cm}p{6.6cm}@{}}
\caption[Repository layout]{Repository layout. Body prose throughout the thesis refers to architectural components by their concept names (memory store, three-tier router, verifier ensemble, retrieval-utility classifier, self-improvement loop, retention guard); this table maps each concept to the file that implements it.}
\label{tab:app-repo-layout}\\
\toprule
\textnormal{\textbf{Path}} & \textbf{Role} & \textbf{Description} \\
\midrule
\endfirsthead
\toprule
\textnormal{\textbf{Path}} & \textbf{Role} & \textbf{Description} \\
\midrule
\endhead
\bottomrule
\endfoot
\multicolumn{3}{@{}l}{\textnormal{\textit{Top-level directories}}} \\
caem/              & package root & Core Python package: pipeline, memory, routing, verification, training. \\
scripts/           & runners      & Cycle runner, calibration scripts, baseline runners, rescore utilities. \\
eval/              & harness      & Shared per-benchmark eval scaffold (EM, em\_llm\_judged, CHM, CES). \\
outputs/           & artefacts    & Run outputs: full\_run/ (main run), cycle\_0/, rerun\_v2/ and baselines\_phase1e/ (baselines), analysis/ (scored summaries). \\
data/              & external     & Passage index, calibration pairs, alias dictionary, retention and RUC datasets. \\
tests/             & test suite   & Pytest suite covering storage gate, deferred buffer, retention probe, SIL, verifier. \\
docs/              & runbook      & PRODUCTION\_RUNBOOK.md: deployment-mode spec and cycle-trigger policies. \\
defense/thesis\_report/ & thesis   & LaTeX chapters, bibliography, figures, main.pdf build target. \\
\midrule
\multicolumn{3}{@{}l}{\textnormal{\textit{Pipeline entry points}}} \\
caem/pipeline.py            & per-query  & Pipeline class: serial driver used in production and demo paths. \\
caem/pipeline\_batch.py     & per-cycle  & BatchPipeline class: batched cycle driver used by the experiment runner. \\
scripts/run\_experiment.py  & runner     & Canonical cycle driver: invokes BatchPipeline, SIL, retention probe, retroverify. \\
\midrule
\multicolumn{3}{@{}l}{\textnormal{\textit{Configuration}}} \\
caem/config.py                       & registry   & Central CAEMConfig dataclass; every threshold, weight, structural toggle. \\
caem/ruc/v2/feature\_spec.json       & RUC config & RUC threshold, regulariser, feature order, model paths. \\
\midrule
\multicolumn{3}{@{}l}{\textnormal{\textit{Memory module}}} \\
caem/memory/store.py     & storage gate    & EpisodicMemoryStore: fixed-threshold admission and FAISS-backed retrieval. \\
caem/memory/deferred.py  & deferred buffer & DeferredBuffer and DeferredEntry: TTL-bounded reconsider queue. \\
caem/memory/encoder.py   & embeddings      & SBERT-side encoder for memory keys and Tier 1 lookup. \\
\midrule
\multicolumn{3}{@{}l}{\textnormal{\textit{Routing}}} \\
caem/routing/router.py & tier dispatch & Tier 1/2/3 router with safety override and combined-score thresholds. \\
caem/routing/ruc.py    & RUC head      & LRRetrievalUtilityClassifier: binary head for Tier 2 versus Tier 3 routing. \\
\midrule
\multicolumn{3}{@{}l}{\textnormal{\textit{Verification}}} \\
caem/verification/verifier.py            & verifier ensemble & UnifiedVerifier: nine-signal composite producing $u_{\text{pre}}$ and $u_{\text{stored}}$. \\
caem/verification/adaptive\_nli\_judge.py & NLI dispatch     & AdaptiveNLIJudge: length-gated dispatcher; in the deployed configuration every hypothesis routes to MiniCheck under a 512-token truncation contract because the long-hypothesis Qwen-3B backend was ablated (\Cref{sec:adj-judge-ablation}). \\
caem/verification/cal\_prob\_composite.py & calibration      & Per-signal isotonic curves plus per-benchmark shrinkage and boost layer. \\
caem/retrieval/rag.py                    & retrieval         & FAISS-backed passage index reader for Tier 3 grounded generation. \\
\midrule
\multicolumn{3}{@{}l}{\textnormal{\textit{Training}}} \\
caem/training/self\_improvement.py & SIL trainer     & SelfImprovementLoop: LoRA fine-tune and retroverify hook. \\
caem/training/retention\_probe.py  & retention guard & Per-cycle multi-modal probe panel and ratio-floor check. \\
caem/training/pool\_reweighting.py & pool balancer   & Per-benchmark temperature-mixed reweighting and share ceiling. \\
caem/training/loop\_filter.py      & SIL admission   & Per-pair filter at $\tau_{\text{train}} = 0.70$ above the storage cut-point. \\
\midrule
\multicolumn{3}{@{}l}{\textnormal{\textit{Evaluation}}} \\
eval/harness.py            & eval scaffold     & Per-benchmark EM, F1, latency, CHM, and CES. \\
eval/benchmarks.py         & benchmark loaders & TriviaQA, FEVER, CommonsenseQA, TruthfulQA, StrategyQA, MMLU loaders. \\
caem/eval/equilibrium.py   & equilibrium       & Equilibrium summary and the automatic early-stop check. \\
\midrule
\multicolumn{3}{@{}l}{\textnormal{\textit{Final results}}} \\
scripts/final\_sources.py      & sources   & The one answer file per system and benchmark that every analysis reads. \\
scripts/make\_final\_results.sh & results   & Rebuilds every number, table and figure of \Cref{ch:evaluation} from the answer files. \\
\end{longtable}
}

\section{Configuration registry}
\label{app:config-registry}

Every constant that governs runtime behaviour lives in the \verb|CAEMConfig| dataclass at \verb|caem/config.py|. The exception is the retrieval-utility classifier, whose configuration lives in \verb|caem/ruc/v2/feature_spec.json| at deployment time (the classifier is loaded once at pipeline construction and is not refit at cycle boundaries; see \Cref{sec:ruc}). The registry below groups every field by concept, records the deployed value, and tags the provenance as \emph{literature} (anchored to a cited paper), \emph{design} (architectural decision documented before any run), \emph{calibrated} (selected empirically at cycle zero), or \emph{pre-specified} (chosen before any data is observed).

{\footnotesize
\setlength{\tabcolsep}{4pt}
\renewcommand{\arraystretch}{1.1}
\begin{longtable}{@{}p{4.6cm}p{2.0cm}p{1.8cm}p{5.8cm}@{}}
\caption[Configuration registry]{Configuration registry. Provenance tags: \emph{literature} (cited paper), \emph{design} (architectural choice documented before any run), \emph{calibrated} (chosen empirically at cycle zero or refit at a cycle boundary), \emph{pre-specified} (registered horizon before any data is observed).}
\label{tab:app-config-registry}\\
\toprule
\textbf{Field} & \textbf{Value} & \textbf{Provenance} & \textbf{Source / rationale} \\
\midrule
\endfirsthead
\toprule
\textbf{Field} & \textbf{Value} & \textbf{Provenance} & \textbf{Source / rationale} \\
\midrule
\endhead
\bottomrule
\endfoot
\multicolumn{4}{@{}l}{\textit{Routing thresholds (\Cref{sec:router})}} \\
\texttt{routing\_lambda}             & $0.70$ & design       & Combined-score weight on similarity. \\
\texttt{tier1\_combined\_threshold}   & $0.90$ & design       & Tier 1 dispatch floor on combined score. \\
\texttt{tier2\_similarity\_threshold} & $0.75$ & design       & Tier 2 dispatch floor on similarity alone. \\
\texttt{safety\_u\_pre\_min}           & $0.38$ & calibrated   & Safety-override floor; lowered from $0.60$ before the main run after an earlier run's first-cycle routing. \\
\texttt{u\_pre\_token\_weight}         & $0.60$ & design       & Token-probability weight in $u_{\text{pre}}$. \\
\texttt{u\_pre\_cconv\_weight}         & $0.40$ & design       & Layer-convergence weight in $u_{\text{pre}}$. \\
\midrule
\multicolumn{4}{@{}l}{\textit{Storage gate thresholds (\Cref{sec:decision})}} \\
\texttt{store\_threshold}            & $0.60$ & design       & STORE if $u_{\text{stored}} \ge$ this. \\
\texttt{defer\_threshold}            & $0.45$ & design       & DEFER if \texttt{defer\_threshold} $\le u_{\text{stored}} < $ \texttt{store\_threshold}. \\
\texttt{abstain\_pground\_ceiling}    & $0.20$ & design       & ABSTAIN if $u_{\text{stored}} < $ \texttt{defer\_threshold} and $p_{\text{ground,max}} <$ this. \\
\texttt{early\_exit\_u\_internal}      & $0.70$ & design       & Confabulation early-exit on internal-uncertainty floor. \\
\texttt{early\_exit\_p\_ground\_max}    & $0.20$ & design       & Confabulation early-exit on grounding ceiling. \\
\midrule
\multicolumn{4}{@{}l}{\textit{Calibration parameters (\Cref{sec:composite})}} \\
\texttt{composite\_mode}             & \texttt{"auto"}  & design     & Auto-switch to calibrated probability composite when the calibration JSON loads. \\
\texttt{composite\_shrinkage\_alpha}  & $0.6$  & calibrated   & Per-benchmark child curve shrinkage toward pooled fit. \\
\texttt{adaptive\_thresholds\_\allowbreak per\_cycle} & \texttt{true} & design   & Refit temperature, isotonic curves, and boost weights at every successful cycle. \\
\texttt{adaptive\_thresholds\_\allowbreak ema\_alpha} & $0.7$ & design       & Exponential-moving-average smoothing on the per-cycle temperature refit. \\
\midrule
\multicolumn{4}{@{}l}{\textit{Retention guard (\Cref{sec:retention})}} \\
\texttt{forgetting\_tolerance}       & $0.93$  & design       & Worst-probe ratio floor; trigger asymmetric rollback if any probe drops below. \\
\texttt{retention\_probes}           & 3-probe panel & design & Multi-modal probe panel \texttt{(mmlu,\allowbreak{} triviaqa\_test,\allowbreak{} commonsense\_qa\_test)}; replaces the single-MMLU probe of v1. \\
\texttt{retention\_probe\_n}          & $200$  & design       & Samples per probe at every cycle. \\
\midrule
\multicolumn{4}{@{}l}{\textit{Deferred buffer (\Cref{sec:deferred-reconsider})}} \\
\texttt{deferred\_buffer\_max\_size}   & $10\,000$ & design     & Bounded FIFO capacity for deferred entries. \\
\texttt{deferred\_buffer\_ttl\_cycles} & $2$     & design       & Reconsider-invocation TTL; entries dropped after two passes. \\
\midrule
\multicolumn{4}{@{}l}{\textit{Retroverify and consolidation (\Cref{sec:retroverify})}} \\
\texttt{retroverify\_prune\_\allowbreak threshold}        & $0.50$ & design & Prune from memory if recalibrated $u_{\text{stored}}$ drops below. \\
\texttt{hit\_counter\_force\_\allowbreak retroverify}      & $10$   & design & Per-entry retrieval-hit threshold forcing reverification next cycle. \\
\texttt{enable\_consolidation}               & \texttt{true} & design & Cycle-boundary SBERT-similarity clustering pass. \\
\texttt{consolidation\_similarity\_\allowbreak threshold} & $0.92$ & design & Minimum cosine similarity for cluster admission. \\
\texttt{consolidation\_search\_k}             & $20$   & design & Per-entry FAISS search depth for neighbours. \\
\texttt{consolidation\_max\_u\_spread}         & $0.15$ & design & Spread guardrail; reject clusters whose $u_{\text{stored}}$ range exceeds this. \\
\midrule
\multicolumn{4}{@{}l}{\textit{Self-improvement training (\Cref{sec:sil})}} \\
\texttt{use\_lora\_training}          & \texttt{true}  & design       & LoRA SIL primary path; replaces the v1 full-fine-tune fallback. \\
\texttt{lora\_r}                     & $32$         & design       & Adapter rank. \\
\texttt{lora\_alpha}                 & $64$         & design       & Scaling factor, set to twice the rank. \\
\texttt{lora\_dropout}               & $0.05$       & design       & Adapter dropout. \\
\texttt{lora\_target\_modules}        & \texttt{(q,k,v,o,\allowbreak{}gate,up,down)\_proj} & literature & All linear projections of the transformer block. \\
\texttt{lora\_learning\_rate}         & $2\!\times\!10^{-4}$ & design & Adapter learning rate. \\
\texttt{use\_8bit\_adamw}             & \texttt{true}  & design       & Applies only to the full fine-tuning fallback; the LoRA path overrides it and uses 32-bit AdamW. \\
\texttt{batch\_size}                 & $4$          & design       & Per-step batch, set for graphics memory. \\
\texttt{grad\_accum\_steps}           & $4$          & design       & Holds the effective batch at $16$. \\
\texttt{epochs\_per\_cycle}           & $3$          & design       & SIL training duration per successful cycle. \\
\texttt{warmup\_steps}               & $500$        & design       & Linear warmup before the standard schedule. \\
\texttt{min\_u\_stored\_for\_training}  & $0.70$       & calibrated   & $\tau_{\text{train}}$; lowered from $0.75$ before the main run. \\
\texttt{pool\_reweighting\_enabled}   & \texttt{true}  & design       & Temperature-mixed reweighting across training benchmarks. \\
\texttt{pool\_reweighting\_\allowbreak temperature} & $2.0$      & design       & Mixture temperature across training benchmarks. \\
\texttt{pool\_reweighting\_\allowbreak upsample\_cap} & $3.0$     & calibrated   & Upper bound on per-benchmark upsampling factor. \\
\texttt{pool\_reweighting\_max\_share} & $0.40$       & design       & Per-benchmark hard share ceiling. \\
\texttt{loop\_distinct4\_threshold}   & $0.25$       & calibrated   & SIL-pool loop filter on token-4-gram distinct ratio. \\
\midrule
\multicolumn{4}{@{}l}{\textit{Verifier ensemble (\Cref{sec:verifier})}} \\
\texttt{verifier\_backend}           & \texttt{"minicheck"} & literature & Short-hypothesis backend; Tang et al.\ 2024. \\
\texttt{minicheck\_model}            & \texttt{lytang/\allowbreak MiniCheck-\allowbreak Flan-T5-\allowbreak Large} & literature & Verifier checkpoint. \\
\texttt{cross\_encoder\_model}        & \texttt{BAAI/\allowbreak bge-reranker-\allowbreak v2-m3} & design & Shared reranker; consumed by grounding and question-answer match. \\
\texttt{verifier\_retrieve\_k}        & $20$  & design & Top-$k$ passages retrieved before reranking. \\
\texttt{verifier\_rerank\_k}          & $3$   & design & Top-$N$ retained after reranking. \\
\texttt{atomic\_min\_tokens}          & $12$  & calibrated & Length gate below which atomic decomposition is bypassed. \\
\texttt{mc\_dropout\_k},\allowbreak{} \texttt{mc\_dropout\_rate} & $5$, $0.1$ & literature & Gal and Ghahramani 2016. \\
\texttt{sc\_chains\_m}                & $3$   & literature & Wang et al.\ 2022 self-consistency. \\
\texttt{se\_samples\_k},\allowbreak{} \texttt{se\_temperature} & $10$, $1.0$ & literature & Farquhar et al.\ 2024 semantic-entropy protocol. \\
\texttt{disable\_h\_norm}             & \texttt{true} & calibrated & Semantic-entropy signal dropped at cycle zero: learned weight near zero, largest verifier cost. \\
\texttt{disable\_p\_contra}           & \texttt{true} & design & Structurally zero under the MiniCheck binary judge. \\
\midrule
\multicolumn{4}{@{}l}{\textit{Generation and retrieval}} \\
\texttt{base\_model\_name}            & \texttt{Qwen/\allowbreak Qwen2.5-3B-\allowbreak Instruct} & design & Apache-licensed three-billion-parameter instruction-tuned backbone. \\
\texttt{rag\_max\_context\_tokens}     & $384$ & design & Sets the input-length limit of the retrieval prompt ($8 \times 384 + 512$ tokens). \\
\texttt{rag\_max\_new\_tokens}         & $512$ & design & Generation budget under retrieval-augmented prompt. \\
\texttt{rag\_do\_sample}              & \texttt{false} & design & Greedy decoding for reproducibility. \\
\texttt{rag\_index\_type}             & \texttt{"ivf\_pq"} & design & FAISS index backend. \\
\texttt{rag\_faiss\_nlist},\allowbreak{} \texttt{nprobe\_tier3},\allowbreak{} \texttt{nprobe\_tier2} & $32768$, $64$, $16$ & design & IVF-PQ density and adaptive probe budget. \\
\texttt{rag\_top\_k}                  & $5$ & literature & Karpukhin et al.\ 2020 DPR standard. \\
\midrule
\multicolumn{4}{@{}l}{\textit{Memory store (\Cref{sec:memory-store})}} \\
\texttt{sbert\_model}                & \texttt{sentence-\allowbreak transformers/\allowbreak all-mpnet-\allowbreak base-v2} & literature & 768-dimensional sentence encoder. \\
\texttt{embedding\_dim}              & $768$ & literature & Must match the encoder output dimension. \\
\texttt{max\_memory\_size}            & $1\,000\,000$ & design & Capacity for thesis-scale and production-scale runs. \\
\texttt{memory\_index\_type}          & \texttt{"ivf\_pq"} & design & FAISS backend from $20{,}000$ entries; below that, and so throughout the run, search is exact. \\
\texttt{novelty\_threshold}          & $0.95$ & design & Skip storage if nearest-neighbour cosine similarity exceeds this. \\
\texttt{pruning\_trigger}            & $0.95$ & design & Trigger pruning when memory reaches this fraction of capacity (never reached). \\
\texttt{pruning\_amount}             & $0.20$ & design & Remove the bottom fraction by value score. \\
\texttt{value\_importance\_weight},\allowbreak{} \texttt{value\_recency\_weight} & $0.70$, $0.30$ & design & Value-score decomposition. \\
\midrule
\multicolumn{4}{@{}l}{\textit{Experiment scope}} \\
\texttt{num\_cycles}                 & $10$  & pre-specified & Registered horizon; the run stopped after cycle five, as planned after the cycle-four rollback. \\
\texttt{calibration\_set\_size}       & $500$ & design & Cycle-zero calibration-fold size per training benchmark. \\
\texttt{purity\_validation\_set\_size} & $500$ & design & Cycle-zero purity-validation slice per training benchmark. \\
\texttt{questions\_per\_cycle}        & $5\,000$ & design & Stream-chunk question budget per cycle. \\
\end{longtable}
}

Only four groups of quantities are refitted at each kept cycle boundary: the pooled and per-benchmark pre-routing temperatures, the per-benchmark safety floors, the per-signal isotonic curves, and the boost weights and intercept (\Cref{tab:calibration-trajectory}); every value in the registry is fixed for the whole run. The model and hardware are given in \Cref{app:repro-hardware,app:repro-versions}, and the cycle-zero validation gate in \Cref{sec:adj-gate}.

The retrieval-utility classifier (\Cref{sec:ruc}) has no fields in \verb|caem/config.py|. Its operating point lives in \verb|caem/ruc/v2/feature_spec.json|: the threshold $\tau_{\text{RUC}} = 0.375$, the regularisation strength $C = 0.5$, the thirty-six-feature ordering, the standard-scaler and logistic-head joblib paths, and the model-type guard \codepath{"logistic\_regression\_scaled"}. The classifier is loaded once at pipeline construction from the on-disk specification and is not refit at any cycle boundary; the architectural rationale is documented in \Cref{sec:ruc}.

\section{Runbook step graph}
\label{app:runbook}

The end-to-end execution sequence partitions into a pre-run setup block, a six-step main cycle loop, and a post-run analysis block. The pre-run block runs once before the first cycle; the main cycle loop repeats per cycle; the post-run block runs once after the final cycle closes. Paths: \verb|archive/scripts_superseded/old_runners/run_phase1a.sh| (the top-level driver used for the run, now archived), \verb|scripts/run_experiment.py| (the in-cycle implementation), and \verb|docs/PRODUCTION_RUNBOOK.md| Section~5 (the production version).

\paragraph{Pre-run setup (Steps 1--6):}
\begin{enumerate}
\item \textbf{Prompt smoke test} (\verb|step_5_9_prompt_smoke|). Renders $5$ samples per benchmark against the registered prompt templates and asserts the answer-line contract. Gate: every sample renders with a non-empty answer line.
\item \textbf{Alias dictionary build} (\verb|step_5_9_5_alias_dict|). Builds \verb|data/alias_dict.json| from the benchmarks' gold aliases. The alias signal is not part of the deployed composite; the dictionary is kept for ablation only.
\item \textbf{Cold-start seeding} (\verb|step_6_reseed| and \verb|step_6_5_em_prune|). Answers questions from each training benchmark's seed pool under a cold-start storage threshold of $0.50$ until about $200$ per benchmark are stored ($603$ answers), then removes the stored answers that are wrong against the gold labels, leaving $431$ entries (\Cref{sec:memory-store}).
\item \textbf{Cycle-zero evaluation and calibration-fold scoring} (\verb|step_7_0_cycle0|). Runs the pristine generator against the held-out evaluation fold and the calibration fold ($500$ samples per training benchmark, $1{,}500$ pooled). Writes per-sample nine-signal vectors and ground-truth labels. Gate: every JSON has the expected sample count and no NaN signals.
\item \textbf{Composite calibration and validation gate} (\codepath{step\_7\_0\_calibrate}, \codepath{step\_7\_0\_2\_5\_rescore\_eval}, \codepath{step\_7\_0\_3\_validate\_weights}). Fits the per-signal isotonic curves, the per-benchmark shrinkage-blended children, the boost-layer logistic regression, the per-benchmark temperature scalars, and the per-benchmark safety floors. Runs the cycle-zero validation gate (\Cref{sec:adj-gate}). Gate: cycle-zero canonical composite locked at \verb|outputs/cycle_0/composite_calibration.json|; never overwritten by later cycles.
\item \textbf{Retention baseline freeze and snapshot upload} (\verb|step_19_2_slice| and \verb|step_gdrive_upload_pre_main|). Records the pristine-model probe accuracies on MMLU, TriviaQA-test, and CommonsenseQA-test as the cycle-zero anchor for the $0.93$ retention floor. Uploads the pre-main snapshot so that a lost machine can be recovered from it (\Cref{app:repro-snapshot}).
\end{enumerate}

\paragraph{Main cycle loop (Steps 1--6, repeating per cycle):}
\begin{enumerate}
\item \textbf{Self-improvement fine-tune:} Draws the SIL training pool from memory entries above $\tau_{\text{train}} = 0.70$, applies the pool-reweighting layer, runs the LoRA fine-tune for three epochs, and probes the retention panel. Gate: every probe ratio against the cycle-zero baseline clears $0.93$. On failure, the adapter is rolled back to the previous cycle's checkpoint, the memory is preserved, and the recalibration and re-verification of step 2 are skipped under the asymmetric-rollback contract.
\item \textbf{Cycle-boundary recalibration:} Re-scores the calibration fold under the post-fine-tune model, refits the per-benchmark temperature scalars, refits the per-signal isotonic curves under the shrinkage prior, and runs the retroactive re-verification pass against the stored memory under the recalibrated verifier. Skipped on aborted cycles.
\item \textbf{Stream-chunk pre-routing pass:} Runs the per-query pipeline on the cycle's stream chunk ($2{,}000$ FEVER, $2{,}000$ TriviaQA, $700$ CommonsenseQA samples), accumulating storage admissions, deferred entries, abstentions, and discards.
\item \textbf{Stream-chunk verification and storage:} Applies the calibrated probability composite, the four-outcome decision tree, and the storage admission. Writes \verb|memory_store_cycle_N.{faiss,meta}| and \verb|deferred_buffer_cycle_N.pkl|.
\item \textbf{Held-out evaluation:} Runs the per-query pipeline on the held-out evaluation fold ($300$ samples per benchmark across the five-benchmark panel) under \verb|store_to_memory=False| so the eval fold does not contaminate the memory pool.
\item \textbf{Checkpoint and halt-trigger evaluation:} Writes the cycle's coverage diagnostic, the per-tier hit-rate roll-up, and the per-cycle composite calibration snapshot. Evaluates the automatic early-stop check.
\end{enumerate}

\paragraph{Post-run analysis:} the baseline re-run under the comparison protocol and the rebuild of every reported number; the two commands are given in \Cref{app:repro-steps}.

\paragraph{Asymmetric-rollback contract:} triggered when any retention probe drops below $0.93$ of its cycle-zero value in step 1 of the main cycle loop. The adapter reverts to the previous kept cycle's checkpoint, but the memory does not: answers stored during the failed cycle are kept. Step 2 (recalibration, retroactive re-verification and deferred reconsideration) is skipped, because the weights it would re-score are unchanged, and the re-verification log records the skip. Steps 3 to 6 still run, so memory grows and the rolled-back model is evaluated. In the run the contract fired once, at cycle four, when the TriviaQA-test probe ratio fell to $0.886$; four of the five fine-tunes were kept.
  % Implementation Reference
% ============================================================================
% Appendix B — Cycle-zero configuration and example answers (final drafting pass 2026-10-19).
% ============================================================================

\chapter{Cycle-zero configuration and example answers}
\label{app:cycle0-catalog}


\section{The $100$-variant composite parameter sweep}
\label{app:sweep-full}

The locked configuration selected from the sweep, with the rationale for each parameter, is given in \Cref{tab:sweep-top}; the locked variant and its neighbours are compared in \Cref{tab:sweep-variants}. The full sweep over the pre-specified $\tau_{\text{store}} \times C \times \alpha_{\text{shrink}}$ grid ($4 \times 5 \times 5 = 100$ variants) is recorded in the cycle-zero sweep archive with one record per variant; no per-variant comparison table is reproduced in the thesis. Per-benchmark sweep diagnostics for any variant in the grid are recoverable by reading the corresponding JSON entry.

% Rationale cells and caption hand-corrected 2026-10.
% Auto-generated by scripts/aggregate_section_5_7.py
% Locked configuration parameters and rationale (replaces stale v1 sweep_top.tex)
% Hand-corrected 2026-10-15 (see ref_check/fixes_applied.md)
\begin{table}[!htbp]
\centering\small
\begin{tabular}{@{}p{4.2cm}p{1.8cm}p{8.5cm}@{}}
\toprule
\textbf{Parameter} & \textbf{Locked value} & \textbf{Selection rationale} \\
\midrule
Storage cut-point $\tau_{\text{store}}$ & $0.60$ & Stores about twice as many answers as $0.65$ at acceptable precision (\Cref{tab:sweep-variants}) \\
Deferral cut-point $\tau_{\text{defer}}$ & $0.45$ & Pre-specified $0.15$ below storage cut \\
Abstention floor $\phi_{\text{pg}}$ & $0.20$ & Early-exit rule \Cref{eq:early-exit-rule} \\
Boost regularisation $C$ & $0.010$ & Strongest regularisation in the grid; a hundredfold tighter penalty than the library default $C = 1.0$ \\
Shrinkage strength $\alpha_{\text{shrink}}$ & $0.6$ & Pulls small per-benchmark fits toward the pooled fit; $\alpha = 1$ stores more correct answers but separates less well on the validation fold (\Cref{tab:sweep-variants}) \\
Retroactive prune $\tau_{\text{retro}}$ & $0.50$ & $0.10$ below storage cut by registration \\
Retention floor $\rho_{\min}$ & $0.93$ & Multi-modal probe panel rollback threshold \\
\bottomrule
\end{tabular}
\caption[Locked configuration carried into every cycle.]{Locked configuration carried into every cycle of the main run, with the reason for each value.}
\label{tab:sweep-top}
\end{table}


\section{Composite weight table}
\label{app:composite-weights}

\Cref{tab:composite-calibration} gives the direction and boost weight of each signal in the cycle-zero composite; the refitted weights of later cycles are in \Cref{tab:calibration-trajectory}. The per-signal isotonic curves are step functions kept in each cycle's composite record, and the storage and deferral cut-points stay fixed throughout.

% Caption shortened by hand in the final drafting pass (2026-10); detail moved to chapter 5 text.
% Auto-generated by scripts/aggregate_section_5_7.py
% Cycle-zero composite calibration (per-signal Pearson rho + boost-layer weights).
% Source: outputs/cycle_0/composite_calibration.json (boost weights, intercept).
% Pearson rho re-computed from outputs/cycle_0/calibration/calibration_fold_samples.json.
\begin{table}[!htbp]
\centering\small
\begin{tabular}{lccr}
\toprule
\textbf{Signal} & \textbf{Pearson $\rho$} & \textbf{Direction} & \textbf{Boost weight} \\
\midrule
$p^{\text{mean}}_{\text{ground}}$ & +0.101 & increasing & +0.184 \\
$p^{\text{max}}_{\text{ground}}$ & +0.140 & increasing & +0.379 \\
$p^{\text{atomic}}_{\text{ground}}$ & +0.101 & increasing & +0.184 \\
$p_{\text{entail}}$ & -0.156 & decreasing & +0.067 \\
$q_{\text{match}}$ & +0.299 & increasing & +0.491 \\
$u_{\text{internal}}$ & +0.425 & increasing & +0.334 \\
$u_{\text{token}}$ & +0.119 & increasing & +0.126 \\
$u_{\text{dropout}}$ & -0.411 & decreasing & +0.190 \\
$s_{\text{avg}}$ & +0.098 & increasing & +0.175 \\
\midrule
\multicolumn{3}{r}{Boost intercept} & +0.265 \\
\bottomrule
\end{tabular}
\caption[The cycle-zero composite.]{Direction of each signal's calibration curve (from its correlation with correctness) and its boost weight in the cycle-zero composite, fitted on the $1{,}500$-question calibration fold.}
\label{tab:composite-calibration}
\end{table}


\section{Example answers under each status tag}
\label{app:production-envelope}

\Cref{tab:production-envelope-examples} shows real answers from the evaluation fold at C5 as the user would see them under each status tag (\Cref{sec:decision}), with whether each answer was right. A verified answer can still be wrong, and a withheld answer can be right; the tag is a probability, not a guarantee (\Cref{tab:status-tags}).

% Generated by scripts/build_final_tables.py from outputs/analysis/chm_judged.json (label em_final); do not edit.
\begin{table}[!htbp]
\centering\footnotesize
\setlength{\tabcolsep}{4pt}
\begin{tabular}{@{}l>{\raggedright\arraybackslash}p{6.4cm}>{\raggedright\arraybackslash}p{2.6cm}rc@{}}
\toprule
\textbf{Tag} & \textbf{Question} & \textbf{Answer shown} & \textbf{Conf.} & \textbf{Correct} \\
\midrule
verified & Libya was a colony of which European country prior to World War II? & Italy & $71\%$ & yes \\
verified & Hippocampus is the Latin name for which emblematic marine creature? & hippocampus & $70\%$ & no \\
provisional & Thomas Hardy's famous novel is 'Far from the Madding (What?)'? & Crowd & $52\%$ & yes \\
provisional & The word barber - meaning a men's/boy's hairdresser - is from Latin... & parrhes & $46\%$ & no \\
conflicting & Mrs Lovett made meat pies from material supplied by whom? & Sweeney Todd & $40\%$ & yes \\
conflicting & Sound engineer Peter Lodge first announced which phrase on the Lond... & "the world is a branch line of the pi... & $45\%$ & no \\
insufficient & What did Ukrainian prime minister Yulia Timoshenko change about her... & (withheld) & $34\%$ & yes \\
insufficient & Who is the nude bow-carrying statue in London's Piccadilly Circus? & (withheld) & $43\%$ & no \\
\bottomrule
\end{tabular}
\caption[What the user sees for each status tag.]{TriviaQA answers at C5 as the user sees them, one right and one wrong per status tag where both occur. Verified answers are shown plainly, provisional ones with a caveat that the system will reconsider them, conflicting ones with a warning, and insufficient-evidence answers are replaced by a message; every answer carries its calibrated probability (Conf.).}
\label{tab:production-envelope-examples}
\end{table}

  % Cycle 0 Evidence Catalog
% ============================================================================
% Appendix C — Per-cycle artifact dictionary
% A flat catalog of the per-cycle artifacts produced by the realised
% six-cycle run. Body prose in Chapter 5 reports trajectories; this
% appendix is the per-cycle index against the on-disk JSON snapshots.
% ============================================================================

\chapter{Per-cycle records}
\label{app:per-cycle}

This appendix holds the per-cycle tables that \Cref{ch:evaluation} refers to but does not print: the memory trajectory with the four mechanisms that change it, and the further evaluation tables. The calibration refit by cycle is in \Cref{tab:calibration-trajectory}, and the fine-tuning and retention record in \Cref{tab:sil-training-trajectory}.

\section{Per-cycle memory size and storage rate}
\label{app:per-cycle-memory}

\Cref{tab:app-per-cycle-memory} reports the memory pool size at each cycle close together with the four mechanisms that move the pool size across the cycle boundary (stream-chunk admissions through the storage gate, retroactive re-verification prunes against the stored pool under the recalibrated verifier, deferred-buffer promotions from the bounded reconsider queue, and consolidation removals of near-duplicate episodes) and the cycle's net change in pool size. The cycle-four row records no retroverify prunes, deferred-buffer promotions, or consolidation removals under the asymmetric-rollback contract that skips these passes on the aborted cycle; the cycle-four stream-chunk admissions and the cycle-four pool growth proceed unchanged because the fresh-data side of the architecture runs unconditionally.

\begin{table}[!htbp]
\centering\footnotesize
\setlength{\tabcolsep}{3pt}
\renewcommand{\arraystretch}{1.1}
\begin{tabular}{@{}lrrrrrr@{}}
\toprule
\textbf{Cycle} & \textbf{Pool size} & \textbf{Stream-chunk} & \textbf{Retroverify} & \textbf{Deferred-buffer} & \textbf{Consolidation} & \textbf{Net} \\
 & \textbf{at close} & \textbf{admissions} & \textbf{prunes} & \textbf{promotions} & \textbf{removals} & \textbf{change} \\
\midrule
C0 & $431$    & $431$    & $0$ & $0$ & $0$ & $+431$ \\
C1 & $2\,818$ & $2\,443$ & $56$  & $0$   & $0$  & $+2\,387$ \\
C2 & $5\,151$ & $2\,477$ & $271$ & $131$ & $3$  & $+2\,333$ \\
C3 & $7\,477$ & $2\,469$ & $352$ & $220$ & $11$ & $+2\,326$ \\
C4$^{\ast}$ & $9\,961$ & $2\,484$ & ---   & ---   & --- & $+2\,484$ \\
C5 & $12\,175$ & $2\,343$ & $382$ & $299$ & $44$ & $+2\,214$ \\
\bottomrule
\end{tabular}
\caption[Per-cycle memory store trajectory]{Per-cycle memory store trajectory. The pool grows from the $431$ cold-start entries to $12{,}175$ at the cycle-five close. The four columns do not add up exactly to the net change, by at most two entries per cycle. The cycle-zero admissions row records the cold-start seed rather than a stream-chunk pass. The asterisk marks cycle four, which was aborted under the retention guard; the retroverify, deferred-reconsider, and consolidation passes are skipped under the asymmetric-rollback contract and shown as dashes. The per-cycle stream-chunk admission count stays in the band $[2\,343, 2\,484]$ across cycles one through five against a stream-chunk candidate budget of $4\,700$ samples per cycle ($2{,}000$ FEVER, $2{,}000$ TriviaQA, $700$ CommonsenseQA), so the storage gate admits approximately half of every candidate stream chunk. The admissions count the entries actually added to the store; they sit slightly below the gate-pass counts $n_{\text{stored}}$ of \Cref{tab:storage-pool-quality} because the memory store's novelty filter rejects near-duplicates of existing entries between the gate and the insert. The cycle-five admission count is the cycle-close pool minus the post-re-verification pool ($12\,175 - 9\,832$). The consolidation-removal column is the episodes-removed count of \Cref{tab:memory-store-trajectory}. Admissions minus prunes plus promotions minus consolidation removals reproduces the net change at every cycle except a residual of one entry at cycle two and two entries at cycle five; the residual is promoted entries dropped as duplicates of entries already in the store (the cycle-two consolidation pass sees $2\,948$ entries against $2\,818 + 131 = 2\,949$). Sources: the per-cycle memory-store metadata side-files supply the cycle-close pool sizes (the length of the metadata list inside the pickled snapshot); the per-cycle retroverification records supply the per-cycle prune counts; per-cycle deferred-buffer promotion counts are read from the deferred-reconsideration log emissions in the experiment run log.}
\label{tab:app-per-cycle-memory}
\end{table}

\section{Additional per-cycle tables from the evaluation}
\label{app:extra-tables}

The tables in this section support statements made in \Cref{ch:evaluation}; each is referred to from the paragraph that uses it.

% Caption shortened by hand in the final drafting pass (2026-10); detail moved to chapter 5 text.
% Auto-generated by scripts/aggregate_ces_caem_trajectory.py
% CAEM-only per-cycle Composite Evaluation Score trajectory, computed under the
% canonical eval/metrics.py:ces_score axis formulas with the v2 multi-modal
% retention rule (worst probe across MMLU + TQA-test + CSQA-test) rather than
% the v1 MMLU-only rule.
% Formulas:
%   ACC = pooled exact match across the 5-benchmark eval panel (TruthfulQA
%         under em_llm_judged Haiku rule).
%   EPI = clamp_{[0,1]}(1 - pooled CHM).
%   RET = min(MMLU ratio, TQA-test ratio, CSQA-test ratio, 1.0); C0 = 1.0 by
%         construction (baseline against itself).
%   CAL = 1 - 2 * min(ECE_after, 0.5).
%   VER = max(2 * BA - 1, 0), with BA = Phi(d / sqrt(2)) the balanced-accuracy
%         proxy of the calibrated probability composite, computed from the
%         per-cycle composite Cohen's d on the training-panel calibration fold.
%   CES = (prod_i clamp_{[eps,1]}(axis_i)) ^ (1/5) with eps = 0.01.
% Cycle four reuses cycle-three CAL and VER because the cycle-four refit was
% skipped under the retention guard's asymmetric-rollback contract.
% Generated by scripts/build_final_tables.py from outputs/analysis/chm_judged.json (label em_final); ACC/EPI/CES columns.
\begin{table}[!htbp]
\centering\small
\setlength{\tabcolsep}{6pt}
\begin{tabular}{@{}lrrrrrr@{}}
\toprule
\textbf{Cycle} & \textbf{ACC} & \textbf{EPI} & \textbf{RET} & \textbf{CAL} & \textbf{VER} & \textbf{CES} \\
\midrule
C0 & $0.517$ & $0.845$ & $1.000$ & $0.898$ & $0.174$ & $0.585$ \\
C1 & $0.533$ & $0.869$ & $1.000$ & $0.912$ & $0.325$ & $0.672$ \\
C2 & $0.535$ & $0.873$ & $0.949$ & $0.901$ & $0.415$ & $0.698$ \\
C3 & $0.516$ & $0.863$ & $0.976$ & $0.919$ & $0.407$ & $0.695$ \\
C4$^{\ast}$ & $0.517$ & $0.864$ & $\mathbf{0.886}$ & $0.919$ & $0.407$ & $\mathbf{0.682}$ \\
\rowcolor{gray!18}
\textbf{C5} & $0.508$ & $0.866$ & $0.958$ & $0.915$ & $\mathbf{0.448}$ & $\mathbf{0.704}$ \\
\bottomrule
\end{tabular}
\par\vspace{2pt}{\footnotesize\parbox{\linewidth}{$^{\ast}$ Cycle four was aborted under the retention guard ($\rho_{\min}=0.93$); CAL and VER reuse cycle three.}\par}
\captionsetup{singlelinecheck=false}
\caption[Composite evaluation score by cycle.]{Composite evaluation score (\Cref{eq:ces}) and its five axes by cycle, for CAEM only: three of the axes do not exist for single-pass baselines.}
\label{tab:ces-caem-trajectory}
\end{table}


% Caption shortened by hand in the final drafting pass (2026-10); detail moved to chapter 5 text.
% Auto-generated from outputs/tab_decision_matrix_trajectory.csv
% Per-bench per-decision (STORE / DEFER / DISCARD) outcome counts and STORE precision
% at C2 (within-trajectory best) and C5 (trajectory-end).
% Hand-corrected 2026-10-15 (see ref_check/fixes_applied.md)
\begin{table}[H]
\centering\footnotesize
\setlength{\tabcolsep}{4pt}
\resizebox{\textwidth}{!}{%
\begin{tabular}{@{}llrrcrrrr@{}}
\toprule
\textbf{Cycle} & \textbf{Benchmark} & \textbf{Store $\checkmark$} & \textbf{Store $\times$} & \textbf{Store $\pi$} & \textbf{Defer $\checkmark$} & \textbf{Defer $\times$} & \textbf{Discard $\checkmark$} & \textbf{Discard $\times$} \\
\midrule
\multicolumn{9}{@{}l}{\textit{Cycle two (the within-trajectory best cycle)}} \\
C2 & FEVER & $803$ & $240$ & $0.770$ & $344$ & $419$ & $58$ & $45$ \\
C2 & TriviaQA & $605$ & $209$ & $0.743$ & $262$ & $329$ & $61$ & $163$ \\
C2 & CommonsenseQA & $465$ & $172$ & $0.730$ & $34$ & $19$ & $1$ & $0$ \\
C2 & TruthfulQA & $33$ & $58$ & $0.363$ & $23$ & $48$ & $11$ & $38$ \\
C2 & StrategyQA & $57$ & $20$ & $0.740$ & $125$ & $73$ & $1$ & $0$ \\
\midrule
\multicolumn{9}{@{}l}{\textit{Cycle five (trajectory-end close)}} \\
C5 & FEVER & $764$ & $210$ & $0.784$ & $474$ & $432$ & $10$ & $2$ \\
C5 & TriviaQA & $532$ & $269$ & $0.664$ & $284$ & $430$ & $8$ & $13$ \\
C5 & CommonsenseQA & $428$ & $177$ & $0.707$ & $32$ & $32$ & $11$ & $6$ \\
C5 & TruthfulQA & $23$ & $52$ & $0.307$ & $38$ & $101$ & $3$ & $13$ \\
C5 & StrategyQA & $45$ & $10$ & $0.818$ & $100$ & $53$ & $0$ & $1$ \\
\bottomrule
\end{tabular}%
}
\caption[Correct and wrong answers per decision at C2 and C5.]{Correct and wrong answers in each decision at C2 and C5, with the precision of the stored class. Training-panel rows come from the stream chunks, transfer-panel rows from the evaluation fold.}
\label{tab:decision-matrix-trajectory}
\end{table}


% Caption shortened by hand in the final drafting pass (2026-10); detail moved to chapter 5 text.
% Auto-generated from outputs/full_run/eval/{bench}_cycle{c}.json
% Per-benchmark abstention rate (eval fold, n=300/bench).
% Abstention = decision == ABSTAIN OR prediction starts with <abstain>.
\begin{table}[H]
\centering\footnotesize
\setlength{\tabcolsep}{6pt}
\begin{tabular}{@{}lrrrrr@{}}
\toprule
\textbf{Cycle} & \textbf{FEVER} & \textbf{TriviaQA} & \textbf{CSQA} & \textbf{TruthfulQA} & \textbf{StrategyQA} \\
\midrule
C0 & $64.3\%$ & $29.0\%$ & $0.0\%$ & $33.3\%$ & $43.0\%$ \\
C1 & $24.3\%$ & $17.7\%$ & $0.3\%$ & $27.3\%$ & $2.3\%$ \\
C2 & $5.7\%$ & $17.3\%$ & $1.0\%$ & $29.3\%$ & $8.0\%$ \\
C3 & $10.7\%$ & $20.0\%$ & $0.0\%$ & $25.0\%$ & $4.0\%$ \\
C4 & $10.0\%$ & $21.0\%$ & $0.0\%$ & $25.3\%$ & $2.3\%$ \\
C5 & $5.3\%$ & $26.0\%$ & $0.3\%$ & $23.0\%$ & $30.3\%$ \\
\bottomrule
\end{tabular}
\caption[Abstention rate by benchmark and cycle.]{Share of evaluation-fold questions on which the system abstained ($300$ per benchmark).}
\label{tab:abstention-per-bench}
\end{table}


% Caption shortened by hand in the final drafting pass (2026-10); detail moved to chapter 5 text.
% Auto-generated from outputs/tab_per_tier_em.csv
% Hand-corrected 2026-10-15 (see ref_check/fixes_applied.md)
% Per-bench per-tier EM at C0, C3, C5 anchors (snapshot of the full 6-cycle CSV).
\begin{table}[H]
\centering\footnotesize
\setlength{\tabcolsep}{4pt}
\begin{tabular}{@{}llrrrrrrr@{}}
\toprule
\textbf{Cycle} & \textbf{Benchmark} & \textbf{T1 EM} & \textbf{T1 n} & \textbf{T2 EM} & \textbf{T2 n} & \textbf{T3 EM} & \textbf{T3 n} & \textbf{Pooled EM} \\
\midrule
\multicolumn{9}{@{}l}{\textit{Cycle zero (cold-start; no Tier 1 firings)}} \\
C0 & FEVER & --- & $0$ & $55.3$ & $132$ & $43.5$ & $168$ & $48.7$ \\
C0 & TriviaQA & --- & $0$ & $46.1$ & $89$ & $46.4$ & $211$ & $46.3$ \\
C0 & CommonsenseQA & --- & $0$ & $78.5$ & $275$ & $52.0$ & $25$ & $76.3$ \\
C0 & TruthfulQA & --- & $0$ & $42.6$ & $197$ & $29.1$ & $103$ & $38.0$ \\
C0 & StrategyQA & --- & $0$ & $68.4$ & $38$ & $59.9$ & $262$ & $61.0$ \\
\midrule
\multicolumn{9}{@{}l}{\textit{Cycle two (the within-trajectory joint EM/CHM best cycle)}} \\
C2 & FEVER & --- & $0$ & $51.6$ & $122$ & $54.5$ & $178$ & $53.3$ \\
C2 & TriviaQA & --- & $0$ & $51.0$ & $96$ & $48.5$ & $204$ & $49.3$ \\
C2 & CommonsenseQA & --- & $0$ & $72.4$ & $275$ & $64.0$ & $25$ & $71.7$ \\
C2 & TruthfulQA & $100.0$ & $1$ & $35.3$ & $204$ & $25.3$ & $95$ & $32.3$ \\
C2 & StrategyQA & --- & $0$ & $66.4$ & $131$ & $64.5$ & $169$ & $65.3$ \\
\midrule
\multicolumn{9}{@{}l}{\textit{Cycle five (trajectory-end close)}} \\
C5 & FEVER & --- & $0$ & $50.0$ & $130$ & $45.9$ & $170$ & $47.7$ \\
C5 & TriviaQA & $100.0$ & $1$ & $42.3$ & $111$ & $46.8$ & $188$ & $45.3$ \\
C5 & CommonsenseQA & --- & $0$ & $74.5$ & $247$ & $58.5$ & $53$ & $71.7$ \\
C5 & TruthfulQA & $100.0$ & $1$ & $32.1$ & $193$ & $25.5$ & $106$ & $30.0$ \\
C5 & StrategyQA & --- & $0$ & $63.9$ & $72$ & $64.0$ & $228$ & $64.0$ \\
\bottomrule
\end{tabular}
\caption[Per-tier EM by benchmark at C0, C2 and C5.]{EM (percent) and number of queries per tier and benchmark at C0, C2 and C5.}
\label{tab:per-tier-em}
\end{table}


% Generated by scripts/build_final_tables.py from outputs/analysis/chm_judged.json (label em_final); do not edit.
\begin{table}[H]
\centering\scriptsize
\setlength{\tabcolsep}{3pt}
\begin{tabular}{@{}l|cc|cc|cc|cc|cc@{}}
\toprule
 & \multicolumn{2}{c|}{$[0.60, 0.65)$} & \multicolumn{2}{c|}{$[0.65, 0.70)$} & \multicolumn{2}{c|}{$[0.70, 0.80)$} & \multicolumn{2}{c|}{$[0.80, 0.90)$} & \multicolumn{2}{c}{$[0.90, 1.00]$} \\
\textbf{Cycle} & $n$ & EM & $n$ & EM & $n$ & EM & $n$ & EM & $n$ & EM \\
\midrule
\multicolumn{11}{@{}l}{\textit{(a) Evaluation fold}} \\
C0 & $174$ & $0.759$ & $104$ & $0.731$ & $46$ & $0.804$ & $10$ & $0.700$ & $2$ & $1.000$ \\
C1 & $139$ & $0.655$ & $145$ & $0.697$ & $170$ & $0.776$ & $59$ & $0.814$ & $2$ & $1.000$ \\
C2 & $113$ & $0.566$ & $155$ & $0.729$ & $225$ & $0.716$ & $35$ & $0.771$ & $6$ & $0.833$ \\
C3 & $85$ & $0.600$ & $201$ & $0.726$ & $244$ & $0.738$ & $34$ & $0.706$ & $1$ & $1.000$ \\
C4 & $91$ & $0.648$ & $203$ & $0.704$ & $231$ & $0.732$ & $40$ & $0.750$ & $1$ & $1.000$ \\
C5 & $164$ & $0.677$ & $159$ & $0.686$ & $164$ & $0.720$ & $14$ & $0.929$ & $3$ & $0.667$ \\
\midrule
\multicolumn{11}{@{}l}{\textit{(b) SIL stream chunks}} \\
C1 & $596$ & $0.666$ & $551$ & $0.693$ & $832$ & $0.798$ & $419$ & $0.845$ & $58$ & $0.914$ \\
C2 & $512$ & $0.635$ & $534$ & $0.702$ & $1\,043$ & $0.785$ & $356$ & $0.871$ & $49$ & $0.898$ \\
C3 & $423$ & $0.572$ & $730$ & $0.675$ & $1\,012$ & $0.768$ & $292$ & $0.880$ & $33$ & $1.000$ \\
C4 & $427$ & $0.578$ & $682$ & $0.705$ & $1\,060$ & $0.776$ & $316$ & $0.870$ & $37$ & $0.973$ \\
C5 & $732$ & $0.637$ & $753$ & $0.716$ & $752$ & $0.793$ & $127$ & $0.858$ & $16$ & $0.875$ \\
\bottomrule
\end{tabular}
\caption[Stored scores and accuracy by score band.]{Stored training-panel answers by score band and the EM within each band, on (a) the evaluation fold and (b) the stream chunks.}
\label{tab:u-stored-distribution}
\end{table}


% Caption shortened by hand in the final drafting pass (2026-10); detail moved to chapter 5 text.
% Auto-generated from outputs/full_run/memory_store_cycle_*.faiss.meta set-difference
% (prune cohort identification) + experiment.log DeferredBuffer.reconsider() per-entry
% trace (promotion cohort identification) + offline EM-label lookup against
% HuggingFace gold (FEVER train+dev, TriviaQA train+val+test, CSQA train+val).
% EM labels are not persisted per entry; this reconstruction is a one-off audit.
% Hand-corrected 2026-10-15 (see ref_check/fixes_applied.md)
\begin{table}[H]
\centering\footnotesize
\setlength{\tabcolsep}{4pt}
\resizebox{\textwidth}{!}{%
\begin{tabular}{@{}l|rrrr|rrrr@{}}
\toprule
 & \multicolumn{4}{c|}{\textbf{Retroverify prunes}} & \multicolumn{4}{c}{\textbf{Deferred-buffer promotions}} \\
\textbf{Cycle} & \textbf{Pruned} & \textbf{em-wrong} & \textbf{em-correct} & \textbf{Wrong-caught} & \textbf{Promoted} & \textbf{em-correct} & \textbf{em-wrong} & \textbf{Correct-recovered} \\
\midrule
C1 & $56$ & $0$ & $56$ & $0.0\%$ & $0$ & $0$ & $0$ & --- \\
C2 & $261$ & $114$ & $147$ & $43.7\%$ & $117$ & $65$ & $52$ & $55.6\%$ \\
C3 & $355$ & $122$ & $233$ & $34.4\%$ & $210$ & $119$ & $91$ & $56.7\%$ \\
C4$^{\ast}$ & --- & --- & --- & --- & --- & --- & --- & --- \\
C5 & $422$ & $184$ & $238$ & $43.6\%$ & $293$ & $161$ & $132$ & $54.9\%$ \\
\midrule
\textbf{Pooled} & $\mathbf{1\,094}$ & $\mathbf{420}$ & $\mathbf{674}$ & $\mathbf{38.4\%}$ & $\mathbf{620}$ & $\mathbf{345}$ & $\mathbf{275}$ & $\mathbf{55.6\%}$ \\
\bottomrule
\multicolumn{9}{@{}l}{\footnotesize $^{\ast}$ Cycle four was aborted under the retention guard; both the retroverify pass and the deferred-reconsider pass were skipped under the asymmetric-rollback contract.} \\
\end{tabular}%
}
\caption[Correctness of pruned and promoted entries.]{Share of wrong entries among those removed by re-verification, and of correct entries among those promoted from the deferred buffer, reconstructed from the memory snapshots.}
\label{tab:em-conditional-breakdown}
\end{table}


% CHM cells re-scored with the Haiku-judged TruthfulQA label by scripts/patch_chm_judged_tables.py (2026-10-17).
% Caption shortened by hand in the final drafting pass (2026-10); detail moved to chapter 5 text.
% Auto-generated by scripts/aggregate_tier_chm.py
% Per-cycle per-tier EM + CHM (cor:tier1-floor empirical receipt)
\begin{table}[!htbp]
\centering
\small
\begin{tabular}{@{}c|cc|cc|cc|cc@{}}
\toprule
\textbf{Cycle} & \multicolumn{2}{c|}{\textbf{Tier 1 (memory)}} & \multicolumn{2}{c|}{\textbf{Tier 2 (zero-shot)}} & \multicolumn{2}{c|}{\textbf{Tier 3 (RAG)}} & \multicolumn{2}{c}{\textbf{Pooled}} \\
  & EM & CHM & EM & CHM & EM & CHM & EM & CHM \\
\midrule
C0 & -- & -- & $0.602$ & $0.109$ & $0.482$ & $0.158$ & $0.541$ & $0.134$ \\
C1 & -- & -- & $0.585$ & $0.104$ & $0.502$ & $0.123$ & $0.543$ & $0.114$ \\
C2 & $1.000$ & $0.000$ & $0.568$ & $0.106$ & $0.514$ & $0.115$ & $0.544$ & $0.110$ \\
C3 & $1.000$ & $0.000$ & $0.566$ & $0.109$ & $0.503$ & $0.126$ & $0.532$ & $0.118$ \\
C4 & $1.000$ & $0.000$ & $0.564$ & $0.110$ & $0.504$ & $0.125$ & $0.532$ & $0.118$ \\
C5 & $1.000$ & $0.000$ & $0.537$ & $0.111$ & $0.497$ & $0.122$ & $0.517$ & $0.117$ \\
\bottomrule
\end{tabular}
\caption[Per-tier CHM by cycle.]{CHM of the answers given by each tier, by cycle, on the evaluation fold.}
\label{tab:tier-chm-trajectory}
\end{table}


% Caption shortened by hand in the final drafting pass (2026-10); detail moved to chapter 5 text.
% Auto-generated by scripts/aggregate_section_5_6.py
% Per-cycle retroactive prune rate (cor:self-correction empirical receipt)
\begin{table}[!htbp]
\centering\small
\resizebox{\textwidth}{!}{%
\begin{tabular}{cccccc}
\toprule
\textbf{Cycle} & \textbf{Pool size at pass} & \textbf{Updated} & \textbf{Pruned} & \textbf{Prune rate} & \textbf{Survivor share} \\
\midrule
C1 & 431 & 375 & 56 & 13.0\% & 87.0\% \\
C2 & 2,945 & 2,674 & 271 & 9.2\% & 90.8\% \\
C3 & 5,358 & 5,006 & 352 & 6.6\% & 93.4\% \\
C5 & 10,214 & 9,832 & 382 & 3.7\% & 96.3\% \\
\midrule
\multicolumn{6}{l}{\textit{Cycle 4: retroactive pass skipped under the cycle-four retention abort (skip-retroverify-on-abort patch).}} \\
\bottomrule
\end{tabular}%
}
\caption[Retroactive re-verification by cycle.]{Entries checked, kept and removed by the retroactive pass at each kept cycle; the pass was skipped at the rolled-back cycle four.}
\label{tab:self-correction-check}
\end{table}


% Caption shortened by hand in the final drafting pass (2026-10); detail moved to chapter 5 text.
% Auto-generated by scripts/aggregate_section_5_7.py
% Per-cycle calibration trajectory: temperature scalar (pooled + per-benchmark),
% expected calibration error after refit, boost-layer intercept, per-benchmark
% safety floor drift, and the full per-cycle boost-layer weight matrix on the
% nine deployed signals with |Delta C0->C5| absolute movement.
% Source artefacts: outputs/cycle_0/composite_calibration.json (C0 row);
% outputs/full_run/calibration/calibrated_config_cycle{N}.json (T + ECE +
% per-bench T + per-bench safety floor for N in {1,2,3,5});
% outputs/full_run/cycle_{N}/composite_calibration.json (per-cycle boost
% intercept + per-signal boost weights). Cycle four is omitted because its
% post-fine-tune adapter was rolled back under the retention guard and the
% cycle-boundary refit was skipped by the asymmetric-rollback contract.
\begin{table}[!htbp]
\centering\scriptsize
\setlength{\tabcolsep}{4pt}
\begin{tabular}{@{}lrrrrrr@{}}
\toprule
\textbf{Parameter} & \textbf{C0} & \textbf{C1} & \textbf{C2} & \textbf{C3} & \textbf{C5} & $\boldsymbol{|\Delta_{\text{C0}\to\text{C5}}|}$ \\
\midrule
\multicolumn{7}{@{}l}{\textit{(a) Pooled calibration discipline}} \\
Temperature scalar $T$ (pooled, post-refit) & $20.086$ & $1.089$ & $0.836$ & $0.740$ & $0.869$ & $19.217$ \\
Expected calibration error $\mathrm{ECE}$ after refit & --- & $0.044$ & $0.049$ & $0.041$ & $0.043$ & --- \\
Boost-layer intercept $b_0$ & $+0.265$ & $-0.613$ & $-0.634$ & $-0.565$ & $-0.520$ & $0.785$ \\
\midrule
\multicolumn{7}{@{}l}{\textit{(b) Per-benchmark temperature scalar (post-refit)}} \\
$T_{\text{FEVER}}$ & --- & $14.092$ & $9.899$ & $6.986$ & $4.972$ & --- \\
$T_{\text{TriviaQA}}$ & --- & $14.123$ & $10.007$ & $7.163$ & $11.040$ & --- \\
$T_{\text{CommonsenseQA}}$ & --- & $1.342$ & $1.209$ & $1.083$ & $1.028$ & --- \\
\midrule
\multicolumn{7}{@{}l}{\textit{(c) Per-benchmark safety floor (drift observed only on CommonsenseQA)}} \\
$\phi_{\text{safe}}^{\text{CommonsenseQA}}$ & $0.380$ & $0.380$ & $0.380$ & $0.491$ & $0.569$ & $0.189$ \\
\midrule
\multicolumn{7}{@{}l}{\textit{(d) Pooled boost-layer weights on the nine deployed signals}} \\
$u_{\text{token}}$ & $+0.126$ & $+0.211$ & $+0.177$ & $+0.060$ & $+0.110$ & $0.016$ \\
$u_{\text{dropout}}$ & $+0.190$ & $+0.132$ & $+0.225$ & $+0.235$ & $+0.204$ & $0.014$ \\
$u_{\text{internal}}$ & $+0.334$ & $+0.257$ & $+0.297$ & $+0.257$ & $+0.321$ & $0.013$ \\
$s_{\text{avg}}$ & $+0.175$ & $+0.220$ & $+0.177$ & $+0.264$ & $+0.242$ & $0.067$ \\
$p_{\text{entail}}$ & $+0.067$ & $+0.368$ & $+0.324$ & $+0.359$ & $+0.285$ & $\mathbf{0.218}$ \\
$p_{\text{ground}}^{\max}$ & $+0.379$ & $+0.296$ & $+0.333$ & $+0.288$ & $+0.356$ & $0.023$ \\
$p_{\text{ground}}^{\mathrm{mean}}$ & $+0.184$ & $+0.262$ & $+0.238$ & $+0.223$ & $+0.230$ & $0.046$ \\
$p_{\text{ground}}^{\text{atomic}}$ & $+0.184$ & $+0.262$ & $+0.238$ & $+0.223$ & $+0.230$ & $0.046$ \\
$q_{a,\text{rel}}$ & $+0.491$ & $+0.360$ & $+0.259$ & $+0.264$ & $+0.265$ & $\mathbf{0.226}$ \\
\bottomrule
\end{tabular}
\captionsetup{singlelinecheck=false}
\caption[Composite refit by cycle.]{What the refit changes at each kept cycle: (a) pooled temperature, calibration error and boost intercept; (b) per-benchmark temperature; (c) per-benchmark safety floor; (d) boost weights.}
\label{tab:calibration-trajectory}
\end{table}


\FloatBarrier
  % Per-Cycle Artifact Dictionary
% ============================================================================
% Appendix D — Significance test details (rewritten 2026-10-19 on the final
% paired tests: outputs/analysis/paired_baseline_tests_final.json, written by
% scripts/paired_baseline_tests.py; label em_final for every system).
% ============================================================================

\chapter{Significance test details}
\label{app:significance}

This appendix gives the procedure behind the paired comparisons of \Cref{sec:significance}: how CAEM's and each baseline's answers are paired, the exact-match and hallucination tests, the bootstrap intervals, the Holm correction and the rule that turns a test into a win or a defeat. Every number below is read from the final test record, which is rebuilt by the results command of \Cref{app:reproducibility}.

\section{Pairing rule and correctness label}
\label{app:sig-pairing}

\textbf{Pairing:} each contrast compares CAEM at one cycle (C0, C2 or C5) with one baseline on one benchmark. The answers of the two systems are matched on the question identifier, restricted to the evaluation-fold identifiers of the dataset-split record, so a contrast uses only questions both systems answered. Under the comparison protocol of \Cref{sec:setup-protocol} every system answered the same $300$ questions per benchmark, and every contrast in this appendix has $n = 300$ pairs.

\textbf{Correctness label:} both sides are scored with the single rule of \Cref{sec:setup-metrics}: the uniform scorer on FEVER, TriviaQA, CommonsenseQA and StrategyQA, and on TruthfulQA the language-model judge applied to the final answer line, the same judge for every system. The earlier ROUGE-based TruthfulQA field is not read.

\textbf{Hallucination value:} the per-question CHM is the share of the seven active failure flags a question raises (\Cref{sec:setup-metrics}). Each baseline answer carries the verifier signals of the comparison protocol, so B8 is tested on this axis as well, and no baseline is charged with a decision it did not take.

\section{Contingency tables}
\label{app:sig-tables}

For exact match, each contrast forms the paired two-by-two table of correct and wrong answers. The two off-diagonal cells are the discordant pairs: $b$ counts questions CAEM answers correctly and the baseline does not, and $c$ the reverse. The McNemar test reads only these two cells, through the continuity-corrected statistic of Edwards~\cite{edwards1948note} with one degree of freedom,
\[
\chi^{2} \;=\; \frac{\bigl(\lvert b - c \rvert - 1\bigr)^{2}}{b + c},
\]
with $p = 1$ when $b + c = 0$. \Cref{tab:app-sig-tables} shows four contrasts at C5: a tie with the zero-shot baseline, CAEM's largest win and its largest defeat, both against vanilla fine-tuning, and a defeat by five-shot chain of thought.

\begin{table}[H]
\centering\footnotesize
\setlength{\tabcolsep}{3pt}
\begin{tabular}{@{}lrrrrrrr@{}}
\toprule
\textbf{Contrast at C5} & \textbf{EM CAEM} & \textbf{EM base} & $\boldsymbol{\Delta_{\text{EM}}}$ & $\boldsymbol{b}$ & $\boldsymbol{c}$ & $\boldsymbol{\chi^{2}}$ & $\boldsymbol{p_{\text{raw}}}$ \\
\midrule
B1 zero-shot $\cdot$ FEVER & $0.477$ & $0.507$ & $-0.030$ & $71$ & $80$ & $0.42$ & $0.515$ \\
B6 vanilla fine-tune $\cdot$ TriviaQA & $0.453$ & $0.290$ & $+0.163$ & $72$ & $23$ & $24.25$ & $8.5 \times 10^{-7}$ \\
B6 vanilla fine-tune $\cdot$ FEVER & $0.477$ & $0.640$ & $-0.163$ & $30$ & $79$ & $21.14$ & $4.3 \times 10^{-6}$ \\
B5 five-shot CoT $\cdot$ CommonsenseQA & $0.717$ & $0.793$ & $-0.077$ & $24$ & $47$ & $6.82$ & $0.009$ \\
\bottomrule
\end{tabular}
\caption[Four McNemar contrasts at C5.]{Four exact-match contrasts at C5 ($n = 300$ each): $b$ counts questions only CAEM answers correctly, $c$ questions only the baseline answers correctly. The two systems disagree on $151$ FEVER questions against B1 but split them almost evenly, so the difference is not significant.}
\label{tab:app-sig-tables}
\end{table}

For the composite hallucination metric, the per-question differences (CAEM minus baseline) are tested with the Wilcoxon signed-rank test~\cite{wilcoxon1945individual}, two-sided, with questions of zero difference dropped.

\section{Bootstrap protocol}
\label{app:sig-bootstrap}

The interval on each difference is the percentile bootstrap over paired questions: $2{,}000$ resamples of the $300$ pairs with replacement, from a generator seeded with $42$, with the $95$ percent interval read from the $2.5$ and $97.5$ percentiles of the resampled means. The same resampling gives the interval on the CHM difference. The percentile method is used rather than the bias-corrected and accelerated method because the per-question exact-match difference takes only the values $-1$, $0$ and $+1$, on which the acceleration estimate is unstable at this sample size.

The interval describes the size of a difference; whether a contrast counts as significant is decided by the Holm-corrected test of the next section. \Cref{tab:app-sig-boot-intervals} shows why the two are reported separately: the interval of the retrieval baseline on TriviaQA excludes zero, but the contrast does not survive the correction for eight comparisons.

\begin{table}[H]
\centering\small
\setlength{\tabcolsep}{6pt}
\begin{tabular}{@{}lrrl@{}}
\toprule
\textbf{Contrast at C5} & $\boldsymbol{\Delta_{\text{EM}}}$ & \textbf{95\% interval} & \textbf{Holm $p$, verdict} \\
\midrule
B1 zero-shot $\cdot$ FEVER & $-0.030$ & $[-0.117,\, +0.050]$ & $1.000$, not significant \\
B3 RAG $\cdot$ TriviaQA & $+0.053$ & $[+0.003,\, +0.107]$ & $0.367$, not significant \\
B6 vanilla fine-tune $\cdot$ TriviaQA & $+0.163$ & $[+0.100,\, +0.220]$ & $<0.001$, CAEM win \\
\bottomrule
\end{tabular}
\caption[Three bootstrap intervals at C5.]{Three bootstrap intervals on the exact-match difference at C5, with the Holm-corrected $p$-value and the verdict.}
\label{tab:app-sig-boot-intervals}
\end{table}

\section{Holm sequence}
\label{app:sig-holm}

The Holm--Bonferroni correction~\cite{holm1979simple} is applied within each benchmark, separately on each axis, over the $m = 8$ baselines. The raw $p$-values are sorted in ascending order, the $k$-th smallest is multiplied by $m - k + 1$, a running maximum keeps the adjusted values in order, and values above one are set to one. The family-wise error rate is held at $0.05$ within each benchmark. \Cref{tab:app-sig-holm-tqa,tab:app-sig-holm-fever} give the exact-match sequences at C5 for TriviaQA, where CAEM's two wins lie, and for CommonsenseQA, the one benchmark with both wins and defeats.

\begin{table}[H]
\centering\footnotesize
\setlength{\tabcolsep}{4pt}
\begin{tabular}{@{}llrrrrl@{}}
\toprule
\textbf{Rank} & \textbf{Baseline} & $\boldsymbol{\Delta_{\text{EM}}}$ & \textbf{Raw $p$} & $\boldsymbol{m-k+1}$ & \textbf{Holm $p$} & \textbf{Verdict} \\
\midrule
1 & B6 vanilla fine-tune & $+0.163$ & $8.5 \times 10^{-7}$ & $8$ & $6.8 \times 10^{-6}$ & CAEM win \\
2 & B7 FLARE & $+0.083$ & $0.0039$ & $7$ & $0.027$ & CAEM win \\
3 & B4 CoT $+$ RAG & $+0.057$ & $0.061$ & $6$ & $0.367$ & -- \\
4 & B3 RAG & $+0.053$ & $0.065$ & $5$ & $0.367$ & -- \\
5 & B2 chain of thought & $+0.033$ & $0.308$ & $4$ & $1.000$ & -- \\
6 & B1 zero-shot & $+0.023$ & $0.494$ & $3$ & $1.000$ & -- \\
7 & B8 semantic entropy & $+0.020$ & $0.556$ & $2$ & $1.000$ & -- \\
8 & B5 five-shot CoT & $+0.013$ & $0.737$ & $1$ & $1.000$ & -- \\
\bottomrule
\end{tabular}
\caption[Holm sequence on TriviaQA at C5.]{Holm sequence for exact match on TriviaQA at C5. CAEM is ahead of every baseline on this benchmark, but only the two largest differences survive the correction. Rank four takes the adjusted value of rank three through the running maximum.}
\label{tab:app-sig-holm-tqa}
\end{table}

\begin{table}[H]
\centering\footnotesize
\setlength{\tabcolsep}{4pt}
\begin{tabular}{@{}llrrrrl@{}}
\toprule
\textbf{Rank} & \textbf{Baseline} & $\boldsymbol{\Delta_{\text{EM}}}$ & \textbf{Raw $p$} & $\boldsymbol{m-k+1}$ & \textbf{Holm $p$} & \textbf{Verdict} \\
\midrule
1 & B3 RAG & $+0.117$ & $8.6 \times 10^{-5}$ & $8$ & $6.9 \times 10^{-4}$ & CAEM win \\
2 & B6 vanilla fine-tune & $-0.090$ & $0.0034$ & $7$ & $0.024$ & baseline win \\
3 & B4 CoT $+$ RAG & $+0.087$ & $0.0037$ & $6$ & $0.024$ & CAEM win \\
4 & B5 five-shot CoT & $-0.077$ & $0.0090$ & $5$ & $0.045$ & baseline win \\
5 & B8 semantic entropy & $-0.050$ & $0.068$ & $4$ & $0.273$ & -- \\
6 & B2 chain of thought & $-0.043$ & $0.137$ & $3$ & $0.410$ & -- \\
7 & B1 zero-shot & $-0.010$ & $0.807$ & $2$ & $1.000$ & -- \\
8 & B7 FLARE & $+0.007$ & $0.903$ & $1$ & $1.000$ & -- \\
\bottomrule
\end{tabular}
\caption[Holm sequence on CommonsenseQA at C5.]{Holm sequence for exact match on CommonsenseQA at C5. CAEM beats the two retrieval baselines and loses to vanilla fine-tuning and five-shot chain of thought; rank three takes the adjusted value of rank two through the running maximum.}
\label{tab:app-sig-holm-fever}
\end{table}

\Cref{tab:sig-full-c5,tab:sig-full-c2} list every contrast at C5 and C2 with its interval; $\ast$ marks a CAEM win and (B) a baseline win under the rule of the next section.

% Auto-generated by scripts/paired_baseline_tests.py (final drafting pass 2026-10).
\begin{table}[!htbp]
\centering\small
\setlength{\tabcolsep}{4pt}
\begin{tabular}{@{}llrcrc@{}}
\toprule
\textbf{Benchmark} & \textbf{Baseline} & $\Delta$\textbf{EM} [95\% CI] & & $\Delta$\textbf{CHM} [95\% CI] & \\
\midrule
FEVER & B1 & $-0.030$ [$-0.12$, $+0.05$] & - & $+0.005$ [$-0.02$, $+0.03$] & - \\
 & B2 & $-0.037$ [$-0.12$, $+0.05$] & - & $+0.006$ [$-0.02$, $+0.03$] & - \\
 & B3 & $+0.040$ [$-0.05$, $+0.13$] & - & $-0.029$ [$-0.06$, $+0.00$] & - \\
 & B4 & $+0.040$ [$-0.05$, $+0.13$] & - & $-0.017$ [$-0.05$, $+0.01$] & - \\
 & B5 & $+0.010$ [$-0.07$, $+0.09$] & - & $+0.012$ [$-0.01$, $+0.03$] & - \\
 & B6 & $-0.163$ [$-0.23$, $-0.10$] & (B) & $+0.068$ [$+0.05$, $+0.09$] & (B) \\
 & B7 & $+0.060$ [$-0.03$, $+0.16$] & - & $-0.065$ [$-0.10$, $-0.03$] & * \\
 & B8 & $-0.037$ [$-0.11$, $+0.04$] & - & $+0.010$ [$-0.01$, $+0.03$] & - \\
\midrule
TriviaQA & B1 & $+0.023$ [$-0.03$, $+0.08$] & - & $+0.003$ [$-0.01$, $+0.02$] & - \\
 & B2 & $+0.033$ [$-0.02$, $+0.09$] & - & $-0.000$ [$-0.02$, $+0.02$] & - \\
 & B3 & $+0.053$ [$+0.00$, $+0.11$] & - & $-0.008$ [$-0.02$, $+0.01$] & - \\
 & B4 & $+0.057$ [$+0.00$, $+0.11$] & - & $-0.002$ [$-0.02$, $+0.01$] & - \\
 & B5 & $+0.013$ [$-0.04$, $+0.07$] & - & $+0.000$ [$-0.02$, $+0.02$] & - \\
 & B6 & $+0.163$ [$+0.10$, $+0.22$] & * & $-0.024$ [$-0.04$, $-0.01$] & - \\
 & B7 & $+0.083$ [$+0.03$, $+0.14$] & * & $-0.075$ [$-0.10$, $-0.06$] & * \\
 & B8 & $+0.020$ [$-0.03$, $+0.08$] & - & $-0.000$ [$-0.02$, $+0.02$] & - \\
\midrule
CommonsenseQA & B1 & $-0.010$ [$-0.07$, $+0.04$] & - & $+0.003$ [$-0.02$, $+0.02$] & - \\
 & B2 & $-0.043$ [$-0.10$, $+0.01$] & - & $+0.012$ [$-0.01$, $+0.03$] & - \\
 & B3 & $+0.117$ [$+0.07$, $+0.17$] & * & $-0.048$ [$-0.07$, $-0.03$] & * \\
 & B4 & $+0.087$ [$+0.03$, $+0.14$] & * & $-0.037$ [$-0.06$, $-0.01$] & * \\
 & B5 & $-0.077$ [$-0.13$, $-0.03$] & (B) & $+0.028$ [$+0.01$, $+0.05$] & - \\
 & B6 & $-0.090$ [$-0.15$, $-0.04$] & (B) & $+0.056$ [$+0.03$, $+0.08$] & (B) \\
 & B7 & $+0.007$ [$-0.05$, $+0.06$] & - & $-0.020$ [$-0.04$, $+0.00$] & - \\
 & B8 & $-0.050$ [$-0.10$, $-0.00$] & - & $+0.017$ [$-0.00$, $+0.04$] & - \\
\midrule
TruthfulQA & B1 & $+0.003$ [$-0.06$, $+0.07$] & - & $+0.021$ [$+0.00$, $+0.04$] & - \\
 & B2 & $-0.027$ [$-0.09$, $+0.03$] & - & $+0.020$ [$+0.00$, $+0.04$] & - \\
 & B3 & $+0.073$ [$+0.01$, $+0.13$] & - & $-0.015$ [$-0.03$, $+0.00$] & - \\
 & B4 & $+0.060$ [$+0.00$, $+0.12$] & - & $-0.010$ [$-0.03$, $+0.01$] & - \\
 & B5 & $-0.093$ [$-0.16$, $-0.03$] & - & $+0.033$ [$+0.02$, $+0.05$] & (B) \\
 & B6 & $+0.090$ [$+0.03$, $+0.15$] & * & $+0.022$ [$+0.01$, $+0.04$] & - \\
 & B7 & $+0.020$ [$-0.05$, $+0.08$] & - & $-0.087$ [$-0.11$, $-0.06$] & * \\
 & B8 & $+0.010$ [$-0.05$, $+0.08$] & - & $-0.005$ [$-0.02$, $+0.01$] & - \\
\midrule
StrategyQA & B1 & $-0.003$ [$-0.08$, $+0.07$] & - & $-0.023$ [$-0.05$, $+0.00$] & - \\
 & B2 & $+0.060$ [$-0.02$, $+0.14$] & - & $-0.051$ [$-0.08$, $-0.02$] & * \\
 & B3 & $+0.030$ [$-0.05$, $+0.11$] & - & $-0.056$ [$-0.08$, $-0.03$] & * \\
 & B4 & $+0.037$ [$-0.04$, $+0.12$] & - & $-0.053$ [$-0.08$, $-0.03$] & * \\
 & B5 & $+0.080$ [$-0.01$, $+0.16$] & - & $-0.048$ [$-0.07$, $-0.02$] & * \\
 & B6 & $-0.020$ [$-0.09$, $+0.05$] & - & $+0.035$ [$+0.02$, $+0.05$] & (B) \\
 & B7 & $+0.040$ [$-0.05$, $+0.12$] & - & $-0.060$ [$-0.09$, $-0.03$] & * \\
 & B8 & $+0.000$ [$-0.07$, $+0.07$] & - & $-0.025$ [$-0.05$, $+0.00$] & - \\
\bottomrule
\end{tabular}
\caption[Per-benchmark paired tests at C5.]{Per-benchmark paired tests of CAEM at C5 against each baseline ($300$ questions per benchmark). Baselines are numbered as in \Cref{sec:comp-baselines}. A positive $\Delta$EM and a negative $\Delta$CHM favour CAEM; $\ast$ marks a significant CAEM win and (B) a significant baseline win.}
\label{tab:sig-full-c5}
\end{table}


% Auto-generated by scripts/paired_baseline_tests.py (final drafting pass 2026-10).
\begin{table}[!htbp]
\centering\small
\setlength{\tabcolsep}{4pt}
\begin{tabular}{@{}llrcrc@{}}
\toprule
\textbf{Benchmark} & \textbf{Baseline} & $\Delta$\textbf{EM} [95\% CI] & & $\Delta$\textbf{CHM} [95\% CI] & \\
\midrule
FEVER & B1 & $+0.027$ [$-0.05$, $+0.10$] & - & $-0.022$ [$-0.05$, $+0.00$] & - \\
 & B2 & $+0.020$ [$-0.06$, $+0.10$] & - & $-0.022$ [$-0.05$, $+0.00$] & - \\
 & B3 & $+0.097$ [$+0.01$, $+0.18$] & - & $-0.056$ [$-0.08$, $-0.03$] & * \\
 & B4 & $+0.097$ [$+0.01$, $+0.18$] & - & $-0.045$ [$-0.07$, $-0.02$] & * \\
 & B5 & $+0.067$ [$-0.01$, $+0.15$] & - & $-0.016$ [$-0.04$, $+0.01$] & - \\
 & B6 & $-0.107$ [$-0.17$, $-0.04$] & (B) & $+0.040$ [$+0.02$, $+0.06$] & (B) \\
 & B7 & $+0.117$ [$+0.03$, $+0.20$] & - & $-0.092$ [$-0.12$, $-0.06$] & * \\
 & B8 & $+0.020$ [$-0.06$, $+0.10$] & - & $-0.018$ [$-0.04$, $+0.01$] & - \\
\midrule
TriviaQA & B1 & $+0.063$ [$+0.01$, $+0.12$] & - & $-0.014$ [$-0.03$, $+0.00$] & - \\
 & B2 & $+0.073$ [$+0.01$, $+0.13$] & - & $-0.018$ [$-0.03$, $-0.00$] & - \\
 & B3 & $+0.093$ [$+0.05$, $+0.14$] & * & $-0.026$ [$-0.04$, $-0.01$] & * \\
 & B4 & $+0.097$ [$+0.05$, $+0.14$] & * & $-0.020$ [$-0.04$, $-0.00$] & * \\
 & B5 & $+0.053$ [$-0.00$, $+0.11$] & - & $-0.017$ [$-0.03$, $-0.00$] & - \\
 & B6 & $+0.203$ [$+0.14$, $+0.26$] & * & $-0.041$ [$-0.06$, $-0.03$] & * \\
 & B7 & $+0.123$ [$+0.07$, $+0.18$] & * & $-0.093$ [$-0.11$, $-0.07$] & * \\
 & B8 & $+0.060$ [$+0.01$, $+0.11$] & - & $-0.018$ [$-0.03$, $-0.00$] & - \\
\midrule
CommonsenseQA & B1 & $-0.010$ [$-0.05$, $+0.03$] & - & $+0.005$ [$-0.01$, $+0.02$] & - \\
 & B2 & $-0.043$ [$-0.09$, $+0.01$] & - & $+0.015$ [$-0.01$, $+0.03$] & - \\
 & B3 & $+0.117$ [$+0.06$, $+0.18$] & * & $-0.046$ [$-0.07$, $-0.02$] & * \\
 & B4 & $+0.087$ [$+0.03$, $+0.14$] & * & $-0.034$ [$-0.06$, $-0.01$] & * \\
 & B5 & $-0.077$ [$-0.12$, $-0.03$] & (B) & $+0.030$ [$+0.01$, $+0.05$] & (B) \\
 & B6 & $-0.090$ [$-0.14$, $-0.04$] & (B) & $+0.059$ [$+0.04$, $+0.08$] & (B) \\
 & B7 & $+0.007$ [$-0.04$, $+0.06$] & - & $-0.018$ [$-0.04$, $+0.01$] & - \\
 & B8 & $-0.050$ [$-0.09$, $-0.01$] & - & $+0.020$ [$+0.00$, $+0.04$] & - \\
\midrule
TruthfulQA & B1 & $+0.027$ [$-0.03$, $+0.09$] & - & $+0.011$ [$-0.00$, $+0.03$] & - \\
 & B2 & $-0.003$ [$-0.06$, $+0.06$] & - & $+0.010$ [$-0.01$, $+0.02$] & - \\
 & B3 & $+0.097$ [$+0.04$, $+0.15$] & * & $-0.025$ [$-0.04$, $-0.01$] & * \\
 & B4 & $+0.083$ [$+0.03$, $+0.14$] & * & $-0.020$ [$-0.04$, $-0.00$] & - \\
 & B5 & $-0.070$ [$-0.13$, $-0.00$] & - & $+0.023$ [$+0.01$, $+0.04$] & (B) \\
 & B6 & $+0.113$ [$+0.06$, $+0.17$] & * & $+0.012$ [$-0.00$, $+0.03$] & - \\
 & B7 & $+0.043$ [$-0.02$, $+0.10$] & - & $-0.097$ [$-0.12$, $-0.07$] & * \\
 & B8 & $+0.033$ [$-0.03$, $+0.10$] & - & $-0.015$ [$-0.03$, $+0.00$] & - \\
\midrule
StrategyQA & B1 & $+0.010$ [$-0.06$, $+0.07$] & - & $-0.006$ [$-0.03$, $+0.02$] & - \\
 & B2 & $+0.073$ [$+0.00$, $+0.15$] & - & $-0.033$ [$-0.06$, $-0.01$] & - \\
 & B3 & $+0.043$ [$-0.02$, $+0.11$] & - & $-0.039$ [$-0.07$, $-0.01$] & * \\
 & B4 & $+0.050$ [$-0.01$, $+0.11$] & - & $-0.035$ [$-0.06$, $-0.01$] & - \\
 & B5 & $+0.093$ [$+0.02$, $+0.16$] & - & $-0.030$ [$-0.06$, $-0.00$] & - \\
 & B6 & $-0.007$ [$-0.07$, $+0.06$] & - & $+0.052$ [$+0.03$, $+0.07$] & (B) \\
 & B7 & $+0.053$ [$-0.02$, $+0.12$] & - & $-0.042$ [$-0.07$, $-0.01$] & * \\
 & B8 & $+0.013$ [$-0.05$, $+0.08$] & - & $-0.007$ [$-0.03$, $+0.02$] & - \\
\bottomrule
\end{tabular}
\caption[Per-benchmark paired tests at C2.]{Per-benchmark paired tests of CAEM at C2 against each baseline ($300$ questions per benchmark). Baselines are numbered as in \Cref{sec:comp-baselines}. A positive $\Delta$EM and a negative $\Delta$CHM favour CAEM; $\ast$ marks a significant CAEM win and (B) a significant baseline win.}
\label{tab:sig-full-c2}
\end{table}


\section{Win and defeat rule}
\label{app:sig-licensing}

A contrast counts as a win or a defeat only when two conditions hold together: the Holm-corrected $p$-value is below $0.05$, and the absolute difference is at least $0.02$ in exact match or $0.01$ in CHM. The direction of the difference decides whether it is a CAEM win or a baseline win (for CHM, lower is better). The second condition matters because with $300$ paired questions a test can detect differences too small to matter; the floors keep such contrasts out of the counts. The floors are the same for every contrast, cycle and baseline.

\textbf{Counts:} over the $40$ contrasts per cycle ($8$ baselines on $5$ benchmarks), exact match gives $9$ CAEM wins and $2$ defeats at C0, $9$ and $3$ at C2, and $5$ and $3$ at C5; CHM gives $6$ wins and $17$ defeats at C0, $13$ and $5$ at C2, and $10$ and $4$ at C5 (\Cref{tab:sig-summary}). Against the prompting baselines and semantic-entropy abstention (B1, B2, B5, B8), no exact-match contrast is a CAEM win at any cycle. The exact-match defeats at C2 and C5 are against vanilla fine-tuning on FEVER and CommonsenseQA and five-shot chain of thought on CommonsenseQA.

\textbf{Verifier-free check:} the same tests on the CHM restricted to the five subtypes that use neither CAEM's composite score nor its decisions give $15$ wins and $4$ defeats at C0, $28$ and $2$ at C2, and $31$ and $3$ at C5. CAEM fares better on this measure than on the full CHM, partly by construction, as discussed in \Cref{sec:disc-threat-verifier}.
  % Significance Test Details
% ============================================================================
% Appendix E — Reproducibility recipe
% The data preparation pipeline, snapshot recovery sequence, step-by-step
% run order with gate verdicts, framework version pins, hardware envelope,
% and a verification checklist a reproducer uses to confirm the headline
% numbers match within tolerance.
% ============================================================================

\chapter{Reproducibility recipe}
\label{app:reproducibility}

\section{Data preparation}
\label{app:repro-data}

The five benchmark loaders that feed the realised evaluation panel are registered in \verb|eval/benchmarks.py| against the Hugging Face dataset hub identifiers. FEVER reads through \verb|lucadiliello/fever| under the development and training splits; TriviaQA reads through \verb|trivia_qa| under the reading-comprehension-no-context configuration; CommonsenseQA reads through \verb|commonsense_qa|; TruthfulQA reads through \verb|truthful_qa| under the generation configuration; StrategyQA reads through \verb|ChilleD/StrategyQA| with documented fallbacks to \verb|wics/strategy-qa| and the Eladsegal train-JSON mirror. Earlier loader registrations for Natural Questions and HotpotQA are preserved in the module for reproducibility of the cycle-zero diagnostics but do not feed the realised five-benchmark panel.

The six-way disjoint partition is constructed once at pool-builder time by \codepath{caem/benchmark\_splits.py::build\_all\_benchmark\_pools}, which is invoked by the experiment runner and writes the per-benchmark identifier manifest to \verb|outputs/full_run/dataset_splits.json|. Each benchmark contributes a $1{,}000$-sample seed pool for cold-start memory population, a $500$-sample calibration fold, a $500$-sample purity-validation slice, ten per-cycle stream-chunk slices (FEVER and TriviaQA at $2{,}000$ samples per cycle, CommonsenseQA at $700$ samples per cycle to fit the smaller available train pool of about $10{,}000$ samples after deduplication), a $300$-sample evaluation fold, and a test fold of up to $500$ samples (clamped to the available remainder where the source split admits fewer than the registered budget). The shuffle is seeded at forty-two on the training side and forty-three on the evaluation side, so the partitions are stable across re-invocations on the same source dataset.

Within-benchmark deduplication runs through the \verb|_dedupe_by_content_id| helper, which keys on the first $500$ characters of the question text under a sixteen-character SHA-256 hash. Cross-benchmark disjointness is verified by the \codepath{assert\_cross\_benchmark\_disjoint} assertion. Both checks run at pool-builder time; the runner re-stamps the same checks at every successful cycle boundary through \codepath{\_dedup\_inplace} and \codepath{assert\_disjoint\_calibration}, so any drift in the partition would be caught at the cycle close rather than at the metric-computation step.

The realised evaluation fold reads $300$ samples per benchmark across the five-benchmark panel, for a total of $1{,}500$ evaluation questions per cycle. The StrategyQA test fold has $387$ samples rather than the registered $500$ because the source split (the ChilleD StrategyQA test split) has only $687$ labelled samples, of which $300$ form the evaluation fold.

\section{Snapshot recovery}
\label{app:repro-snapshot}

The pre-main snapshot is the state of the run just before the main cycle loop. It is held in the private Hugging Face dataset repository \verb|aksaN000/caem-passage-index-21m| (access on request), which also holds the $64$-gigabyte Wikipedia passage index at its root, with a copy in the Drive folder's \texttt{9\_archive/pre\_main\_snapshot/}. Each bundle sits under \verb|pre_main_snapshot/<git-commit>/|: the anchor before the main run is \verb|c5b8066|, and the anchor after the retrieval-utility classifier was added is \verb|700b9e3|. A manifest in each bundle records the upload time, the source machine, the commit and a recovery hint.

The three-command recovery sequence runs from a fresh clone with the Hugging Face command-line interface authenticated to the private repository:
\begin{verbatim}
# 1. Fetch the index and the pre-main snapshot bundles
hf download aksaN000/caem-passage-index-21m \
    --include 'passages.*' 'pre_main_snapshot/*'

# 2. Stage the chosen anchor into outputs/
cp -r pre_main_snapshot/c5b8066/* outputs/

# 3. Pin the code to the matching commit
git checkout c5b8066    # or 700b9e3 for the RUC-era resume point
\end{verbatim}
After the third command the working tree holds the cold-start memory store, the cycle-zero memory and deferred-buffer snapshots, the cycle-zero calibration fold (per-sample nine-signal vectors and gold labels), the cycle-zero composite calibration JSON, the calibrated configuration JSON, the cycle-zero validation-gate verdict, the cycle-zero evaluation fold under the pristine generator, the post-calibration rescored evaluation fold, the dataset-splits manifest, and the run logs. The four artefacts that must be regenerated rather than pulled from the snapshot are the Qwen-2.5-3B-Instruct base weights (recovered from the Hugging Face model hub at deployment time), the per-cycle LoRA adapters at cycles one through five, the per-cycle self-improvement training artefacts, and the post-cycle-zero composite refits.

\section{Step-by-step run sequence}
\label{app:repro-steps}

The commands below reproduce the run from a fresh clone. What each stage does, and the gate it must pass, is described once in the stage graph of \Cref{app:runbook}; the run itself was driven by the top-level script that is now archived (\verb|archive/scripts_superseded/old_runners/run_phase1a.sh|), whose stages call these commands.

\paragraph{Pre-run setup:}
\begin{verbatim}
python scripts/prompt_smoke_test.py \
    --benchmarks fever triviaqa commonsense_qa
python scripts/build_alias_dict_from_benchmarks.py
python -m scripts.seed_cold_start
python -m scripts.run_experiment --output_dir outputs/cycle_0 \
    --num_cycles 0 --n_questions 500 --n_eval_questions 500 \
    --benchmarks fever triviaqa commonsense_qa truthfulqa strategyqa \
    --passage_index data/passage_index \
    --cold_start_memory outputs/cold_start_memory/memory_store
python scripts/fit_composite_calibration.py --shrinkage_alpha 0.6 \
    --cherian_boost --boost_C 0.01
python scripts/validate_composite_weights.py
\end{verbatim}
The retention-probe baseline and the snapshot upload are the stages \verb|step_19_2_slice| and \verb|step_gdrive_upload_pre_main| of the archived driver.

\paragraph{Main cycle loop:} A single invocation drives all six per-cycle sub-steps and resumes from the highest-cycle artefact set on disk:
\begin{verbatim}
python -m scripts.run_experiment \
    --output_dir outputs/full_run \
    --num_cycles 10 \
    --n_questions 5000 \
    --n_eval_questions 300 \
    --benchmarks fever triviaqa commonsense_qa truthfulqa strategyqa \
    --passage_index data/passage_index \
    --cold_start_memory outputs/cold_start_memory/memory_store \
    --verifier_backend minicheck \
    --resume_from_cycle <N>
\end{verbatim}
The six steps of each cycle and the rollback contract are described in \Cref{app:runbook}; the time per cycle is given in \Cref{app:repro-hardware}.

\paragraph{Post-run analysis:} The baselines and every reported number are produced by two commands.
\begin{description}[style=nextline,leftmargin=0pt]
\item[Baseline panel under the comparison protocol:] The baseline re-run of \Cref{sec:setup-protocol} on one rented $32$-gigabyte graphics processor: retrieval of the evaluation passages by exact search, B1, B2, B5 and B8 on all five benchmarks, B3 and B4 on TruthfulQA, FLARE, the evaluation of the vanilla fine-tune, and the verifier pass over every baseline answer with the cycle-three adapter and composite. Each stage can be run on its own.
\begin{verbatim}
bash scripts/rental/final_version.sh all
\end{verbatim}
\item[Final numbers, tables and figures:] On CPU in about two minutes: writes the final correctness label into every answer file, recomputes CHM in the model and delivered views and the status-tag accuracy, runs the paired tests, and writes every chapter-five table and figure and the result tables of the appendices. The option \texttt{--judge} first re-judges every TruthfulQA answer with the language-model judge (requires \texttt{ANTHROPIC\_API\_KEY}); the option \texttt{--check} fails if any committed output changes.
\begin{verbatim}
bash scripts/make_final_results.sh --check
\end{verbatim}
\end{description}

\section{Version pinning}
\label{app:repro-versions}

The environment is pinned in \verb|requirements.txt|, with PyTorch installed first from the CUDA wheel index; the exact versions of the baseline re-run environment are frozen in \verb|requirements-lock-rental-2026-10-18.txt|. The main run used the versions below.

\begin{table}[!ht]
\centering\footnotesize
\setlength{\tabcolsep}{5pt}
\renewcommand{\arraystretch}{1.1}
\rowcolors{2}{gray!8}{white}
\begin{tabular}{@{}p{5.0cm}p{3.0cm}p{5.5cm}@{}}
\toprule
\textbf{Package} & \textbf{Version} & \textbf{Role} \\
\midrule
Python                & $3.12$                   & interpreter \\
CUDA toolkit          & $13.0$ (driver $13.1$)   & RTX 5090 Blackwell compute capability \texttt{sm\_120} \\
\midrule
torch                 & $2.10.0\text{+cu130}$    & tensor backend, bf16 forward path, MC-dropout sampler \\
transformers          & $4.57.6$                 & Qwen-2.5-3B-Instruct loader, generate, MiniCheck encoder \\
peft                  & $0.19+$                  & LoRA wrapper for the self-improvement fine-tune \\
accelerate            & $0.34+$                  & device placement during the fine-tune step \\
\midrule
datasets              & $3.0+$                   & benchmark loaders and retention-probe shuffles \\
faiss-cpu             & $1.8+$                   & passage index and episodic memory store \\
sentence-transformers & $4.1.0$                  & SBERT memory encoder and cross-encoder reranker \\
\midrule
scipy                 & $1.11+$                  & Wilcoxon test, temperature fit (L-BFGS) \\
scikit-learn          & $1.4+$                   & per-signal isotonic regression and boost-layer logistic \\
\midrule
bitsandbytes          & $0.43+$                  & eight-bit AdamW optimiser path \\
anthropic             & $0.108$ to below $1$     & TruthfulQA judge client \\
numpy, pandas, tqdm   & latest stable            & arithmetic, table aggregation, progress bars \\
\midrule
biber + biblatex      & TeX Live 2024+           & thesis bibliography compile \\
\bottomrule
\end{tabular}
\caption[Framework version pinning]{Framework version pinning. Lower bounds follow the requirements file; the lock file gives the exact versions of the re-run.}
\label{tab:app-repro-versions}
\end{table}

Seeding is partial, and the numbers are reproducible within the tolerances of \Cref{app:repro-verify} rather than bit for bit. The seed $42$ fixes the benchmark partitions (\Cref{app:repro-data}) and the order of each fine-tune; generation is greedy, so answers repeat given the same weights and inputs. The signals computed from sampled text (dropout agreement, chain similarity and chain entailment) are random draws by design, PyTorch's deterministic-algorithm mode is left off because the attention kernels are not bitwise deterministic, and TensorFloat-32 matrix multiplication is left on.

The two required environment variables are \verb|ANTHROPIC_API_KEY| for the TruthfulQA judge and \verb|HF_TOKEN| (equivalently \verb|HUGGINGFACE_HUB_TOKEN|) for the pull from the private pre-main snapshot repository. Optional environment variables \verb|HUGGINGFACE_HUB_CACHE|, \verb|CAEM_DEVICE|, and \verb|CUDA_VISIBLE_DEVICES| override the cache location, the compute device, and the GPU index respectively.

\section{Hardware envelope}
\label{app:repro-hardware}

\paragraph{Production envelope:} the run used a rented RTX 5090 with $32$ gigabytes of graphics memory, about $150$ gigabytes of disk and about $90$ gigabytes of host memory, at about $0.70$ US dollars per hour. Cycles zero to five of the main run took $6.0$, $29.0$, $35.4$, $43.2$, $24.3$ and $54.6$ hours according to the run log, about $193$ hours in all (\Cref{sec:env-experiment}); the growth comes mostly from the retroactive re-verification pass, which grows with the memory (about $38$ minutes for $431$ entries at cycle one, about $16$ hours for $10{,}214$ entries at cycle five). The verifier, which re-scores every answer, takes most of the time; the fine-tune itself is a small share. With the calibration phase, the earlier rejected runs and the analysis, the programme used several times the compute of the main run (\Cref{sec:env-experiment}). The baseline re-run of \Cref{sec:setup-protocol} ran later on a second rented machine of the same class.

\paragraph{Writeup envelope:} the thesis was written on a Windows Subsystem for Linux 2 environment with an RTX 3060 ($12$ gigabytes of graphics memory). The $64$-gigabyte passage index does not fit in its memory; all retrieval-side numbers in the thesis are read from cached evaluation JSON files under \verb|outputs/full_run/| rather than recomputed locally. The local writeup envelope hosts the thesis source, the table and figure compilation, and the signal-level audits over saved verifier traces, but does not host the live retrieval substrate.

\section{Verification of reproduced numbers}
\label{app:repro-verify}

A reproducer who has executed the full trajectory uses the ten-item checklist below to confirm that the headline numbers match the thesis within the stated tolerance. The per-benchmark exact match on $300$ questions carries a stochastic floor of about $\pm 0.5$ percentage points from the non-deterministic attention path; the pooled exact match over the five benchmarks tightens to about $\pm 0.3$ points; the composite hallucination metric and the storage precision tolerate about $\pm 1.0$ point; integer counts (commit pattern, win and defeat tallies) must match exactly when the results command is run on the same answer files. Items one to three, eight and ten are rebuilt from the saved answer files by the results command above, which reproduces them exactly.

\begin{description}[style=nextline,leftmargin=0pt]
\item[Pooled CAEM exact match:]
\begin{itemize}
\item Receipt: the final correctness label of every answer in \codepath{outputs/full\_run/eval/}, pooled over the five benchmarks (\Cref{tab:per-bench-trajectory}).
\item Expected: $0.517$ at C0, $0.535$ at C2, $0.508$ at C5; B1 $0.511$.
\item Tolerance: $\pm 0.3$ percentage points.
\end{itemize}
\item[Pooled composite hallucination metric:]
\begin{itemize}
\item Receipt: \codepath{outputs/analysis/chm\_judged.json}, model view (mean of the seven active failure flags).
\item Expected: $0.155$ at C0, $0.127$ at C2, $0.134$ at C5; B1 $0.132$.
\item Tolerance: $\pm 1.0$ percentage point.
\end{itemize}
\item[Per-benchmark exact match and CHM by cycle:]
\begin{itemize}
\item Receipt: the same record, per benchmark.
\item Expected: each cell of \Cref{tab:per-bench-trajectory}.
\item Tolerance: $\pm 0.5$ percentage points.
\end{itemize}
\item[Calibration-fold pooled storage precision band:]
\begin{itemize}
\item Receipt: \codepath{outputs/full\_run/pi\_store\_trajectory.csv}, column \texttt{cal\_fold\_pi\_store} across the committed-cycle rows.
\item Expected: band $[0.732, 0.746]$.
\item Tolerance: $\pm 1.0$ percentage point on each endpoint.
\end{itemize}
\item[Cycle-zero validation-gate composite Cohen's $d$:]
\begin{itemize}
\item Receipt: \codepath{outputs/analysis/validation\_gate\_locked.json}, field \texttt{d\_pooled\_train} (the locked composite on the validation fold, \Cref{sec:adj-gate}).
\item Expected: $0.717$.
\item Tolerance: $\pm 0.02$.
\end{itemize}
\item[Memory pool size at cycle five close:]
\begin{itemize}
\item Receipt: \codepath{outputs/full\_run/memory\_store\_cycle\_5.meta}, entry count.
\item Expected: $12\,175$ entries.
\item Tolerance: $\pm 5$ for content-hash collision drift.
\end{itemize}
\item[Realised commit pattern and abort cycle:]
\begin{itemize}
\item Receipt: \codepath{outputs/full\_run/stochastic\_equilibrium\_trajectory.csv}, columns \texttt{verdict} and \texttt{committed}.
\item Expected: commit pattern $(1, 1, 1, 0, 1)$ across C1 to C5, with cycle four aborted by the TriviaQA-test retention probe.
\item Tolerance: exact match required.
\end{itemize}
\item[Retrieval-utility classifier ablation:]
\begin{itemize}
\item Receipt: the classifier-off answers at cycle three in \codepath{outputs/phase1d\_eval/} against the main run at cycle three (\Cref{tab:ruc-ablation}).
\item Expected: pooled exact match $0.467$ off and $0.516$ on, a lift of $4.9$ percentage points.
\item Tolerance: $\pm 0.5$ percentage points.
\end{itemize}
\item[Status-tag accuracy:]
\begin{itemize}
\item Receipt: \codepath{outputs/analysis/chm\_judged.json}, field \texttt{tags} of each CAEM cycle (\Cref{tab:status-tags}).
\item Expected: verified answers correct $0.746$, $0.655$ and $0.667$ of the time at C0, C2 and C5.
\item Tolerance: $\pm 1.0$ percentage point.
\end{itemize}
\item[Significance win and defeat tallies:]
\begin{itemize}
\item Receipt: \codepath{outputs/analysis/paired\_baseline\_tests\_final.json} (\Cref{tab:sig-summary}).
\item Expected: exact match $9$ wins and $2$ defeats at C0, $9$ and $3$ at C2, $5$ and $3$ at C5; CHM $6$ and $17$ at C0, $13$ and $5$ at C2, $10$ and $4$ at C5, each out of $40$ contrasts.
\item Tolerance: exact match required.
\end{itemize}
\end{description}

A reproducer whose readings clear every item within the stated tolerance has matched the thesis to within the determinism floor of the compute envelope. An item that fails by a larger margin points to a determinism gap the runner did not absorb (the seed-everything helper is the first place to inspect), a version mismatch against \Cref{tab:app-repro-versions}, a hardware change not registered in the runner's overrides, or a genuine correction that the empirical chapter should adopt.
  % Reproducibility Recipe
% ============================================================================
% Appendix F — Glossary
% Discipline: only terms that appear in body prose more than once and
% whose meaning a typical reader might not infer go here. Concept
% words always; implementation names never.
% ============================================================================

\chapter{Glossary}
\label{app:glossary}

The glossary collects every concept word that recurs in body prose and whose meaning a reader who has read \Cref{ch:introduction} alone might not infer. Each entry pairs the term with a one-sentence concept-word definition and a cross-reference to the chapter subsection that introduces the term. Acronyms used in body prose are listed together with their expansions; concept terms are listed in alphabetical order. Implementation identifiers (module paths, class names, script names) do not appear here; they live in \Cref{app:implementation}.

\section{Concept terms}

\begin{description}[style=nextline, leftmargin=0pt, labelindent=0pt, itemsep=6pt, parsep=0pt, font=\bfseries, beginpenalty=10000]

\item[Asymmetric rollback.] The retention guard's response on a cycle whose worst-probe ratio falls below the registered floor: the adapter weights revert to the previous successful cycle's checkpoint, but the memory store does not revert, so the storage admissions accumulated during the failed cycle's stream-chunk pass are preserved. The retroactive re-verification pass and the deferred-buffer reconsider pass are skipped on the aborted cycle because both would be provable no-ops under unchanged weights. See \Cref{sec:retention}.

\item[Boost layer.] The L2-regularised logistic regression that aggregates the nine per-signal calibrated probabilities produced by the per-signal isotonic regression stage into the final calibrated probability composite. Introduced at \Cref{sec:composite}.

\item[Calibration fold.] The five-hundred-sample-per-training-benchmark labelled slice that fits the per-signal isotonic regression, the per-benchmark shrinkage prior, and the boost-layer logistic regression at every successful cycle boundary. Held content-hash disjoint from the seed pool, the purity-validation slice, the stream chunks, the evaluation fold, and the test fold. See \Cref{sec:adj-discipline}.

\item[Calibrated-probability composite.] The two-stage scoring head that maps the nine verifier signals into a single calibrated correctness probability $u_{\text{stored}} \in [0, 1]$ through per-signal isotonic regression followed by the boost layer. The fixed-threshold storage gate reads against this output. See \Cref{sec:composite}.

\item[Combined dispatch score.] The router's weighted blend of the cosine similarity to the nearest stored entry and the stored confidence of that neighbour, against which the memory-direct tier threshold is read. See \Cref{sec:router}.

\item[Committed cycle.] A cycle whose self-improvement fine-tune cleared the retention guard and whose calibration-fold refit was applied at the cycle boundary. The realised trajectory records four committed cycles (C1, C2, C3, C5) against the cycle-four abort. See \Cref{sec:retention}.

\item[Composite evaluation score (CES).] The five-axis geometric-mean roll-up that combines accuracy, episodic contribution, retention, calibration, and verifier reliability into a single per-cycle scalar. The CES peak across the trajectory falls at cycle five. See \Cref{sec:headline-trajectory}.

\item[Composite hallucination metric (CHM).]\nopagebreak[4] The eight-axis equal-weighted mean over the post-generation hallucination subtypes (confident confabulation, false refusal, factual fabrication, logical fabrication, off-topic, template leak, defensive evasion, over-long generation). The architecture-vs-baseline CHM contrast is the headline hallucination-reduction axis. See \Cref{sec:headline-trajectory}.

\item[Cycle.]\nopagebreak[4] A single iteration of the self-improvement loop: a stream-chunk pre-routing pass that accumulates storage admissions, a cycle-boundary recalibration that refits the per-signal isotonic curves and the per-benchmark shrinkage and boost-layer weights, a retroactive re-verification pass that updates the stored confidences against the recalibrated verifier, and a self-improvement fine-tune that distils verified retrieval-augmented reasoning into the generator's parametric capacity. See \Cref{sec:overview-cycle}.

\item[Deferred buffer.] The bounded-capacity reconsider queue that holds outputs whose calibrated probability composite cleared the deferral threshold but fell short of the storage threshold. Each entry carries an age counter, is rescored under the recalibrated verifier at every successful cycle boundary, and is dropped from the buffer when the counter reaches the registered time-to-live of two reconsider invocations without a promotion. See \Cref{sec:deferred-reconsider}.

\item[Empirical-precision anchor.] The realised storage precision on the calibration fold at the locked storage cut-point, audited once per cycle and reported as the headline diagnostic that licences the calibration-consistent precision floor of the pool-purity theorem. See \Cref{sec:adj-cal-eval-gap}.

\item[Equilibrium-saturation gate.]\nopagebreak[4] The two-of-three composite stop trigger that closes the trajectory when at least two of the three signals fire simultaneously: storage saturation under repeated zero-admission cycles, MMLU-floor breach against the cycle-zero pristine baseline, and composite-evaluation-score plateau across the late-cycle subsequence. The realised trajectory closes at cycle five under the storage-saturation and MMLU-floor signals. See \Cref{sec:check-convergence}.

\item[Evaluation fold.] The three-hundred-sample-per-benchmark held-out slice on which every per-cycle headline metric is computed. Held content-hash disjoint from the calibration fold, the seed pool, the purity-validation slice, and the stream chunks. See \Cref{sec:setup}.

\item[Fixed-threshold storage gate.] The architecturally frozen storage admission rule that admits a generated output into the memory pool when its calibrated probability composite clears the registered cut-point $\tau_{\text{store}} = 0.60$. The cut-point does not move across the trajectory; the per-cycle composite refit moves the signal-to-probability mapping the cut-point operates on. See \Cref{sec:decision}.

\item[Four-outcome decision tree.] The per-query post-verification routing rule that maps each generated output to one of the four outcomes \textsc{store}, \textsc{defer}, \textsc{abstain}, or \textsc{discard} on the basis of the calibrated probability composite and the grounding signal. See \Cref{sec:decision}.

\item[Matched-protocol contract.] The pairing rule that every per-benchmark contrast between the architecture and an external baseline is computed sample-by-sample on the identical evaluation-fold sample identifiers under the identical prompt template family and the identical scoring rule. See \Cref{sec:sig-pairing}.

\item[Memory-direct tier (Tier 1).] The serving tier that returns the nearest stored entry's answer directly when the combined dispatch score clears the registered Tier 1 threshold. The cheapest of the three tiers per query because it neither runs the generator nor retrieves passages. See \Cref{sec:router}.

\item[Memory-poisoning rate.] The share of stored entries whose stored answer is wrong against the ground-truth label, audited at every cycle boundary as the complement of the storage precision. The validation-gate rule registers this rate at a fifty-percent ceiling on every training benchmark at cycle zero. See \Cref{sec:adj-gate}.

\item[Per-cycle composite refit.] The cycle-boundary recalibration of the per-signal isotonic curves, the per-benchmark shrinkage prior, the per-benchmark temperature scalar, the per-benchmark safety floor, and the boost-layer weights and intercept. The cut-points themselves are architecturally frozen; the refit moves the signal-to-probability mapping under which the cut-points operate. See \Cref{sec:adj-trajectory}.

\item[Pool purity.] The empirical correctness rate of the stored memory pool against the ground-truth label, reported as $\pi_{\text{store}}$ on the calibration fold (the theorem-anchored reading) and on the evaluation fold (the deployment-realised reading). The calibration-fold band across the committed-cycle subsequence is $[0.732, 0.746]$. See \Cref{sec:check-purity}.

\item[Pre-routing confidence (\boldmath$u_{\text{pre}}$).] A cheap confidence score computed before an answer is generated: it combines the token confidence of a $32$-token greedy decode with a decoder layer-convergence ratio from one forward pass, calibrated by temperature, and drives the safety override. See \Cref{sec:pre-routing}.

\item[Production envelope.] The per-query four-outcome rendering on the deployment surface: every served answer carries the calibrated correctness probability as a continuous percentage alongside the categorical decision tag (store-and-serve, defer, abstain, retrieve-and-serve). See \Cref{sec:func-reqs}.

\item[Purity-validation slice.] The five-hundred-sample-per-training-benchmark slice held disjoint from the calibration fold and the evaluation fold, against which the empirical-precision audit is run once per cycle to license the calibration-consistent precision floor. See \Cref{sec:adj-discipline}.

\item[Retention guard.] The cycle-boundary mechanism that measures the post-fine-tune generator on the three-probe multi-modal panel (MMLU, TriviaQA-test, CommonsenseQA-test) and triggers the asymmetric rollback when the worst-probe ratio against the cycle-zero pristine baseline falls below the registered floor of $0.93$. See \Cref{sec:retention}.

\item[Retrieval-augmented tier (Tier 3).] The serving tier that retrieves the top-$k$ passages from the dense passage index, reranks them through a cross-encoder, and runs a grounded generation against the reranked context. The architectural fallback when neither the memory-direct nor the zero-shot tier earns the right to answer. See \Cref{sec:router}.

\item[Retrieval-utility classifier (RUC).] The binary classifier that intercepts the safety-veto fall-through and the score-formula fall-through of the three-tier router and decides whether to dispatch the query to the retrieval-augmented tier or to the zero-shot tier on the basis of a learned per-query retrieval-utility prediction. Closes the retrieval-amplification regression on misconception-bait benchmarks. See \Cref{sec:ruc}.

\item[Retroactive re-verification.]\nopagebreak[4] The cycle-boundary pass that re-scores every stored entry under the recalibrated verifier, prunes entries whose recalibrated stored confidence drops below the retention threshold $\tau_{\text{retro}} = 0.50$, and writes the recalibrated confidence back to the memory store. Skipped on aborted cycles under the asymmetric-rollback contract. See \Cref{sec:retroverify}.

\item[Safety override.] The router rule that dispatches the query off the score-formula path whenever the pre-routing confidence falls below the registered safety floor $\phi_{\text{safe}} = 0.38$. Post-classifier, the retrieval-utility classifier intercepts the safety-veto fall-through and chooses between the zero-shot and retrieval-augmented tiers. See \Cref{sec:router}.

\item[Seed pool.] The one-thousand-sample-per-training-benchmark slice used to populate the cold-start memory through the cold-start seeding pass at the registered cold-start storage threshold of $0.50$. See \Cref{sec:setup}.

\item[Self-improvement loop.]\nopagebreak[4] The cycle-boundary fine-tune that draws its training pool from stored memory entries whose calibrated probability composite clears the SIL admission threshold $\tau_{\text{train}} = 0.70$, applies the parameter-efficient update through the bounded low-rank adapter envelope on the frozen backbone, and re-tests the post-fine-tune model against the retention guard. See \Cref{sec:sil}.

\item[Stochastic equilibrium.] The calendar-timeline refinement of the bounded-band stationarity result that reads the realised commit fraction $\pi_{\mathrm{commit}}$ off the per-cycle commit pattern. The realised pattern $(1, 1, 1, 0, 1)$ yields $\pi_{\mathrm{commit}} = 4/5$. See \Cref{cor:stochastic-equilibrium}.

\item[Stored confidence (\boldmath$u_{\text{stored}}$).] The calibrated probability composite's output for a generated answer: the per-signal isotonic curves and the boost-layer logistic regression aggregate the nine verifier signals into a single scalar in the unit interval, against which the storage cut-point is read. See \Cref{sec:composite}.

\item[Stream chunk.] The per-cycle pre-routing pass over a fresh slice of training-benchmark samples (two thousand FEVER, two thousand TriviaQA, seven hundred CommonsenseQA per cycle) that accumulates the per-cycle storage admissions, deferred entries, abstentions, and discards. See \Cref{sec:setup}.

\item[Test fold.] The held-out slice reserved for the post-trajectory comparative panel. Distinct from the evaluation fold against which the per-cycle headline metrics are reported. See \Cref{sec:setup}.

\item[Training-pool admission threshold (\boldmath$\tau_{\text{train}}$).] The second operating point above the storage cut-point at which a stored memory entry is admitted into the self-improvement training pool. Set at $0.70$ and held strictly greater than the storage cut-point $\tau_{\text{store}} = 0.60$ so the SIL pool is a strict refinement of the storage pool. See \Cref{sec:sil}.

\item[Verifier-balanced-accuracy precondition.] The pre-specified floor on the verifier's balanced accuracy at the locked storage cut-point: the precondition of the purification condition (\Cref{thm:bayes-purity}) on which the calibration-consistent precision floor of \Cref{thm:purity} relies, audited at cycle zero through the composite Cohen's $d$ gate at the registered floor of $0.20$. See \Cref{sec:adj-gate}.

\item[Verifier ensemble.] The nine-signal post-generation block that converts a generated answer into the four-family signal vector (internal calibration, sample-set agreement, external grounding, question-answer relevance) feeding the calibrated probability composite. See \Cref{sec:verifier}.

\item[Zero-shot tier (Tier 2).] The serving tier that runs a generation on the fine-tuned backbone without retrieving passages, dispatched when the similarity to the nearest stored entry clears the registered Tier 2 threshold but the combined dispatch score does not clear the Tier 1 threshold. See \Cref{sec:router}.

\end{description}

\section{Acronyms}

\begin{description}[style=nextline, leftmargin=0pt, labelindent=0pt, itemsep=6pt, parsep=0pt, font=\bfseries, beginpenalty=10000]

\item[Bootstrap CI.] Bootstrap confidence interval. The percentile-method confidence interval on the per-sample exact-match-difference vector at the matched-paired sample size of three hundred. See \Cref{sec:sig-bootstrap}.

\item[CES.] Composite evaluation score. See the concept entry above.

\item[CHM.] Composite hallucination metric. See the concept entry above.

\item[Cohen's $d$.] An effect-size measure for the standardised mean difference between two distributions. The cycle-zero validation gate reads the composite Cohen's $d$ between exact-match-correct and exact-match-incorrect outputs against the registered precondition floor of $0.20$. See \Cref{sec:adj-gate}.

\item[EM.] Exact match. The registered correctness metric. On TruthfulQA the exact-match rule is the Haiku-judged correctness rule throughout the empirical chapter; on every other benchmark the rule is the standard span match. See \Cref{sec:setup-metrics}.

\item[FAISS.] The dense vector similarity-search index used to host the memory store, the Wikipedia passage substrate, and the cross-encoder pair index. See \Cref{sec:memory-store}.

\item[Haiku-judged.]\nopagebreak[4] The Claude Haiku 4.5 large-language-model judge used as the truthfulness adjudicator on the TruthfulQA evaluation fold. The judge reads the question, the model's answer, the correct reference answers, and the incorrect reference answers, and returns a binary truthful-or-not verdict. The uniform Haiku-judged rule is applied to every TruthfulQA exact-match cell in the empirical chapter, including every baseline. See \Cref{sec:setup-metrics}.

\item[Holm-Bonferroni.] The step-down multiple-comparison correction applied within each benchmark family of size $m = 8$ on each metric axis. Sorts the raw $p$-values ascending and multiplies the $k$-th smallest by $m - k + 1$ under a running-maximum monotonicity constraint, clipping at one. See \Cref{sec:sig-holm}.

\item[LoRA.] Low-rank adaptation. The parameter-efficient fine-tuning technique that injects trainable low-rank matrices into the seven linear projections of every transformer block of the frozen base generator at rank $32$ and alpha $64$. See \Cref{sec:sil}.

\item[McNemar test.] The paired significance test for binary outcomes, applied on the exact-match axis. Reads only the two discordant cells of the matched-pair contingency table under Edwards' continuity correction. See \Cref{sec:sig-mcnemar}.

\item[MMLU.] Massive multitask language understanding. The fifty-seven-subject multiple-choice benchmark used as one of the three probes in the multi-modal retention probe panel. See \Cref{sec:retention}.

\item[NLI.] Natural language inference. The canonical three-class textual-entailment task (entailment, contradiction, neutral) in the literature; the verifier's grounding signal family consumes the binary supported-versus-unsupported MiniCheck variant of this task, with the contradiction class structurally unavailable under the deployed backend. See \Cref{sec:verifier}.

\item[RUC.] Retrieval-utility classifier. See the concept entry above.

\item[SBERT.] Sentence-BERT. The siamese-network-fine-tuned sentence encoder that produces the 768-dimensional embeddings consumed by the memory-store lookup, the cycle-boundary consolidation pass, and the cold-start seeding pass. See \Cref{sec:memory-store}.

\item[SIL.] Self-improvement loop. See the concept entry above.

\item[Wilcoxon signed-rank.] The paired significance test for continuous outcomes, applied on the composite-hallucination-metric axis. Reads the signed-rank sum of the per-sample architecture-minus-baseline composite-hallucination-metric difference vector. See \Cref{sec:sig-mcnemar}.

\end{description}
  % Glossary
% ============================================================================
% Appendix G — Architecture positioning: frequently asked questions
% Status: COMPLETE.
%
% Scope: this appendix answers the comparative-positioning question that
% reduces the entire architecture to its simplest skeptical form: "if we
% are willing to spend an exact-match-labelled batch every cycle anyway,
% what does CAEM contribute that plain fine-tuning on the same batch
% would not?". The body of the thesis carries the full architectural
% specification in formal academic register; this appendix lays out the
% six load-bearing differentiators that justify the architecture even
% under that simplest skeptical framing.
% ============================================================================

\chapter[Architecture positioning]{Architecture positioning, frequently asked questions}
\label{app:positioning-qa}

The body chapters specify what CAEM does and how. This appendix collects the comparative-positioning and design-justification questions a reviewer in an adjacent area would ask, and answers each by naming the specific mechanism CAEM contributes that the alternative would lack. The first part answers \emph{why CAEM rather than the simpler alternative of plain exact-match-labelled fine-tuning at the same labelling budget} along the six load-bearing differentiators. The second part answers \emph{why the calibration discipline and cycle-boundary ordering chosen, rather than the straightforward alternatives}.

\part*{Part I: Positioning against plain fine-tuning}

% ---------------------------------------------------------------------------- %

\needspace{4\baselineskip}
\section[Per-query precision contract]{What does the per-query precision contract give a user that plain fine-tuning cannot?}
\label{app:qa-precision-contract}

Plain fine-tuning at the same labelling budget produces a model with an average exact-match accuracy on the held-out evaluation set; the deployed system serves whatever the model generates, with no per-query confidence axis. CAEM's calibrated-probability composite paired with the fixed-threshold storage gate produces a different deliverable: at every query the system outputs a calibrated probability $\ustored$ that the served answer is correct, and at the storage threshold $\tau_{\text{store}} = 0.60$ the cycle-zero calibration fold's empirical precision projects forward as a one-sided lower bound on the next cycle's storage precision under exchangeability between the calibration fold and the cycle-$t+1$ candidate pool; the per-cycle composite refit absorbs distribution drift into the signal-to-probability mapping so that the same threshold continues to address the same precision target as the generator improves.

The user-facing consequence is the abstain class. A query whose composite is below the deferral threshold and whose grounding is below the abstention floor returns an explicit refusal rather than a guessed answer. Plain fine-tuning has no analogous behaviour; the model serves an answer to every query because it has no machinery for deciding when not to answer. For deployments where confident-but-wrong answers carry a higher cost than refusals (clinical question answering, legal document review, customer-facing knowledge serving), the abstain class is the architectural deliverable that plain fine-tuning structurally cannot produce regardless of how much the model is fine-tuned. Chatbots have been documented fabricating every reference offered for a clinical explanation \cite{alkaissi2023artificial}, up to 91.4\% of chatbot-generated references in systematic-review searches are hallucinated \cite{info:doi/10.2196/53164}, and TruthfulQA's inverse-scaling result \cite{lin-etal-2022-truthfulqa} establishes that scaling alone does not close this regime; only an architectural intervention with a per-query confidence axis does.

% ---------------------------------------------------------------------------- %

\needspace{4\baselineskip}
\section[Memory-direct tier at scale]{The memory-direct tier saves inference cost; does that saving survive at small deployment scale?}
\label{app:qa-memory-direct}

A memory-direct query serves an answer through a sentence-embedding cosine search over the FAISS-backed episodic memory and returns the stored answer at retrieval-only cost. A zero-shot generation through the same generator pays the full forward-pass cost. A retrieval-augmented generation pays the FAISS lookup cost plus the cross-encoder reranker cost plus the forward-pass cost over a longer context. The per-tier cost ratio at the registered hardware envelope is reported in \Cref{sec:tier-cost}; the memory-direct path is meaningfully cheaper than either generation path.

The cost saving is not constant across scale. At one query per hour the per-query saving is operationally invisible; at the deployment scale projected in \Cref{sec:econ-deployment} the saving compounds across millions of queries. The architectural value of the memory-direct tier is therefore conditional on deployment scale, but it is also conditional on the fraction of queries the memory-direct tier can serve confidently, a fraction that grows monotonically across cycles as the verifier admits more verified episodes (\Cref{thm:purity}). Plain fine-tuning has no equivalent mechanism: every query pays the full forward-pass cost regardless of how many times the question has been seen, because the knowledge is baked into the weights rather than retrievable as a separable episode. At any scale where the memory-direct hit rate exceeds the architectural overhead of the verifier ensemble, the cost saving is binding; the locked configuration is engineered to enter that regime by approximately Cycle 5.

% ---------------------------------------------------------------------------- %

\needspace{4\baselineskip}
\section[Asymmetric rollback vs manual backup]{What does asymmetric rollback give the operator that a manual checkpoint backup does not?}
\label{app:qa-asymmetric-rollback}

The retention guard at every cycle boundary measures the post-fine-tune model on a three-probe multi-modal panel (MMLU, TriviaQA-test, CommonsenseQA-test) and computes the per-probe retention ratio $\rho_{k} = \mathrm{score}_{k,\text{post}} / \mathrm{score}_{k,\text{pre}}$ for each probe $k$. When any worst-probe ratio falls below the registered floor of $0.93$ the cycle is declared aborted: the weights revert to the previous cycle's parameters, and the new memory written during the cycle is preserved unchanged. The architectural property is that the model is replaceable but the memory is not.

Plain fine-tuning has no equivalent guarantee even when the operator manually keeps the previous checkpoint. The reason is that plain fine-tuning at any cycle commits both the weights and the gradient signal that produced them; rolling the weights back rolls back the cycle's accumulated training experience entirely. CAEM's separation of weights from memory means a failed cycle costs only the cycle's fine-tune compute but does not lose the verified episodes that the cycle accumulated. The next cycle inherits the previous cycle's weights together with the larger memory, and the retroactive re-verification pass at the next cycle boundary refreshes the stored composites under the rolled-back generator. This asymmetric-rollback property is the architectural anchor for the memory-growth invariant established in \Cref{cor:stochastic-equilibrium}: memory accumulates strictly across cycles even when weights do not. A weights-only architecture has no monotonicity property to anchor.

% ---------------------------------------------------------------------------- %

\needspace{4\baselineskip}
\section[Sample efficiency]{If both architectures need exact-match-labelled data at the cycle boundary, how does CAEM use ten to a hundred times fewer labels?}
\label{app:qa-sample-efficiency}

Plain fine-tuning consumes exact-match-labelled samples as the gradient signal. The number of labelled samples drives the size of the supervised-learning fine-tune; ten thousand to one hundred thousand labelled samples per fine-tune is the conventional supervised-learning budget. CAEM consumes exact-match-labelled samples for a different purpose: the labelled batch at every cycle boundary validates the empirical-precision anchor on the storage class. Approximately fifteen hundred labelled samples per cycle (five hundred per training benchmark) is sufficient for stable per-signal isotonic curves, stable per-benchmark composite boost weights under the shrinkage prior, and a stable empirical-precision reading at the registered fixed cut-points; the labels are not driving the gradient signal but anchoring the precision-floor diagnostic.

The mechanism that produces the order-of-magnitude difference is the verifier ensemble. The verifier runs label-free on every query and admits to memory only the high-confidence subset of traffic; in a production deployment, the cycle boundary's labelled batch is a stratified sample of approximately three thousand entries (about seven hundred and fifty per decision class) drawn from the memory-admission distribution rather than from the unlabelled-traffic distribution, against the fifteen hundred labelled samples per cycle used in the research run. The architecture spends labels validating the verifier rather than training the generator. Plain fine-tuning at the same three-thousand-sample budget would produce a much weaker fine-tune because three thousand labels carrying the gradient signal directly under-trains a three-billion-parameter model. The architecture's labelling budget is therefore not the same instrument as the plain-fine-tuning labelling budget; it is a verifier-precision-validation instrument that is informationally efficient by a factor of ten to a hundred.

% ---------------------------------------------------------------------------- %

\needspace{4\baselineskip}
\section[What the composite catches]{What does the composite catch at inference time that a fine-tune cannot fix?}
\label{app:qa-composite-catches}

A fine-tune updates the generator's weights so that the model is more likely to produce the correct answer on questions it has seen. The fine-tune does not alter the model's behaviour on a specific in-flight query that is hallucinating because of a property the training data did not cover: a fresh trivia question on a recently-changed world fact, an off-topic-but-grounded answer that drifts from the question's intent, a confidently-wrong output where the model has a strong prior toward a misconception. These failure modes are not closeable by parametric updates because they are first-order properties of the model's prior, not properties the training data instances reveal.

The calibrated-probability composite catches each of these failure modes at inference time through a different signal family. Confidently-wrong outputs trip the early-exit conjunction $(u_{\text{internal}} \ge 0.70) \wedge\allowbreak (p_{\text{ground,max}} \le 0.20)$, route to retrieval-augmented regeneration, and avoid serving the hallucinated answer. Off-topic-but-grounded outputs are caught by the question-answer relevance signal $q_{\text{a,rel}}$, which scores the served answer against the question through a separate cross-encoder that the generator does not influence. Common-misconception outputs that match the model's prior are caught by the sample-set agreement family: the chain-to-answer entailment probability $p_{\text{entail}}$, the average pairwise similarity $s_{\text{avg}}$ over $M=3$ resampled chains, and the MC-dropout dispersion signal $u_{\text{dropout}}$ over $K=5$ stochastic forward passes (the originally-registered semantic-entropy signal $h_{\text{norm}}$ was dropped before the main run because its learned weight in the cycle-zero composite fit was effectively zero while its ten extra samples were the verifier's largest cost). The verifier's value at inference time is structurally orthogonal to the value that fine-tuning provides; once the fine-tune saturates against the parametric capacity of the generator, the architecture continues to add value because every query is still re-checked through the verifier ensemble before serving. \cite{xiong2023can} report that verbalised confidence distinguishes correct from incorrect outputs only weakly (an area under the curve of about $0.6$ for GPT-4) and that no single confidence-elicitation method is consistently best across tasks; the composite therefore aggregates nine signals across four families through per-signal isotonic regression and a regularised logistic boost layer to produce a calibrated correctness probability that does not depend on any one signal.

% ---------------------------------------------------------------------------- %

\needspace{4\baselineskip}
\section[Tier-2 distillation]{What does the zero-shot tier gain from the cycle loop that plain fine-tuning cannot match?}
\label{app:qa-tier2-distillation}

A plain fine-tune produces a single static checkpoint whose zero-shot output distribution does not change between deployments. The cycle loop produces a sequence of checkpoints whose zero-shot output distribution evolves under a specific architectural pressure: at each cycle boundary the self-improvement loop draws its training pool from stored episodes whose answers were originally produced by the retrieval-augmented tier and were certified by the nine-signal verifier, so the fine-tune folds verified retrieval-augmented reasoning into the generator's parametric capacity. The next cycle is therefore able to answer a measurable share of the previously retrieval-dependent queries from the zero-shot tier alone, without paying the retrieval cost. The deployment-side reading of this mechanism is a tier-share migration over cycles: the zero-shot tier's share of the served stream grows, the retrieval-augmented tier's share contracts, and the per-query cost of the workload falls accordingly.

Plain fine-tuning has no analogous mechanism because it has no separate retrieval-augmented tier whose certified outputs to distil from; the training data for a plain fine-tune is either pre-collected and static or it accumulates without the verifier-gated quality filter that the CAEM cycle imposes. The empirical receipt across the realised trajectory in \Cref{ch:evaluation} is the per-tier hit-rate trajectory over the four committed cycles (C1, C2, C3, C5; C4 aborted under the retention guard), reported in \Cref{sec:tier-hit} alongside the retrieval-utility classifier ablation of \Cref{sec:ruc-ablation} that decomposes the share migration into the self-improvement distillation contribution and the per-query routing refinement contribution.

% ---------------------------------------------------------------------------- %

\needspace{4\baselineskip}
\section[When not to deploy]{When would CAEM not be worth deploying?}
\label{app:qa-when-not}

The architecture is not a free contribution. Each of the six differentiators above corresponds to an architectural component whose cost is non-zero: the verifier ensemble pays forward passes through MiniCheck, the cross-encoder, and the sentence encoder at every query; the labelled cycle-boundary batch requires ongoing labelling effort; the cycle policy requires operational discipline. The architecture is worth deploying exactly when at least one of the six differentiators is binding for the deployment, and is not worth deploying when all six relax simultaneously:

\begin{description}
\item[No abstain class.] Deployments with no requirement for an abstain class.
\item[No scale compounding.] Deployments where memory-direct cost savings do not compound (very low traffic, or a product requirement that every query traverse the same path for uniformity).
\item[Distillation traction absent.] Deployments where the zero-shot tier's per-cycle improvement does not bind because the deployed workload distribution is too narrow for the SIL-driven distillation to gain traction, or too broad for a 3B-class generator to absorb at any cycle horizon.
\item[Manual checkpoint sufficient.] Deployments where a manual previous-checkpoint backup is sufficient and an external snapshot pipeline already solves the catastrophic-forgetting recovery problem.
\item[Labelling abundant.] Deployments with a labelling budget large enough to support frequent direct fine-tuning, where the verifier-as-sample-efficiency-multiplier value collapses because labelling cost is not the binding constraint.
\item[Below parametric ceiling.] Deployments below the parametric-capacity ceiling of the underlying generator, where direct fine-tuning still produces meaningful accuracy gains so the inference-time-error-catching value of the composite is dominated by what fine-tuning would also provide.
\end{description}

A deployment that satisfies all six conditions simultaneously is one where plain fine-tuning at the same labelling budget would deliver comparable user-facing quality with simpler implementation. A deployment that fails any one of these conditions is one where the corresponding architectural component is binding. The thesis's value proposition is the regime where at least one of the six conditions binds; the empirical receipts in \Cref{ch:evaluation} measure the contribution of each component on the registered five-benchmark workload, where every component is binding by construction of the failure-mode taxonomy in \Cref{sec:hallucination-taxonomy}.

\part*{Part II: Design-justification: calibration discipline and cycle-boundary ordering}

The architecture makes several non-obvious design choices in how the calibration fold is constructed, how the cycle boundary is sequenced, and how the deduplication discipline guards the disjointness of the data partition. Part~II answers each as a direct ``is this correct, and why not the simpler alternative'' question.

% ---------------------------------------------------------------------------- %

\needspace{4\baselineskip}
\section[Same calibration fold across cycles]{Why use the same calibration fold across all cycles of the registered horizon rather than rotating folds?}
\label{app:qa-same-fold}

The calibration fold is a fixed thousand-five-hundred-sample partition (five hundred samples drawn from each of the three training benchmarks) constructed once at pool-builder time and reused at every cycle boundary. The natural alternative would be to rotate the fold across cycles, drawing a fresh fifteen-hundred-sample slice each cycle from a larger reservoir. The same-fold protocol is preferred for four reasons.

First, the cycle-boundary recalibration is itself a measurement. The labelled fold is the reference distribution against which the verifier's calibration is read; if the reference distribution varies between cycles, then a cycle-to-cycle change in the fitted curves entangles two sources of variation (the verifier improving under self-improvement, and the reference set drifting) and the trajectory analysis cannot decompose them. A fixed reference fold is the variance-isolation lever that lets the per-cycle isotonic and composite-calibration readings be attributed cleanly to verifier change.

Second, the gold labels on the fold are constant across cycles. The model's verifier signals on a given fold question shift cycle-to-cycle (because the post-fine-tune model produces different token probabilities, dropout variances, semantic-entropy values, and so on), but the question's correct answer does not. Re-fitting the per-signal isotonic curves on the same (variable signals, constant labels) data keeps the calibrator matched to the current model, which is the regime in which the post-hoc calibration literature finds such calibration reliable; calibration fitted before a model or data shift degrades after it \cite{ovadia2019trust}. Rotating the fold would replace the constant-labels property with a moving-target labelling distribution that confounds the fit.

Third, the fixed-threshold empirical-precision anchor holds against the same-fold protocol. The anchor requires that the calibration-fold composite distribution at the registered cut-point projects forward as a one-sided lower bound on the cycle-$t+1$ candidate pool's storage precision under exchangeability between those two distributions. The protocol does \emph{not} require exchangeability between cycle $N$'s calibration fold and cycle $N+1$'s calibration fold; the same-fold choice across cycles is therefore not a contract violation, because the per-cycle composite refit absorbs any drift in the underlying signal distribution into the signal-to-probability mapping rather than into the threshold.

Fourth, the storage cut-points themselves are held constant across cycles at $\tau_{\text{store}} = 0.60$ and $\tau_{\text{defer}} = 0.45$. With the same fold every cycle and the same cut-points across cycles, any cycle-to-cycle variation in the realised storage precision at the cut-point is attributable to the per-cycle composite refit absorbing distribution drift, rather than to fold-resampling variance or to the cut-points themselves moving.

A defensible alternative for a substantially longer trajectory (say fifty or more cycles) would be a rotating-fold protocol drawn from a five-thousand-sample reservoir partitioned into ten overlapping subsets, where each cycle's fold differs but is still an honest sample from the labelled population. The variance-isolation argument weakens at long horizons because the same-fold protocol risks calibration over-fitting to fold-specific noise. The realised six-cycle horizon (the cycle-zero calibration plus cycles one through five, with cycle four aborted under the retention guard and cycle five closing the trajectory under the equilibrium-saturation gate) is short enough that variance isolation dominates, and the rotating-fold extension is registered as a future-work item should the trajectory horizon extend.

% ---------------------------------------------------------------------------- %

\needspace{4\baselineskip}
\section[Production vs research fold]{In production each calibration split would necessarily differ; how is this reconciled with the same-fold research protocol?}
\label{app:qa-production-fold}

In production there is no pre-built calibration partition. The labelled batch at every production cycle boundary is drawn from accumulated serving traffic via stratified sampling (per the production runbook), so each cycle's fold is a fresh sample of three thousand or so labelled queries with content that varies cycle-to-cycle. This is a different operational mode from the research protocol but architecturally compatible with it.

Three points reconcile the two modes. First, the fitting procedure is fold-agnostic. The per-signal isotonic regression is a deterministic monotone fit on whatever (signal, label) pairs are supplied; sklearn's isotonic regression does not care whether the questions were the same as in a previous fit or different. The cycle-zero cut-point selection walks the labelled fold sorted by composite and picks the threshold at the registered precision target; this is run once at cycle zero and the resulting cut-points are held constant across cycles. Both procedures work identically on a fixed fold (research) or a rotating fold (production).

Second, production loses one variance-isolation property and gains a different fidelity property. The research protocol uses a fixed fold to isolate verifier-distribution change as the only source of cycle-to-cycle calibration variation. Production cannot use the fixed fold because the fold material does not exist. Instead, production gains a different property: the calibration fold matches the live deployment distribution exactly, because it is drawn from live deployment traffic. Research's fixed fold is a proxy for deployment (drawn from training-benchmark dev splits); production's rotating fold is the real thing. Each protocol is correct under its own labelling-availability constraint.

Third, the EMA smoothing on the thresholds carries forward the previous cycle's threshold value into the next cycle's fit at weight $\alpha_{\text{ema}} = 0.7$. In production this means a single cycle's traffic anomaly (an unusual spike of out-of-domain queries, for example) cannot drag the threshold to a degenerate value, because the next cycle's labelled batch will return the threshold trajectory toward the running average. The same EMA mechanism that prevents same-fold over-fitting in research prevents rotating-fold instability in production.

The practical consequence: an operator who wants production calibration to behave as close as possible to research-mode discipline can stratify the production sample to be representative of the five-benchmark workload composition and can hold the labelled batch size near the research-mode fifteen-hundred-sample point. An operator whose live distribution differs sharply from research can let the production fold reflect the live distribution and accept that the calibration trajectory will diverge from the research trajectory by exactly the distribution-shift magnitude. Both behaviours are correct; the choice is between distribution fidelity and protocol parity.

% ---------------------------------------------------------------------------- %

\needspace{4\baselineskip}
\section[Cycle-boundary ordering]{Why apply the cycle-boundary recalibration before retroactive re-verification rather than after?}
\label{app:qa-cycle-order}

The cycle-boundary order is: self-improvement fine-tune, then score the calibration fold under the now-updated generator, then re-fit the temperature scalar and per-signal isotonic curves and per-benchmark composite boost weights on the freshly-scored fold (the storage cut-points themselves are held constant across cycles), then reload the verifier from the new calibration files, then run retroactive re-verification against the now-recalibrated verifier. The natural alternative would be to run retroactive re-verification first under the previous cycle's calibration and recalibrate afterwards. The current order is the architecturally correct one for two reasons.

First, the post-fine-tune signal distribution falls outside the previous cycle's isotonic fit envelope on a non-trivial fraction of stored episodes. The fine-tune saturates token probability and collapses semantic entropy on chains the model trained on; the previous cycle's isotonic curves were fitted on the previous cycle's signal distribution and do not extrapolate cleanly past the upper edge of their training-input range. Running retroactive re-verification under the previous cycle's calibration produces calibrated probabilities on memorised chains that are artefactually low, and the resulting composite drops below the prune threshold for many EM-correct stored entries. The empirical evidence for this artefact is documented in the original cycle-one prune of a large share of the cold-seed memory under the stale-calibration ordering; the corrected order has the recalibration absorb the post-fine-tune signal-distribution shift before the prune rule fires.

Second, the corrected order preserves the architectural intent of retroactive re-verification. The re-verification pass is meant to upgrade the stored composite of episodes whose underlying correctness has improved under the new generator and to prune episodes whose correctness has genuinely degraded. The pass cannot do either if the calibration it reads through has not yet been updated to the new model. Under the corrected order, retroactive re-verification reads through fresh per-signal isotonic curves and fresh per-benchmark boost weights at the registered fixed cut-points; an entry whose correctness genuinely improves under the new model maps to a higher calibrated probability and stays in memory with an upgraded $u_{\text{stored}}$, an entry whose correctness genuinely degrades maps to a lower probability and gets pruned. The order is therefore not an implementation detail but a precondition on the theorem-licensing receipt for pool purity in \Cref{thm:purity}.

% ---------------------------------------------------------------------------- %

\needspace{4\baselineskip}
\section[Early-exit disabled at retro-verify]{Why disable the confabulation early-exit during retroactive re-verification when it remains active at query time?}
\label{app:qa-early-exit-retro}

The early-exit rule fires when the internal-confidence composite is high (above $\eta_u = 0.70$) and the top-passage entailment is low (below $\eta_g = 0.20$). At query time this conjunction is the architectural defence against confabulation: a model that is internally confident on a fresh question but has no supporting passage is generating a confident hallucination, and the rule routes the answer to retrieval-augmented regeneration rather than serving the hallucination. At retroactive re-verification time the same conjunction reads differently. The fine-tune just memorised the stored chains, so the internal-confidence composite is automatically near one for every entry the self-improvement loop touched; the first conjunct is no longer informative. Letting the rule fire under this regime collapses the conjunction into ``low top-passage entailment, regardless of confidence,'' which is a much wider catchment than the registered confabulation rule and over-prunes EM-correct stored entries whose chain happened to have thin passage support.

The architectural fix is a query-time-only flag on the rule. The verifier's $\mathrm{verify}$ method takes an $\mathrm{is\_query\_time}$ parameter that defaults to true (preserving query-time inference and deferred-buffer reconsideration behaviour) and is set to false in the retroactive re-verification closure. Under the false setting the early-exit conjunction is short-circuited and the verifier falls through to the standard composite-and-threshold path, which still prunes any entry whose recalibrated $u_{\text{stored}}$ falls below $\tau_{\text{retro}} = 0.50$. Confidently-wrong stored entries are still removed; EM-correct stored entries with thin passage support are no longer falsely pruned. The rule's intended semantics (catching query-time confabulation) is preserved; its inheritance into retroverify-time is corrected.

% ---------------------------------------------------------------------------- %

\needspace{4\baselineskip}
\section[Training vs storage threshold]{Why is the training threshold strictly higher than the storage threshold rather than equal?}
\label{app:qa-tau-train-vs-store}

The architecture registers two different thresholds on the calibrated composite. The storage threshold $\tau_{\text{store}}$ governs admission to the user-facing memory store and to the retrieval tier. The training threshold $\tau_{\text{train}}$ governs admission to the next cycle's self-improvement fine-tune pool; only entries whose composite clears $\tau_{\text{train}}$ enter the gradient signal. The convention is $\tau_{\text{train}} > \tau_{\text{store}}$, with the locked configuration setting $\tau_{\text{store}} = 0.60$ and $\tau_{\text{train}} = 0.70$ (held constant across cycles; the training-pool threshold was lowered from $0.75$ to $0.70$ on the trajectory-launch date to restore a non-empty verified-episode share while preserving the strictly-stricter invariant). The asymmetry is deliberate.

The reason is that the training pool contributes to all subsequent cycles through the generator's weight update, while the storage pool contributes only to the user-facing retrieval tier at the current cycle. A noisy entry in the storage pool serves at most a single Tier-1 retrieval and is re-scored at the next cycle's retroactive re-verification; the architectural cost of a borderline storage admission is bounded. A noisy entry in the training pool propagates through the gradient signal into every subsequent cycle's generator and can compound the verifier's bias over the trajectory. The noise-tolerant-learning literature \cite{natarajan2013learning} bounds the corruption rate at which training under noisy labels still improves a model rather than degrading it; the architecture spends the strict-training-threshold budget to keep the training pool's empirical purity above that rate by construction.

A defensible alternative would be a single threshold $\tau = \tau_{\text{store}} = \tau_{\text{train}}$ that admits the same entries to both purposes. The simpler alternative loses the asymmetry between storage-class consequences (bounded, single-cycle) and training-class consequences (compounding, multi-cycle), and the noise-tolerant-learning bound on training-pool corruption is no longer enforced architecturally. The two-threshold protocol is the registered defence; the empirical purity of the resulting training pool is reported in \Cref{ch:evaluation}.

% ---------------------------------------------------------------------------- %

\needspace{4\baselineskip}
\section[Three-layer dedup]{Why deduplicate the calibration fold at three layers rather than relying on the dataset's native unique identifiers?}
\label{app:qa-three-layer-dedup}

The deduplication discipline runs at three layers in the pool-builder. The first layer hashes every sample's question text into a sixteen-character content identifier and dedupes within each loaded dataset split, dropping ten thousand or more duplicate records on average per training-benchmark train split. The second layer dedupes within each individual partition slice (calibration, purity, training, evaluation) to catch any residual within-slice duplicates. The third layer cross-checks the calibration and purity slices against the evaluation slice and drops any sample whose content hash collides with an evaluation-slice sample. The natural alternative would be to rely on the dataset's native sample identifiers and assume disjointness by construction.

The dataset's native identifiers are unreliable for two reasons. The lucadiliello FEVER dump reuses integer claim identifiers across the train and dev splits, so a sample with id $13957$ in the train split is content-distinct from a different sample with id $13957$ in the dev split, and a join on native id produces spurious overlap warnings. The TriviaQA dump reuses sfq-prefixed and qw-prefixed and qz-prefixed identifiers within a single split, so two content-identical questions can both appear in the same split with the same id. Relying on native identifiers therefore both over-flags overlap (false positives on different-content same-id pairs) and under-flags overlap (false negatives on same-content different-id pairs). Content-hashing on the question text resolves both regimes by definition: identical questions hash to the same identifier regardless of native id; different questions hash to different identifiers regardless of native id collision.

The three-layer structure is the minimum sufficient cover. A single dataset-level dedup pass leaves intra-slice duplicates untouched. A single intra-slice dedup pass leaves cross-slice content collisions untouched (a calibration question that happens to also appear in the evaluation split). A single cross-slice pass relies on the per-slice slice having been itself dedup-clean; without the within-slice dedup, the cross-slice check sees inflated counts and miscomputes the disjointness invariant. The empirical evidence for the three-layer requirement is reported in the loading-time logs of every cycle: each layer drops a non-trivial count, and the fold size after all three layers matches the registered fifteen-hundred-sample target rather than a stale-protocol nineteen-hundred-or-more sample inflated count.
  % Architecture positioning -- FAQ
% Appendix H (Proofs) reserved for any individual proof that overflows the
% half-page-in-Chapter-4 budget set during the rewrite. None have overflowed
% in Phase 1a; all proofs are inline in Chapter 4 §4.3.
% \input{appendix/appendix_h.tex}  % Proofs (created only if overflow occurs)

\end{document}
